\documentclass[UTF8,zihao=-4,openany,fontset=fandol]{ctexbook}

\usepackage[a4paper,top=2.35cm,bottom=2.35cm,left=2.55cm,right=2.55cm]{geometry}
\usepackage{amsmath,amssymb,mathtools}
\usepackage{booktabs,longtable,tabularx,array}
\usepackage{enumitem}
\usepackage{tikz}
\usetikzlibrary{arrows.meta,calc,positioning,shapes.geometric}
\usepackage{xcolor}
\usepackage{listings}
\usepackage{hyperref}
\usepackage{bookmark}

\hypersetup{
  colorlinks=true,
  breaklinks=true,
  linkcolor=blue!55!black,
  urlcolor=blue!65!black,
  pdftitle={大规模并行处理器编程（第5版）习题详解},
  pdfauthor={习题详解编写组},
}

\setcounter{tocdepth}{1}
\setlist{itemsep=0.25em,topsep=0.35em}
\setlength{\parskip}{0.35em}
\setlength{\parindent}{2em}
\setlength{\emergencystretch}{2em}
\raggedbottom

\definecolor{solutionblue}{RGB}{35,85,135}
\definecolor{solutiongray}{RGB}{245,247,249}
\lstset{
  basicstyle=\ttfamily\small,
  columns=fullflexible,
  keepspaces=true,
  breaklines=true,
  breakatwhitespace=false,
  showstringspaces=false,
  commentstyle=\color{green!40!black}\upshape\ttfamily,
  keywordstyle=\color{blue!60!black}\bfseries,
  stringstyle=\color{red!55!black},
  frame=single,
  rulecolor=\color{black!25},
  numbers=left,
  numberstyle=\scriptsize\color{gray},
  numbersep=6pt,
  xleftmargin=2em,
  framexleftmargin=2em,
  aboveskip=0.7em,
  belowskip=0.7em,
}

\newcommand{\code}[1]{\allowbreak\texttt{\detokenize{#1}}\allowbreak}
\newenvironment{solutionexercise}[1]{%
  \section{第 #1 题}\label{sol:\thechapter:#1}%
}{\par\medskip}
\newcommand{\problemheading}{\par\noindent\textbf{题目}\quad}
\newcommand{\approachheading}{\par\noindent\textbf{解题思路}\quad}
\newcommand{\answerheading}{\par\noindent\textbf{详细解答}\quad}
\newcommand{\conclusionheading}{\par\noindent\textbf{结论}\quad}
\newcommand{\discussionheading}{\par\medskip\noindent\textbf{进一步讨论}\par\nopagebreak[3]}
\newcommand{\discussionpoint}[1]{\par\smallskip\noindent\textbf{#1}\quad}
\newcommand{\variantheading}{\par\noindent\textbf{变体练习与答案要点}\quad}
\newcommand{\practiceheading}{\par\noindent\textbf{工程实践}\quad}
\newcommand{\pitfallheading}{\par\noindent\textbf{常见错误}\quad}
\newcommand{\verificationheading}{\par\noindent\textbf{联网核验与对抗审查}\quad}
\newcommand{\sourcecheck}[2]{\verificationheading\href{#1}{[来源]}\ #2。}
\newcommand{\editorialnote}{\par\noindent\textbf{题解编者注：}}
\newcommand{\complexity}[3]{%
  \par\noindent\textbf{复杂度与并行性}\par
  \noindent\textbf{时间复杂度：}#1。\par
  \noindent\textbf{空间复杂度：}#2。\par
  \noindent\textbf{并行性：}#3。}
\newcommand{\noexercises}{\par\bigskip\begin{center}\fbox{\begin{minipage}{0.82\textwidth}\centering\vspace{0.7em}\textbf{本章无习题}\vspace{0.7em}\end{minipage}}\end{center}}

\begin{document}

\frontmatter
\begin{titlepage}
  \centering
  \vspace*{2.6cm}
  {\Huge\bfseries 大规模并行处理器编程（第 5 版）\par}
  \vspace{0.9cm}
  {\LARGE\bfseries 习题详解\par}
  \vspace{1.5cm}
  {\large 依据中文学习译稿与英文底本逐题编写\par}
  \vfill
  {\small 独立学习资料；不替代原书正文\par}
  \vspace{0.6cm}
  {\small \today\par}
\end{titlepage}

\chapter*{使用说明}
\addcontentsline{toc}{chapter}{使用说明}
本题解按原书第 1--24 章及附录 A--C 的顺序编排。每题均给出题目、解题思路、详细解答、
结论、进一步讨论、变体练习与答案要点以及工程实践；算法题另说明复杂度与并行性。题目
文字以仓库中的中文译稿为主要来源，并与英文底本练习页逐项核对。底本或译稿中需要修正
之处以“题解编者注”明确说明。

CUDA 代码以 CUDA C++ 为准。代码中的 CUDA API 错误检查、内核启动错误检查、同步和 CPU
参考验证用于说明完整程序结构；具体性能数字仍依赖 GPU 架构、CUDA 版本和编译选项。
可执行版本位于 \code{solutions/code/}，与本 PDF 独立构建。每题的联网来源优先采用官方
文档、作者公开版和机构仓储；题解只作短引或独立转述，不以网络材料代替推导。

\tableofcontents

\mainmatter
\chapter{引言}
\noexercises


\part{基础概念}
\chapter{异构数据并行计算}

\begin{solutionexercise}{1}
\problemheading 如果希望网格中的每个线程计算向量加法的一个输出元素，那么在内核函数中，
把线程/块索引映射到数据索引 $i$ 的表达式是什么？
\begin{enumerate}[label=\alph*.,leftmargin=2.2em]
  \item \code{i = threadIdx.x + threadIdx.y;}
  \item \code{i = blockIdx.x + threadIdx.x;}
  \item \code{i = blockIdx.x * blockDim.x + threadIdx.x;}
  \item \code{i = blockIdx.x * threadIdx.x;}
\end{enumerate}

\approachheading 全局线程号由“前面完整线程块中的线程数”与“块内线程号”相加得到。

\answerheading 已知一维网格、一维线程块，\code{blockIdx.x} 是块号，\code{blockDim.x}
是每块线程数，\code{threadIdx.x} 是块内线程号。因此
\[
 i=\code{blockIdx.x}\,\code{blockDim.x}+\code{threadIdx.x}.
\]
选项 c 正确。该映射在网格内唯一且连续；实际内核还应检查 $i<N$。

\conclusionheading 选 c。
\discussionheading
\discussionpoint{全局线程号是一种混合进位表示。}
为什么块号不能直接与线程号相加？因为块 $b$ 之前已经排着 $b$ 个完整线程块，每块各含 $T$ 个位置，所以它们共同占去 $bT$ 个线性编号，当前线程再在此基础上加块内偏移 $t$。这和把二维坐标 $(r,c)$ 展平为 $rW+c$、把 MPI 进程坐标转换为 rank 是同一种混合进位思想：先确定每一位的容量，再累加高位所跨过的完整区间。推广到二维或三维网格时，还必须先声明哪一维变化最快，否则同一组坐标会被解释成不同的线性次序。索引公式因而不仅是算术技巧，也是线程布局、数据布局和调试输出共同遵守的一份契约。
\discussionpoint{网格步长循环把启动规模与数据规模解耦。}
若数组有数十亿项，既没有必要也未必能用一次网格为每项创建唯一线程；更常见的做法是先算本题的首索引，再反复增加整个网格的线程总数 $G T$。线程 $i$ 依次处理 $i,i+GT,i+2GT,\ldots$，只要每轮检查上界，所有余数类便共同覆盖数组且互不重叠。这样可以把网格限制在足以占满设备的规模，并让常驻线程摊薄启动和地址计算成本，归约前处理与持久化内核经常采用这种形式。不过循环增加了每线程工作量和寄存器生存期，若任务耗时不均，还可能把同一余数类中的慢项集中到一个线程，不能只凭覆盖正确就认定配置合理。
\discussionpoint{正确索引仍需要完整的边界契约。}
向上取整块数时，最后一个块通常包含一些没有对应数据的线程；因此 $i=bT+t$ 即使唯一且连续，也只证明映射没有重号，并不证明 $i$ 一定小于 $N$。边界谓词必须支配由 $i$ 派生的每一次读取和写入，例如先读取 \code{a[i]}、最后才检查写回仍然已经越界。规模很大时还要审查类型：若 $bT$ 先按 32 位整数计算，即使随后赋给 \code{size_t}，溢出也已经发生，安全写法应让乘法的操作数本身足够宽。这个区别在二维展平时更重要，因为 $rW+c$ 比单维索引更早越过 32 位范围。
\discussionpoint{连续编号只是良好访存的起点。}
当一个 warp 的 lane 0--31 分别访问 $i,i+1,\ldots,i+31$ 时，请求字节通常能落入少量连续内存扇区，这正是该映射常被采用的性能原因。但若数组元素变宽、基址带非对齐偏移，或逻辑相邻元素在结构体中实际相隔很远，线程号连续也不能保证事务利用率理想。验证时可以让一个小内核把每个线程算出的秩写回，与 CPU 枚举逐项对照，从而先排除覆盖错误；随后再比较请求字节与实际传输字节，判断布局是否浪费带宽。把正确性验证和事务效率分开，能避免把“答案数值正确”误当成“线程映射已经高效”。
\variantheading 若每块 256 个线程且处理 $N=1000$ 个元素，则需 4 个块、启动 1024 个线程，最后 24 个线程应被 $i<N$ 屏蔽。
\practiceheading 生产内核可用 Nsight Compute 的 Launch Statistics 核对实际网格、块维度和越界线程比例，并用 Compute Sanitizer memcheck 捕获漏写边界判断造成的越界。
\sourcecheck{https://docs.nvidia.com/cuda/cuda-c-programming-guide/}{NVIDIA CUDA C++ Programming Guide 的线程层次定义支持上述映射；高置信度}
\end{solutionexercise}

\begin{solutionexercise}{2}
\problemheading 假设希望每个线程计算向量加法中相邻的两个元素。把线程/块索引映射到线程
要处理的第一个元素的数据索引 $i$ 的表达式是什么？
\begin{enumerate}[label=\alph*.,leftmargin=2.2em]
  \item \code{i = blockIdx.x * blockDim.x + threadIdx.x + 2;}
  \item \code{i = blockIdx.x * threadIdx.x * 2;}
  \item \code{i = (blockIdx.x * blockDim.x + threadIdx.x) * 2;}
  \item \code{i = blockIdx.x * blockDim.x * 2 + threadIdx.x;}
\end{enumerate}

\approachheading 先求线程的线性全局编号 $t$，再把每个线程负责的元素数 2 乘进去。

\answerheading $t=\code{blockIdx.x*blockDim.x+threadIdx.x}$，故第一个元素
$i=2t$，第二个元素为 $i+1$。即
\[
 i=(\code{blockIdx.x*blockDim.x+threadIdx.x})*2.
\]
选项 c 正确。边界条件分别是 $i<N$ 和 $i+1<N$。

\conclusionheading 选 c；不能简单在索引末尾加 2。
\discussionheading
\discussionpoint{线程粗化是在并行度与单线程工作量之间换挡。}
每线程只算一个元素时，索引和边界判断几乎与一次加法同样显眼；让线程连续算两项，可以让这些固定指令服务更多有效工作，这就是线程粗化。它看起来像 CPU 的循环展开，却多出 GPU 特有的代价：逻辑线程数减半，一个 warp 覆盖的地址集合发生变化，每线程同时存活的输入和中间值也可能增加。对于很短的逐元素算子，减少控制指令有时能带来收益；对于本来就缺少并行任务的小数组，减半线程数反而可能让设备更空闲。因此粗化不是“每线程多做事必然更快”，而是把调度开销、并行度和局部状态重新分配。
\discussionpoint{相邻两项与分段两项具有不同的访存几何。}
本题映射让线程 $t$ 负责 $2t$ 与 $2t+1$，最直观的实现是一次读取一个对齐的 \code{float2}，此时相邻线程的向量对象在内存中仍首尾相接。若编译器生成两条标量指令，第一条却会让 lane 依次访问 $0,2,4,\ldots$，第二条访问 $1,3,5,\ldots$，每条指令的地址集合都有步长 2，事务形态便不同。分段式方案则让所有 lane 第一轮访问连续区间、第二轮整体平移一个块宽，所以不依赖向量指令也能保持每轮单位步长。两种布局数学工作相同，但其性能取决于对齐、生成的机器指令和缓存层次，不能仅凭源码中“元素相邻”四个字判断。
\discussionpoint{奇数长度暴露了粗化边界的真实难点。}
设 $N=9$，线程 4 的首索引是 8，第一项合法而第二项 9 已越界；如果用 \code{if (i+1<n)} 包住两次运算，就会错误漏掉元素 8，若只检查 \code{i<n} 又会非法访问元素 9。正确结构是分别保护两项，或者先处理所有完整二元组，再由唯一线程或独立尾部路径处理余数。这个例子虽只有一个尾元素，却代表任意粗化因子 $C$ 的共同规则：必须覆盖 $N\bmod C$ 的每个余数类，读取本身也要在谓词之内。随着 $C$ 增大，尾部失衡和寄存器占用同时增加，边界测试因而是资源调优之前不可省略的正确性门槛。
\discussionpoint{生产级粗化通常服务于减少数据搬运。}
在实际框架中，单纯省下一次索引乘法往往不值得复杂化代码；更有价值的场景是把加法、缩放和激活融合到同一线程，使刚读入的值在寄存器中完成多步处理，而不把中间数组写回全局内存。相邻式粗化还可与对齐向量加载结合，网格步长循环则让固定数量线程持续消费多个向量包，但每一种组合都必须保留标量前缀和尾部路径。评估时应先用 $N\bmod2=0,1$ 以及非对齐子数组证明结果和边界，再比较端到端字节流量与吞吐，而不是只比较内核指令数。只有数据搬运确实减少且资源占用没有抵消收益，粗化才成为可维护的工程优化。
\variantheading 若 $N=9$ 且启动 5 个逻辑线程，则线程 4 的首索引为 8，只处理元素 8；其第二个索引 9 越界。
\practiceheading 可编译运行 \code{solutions/code/ch02-09/ch02_ex02_vector_add.cu}，并在 Compute Sanitizer memcheck 下覆盖 $N\bmod2=0$ 和 $N\bmod2=1$ 两类尾部。
\sourcecheck{https://docs.nvidia.com/cuda/cuda-c-programming-guide/}{CUDA 线程索引规则；高置信度}
\end{solutionexercise}

\begin{solutionexercise}{3}
\problemheading 希望每个线程计算两个元素。每个线程块处理由两个区段组成的、连续的
$2\times\code{blockDim.x}$ 个元素。块中的所有线程先处理第一个区段，每个线程
处理一个元素；然后它们一起移动到下一个区段，每个线程再处理一个元素。假设
变量 $i$ 是线程要处理的第一个元素的索引，那么把线程/块索引映射到第一个
元素数据索引的表达式是什么？
\begin{enumerate}[label=\alph*.,leftmargin=2.2em]
  \item \code{i = blockIdx.x * blockDim.x + threadIdx.x + 2;}
  \item \code{i = blockIdx.x * threadIdx.x * 2;}
  \item \code{i = (blockIdx.x * blockDim.x + threadIdx.x) * 2;}
  \item \code{i = blockIdx.x * blockDim.x * 2 + threadIdx.x;}
\end{enumerate}

\approachheading 块 $b$ 之前已有 $2b$ 个完整区段；块内线程 $t$ 在当前区段取第 $t$ 个元素。

\answerheading
\[
 i=2\,\code{blockIdx.x}\,\code{blockDim.x}+\code{threadIdx.x}.
\]
第二个元素位于 $i+\code{blockDim.x}$。因此选项 d 正确。与第 2 题不同，这里同一块的
线程先协同处理整个区段，连续线程仍访问连续地址，利于访问合并。

\conclusionheading 选 d。
\discussionheading
\discussionpoint{分段式粗化为什么仍然合并。}
乍看之下，一个线程承担两个元素似乎会把相邻线程的地址拉开，然而关键不在“某个线程最终碰过哪些地址”，而在一个 warp 执行\emph{同一条}访存指令时形成的地址集合。设块号为 $b$、块宽为 $T$、lane 为 $t$，第一轮地址是 $2bT+t$，32 个 lane 仍覆盖连续元素；第二轮只是整体平移 $T$，地址变成 $2bT+T+t$，仍然是单位步长。两轮之间相隔多远不会破坏各轮的合并性，这与逐线程相邻映射的标量首轮地址 $2t$ 形成鲜明对比。由此可见，判断 GPU 数据布局必须沿动态指令横向观察整个 warp，而不能沿单线程时间线纵向观察。
\discussionpoint{从两个区段推广到任意粗化因子。}
若每线程处理 $C$ 个区段，块 $b$ 应先跳过前面 $b$ 个块覆盖的 $CbT$ 个元素，第 $k$ 轮再加区段偏移 $kT$ 和线程偏移 $t$，所以统一公式是 $i_{b,k,t}=CbT+kT+t$，其中 $0\le k<C$。固定 $b,k$ 后，$t$ 每增加 1，地址也只增加 1，因此任意一轮都保持连续；固定 $b,t$ 后改变 $k$，才表现为同一线程跨区段工作。这个推导可直接迁移到 SAXPY、像素通道变换或规则网格更新，因为它把“任务属于哪个块、哪个轮次、哪个 lane”三层含义明确分开。若 $C$ 是编译期常量，循环还可展开，但展开只是减少控制指令，并不改变上述覆盖证明。
\discussionpoint{尾块为何必须逐轮保护。}
危险通常藏在最后一个块：第一段存在有效元素，并不能推出第二段或更后的区段也有效。以 $T=128,C=2,N=850$ 为例，块 3 的第一轮覆盖 $768\ldots895$，只有前 82 个 lane 有效，而第二轮从 896 开始，所有 lane 都越界；若只在循环外检查一次首地址，第二轮仍会非法读取。稳妥写法是在每个 $k$ 上计算 $i_{b,k,t}$ 后独立判断 $i<N$，并让读取与写回都受这一谓词支配。对于每个元素计算量还随数据变化的任务，粗化会把多个“重元素”固定交给同一线程，尾部问题便从越界进一步演化成负载不均，这时动态工作队列或重新分桶可能比继续增大 $C$ 更合适。
\discussionpoint{粗化因子不是越大越好。}
$C$ 增大时，索引计算和块调度成本可由更多有效工作摊薄，同一线程也可能同时准备多个独立加载，从而增加指令级并行；但每多保留一个待处理值，寄存器生存期往往随之延长。寄存器过多会减少可驻留 warp，极端时还会溢出到 local memory，使原本想节省的流量反而增加，而线程总数下降也可能让小规模输入无法填满设备。实际选择可先保证 $C=1,2,4,8$ 都通过余数类与边界测试，再在目标尺寸分布上比较吞吐、实际传输字节和编译器报告的资源用量。这样的实验不仅告诉我们哪个 $C$ 更快，也能解释为何同一取值在另一代 GPU 或另一个问题规模上会反转。
\variantheading 若每块 128 线程、块号为 3，则第一区段首索引为 $2\times3\times128=768$，线程 7 处理 775 和 903。
\practiceheading 可用 Nsight Compute 的 Global Memory Load/Store Efficiency 比较分段式与逐线程连续式粗化，并同时检查两种布局在实际缓存层次上的差异。
\sourcecheck{https://docs.nvidia.com/cuda/cuda-c-best-practices-guide/}{CUDA Best Practices Guide 的合并访问原则；高置信度}
\end{solutionexercise}

\begin{solutionexercise}{4}
\problemheading 对于一次向量加法，向量长度为 8000，每个线程计算一个输出元素，线程块
大小为 1024。程序员把内核调用配置为覆盖所有输出元素所需的最少线程块数。
网格中有多少个线程？
\begin{enumerate}[label=\alph*.,leftmargin=2.2em]
  \item 8000
  \item 8196
  \item 8192
  \item 8200
\end{enumerate}

\approachheading 使用向上取整块数 $B=\lceil N/T\rceil$，再算 $BT$。

\answerheading
\[
 B=\left\lceil\frac{8000}{1024}\right\rceil=8,\qquad
 N_{\rm thread}=8\times1024=8192.
\]
其中 8000 个线程处理有效元素，192 个线程由边界检查屏蔽。

\conclusionheading 选 c，网格共 8192 个线程。
\discussionheading
\discussionpoint{向上取整是离散执行资源的覆盖证明。}
数据规模 8000 不是块容量 1024 的整数倍，因此块数不能用整数除法截断为 7；那样只能覆盖前 7168 个元素，剩余 832 项永远没有线程。写成 $\lceil N/T\rceil$ 的意义，是把连续编号区间 $[0,N)$ 分割成若干容量为 $T$ 的离散容器，最后一个容器允许不满但不能缺席。相同思想出现在缓存瓦片、文件分片和分布式批处理中，只是 GPU 还要求多出来的线程通过边界谓词保持无副作用。到了二维或三维问题，每一维都必须独立向上取整，遗漏任一维的尾块都会漏掉整条边或整张切片，而不只是少算几个点。
\discussionpoint{最大合法线程块不等于最佳线程块。}
1024 常是设备允许的单块线程上限，但“能启动”只回答合法性，没有回答吞吐量。一个大块会一次占用较多寄存器和线程槽，某些内核因而只能让一个块驻留在 SM 上；若该块在屏障处等待，调度器可选择的其他独立工作就很少。较小块提供更细的调度粒度，也较容易适配短尾部，却可能削弱共享内存协作范围并增加块级固定开销。选择 128、256、512 或 1024 时，应先同时核对线程、寄存器和共享内存限制，再让代表性输入决定哪个配置真正缩短时间，而不是把硬件上限当作推荐值。
\discussionpoint{块内尾部与网格波次是两种不同浪费。}
本题最后一个块有 832 个有效线程和 192 个无效线程，后者仍会执行索引与条件，但组成的是 6 个完整越界 warp，并不会在边界分支内部产生半 warp 分歧。另一方面，网格只有 8 个块，若设备一次能并发远多于 8 个块，整个网格从一开始就无法填满全部 SM；这属于网格规模与并发槽位之间的“波次”量化，不是 192 个线程直接造成的。长向量含有很多完整波次时，一个尾块的成本会被摊薄；大量短向量则会反复支付启动和空槽成本，常适合批处理或融合。区分这两层尾效应，才能知道应调整块宽、合并任务，还是接受不可避免的最后余数。
\discussionpoint{占用率只描述潜在并发，不描述完成速度。}
占用率回答“资源最多允许多少 warp 同时驻留”，它有助于判断隐藏内存延迟的余地，却没有说明这些 warp 是否持续发出有用指令。一个带宽受限的向量内核可能在中等占用率就把 DRAM 用满，再增加驻留 warp 也不会提高传输上限；相反，依赖链很长的内核可能需要更多就绪 warp 才能填补停顿。因而块宽实验要同时观察总运行时间、实际带宽、寄存器用量和完整波次数，并覆盖实际业务中的长度分布。这样得到的选择才具有跨输入解释力，而不是只对 8000 这个单点恰好最优。
\variantheading 若仍有 8000 个元素但每块 256 线程，则需 $\lceil8000/256\rceil=32$ 块，共 8192 个线程，同样有 192 个越界线程。
\practiceheading 启动前可用 \code{cudaOccupancyMaxPotentialBlockSize} 比较候选块大小，再以 Nsight Compute 的 Achieved Occupancy 和尾块效率验证选择。
\sourcecheck{https://docs.nvidia.com/cuda/cuda-c-programming-guide/}{CUDA 执行配置与线程块定义；数值另行复算，置信度高}
\end{solutionexercise}

\begin{solutionexercise}{5}
\problemheading 要为包含 $v$ 个整数元素的数组分配设备全局内存。
\code{cudaMalloc} 的第二个参数应使用哪个表达式？
\begin{enumerate}[label=\alph*.,leftmargin=2.2em]
  \item \code{n}
  \item \code{v}
  \item \code{n * sizeof(int)}
  \item \code{v * sizeof(int)}
\end{enumerate}

\approachheading \code{cudaMalloc} 的大小参数以字节为单位。

\answerheading 数组元素数乘每个元素的字节数，得到
\code{v * sizeof(int)}。使用 \code{sizeof} 比假定 \code{int} 固定为 4 字节更可移植。

\conclusionheading 选 d。
\discussionheading
\discussionpoint{分配器接收字节，而程序员思考元素。}
题目给出的 $v$ 是整数元素个数，但 \code{cudaMalloc} 并不知道返回的区域以后会被解释为 \code{int}、\code{float} 还是某个结构体；它只接收一段无类型存储需要多少字节。这就是为什么必须写 $v\,\code{sizeof(int)}$：\code{sizeof} 把类型的对象表示宽度显式带入容量计算，也避免把“某台机器上 int 恰为 4 B”误写成接口契约。这个模型与 C 的 \code{malloc} 相同，类型信息存在于调用方指针和后续访问中，而不附着在设备分配本身。把元素数、元素类型和字节容量分开记录，还能让模板缓冲区在类型改变时自动重算所需空间。
\discussionpoint{容量乘法本身也可能越界。}
若 $v$ 来自文件或网络，直接计算 $v\times\code{sizeof(int)}$ 可能在 \code{size_t} 中回绕成一个较小值，分配调用甚至会成功，随后内核却按原始 $v$ 大量越界。安全前置条件是 $v\le\code{SIZE_MAX}/\code{sizeof(int)}$，并且在任何有符号值转换为 \code{size_t} 之前先排除负数。零长度也应作为独立状态处理，因为不同分配接口对零字节返回值的细节不适合作为可解引用数组的依据。类似检查不仅用于 GPU，图像尺寸乘通道数、矩阵行列乘积和序列批量容量都需要在分配之前逐层验证。
\discussionpoint{获得存储不等于获得已初始化对象。}
\code{cudaMalloc} 成功只说明地址范围存在，并不保证其中位为零；内核在任何写入前读取它，会得到不确定内容，数值测试有时却因旧页恰好为零而偶然通过。若算法确实需要全零\emph{位模式}，可以调用 \code{cudaMemset}，但它按字节重复一个值，不能用来构造任意非零整数或浮点数组。更复杂的 C++ 类型还涉及构造、析构和对象生命周期，取得原始设备存储并不会自动执行构造函数。工程接口因此应把“分配”“初始化”和“对象建立”作为三个可检查阶段，而不是让一个裸指针暗示它们都已经完成。
\discussionpoint{高频分配会把内存管理变成运行时瓶颈。}
若训练循环每一步都分配并释放多个临时数组，即使每次字节数正确，分配器的同步、元数据维护和地址空间操作也可能淹没短内核的计算时间。长寿命程序通常复用 RAII 缓冲区，或使用 stream-ordered 内存池，让释放与后续复用按流依赖排序；代价是需要管理池的高水位和跨 stream 生命周期。容量规划也不能只把永久数组相加，还要考虑库工作区、并发批次和峰值临时区，因为它们可能在时间线上重叠。将分配失败和缩小批次路径纳入测试，能把一次偶然的显存不足变成可诊断、可恢复的系统行为。
\variantheading 若分配 1500 个 \code{double}，第二个参数为元素数 1500 乘
\code{sizeof(double)}。常见平台的 \code{double} 为 8 B，因此总计 12,000 B。
\practiceheading 生产封装应检查 \code{cudaMalloc} 返回值并使用 \code{cudaMemGetInfo} 或内存池指标监控容量；需要确定初值时另行调用 \code{cudaMemset} 或初始化内核。
\sourcecheck{https://docs.nvidia.com/cuda/cuda-runtime-api/group__CUDART__MEMORY.html}{CUDA Runtime API 的 cudaMalloc 参数定义；高置信度}
\end{solutionexercise}

\begin{solutionexercise}{6}
\problemheading 如果希望分配一个包含 $n$ 个浮点元素的数组，并让浮点指针变量
\code{A_d} 指向所分配的内存，\code{cudaMalloc} 调用的第一个参数应使用什么
表达式？
\begin{enumerate}[label=\alph*.,leftmargin=2.2em]
  \item \code{n}
  \item \code{(void *) A_d}
  \item \code{*A_d}
  \item \code{(void **) &A_d}
\end{enumerate}

\approachheading API 要修改调用者的指针变量，因此必须传指针变量的地址，并转换成
\code{void **}。

\answerheading 正确形式为
\begin{lstlisting}[language=C++,numbers=none]
cudaMalloc((void **)&A_d, n * sizeof(float));
\end{lstlisting}
故选 d。现代 C++ 的 CUDA 头文件也提供模板重载；题目所列的是传统 C API 形式，
其第一个实参是转换成 \code{void **} 的指针变量地址。

\conclusionheading 选 d。
\discussionheading
\discussionpoint{二级指针来自“返回一个指针”的需要。}
\code{cudaMalloc} 要把新地址写进调用者的变量 \code{A_d}，若只传 \code{A_d} 的当前值，函数得到的只是一个副本，无法改变调用方保存的地址。传入 \code{&A_d} 后，函数拿到“指针变量本身的地址”，于是参数类型自然多出一级，传统 C 接口再把它统一表示为 \code{void**}。这与用输出参数返回文件句柄、树根指针或 MPI 对象是同一种接口设计，而强制转换只处理类型匹配，并不会替调用者证明字节数正确。现代 CUDA C++ 头文件可提供类型化重载来隐藏转换，但底层的写回语义没有改变，理解指针层次仍是排查错误地址的基础。
\discussionpoint{指针数值相同不代表处理器都能访问。}
统一虚拟寻址让运行时能从地址识别它属于哪类内存，也让支持对等访问的 GPU 在适当配置下共享地址值；这并不意味着 CPU 可以用普通加载指令解引用纯设备分配。托管内存会在处理器间迁移或建立映射，映射的页锁定主机内存则由设备经互连访问，它们都与本题的设备局部分配具有不同的可见性和成本。设计接口时应把“地址如何表示”“哪些执行主体可以访问”以及“何时完成迁移或同步”分别说明。否则一个看似普通的 \code{float*} 很容易越过地址空间边界，在编译期完全合法、运行时才失败。
\discussionpoint{所有权还必须包含异步时间线。}
分配成功后，地址只在其分配生命周期内有效；释放后即使变量仍保留旧数值，也不能再次读取、重复释放或交给新内核。传统 \code{cudaFree} 对普通分配可能引入隐式同步，而 stream-ordered 分配得到的指针交给 \code{cudaFree} 时并不保证替调用者等待所有访问，\code{cudaFreeAsync} 又必须放在正确的流依赖之后，因此不能把“调用了释放函数”等同于生命周期自动安全。RAII 包装除了保存地址和大小，还应记录设备归属、分配方式及最后使用它的 stream 或事件。这样析构策略才能既避免 use-after-free，又不会无意中用全设备同步破坏原本的并发流水线。
\discussionpoint{类型化缓冲区把易错契约提升到接口层。}
裸 \code{void**} 把元素类型、元素数、设备和释放方式都留给调用点记忆，错误路径中很容易遗漏释放，复制对象时也可能制造两个所有者。专门的 \code{device_buffer<T>} 可以在构造时检查 $n\,\code{sizeof(T)}$ 是否溢出，在移动时转移唯一所有权，并提供显式的设备号与 stream-aware 释放策略。多 GPU 对等访问、进程间 IPC 和 CUDA Graph 会继续增加句柄导出、捕获期限制与上下文归属，因而更需要禁止无意的浅复制。这样的封装并非为了隐藏 CUDA，而是把本题揭示的二级指针契约变成编译器可以协助维护的不变量。
\variantheading 若变量为 \code{double *D_d}，传统接口应写成 \code{cudaMalloc((void **)&D_d, n*sizeof(double))}，成功后 \code{D_d} 指向设备分配区。
\practiceheading C++ 项目宜用 RAII 包装 \code{cudaFree}，并用 Compute Sanitizer initcheck/racecheck 配合单元测试排查未初始化读取和错误生命周期。
\sourcecheck{https://docs.nvidia.com/cuda/cuda-runtime-api/group__CUDART__MEMORY.html}{CUDA Runtime API；高置信度}
\end{solutionexercise}

\begin{solutionexercise}{7}
\problemheading 如果希望把 3000 字节数据从主机数组 \code{A_h}（\code{A_h} 指向源数组
的第 0 个元素）复制到设备数组 \code{A_d}（\code{A_d} 指向目标数组的第 0 个
元素），CUDA 中适合进行该复制的 API 调用是什么？
\begin{enumerate}[label=\alph*.,leftmargin=2.2em]
  \item \code{cudaMemcpy(3000, A_h, A_d, cudaMemcpyHostToDevice);}
  \item \code{cudaMemcpy(A_h, A_d, 3000, cudaMemcpyDeviceToHost);}
  \item \code{cudaMemcpy(A_d, A_h, 3000, cudaMemcpyHostToDevice);}
  \item \code{cudaMemcpy(3000, A_d, A_h, cudaMemcpyHostToDevice);}
\end{enumerate}

\approachheading \code{cudaMemcpy} 参数次序是目标、源、字节数、方向。

\answerheading 应调用
\begin{lstlisting}[language=C++]
cudaMemcpy(A_d, A_h, 3000, cudaMemcpyHostToDevice);
\end{lstlisting}
即选 c。生产代码还要检查返回值；源和目标至少都要覆盖 3000 字节的有效区间。

\conclusionheading 选 c。
\discussionheading
\discussionpoint{四个参数共同构成一次原始字节复制。}
把 \code{cudaMemcpy} 记成“复制数组”的函数容易掩盖两个事实：第一个参数永远是目标，第二个永远是源，第三个单位是字节而不是元素。方向枚举还必须与两个地址实际所属空间一致，主机到设备不能通过交换参数后仍保留原枚举来补救。函数不会把 3000 B 的整数转换成浮点，也不知道源和目标是否具有相同对象布局；它只搬运位模式。工程封装最好从元素类型和计数推导字节数，并在调试构建中核对指针属性，这样才能同时防住参数次序、容量和地址空间三类彼此独立的错误。
\discussionpoint{“同步”和“异步”不能只看函数名字。}
CUDA 运行时明确说明，复制是否相对主机同步还取决于方向及主机内存是否 pageable。同步调用从 pageable 主机内存执行 H2D 时，运行时可先把源数据复制到内部 staging buffer，待这一步完成并且调用返回后，原 pageable 源缓冲便可修改或复用，即使 staging 到设备的 DMA 尚未结束。页锁定源配合 \code{cudaMemcpyAsync} 时则不同：在所属 stream 的复制完成前，应用不能改写或释放该源缓冲；设备目标同样要由同一 stream 的顺序、事件或显式同步保护，之后才能被无依赖地消费、覆盖或释放。D2D 复制也未必阻塞主机，而 pageable 的 \code{Async} 传输反过来可能因暂存而阻塞，因此应依据具体内存类型和完成事件管理生命周期，不能只看函数后缀。
\discussionpoint{复制带宽属于互连层而非设备内存层。}
若 GPU 的 DRAM 峰值达到每秒数百 GB，也不能推出主机到设备复制具有同样速度，因为数据还要经过 PCIe、CPU 内存控制器或特定平台上的 NVLink 路径。多 GPU 系统又要区分直接点对点传输与经主机内存中转，两条路径的拓扑距离和并发能力可能完全不同。统一内存省去了显式复制语句，却没有消灭数据移动，只是把它转换成页迁移、缺页处理和一致性开销。容量规划和性能报告因此应分别列出主机互连、GPU 间互连和设备 DRAM 三个层级，避免把不同屋顶混成一个“内存带宽”。
\discussionpoint{流水线成立需要数据依赖与硬件能力同时满足。}
双缓冲的理想时间线是：一块缓冲执行 H2D 时，另一块运行内核，第三阶段把更早结果传回主机；stream 和事件负责表达它们之间真正的数据依赖。要实现重叠，主机缓冲通常需要页锁定，设备需要相应复制引擎，内核资源还不能独占所有可并发能力，所以代码中出现 \code{cudaMemcpyAsync} 只是必要条件之一。验证时应使用非整块长度和带偏移子区间检查正确性，再观察时间线确认三阶段确实重叠且没有默认流或分配调用引入隐式同步。最终数组相等只能证明数据大致到了目的地，无法证明流水线按预期并发。
\variantheading 若把 4096 B 从 \code{A_d} 回传 \code{A_h}，调用为 \code{cudaMemcpy(A_h, A_d, 4096, cudaMemcpyDeviceToHost)}。
\practiceheading 流水化程序可改用页锁定主机内存和 \code{cudaMemcpyAsync}，再用 Nsight Systems 检查复制与内核是否真正重叠及是否发生隐式同步。
\sourcecheck{https://docs.nvidia.com/cuda/cuda-runtime-api/api-sync-behavior.html}{CUDA Runtime API 对同步/异步 memcpy 按传输方向及 pageable/pinned 内存分类的语义；对抗检查：不能仅凭函数名推断主机完成边界}
\end{solutionexercise}

\begin{solutionexercise}{8}
\problemheading 应如何声明变量 \code{err}，才能适当地接收 CUDA API 调用返回的值？
\begin{enumerate}[label=\alph*.,leftmargin=2.2em]
  \item \code{int err;}
  \item \code{cudaError err;}
  \item \code{cudaError_t err;}
  \item \code{cudaSuccess_t err;}
\end{enumerate}

\approachheading 查 API 的统一错误类型。

\answerheading Runtime API 的返回类型是 \code{cudaError_t}，故选 c。声明为：
\begin{lstlisting}[language=C++,numbers=none]
cudaError_t err;
\end{lstlisting}
\code{cudaSuccess} 是这个枚举类型的成功值，不是类型名。

\conclusionheading 选 c。
\discussionheading
\discussionpoint{CUDA 错误沿异步时间线分层出现。}
主机调用若参数非法或资源申请立即失败，通常能在该 API 的返回值中看到错误；内核里的越界访问却发生在设备稍后执行时，启动语句返回时可能还什么都没做。于是某个同步或复制 API 报错，不一定是它自身出错，也可能是在替之前排队的工作传递状态。理解这一点要记录“最后一次已知成功的完成边界”，再向后定位最早可能失败的异步操作。错误值 \code{cudaError_t} 因而不仅是一个枚举结果，还需要与 stream、调用顺序和同步位置共同解释。多 stream 程序还应把事件依赖纳入时间线，否则“之前”本身就未必等于主机代码的文本顺序。
\discussionpoint{启动检查和执行检查回答不同问题。}
内核启动后立即调用 \code{cudaGetLastError}，可以发现块维度非法、参数设置或启动资源超限等问题，但设备代码尚未完成时，它无法证明数组索引和同步语义正确。随后必须在真正依赖结果的位置等待事件、同步对应 stream 或同步设备，才能暴露非法访存和断言失败等执行期错误。开发阶段可在每个可疑内核后临时同步，以牺牲性能换取更短的故障区间；生产路径则应把检查放在已有依赖边界，避免把流水线强行串行化。两层检查配合，才能避免“启动成功”被误读成“计算成功”。
\discussionpoint{不同错误具有不同的恢复边界。}
显存不足可能通过释放缓存、减小批次或换用分块算法恢复，参数错误通常应在调用方修正后重试；非法设备访问或严重启动失败却可能使当前上下文处于不再可靠的状态。若程序统一打印一行字符串后继续，后续 API 可能接连报告派生错误，最终覆盖最初原因并把损坏结果交给下一阶段。服务程序应预先区分可重试、需要重建上下文和必须终止进程的类别，并为每类设计资源清理与请求失败语义。恢复策略也是正确性的一部分，因为“进程没有退出”并不意味着 GPU 状态仍可安全使用。请求标识与首个错误位置也应保留，便于把设备故障追溯到原始输入。
\discussionpoint{错误封装应保留第一现场。}
一个有用的检查宏不应只输出错误数字，还要保存失败表达式、源文件行号、设备号、stream 以及 \code{cudaGetErrorName} 和 \code{cudaGetErrorString}。清理阶段本身也可能失败，若每次都覆盖同一个变量，最有价值的首个错误就会被最后一次 \code{cudaFree} 的状态替代。对于异步故障，可在最小输入上逐步增加同步点，再借助内存与同步检查工具定位实际设备指令，而不是责怪最先观察到错误的同步函数。良好的可观测性把一次难以复现的“稍后失败”转化为带时间边界和上下文的证据链。
\variantheading 若 \code{cudaMalloc} 失败，应保存其 \code{cudaError_t}，与 \code{cudaSuccess} 比较，并用 \code{cudaGetErrorString(err)} 输出诊断。
\practiceheading 生产代码应统一封装每个 Runtime API 返回值，并在内核启动后调用 \code{cudaGetLastError}、在同步边界检查异步执行错误。
\sourcecheck{https://docs.nvidia.com/cuda/cuda-runtime-api/group__CUDART__TYPES.html}{CUDA Runtime API 数据类型；高置信度}
\end{solutionexercise}

\begin{solutionexercise}{9}
\problemheading 考虑下面的 CUDA 内核及调用它的主机函数：
\begin{lstlisting}[language=C++,firstnumber=1]
__global__ void foo_kernel(float* a, float* b, unsigned int N) {
    unsigned int i = blockIdx.x * blockDim.x + threadIdx.x;
    if (i < N) {
        b[i] = 2.7f * a[i] - 4.3f;
    }
}
void foo(float* a_d, float* b_d) {
    unsigned int N = 200000;
    foo_kernel<<<(N + 128 - 1)/128, 128>>>(a_d, b_d, N);
}
\end{lstlisting}
\begin{enumerate}[label=\alph*.,leftmargin=2.2em]
  \item 每个块中的线程数是多少？
  \item 网格中的线程总数是多少？
  \item 网格中的块数是多少？
  \item 执行第 02 行代码的线程数是多少？
  \item 执行第 04 行代码的线程数是多少？
\end{enumerate}

\approachheading 先向上取整网格大小，再区分所有启动线程与通过 $i<N$ 检查的线程。

\answerheading
\begin{enumerate}[label=\alph*.,leftmargin=2.2em]
  \item 每块线程数为 128。
  \item 块数 $B=\lceil200000/128\rceil=1563$，故总线程数为
  $1563\times128=200064$。
  \item 网格含 1563 个块。
  \item 第 2 行索引计算在所有已启动线程中执行，共 200064 个线程。
  \item 第 4 行赋值受 $i<N$ 保护，恰由 200000 个线程执行。
\end{enumerate}
最后 64 个线程只计算索引并判断条件，不访问数组。

\conclusionheading 依次为 128、200064、1563、200064、200000。
\discussionheading
\discussionpoint{逻辑任务数与启动线程数回答不同问题。}
数组有 200000 个待计算元素，这是算法层的逻辑任务数；网格却只能按 128 线程的整块分配，所以创建了 200064 个线程上下文。多出的 64 个线程不是计算错误，而是离散调度粒度为覆盖任意长度付出的冗余，只要边界谓词让它们不触碰数组即可。相同区分在稀疏计算中更明显：线程即使索引合法，也可能发现该行没有非零元或该顶点没有边，启动数量仍不等于有用运算数量。报告性能时若只写线程数，会把调度资源、有效工作和 FLOP 三个层次混在一起。至少还应给出有效线程比例，才能知道空槽在当前规模上是否值得优化。
\discussionpoint{部分 warp 与全空 warp 的行为并不相同。}
本题 $200000\bmod32=0$，有效区间恰好落在 warp 边界，所以最后 64 个越界线程组成两个全空 warp，所有 lane 对 \code{i<N} 得到相同结果，不发生分支内部分歧。若 $N$ 增加 1，则下一个 warp 只有一个 lane 有效，硬件会用活动掩码执行主体，其余 31 个 lane 在该路径上闲置。无论是否分歧，空闲线程仍需要完成地址计算和条件判断，并占据已启动 warp 的调度位置。这个反例说明分支效率与有效工作比例是两种指标：一个全空 warp 可以没有分歧，却没有产生任何输出。
\discussionpoint{向上取整公式也需要防整数溢出。}
常见写法 $(N+T-1)/T$ 简洁，但当 $N$ 已接近整数类型上限时，加上 $T-1$ 会先回绕，随后得到完全错误的块数。等价形式 $N/T+(N\bmod T\ne0)$ 不需要这次加法，更适合处理外部输入和超大数组；全局索引中的 $bT+t$ 也应在足够宽的类型中计算。二维矩阵的 $rW+c$、批量张量的多维乘积更容易超过 32 位，所以只用小型单元测试无法覆盖这种故障。容量检查、网格计算和内核索引必须使用一致宽度，否则主机认为已分配足够空间，设备仍可能以截断索引写错位置。
\discussionpoint{独立枚举能把覆盖问题变成可检验计数。}
对小规模配置，可以在 CPU 上遍历每个块和线程，记录总启动数、有效数、部分 warp 数以及全空 warp 数，再与内核回写的索引集合比较。这样不仅能发现漏算和重复，还能把“边界造成多少浪费”从模糊印象变成精确计数；随后才需要用设备测量解释调度和带宽后果。对于在线服务，还应测试块数刚好等于、略小于和略大于设备一次可并发块数的情形，因为最后一波只剩少量块时会影响尾延迟。覆盖枚举提供算法证据，时间分布提供系统证据，两者结合比只检查前几个输出更可靠。索引集合若还能验证恰等于 $[0,N)$，就同时证明了不重与不漏。
\variantheading 若 $N=200001$，块数仍为 1563、总线程仍为 200064，但赋值线程变为 200001，尾部仅 63 个线程被屏蔽。
\practiceheading 可用 Nsight Compute 的线程/warp 启动统计和分支效率验证尾块行为，并用 Compute Sanitizer memcheck 确认保护条件覆盖每次数组访问。
\sourcecheck{https://docs.nvidia.com/cuda/cuda-c-programming-guide/}{CUDA SIMT 执行与线程索引定义；算术独立复核，置信度高}
\end{solutionexercise}

\begin{solutionexercise}{10}
\problemheading 一位新的 CUDA 程序员对 CUDA 感到沮丧。他抱怨 CUDA 太繁琐：自己计划
在主机和设备上都执行的许多函数，必须分别声明两次，一次声明为主机函数，一次
声明为设备函数。你会如何帮助这位新程序员？

\approachheading CUDA 允许把执行空间限定符组合使用。

\answerheading 把函数同时声明为 \code{__host__ __device__}，编译器会分别生成主机版本
和设备版本，例如：
\begin{lstlisting}[language=C++]
__host__ __device__ float affine(float x) {
    return 2.7f * x - 4.3f;
}
\end{lstlisting}
函数体中只能无条件使用两端都可用的语言或库功能。若需分支，可用预处理条件隔离设备路径：
\begin{lstlisting}[language=C++,numbers=none]
#if defined(__CUDA_ARCH__)
// device path
#endif
\end{lstlisting}
双重限定不会让任意主机 API 自动变成设备可调用。

\conclusionheading 使用组合限定符共享源代码，同时审查两种编译目标的可用功能。
\discussionheading
\discussionpoint{一份函数体会生成两个执行世界中的版本。}
\code{__host__ __device__} 并不是让同一段机器码在 CPU 和 GPU 间移动，而是要求编译器分别为两个目标编译函数体。于是一个普通加法在两端都可成立，读取 \code{threadIdx.x}、调用主机文件接口或依赖设备地址空间的操作却只属于其中一端。这个模型类似跨后端的泛型算法：数学接口可以共享，可调用的库、异常机制和存储语义仍受目标约束。理解“双重生成”后，新程序员就不会期待加上限定符能自动把任意主机函数变成设备函数，也更容易读懂为何错误有时只在一端编译时出现。
\discussionpoint{条件编译应隔离能力差异而非复制算法。}
当设备版本需要使用内建指令、主机版本要提供参考路径时，可以用 \code{__CUDA_ARCH__} 选择局部实现；问题在于大量条件分支会逐渐形成两套难以同步的程序。更稳妥的结构是把无副作用的纯计算核心保持共同，把分配、访存和目标专用操作放进薄适配层，并让两端调用相同的数学接口。这样既能利用设备特性，也能让 CPU 版本成为容易执行的正确性基线。条件编译的范围越小，未来修改公式时遗漏某一目标的风险越低，这一原则同样适用于 SIMD、HIP 或其他异构后端。
\discussionpoint{源码相同不保证位级结果和成本相同。}
双重限定会产生两个函数版本，模板与多种数据类型还会成倍增加实例，因而可能扩大编译时间和二进制体积。CPU 与 GPU 可采用不同的 FMA 融合、舍入模式、数学库近似和次正规数处理，即使 C++ 表达式逐字相同，浮点结果也未必逐位一致；设备端的内联与寄存器分配还会改变性能。测试应依据算法误差建立绝对与相对容差，而不是把主机位型当作唯一正确答案。若业务确实需要可复现位序列，则必须明确限制编译选项和运算路径，并接受由此带来的吞吐代价。跨后端参考还应固定输入与求和次序，避免把两个变量同时改变后无法归因。
\discussionpoint{双端函数最适合承载可验证的小型纯算法。}
几何坐标变换、边界谓词、哈希步骤和短小数值算子常被声明为主机与设备均可调用，因为它们容易保持无副作用，也便于直接在 CPU 单元测试中覆盖大量输入。构建系统仍要确保两个版本都被实际实例化：只调用主机路径会让设备端不支持的函数长期隐藏，只编译 GPU 又失去快速参考验证。一个可靠流程是在主机执行穷举或随机属性测试，再至少为部署目标编译设备版本，并在有 GPU 时比较结果与容差。这样，代码共享带来的主要价值不是少写一份函数，而是让同一数学契约在两个后端接受独立验证。
\variantheading 若函数需要设备端 \code{threadIdx.x}、主机端固定返回 0，可用 \code{#if defined(__CUDA_ARCH__)} 分支，设备版本读取线程号而主机版本返回 0。
\practiceheading 构建系统应同时编译并测试主机路径和至少一个目标 SM 的设备路径，借助 \code{nvdisasm} 或 \code{cuobjdump} 确认预期架构代码确已生成。
\sourcecheck{https://docs.nvidia.com/cuda/cuda-c-programming-guide/}{CUDA C++ 函数执行空间限定符；高置信度}
\end{solutionexercise}

\chapter{多维网格与数据}

\begin{solutionexercise}{1}
\problemheading 本章实现了一个矩阵乘法内核，每个线程生成一个输出矩阵元素。本题要求
实现不同的矩阵--矩阵乘法内核并比较它们。
\begin{enumerate}[label=\alph*.,leftmargin=2.2em]
  \item 编写一个让每个线程生成一个输出矩阵行的内核，并为该设计填写执行
  配置参数。
  \item 编写一个让每个线程生成一个输出矩阵列的内核，并为该设计填写执行
  配置参数。
  \item 分析上述两种内核设计各自的优点和缺点。
\end{enumerate}

\approachheading 对 $P=MN$，一行线程固定 $i$、遍历全部 $j,k$；一列线程固定 $j$、
遍历全部 $i,k$。两者都要检查线程对应的行或列是否越界。

\answerheading 可编译的核心实现如下（矩阵按行主序）：
\begin{lstlisting}[language=C++]
__global__ void matmul_row(const float* M, const float* N,
                           float* P, int n) {
    int row = blockIdx.x * blockDim.x + threadIdx.x;
    if (row >= n) return;
    for (int col = 0; col < n; ++col) {
        float sum = 0.0f;
        for (int k = 0; k < n; ++k)
            sum += M[row*n + k] * N[k*n + col];
        P[row*n + col] = sum;
    }
}
__global__ void matmul_col(const float* M, const float* N,
                           float* P, int n) {
    int col = blockIdx.x * blockDim.x + threadIdx.x;
    if (col >= n) return;
    for (int row = 0; row < n; ++row) {
        float sum = 0.0f;
        for (int k = 0; k < n; ++k)
            sum += M[row*n + k] * N[k*n + col];
        P[row*n + col] = sum;
    }
}
// dim3 block(256), grid((n + block.x - 1) / block.x);
\end{lstlisting}
行方案中单线程对 $M$ 的行访问连续，但同一 warp 在同一时刻访问不同的 $M$ 行；列方案中
同一 warp 的线程在给定 $k$ 时读取 $N$ 的连续列，合并性较好，但各线程写不同列也连续。
两者每线程都做 $O(n^2)$ 工作，暴露的并行任务仅 $n$ 个，远少于每元素一线程的 $n^2$ 个；
且没有共享内存分块复用，所以通常性能较差。

\complexity{$O(n^3)$ 总工作、每线程 $O(n^2)$}{$O(1)$ 每线程额外空间}{仅 $n$ 个并行任务}
\conclusionheading 两种实现均正确，但并行度低；列方案通常更利于 $N$ 的合并读取，最终应以
测量结果判断。
\discussionheading
\discussionpoint{同样的总工作量可以暴露完全不同的并行度。}
三种映射最终都执行约 $2n^3$ 次浮点运算，但每元素一线程会产生 $n^2$ 个独立点积，而每行或每列一线程只留下 $n$ 个长任务。以 $n=1024$ 为例，后者仅有 1024 个线程，每线程顺序计算 1024 个输出且每个输出还含 1024 项累加；在大 GPU 上，这点线程可能不足以持续隐藏访存延迟。长线程也可能复用某些地址计算并提高局部指令连续性，所以它并非逻辑上荒谬，只是把大量线程级并行换成了线程内顺序工作。这个例子说明复杂度相同只约束总功，不决定关键路径和可调度任务数量。
\discussionpoint{行线程与列线程的好坏由 warp 地址集合决定。}
行主序常让人直觉认为“每线程算一行”必然更好，因为单个线程写回一行时地址连续；然而同一时刻一个 warp 的 32 个行线程可能分别落在相隔 $n$ 的 32 行上，形成大步长请求。列线程在给定 $k$ 时读取 $N[k,n\text{的连续列}]$，相邻 lane 往往访问相邻元素，输出列的写回也可呈连续地址。反过来，leading dimension、子矩阵偏移和缓存复用都可能改变物理事务，所以不能只凭线程名字判断。正确分析应写出 lane $t$ 在某条动态加载上的地址公式，再观察 32 个地址如何映射到内存扇区。
\discussionpoint{朴素映射可以沿协作层次逐步演化。}
若一条点积太长，可以让一个 warp 的 lane 分别累加不同 $k$ 项，再用 shuffle 归约；若输入复用不足，则让线程块生成一个输出瓦片，把 $M$ 与 $N$ 的子块载入共享内存。继续增加寄存器微瓦片和双缓冲，便接近高性能 GEMM 的多级分块结构，低精度场景还可把合适的片段交给 Tensor Core。每一步都不是凭空更换算法，而是把顺序循环中的独立工作和数据复用暴露给更宽的硬件层级。成熟库会按矩阵形状、转置、类型和架构选择不同微内核，正因为没有一个固定映射能覆盖所有这些边界。
\discussionpoint{实验必须把正确性基线和性能基线分开。}
CPU 双精度实现适合检查数值，因为它提供独立且较高精度的参考；cuBLAS 更适合充当性能基线，但库与朴素代码的累加次序不同，不能要求结果逐位一致。比较三种内核时应使用相同矩阵、预热和计时区间，并依据 $n$ 与数据幅值建立混合误差容差，再报告有效 GFLOP/s 和实际数据流量。若还记录线程数量、寄存器和事务效率，就能解释某个尺寸上的交叉点来自并行度不足、访存布局还是资源压力。只有一条“加速了几倍”的曲线无法告诉读者优化是否能迁移到另一尺寸或设备。
\variantheading 对 $n=1024$、每块 256 线程，两种方案都启动 4 个块、1024 个工作线程；每个线程计算 1024 个输出元素并执行约 $1024^2$ 次乘加。
\practiceheading 用 cuBLAS \code{cublasSgemm} 作为正确性和性能基线，并在 Nsight Compute 中比较 Global Load Efficiency、Store Efficiency 与每线程寄存器压力。
\sourcecheck{https://docs.nvidia.com/cuda/cuda-c-best-practices-guide/}{CUDA Best Practices Guide 的合并访问原则；高置信度}
\end{solutionexercise}

\begin{solutionexercise}{2}
\problemheading 矩阵--向量乘法接收输入矩阵 $B$ 和向量 $\mathbf{c}$，生成输出向量
$\mathbf{a}$。输出向量 $\mathbf{a}$ 的每个元素都是输入矩阵 $B$ 的一行与
$\mathbf{c}$ 的点积，即
\[
  a_i=\sum_j B_{i,j}\times c_j.
\]
为简单起见，只处理元素为单精度浮点数的方阵。编写矩阵--向量乘法内核以及
可以调用它的主机存根函数。该存根函数接收四个参数：输出向量 $\mathbf{a}$ 的
指针、输入矩阵 $B$ 的指针、输入向量 $\mathbf{c}$ 的指针，以及每个维度的
元素数量 $N$。使用一个线程计算输出向量的一个元素。

\approachheading 把线程映射到行 $i$，检查 $i<N$，再顺序累加该行点积。主机端完成分配、
复制、启动、同步、错误检查、回传与释放。

\answerheading
\begin{lstlisting}[language=C++]
#include <cuda_runtime.h>
#include <cstdio>
#include <cstdlib>
#define CHECK(x) do { cudaError_t e=(x); if(e!=cudaSuccess){ \
  std::fprintf(stderr,"CUDA: %s\n",cudaGetErrorString(e)); std::exit(1); } } while(0)

__global__ void matvec(const float* B, const float* c, float* a, int n) {
    int i = blockIdx.x * blockDim.x + threadIdx.x;
    if (i >= n) return;
    float sum = 0.0f;
    for (int j = 0; j < n; ++j) sum += B[i*n + j] * c[j];
    a[i] = sum;
}
void matvec_host(float* a, const float* B, const float* c, int n) {
    float *dB=nullptr, *dc=nullptr, *da=nullptr;
    size_t mb=size_t(n)*n*sizeof(float), vb=size_t(n)*sizeof(float);
    CHECK(cudaMalloc(&dB, mb)); CHECK(cudaMalloc(&dc, vb));
    CHECK(cudaMalloc(&da, vb));
    CHECK(cudaMemcpy(dB, B, mb, cudaMemcpyHostToDevice));
    CHECK(cudaMemcpy(dc, c, vb, cudaMemcpyHostToDevice));
    matvec<<<(n+255)/256,256>>>(dB,dc,da,n);
    CHECK(cudaGetLastError()); CHECK(cudaDeviceSynchronize());
    CHECK(cudaMemcpy(a, da, vb, cudaMemcpyDeviceToHost));
    CHECK(cudaFree(da)); CHECK(cudaFree(dc)); CHECK(cudaFree(dB));
}
\end{lstlisting}
CPU 参考式就是题给求和；逐元素比较时应使用绝对/相对混合容差，因为浮点累加次序可能不同。

\complexity{$O(N^2)$ 总工作、$O(N)$ 理想并行时间}{$O(N^2)$ 输入与 $O(N)$ 输出设备空间}{$N$ 个独立输出}
\conclusionheading 索引、边界、生命周期和错误检查均完整；该基础实现没有缓存向量或分块优化。
\discussionheading
\discussionpoint{GEMV 缺少 GEMM 的第二个复用维度。}
矩阵--向量乘的每个输出是一行与同一向量的点积，因此有 $N$ 个独立输出，却必须流过约 $N^2$ 个矩阵元素。矩阵--矩阵乘还能让一个输入瓦片服务二维输出区域，而 GEMV 中每个矩阵元素通常只贡献一次，输出写回也仅有 $N$ 项，所以算术强度较低。即使 GPU 具有很高的乘加峰值，性能也常先被读取矩阵的带宽限制；增加计算单元不能突破这个数据移动屋顶。这个结构解释了为什么看似简单的 GEMV 很难达到 GEMM 的 GFLOP/s，也提示优化应优先考虑减少流量或与相邻算子融合。
\discussionpoint{共享向量与协作点积代表两种分解方向。}
基础实现让一个线程串行累加整行，所有线程都会读取公共向量 $c$；可按片段把 $c$ 载入共享内存，使块内多个行线程复用同一片段，但每个片段需要生产者与消费者同步。另一方案让一个 warp 共同计算一行，每个 lane 负责若干列，再用 shuffle 或树形归约合并部分和，这会增加单行内部并行，却减少同时处理的行数。前者围绕输入复用组织线程，后者围绕点积关键路径组织线程，两者适合的 $N$ 和批量规模不同。理解这种双向选择，也有助于分析稀疏矩阵向量乘中“每行一个线程”与“每行一个 warp”的经典取舍。
\discussionpoint{点积的归约树也决定数值误差。}
浮点加法不满足结合律，串行顺序 $(((x_0+x_1)+x_2)+\cdots)$ 与平衡树会在不同位置舍入，所以 GPU 协作归约不应与 CPU 单精度结果逐位比较。对 $N$ 项数量级相近的和，成对归约通常把误差增长从与 $N$ 相关的长链缩短到与 $\log N$ 相关的层数，但数据抵消严重时仍可能丢失有效位。双精度累加或补偿求和能改善准确度，却会增加指令和状态，某些设备上的吞吐代价很大。工程接口应先规定允许的误差和输入范围，再选择归约结构，而不是先追求最快实现后才为差异寻找容差。
\discussionpoint{小批量场景往往被启动与权重流量主导。}
GEMV 出现在小批量全连接推理、迭代线性求解器和图排序中，这些场景的共同点是每次矩阵只乘很少的向量。若向量很短，单次内核启动时间可能接近计算本身；若矩阵很大，权重又会在每次请求中反复从显存读取。批量接口可以把多个向量提升为小型 GEMM，融合后续偏置或激活则避免中间写回，两者都从端到端数据流而非单条点积寻找收益。评估时应同时比较库实现、融合实现和真实批量分布，并用算术强度与有效带宽解释结果，避免只在一个大矩阵上给出孤立加速比。在线系统还要报告批处理等待时间，否则吞吐改善可能掩盖请求延迟恶化。
\variantheading 若 $N=513$ 且每块 256 线程，网格需 3 个块、768 个线程，其中 255 个线程由边界判断提前返回。
\practiceheading 以 cuBLAS \code{cublasSgemv} 和 CPU 双精度参考结果作基线，用 Compute Sanitizer 检查边界，并用 Nsight Compute 判断瓶颈是带宽还是低占用率。
\sourcecheck{https://docs.nvidia.com/cuda/cuda-c-programming-guide/}{CUDA 编程模型与 Runtime API；高置信度}
\end{solutionexercise}

\begin{solutionexercise}{3}
\problemheading 考虑下面的 CUDA 内核以及调用它的主机函数：
\begin{lstlisting}[language=C++,firstnumber=1]
__global__ void foo_kernel(float* a, float* b, unsigned int M, unsigned int N) {
    unsigned int row = blockIdx.y*blockDim.y + threadIdx.y;
    unsigned int col = blockIdx.x*blockDim.x + threadIdx.x;
    if (row < M && col < N) {
        b[row*N + col] = a[row*N + col]/2.1f + 4.8f;
    }
}
void foo(float* a_d, float* b_d) {
    unsigned int M = 150;
    unsigned int N = 300;
    dim3 bd(16, 32);
    dim3 gd((N - 1)/16 + 1, (M - 1)/32 + 1);
    foo_kernel<<<gd, bd>>>(a_d, b_d, M, N);
}
\end{lstlisting}
\begin{enumerate}[label=\alph*.,leftmargin=2.2em]
  \item 每个线程块中的线程数是多少？
  \item 网格中的线程数是多少？
  \item 网格中的线程块数是多少？
  \item 执行第 05 行代码的线程数是多少？
\end{enumerate}

\approachheading $x$ 方向覆盖列，$y$ 方向覆盖行；分别向上取整。

\answerheading 每块线程数 $16\times32=512$。网格维度为
\[
 G_x=\lceil300/16\rceil=19,\qquad G_y=\lceil150/32\rceil=5.
\]
因此块数 $19\times5=95$，线程总数 $95\times512=48640$。第 5 行条件体仅由有效
$150\times300=45000$ 个线程执行；3640 个边界线程被屏蔽。

\conclusionheading (a)--(d) 分别为 512、48640、95、45000。
\discussionheading
\discussionpoint{二维坐标最终仍按 x 维优先组成 warp。}
线程块写成 $(16,32)$ 并不意味着硬件先把同一列的 32 个线程组成 warp；CUDA 线性化线程号时让 x 维变化最快，所以前 16 个线程来自 $y=0$，后 16 个来自 $y=1$，一个 warp 横跨两行。若相邻 x 对应矩阵连续列，这种组织通常仍能产生每半 warp 连续的地址片段，但交换 x、y 映射可能把 lane 间距放大为整行宽度。坐标公式只决定“算哪个元素”，线性线程次序还决定“哪些元素同时发出请求”。因此设计二维块时应同时画出坐标覆盖和前两个 warp 的 lane 分布，而不能只检查网格面积。
\discussionpoint{合法启动需要同时通过总量和逐维限制。}
块内线程总数 $16\times32=512$ 必须不超过 \code{maxThreadsPerBlock}，但这不是唯一条件；x、y、z 各维还分别受 \code{maxThreadsDim} 约束，网格维度又有自己的范围。一个乘积不大的瘦长块仍可能因为某一维过大而非法，反过来各维单独合法也不能保证乘积合法。题目提供的配置在典型目标上可用，生产代码却应针对部署设备查询属性，并紧接启动检查返回状态。把硬约束验证放在调优之前，可以避免把“没有输出”误诊成访存或占用率问题。
\discussionpoint{二维尾块有边、角和不同活动掩码。}
x 方向 300 列需要 19 个块，最右块只有 12 列有效；y 方向 150 行需要 5 个块，最下块只有 22 行有效。于是普通内部块全有效，右边块、下边块和右下角块分别具有不同的活动线程比例，内核必须让行、列两个条件共同支配读写。把矩阵展平成 45000 个一维任务可以减少二维角落的空闲组合，却要用除法和取模还原坐标，也可能改变行局部性。选择二维还是一维映射，本质上是在边界利用率、地址计算和后续邻域结构之间权衡，而不是简单追求更少的越界线程。若后续算子按行取邻域，保留二维坐标往往还能简化共享瓦片与 halo 处理。
\discussionpoint{非方形哨兵数据能同时揭露多类索引错误。}
若测试使用方阵或全 1，交换行列、写错步长甚至局部转置都可能得到看似合理的结果。给位置 $(r,c)$ 写入唯一编码 $1000r+c$，再选择长宽均不整除块维的矩阵，就能从错误值的十进制结构看出是行号、列号还是边界发生偏移。还应显式检查四角、跨行位置和右下尾块，并把手算的 95 个块、48640 个启动线程与运行时配置对照。这样的测试先建立可解释的正确性证据，再使用性能工具观察活动线程与事务，能避免在索引尚未可靠时过早调优。CPU 参考若采用相同哨兵公式，还能自动定位首个错误坐标而非只返回失败布尔值。
\variantheading 若块改为 $(32,8)$，网格为 $(10,19)$，共 190 块、48640 个启动线程，执行第 05 行的仍是 45000 个线程。
\practiceheading 启动前查询 \code{cudaDeviceProp} 的线程维度限制，并在 Nsight Compute Launch Statistics 中核对块、网格和尾部未满 warp。
\sourcecheck{https://docs.nvidia.com/cuda/cuda-c-programming-guide/}{CUDA 多维线程索引；数值独立复核，置信度高}
\end{solutionexercise}

\begin{solutionexercise}{4}
\problemheading 考虑一个宽度为 400、高度为 500 的二维矩阵。该矩阵以一维数组形式存储。
请分别指定矩阵第 20 行、第 10 列元素的数组索引：
\begin{enumerate}[label=\alph*.,leftmargin=2.2em]
  \item 矩阵按行主序存储时；
  \item 矩阵按列主序存储时。
\end{enumerate}

\approachheading 使用从 0 开始的行、列下标；行主序步长为宽度，列主序步长为高度。

\answerheading 行主序：$20\times400+10=8010$。列主序：$10\times500+20=5020$。
若题意采用从 1 开始的自然序号，需先减 1；但全书 CUDA 下标约定为从 0 开始。

\conclusionheading (a) 8010；(b) 5020。
\discussionheading
\discussionpoint{行主序和列主序是存储契约，不是矩阵属性。}
数学上的矩阵只有行列关系，并没有规定内存中哪个方向相邻；行主序选择列坐标为低位，故索引是 $rW+c$，列主序则选择行坐标为低位，故索引是 $cH+r$。若文件读取器、CPU 库和 CUDA 内核对契约理解不同，访问仍可能全部落在合法范围内，却把数据解释成转置或更隐蔽的错位排列。错误因此不一定触发越界工具，甚至在方阵上看起来维度完全正确。接口应把布局与尺寸一起记录，而不是仅传一段裸指针后靠调用方猜测。批次跨距与子矩阵偏移也属于同一契约，不能由宽高两个整数完整表达。
\discussionpoint{主导维描述物理跨距而非逻辑宽度。}
实际矩阵可能在每行末尾填充若干元素以满足对齐，也可能只是大矩阵中的一个矩形视图，此时相邻逻辑行在内存中的距离不等于视图宽度。行主序公式应写成 $r\,ld+c$，其中 $ld$ 是 leading dimension；列主序同理以物理列跨距为准。cuBLAS 的 \code{lda}、\code{ldb} 和 \code{ldc} 正是把这种信息纳入运算契约，否则对子矩阵的第二行开始就会寻址错误。紧凑方阵会让 $ld=W$ 恰好成立，因此测试必须包含填充和非方形视图，才能证明代码真正使用了物理步长。
\discussionpoint{连续维应与 warp 的主要变化方向配合。}
若一个 warp 的 lane 依次改变列号，行主序能让它们请求相邻元素；若 lane 依次改变行号，列主序才提供同样的连续性。转置、结构数组到数组结构转换以及张量维度置换，本质上都在选择哪个逻辑维成为物理连续维，以适应后续算子的主要访问。转换本身会产生额外读写，所以不能为了一个短内核盲目重排整个数据集；若许多后续阶段共享新布局，成本才可能被摊薄。布局选择因而是跨算子的数据流决策，而不只是单个索引公式的局部优化。端到端模型应把一次重排成本按后续复用次数摊销，再与不重排的跨距事务比较。
\discussionpoint{跨库互操作需要同时转换布局、尺寸和步长。}
与列主序 BLAS 互操作时，可以真实转置缓冲区，也可以交换操作数并调整转置标志，从代数上重新解释同一位模式；后一方案节省搬运，却要求行列尺寸和 leading dimension 一起修正。若只切换一个转置标志，方阵或单位矩阵常会掩盖错误，因为转置后形状不变或数值仍相同。更有诊断力的测试使用 $2\times3$ 等非方形矩阵，填入不对称递增值，并检查四角与跨行位置。这样既能验证公式，也能证明库调用的元数据与实际内存布局一致。若调用处理子矩阵，还应加入大于逻辑行列数的 leading dimension，专门测试填充区域未被误读。
\variantheading 对第 499 行、第 399 列，行主序索引为 $499\times400+399=199999$，列主序同样为 $399\times500+499=199999$。
\practiceheading 与 cuBLAS 交互时应显式记录 leading dimension 和转置标志，并用小型非方阵测试防止行列主序在方阵上偶然掩盖错误。
\sourcecheck{https://docs.nvidia.com/cuda/cuda-c-programming-guide/}{CUDA C++ 数组采用 C/C++ 的零基下标；高置信度}
\end{solutionexercise}

\begin{solutionexercise}{5}
\problemheading 考虑一个宽度为 400、高度为 500、深度为 300 的三维张量。该张量按行主序
存储为一维数组。请指定 $x=10$、$y=20$、$z=5$ 的张量元素的数组索引。

\approachheading 令 $x$ 为最快变化维，先线性化一个 $xy$ 平面，再加第 $z$ 个平面偏移。

\answerheading
\[
 i=x+yW+zWH=10+20\times400+5\times400\times500=1{,}008{,}010.
\]
边界检查为 $0\le x<400,0\le y<500,0\le z<300$，给定坐标满足条件。

\conclusionheading 数组索引为 1,008,010。
\discussionheading
\discussionpoint{三维线性化仍是逐层跨过完整低维区域。}
要到达切片 $z$，必须先跨过 $z$ 个完整的 $W\times H$ 平面；进入该切片后，到达行 $y$ 又要跨过 $y$ 个宽度为 $W$ 的行，最后才加列偏移 $x$。于是公式自然成为 $zWH+yW+x$，它把 $(z,y,x)$ 看成一个 x 维最快的混合进位数。更高维张量只需继续在外层乘上所有更快维度的容量，但维度乘积应使用足够宽的整数类型，否则体数据还未耗尽显存，32 位偏移就可能先溢出。线性化顺序也是数据格式的一部分，改变最快维意味着所有生产者与消费者都要同步改变。
\discussionpoint{三维 pitched memory 只填充行，切片跨距仍用逻辑高度。}
\code{cudaMalloc3D} 可把每行字节宽度向上填充为 \code{pitch}，所以地址应写成 $z\,\text{slicePitch}+y\,\text{pitch}+x\,\code{sizeof(T)}$。返回的 \code{cudaPitchedPtr.ysize} 表示一个完整二维切片的逻辑行数，官方典型公式正是 $\text{slicePitch}=\text{pitch}\times\text{ysize}$；“物理高度”并不是另一个被自动填充的字段。若参数只描述所属分配中的子视图，跨到下一 z 切片时仍须使用完整分配的 slice pitch，不能改乘子视图高度。这个边界说明 pitch 不是抽象宽度的替代词，而是必须随分配元数据一起传递的物理字节跨距。
\discussionpoint{张量布局是在为未来访问选择连续方向。}
题目让 x 连续，但图像批次可能采用 NCHW 或 NHWC，推理内核还可能把通道分成固定宽度的 blocked layout，以便向量指令或矩阵单元消费。某种布局对卷积读取友好，未必对归一化、转置或注意力友好，框架因此有时插入显式转换，有时选择能直接读取现有布局的内核。转换会遍历整个张量并产生额外流量，只有后续多个算子共享收益时才值得。理解 $zWH+yW+x$ 的真正价值，是能够针对任意维序和 stride 写出地址，而不是把一个公式硬套到所有张量格式。
\discussionpoint{体数据测试必须覆盖面、棱、角和跨切片位置。}
医学成像、流体网格和三维卷积都使用体数组，但它们还会叠加各向异性体素、halo 与不同边界条件，索引错一维往往仍生成平滑图像而不易肉眼发现。可给每点编码 $10^6z+10^3y+x$，检查八个角、相邻切片和带 padding 的非立方 extent，从数值中直接识别哪个坐标被写错。只验证最后一个紧凑元素不足以证明 pitch 正确，因为首行与末尾有时恰好避开填充。再配合内存边界工具检查字节地址，才能同时证明坐标映射、物理跨距和分配范围三份契约一致。
\variantheading 对同一张量的 $(x,y,z)=(399,499,299)$，索引为 $399+499\times400+299\times400\times500=59{,}999{,}999$。
\practiceheading 使用 \code{cudaMalloc3D} 时应按返回的 pitch 和 slice pitch 寻址，并以端点坐标单元测试及 Compute Sanitizer memcheck 验证字节偏移。
\sourcecheck{https://docs.nvidia.com/cuda/cuda-c-programming-guide/}{CUDA 多维数据的线性寻址示例；高置信度}
\end{solutionexercise}

\chapter{计算架构与调度}

\begin{solutionexercise}{1}
\problemheading 考虑下面的 CUDA 内核及调用它的主机函数：
\begin{lstlisting}[language=C++,firstnumber=1]
__global__ void foo_kernel(int* a, int* b) {
    unsigned int i = blockIdx.x*blockDim.x + threadIdx.x;
    if (threadIdx.x < 40 || threadIdx.x >= 104) {
        b[i] = a[i] + 1;
    }
    if (i % 2 == 0) {
        a[i] = b[i] * 2;
    }
    for (unsigned int j = 0; j < 5 - (i % 3); ++j) {
        b[i] += j;
    }
}
void foo(int* a_d, int* b_d) {
    unsigned int N = 1024;
    foo_kernel<<<(N + 128 - 1)/128, 128>>>(a_d, b_d);
}
\end{lstlisting}
\begin{enumerate}[label=\alph*.,leftmargin=2.2em]
  \item 每个线程块中有多少个 warp？
  \item 网格中有多少个 warp？
  \item 对第 04 行语句：
    \begin{enumerate}[label=\roman*.,leftmargin=2em]
      \item 网格中有多少个活动 warp？
      \item 网格中有多少个分歧 warp？
      \item 块 0 的 warp 0 的 SIMD 效率（百分比）是多少？
      \item 块 0 的 warp 1 的 SIMD 效率（百分比）是多少？
      \item 块 0 的 warp 3 的 SIMD 效率（百分比）是多少？
    \end{enumerate}
  \item 对第 07 行语句：
    \begin{enumerate}[label=\roman*.,leftmargin=2em]
      \item 网格中有多少个活动 warp？
      \item 网格中有多少个分歧 warp？
      \item 块 0 的 warp 0 的 SIMD 效率（百分比）是多少？
    \end{enumerate}
  \item 对第 09 行循环：
    \begin{enumerate}[label=\roman*.,leftmargin=2em]
      \item 有多少次迭代没有分歧？
      \item 有多少次迭代产生分歧？
    \end{enumerate}
\end{enumerate}

\approachheading warp 大小为 32。先把每个 128 线程块分成四个 warp，再对每条语句
逐 warp 写出参与线程集合；SIMD 效率按“执行该语句的活动 lane 数/32”计算。

\answerheading 网格含 $\lceil1024/128\rceil=8$ 个块，每块 $128/32=4$ 个 warp，
全网格共有 $8\times4=32$ 个 warp。

对第 04 行，块内条件为 $t<40$ 或 $t\ge104$。warp 0 的 32 个线程全参与；warp 1
只有 $t=32,\ldots,39$ 的 8 个线程参与；warp 2 无线程参与；warp 3 只有
$t=104,\ldots,127$ 的 24 个线程参与。因此每块有 3 个活动 warp、2 个分歧 warp，
全网格分别为 24 和 16 个。块 0 的 warp 0、1、3 的 SIMD 效率依次为
$100\%$、$8/32=25\%$、$24/32=75\%$。

对第 07 行，连续整数索引中偶数线程执行、奇数线程不执行。32 个 warp 全都活动且
全都分歧；块 0 的 warp 0 有 16 个活动 lane，效率为 $50\%$。

第 09 行循环的迭代次数由 $i\bmod3$ 决定，分别为 5、4、3。$j=0,1,2$ 时所有线程
仍在循环中，共 3 次无分歧迭代；$j=3$ 和 $j=4$ 时只有部分线程继续，共 2 次分歧
迭代。

\conclusionheading (a) 4；(b) 32；(c) 24、16、$100\%$、$25\%$、$75\%$；
(d) 32、32、$50\%$；(e) 3 次无分歧、2 次分歧。
\discussionheading
\discussionpoint{分歧发生在动态 warp，而不是整个线程块。}
若块内 warp 0 全走真分支、warp 1 全走假分支，两组线程虽然选择不同，硬件仍可分别执行，没有 warp 内分歧；只有同一 warp 的 lane 在同一动态分支上分成两组，才需要用不同活动掩码依次处理路径。本题的区间边界恰好切过 warp 1 和 warp 3，所以它们分歧，完整落在区间内外的 warp 则一致。循环还会在每次回边重新判断参与者，静态源码只有一个 \code{for}，动态执行却产生多个不同掩码。这个层次区别也适用于变长序列和稀疏行：数据类别必须按 warp 分布分析，不能只给出全块比例。
\discussionpoint{短分支可能被谓词化，但闲置 lane 仍然存在。}
编译器有时不用显式跳转，而给几条短指令附上谓词，让 warp 发射两侧指令，仅允许满足条件的 lane 提交结果。源码级分支看似消失，活动掩码的成本却没有凭空消失：不满足谓词的 lane 仍占据发射槽，只是不改变架构状态。较长分支、嵌套循环和独立线程调度还会引入不同的重汇合行为，不能一律用“两条路径时间相加”预测。理解谓词化的意义，是知道源码分支数与机器分支数可能不同，因此最终诊断要结合生成指令和活动线程，而不是只数 \code{if}。比较编译选项前还应检查 SASS，确认谓词化边界是否真的发生变化。
\discussionpoint{重排数据能减少分歧，却可能付出更大代价。}
若条件由数据类别决定，可以先扫描并按类别分桶，让每个 warp 尽量处理同一类对象；这在一条路径包含大量计算、类别在多个阶段复用时可能值得。重排本身要读取、写入索引并可能进行全局同步，若分支体只有一次加法，它的成本往往比原分歧更高。类别数量多或频繁变化时，分桶还会破坏原有空间局部性和稳定顺序，需要额外元数据恢复结果。因而优化问题不是“能否消灭分歧”，而是“节省的被屏蔽指令是否超过分类、搬运与顺序恢复的总成本”。输入分布改变时还要重新评估，因为稳定训练集上的桶比例未必代表线上请求。
\discussionpoint{诊断需要构造具有不同活动掩码的输入。}
只用均匀随机数据测一次，很难判断慢是由分支、访存还是尾部造成。更有解释力的实验让条件分别全真、全假、lane 交替和按半 warp 分组，再比较每条源码的动态实例、平均活动 lane 和总时间；这些输入具有相同规模，却形成截然不同的掩码。若全真与全假都快而交替明显变慢，分歧解释才获得支持；若四者接近，瓶颈可能在内存或另一阶段。测量应围绕可证伪的假设设计，而不是堆叠一串指标后凭相关性下结论。还应固定每条路径的数据地址，避免掩码实验同时改变缓存命中而混淆因果。
\variantheading 若第 3 行条件改为 $t<32$ 或 $t\ge96$，每块有 2 个活动 warp且都不分歧，全网格为 16 个活动 warp、0 个分歧 warp。
\practiceheading 用 Nsight Compute 的 Branch Efficiency、Warp State Statistics 和源码关联视图验证动态分支，而不要仅从源码推断实际指令。
\sourcecheck{https://docs.nvidia.com/cuda/cuda-c-programming-guide/}{NVIDIA CUDA C++ Programming Guide 对 32 线程 warp、活动线程与分歧的定义；逐 lane 枚举复核，置信度高}
\end{solutionexercise}

\begin{solutionexercise}{2}
\problemheading 对向量加法，假设向量长度为 2000，每个线程计算一个输出元素，线程块
大小为 512。网格中有多少个线程？

\approachheading 块数必须向上取整，实际启动线程数等于块数乘块大小。

\answerheading
\[
 B=\left\lceil\frac{2000}{512}\right\rceil=4,
 \qquad T_{\rm launch}=4\times512=2048\ \text{个线程}.
\]
其中 2000 个线程处理元素，48 个线程由边界检查屏蔽。

\conclusionheading 网格启动 2048 个线程（4 个线程块）。
\discussionheading
\discussionpoint{任意数组长度必须适配固定的执行粒度。}
数组可以有 2000 个元素，硬件却按完整线程块分派、按完整 warp 发射，因此四个 512 线程块覆盖到索引 2047 是正常的离散化结果。边界谓词只保证最后 48 个线程不产生数组副作用，并不会让硬件把它们从块中删掉；它们仍执行索引和条件判断。这个现象与存储系统按整页分配、网络按整包传输相似：逻辑有效载荷不必填满物理容器。评估冗余时应区分必须支付的容器开销和可以通过改变任务组织减少的重复工作。这里的 48 个空线程只占一个块尾部，不能据此推断整个网格都以同样比例浪费。
\discussionpoint{专用尾内核只有在尾部足够重要时才划算。}
理论上可以让主体内核只处理完整块，再启动一个更小的内核计算最后 464 项，从而减少 48 个空线程；但第二次启动、参数设置和代码维护通常远贵于几十次条件判断。若程序连续处理数百万个短向量，或每个尾元素执行昂贵且分歧严重的算法，把尾部批量收集到专门内核才可能收回成本。另一选择是融合多个小数组，让一个块处理一批对象，同时减少启动次数与空槽。优化应比较端到端时间，而不是以“没有空线程”为目标，因为消除局部浪费可能引入更大的全局开销。专用路径还会扩大测试矩阵，必须覆盖主体与尾部交界处的不重不漏。
\discussionpoint{块内尾线程与网格最后一波是两个独立问题。}
48 个无效线程位于最后一个块内部，它们影响该块的活动 lane；“最后一波不足”则取决于整个网格有多少块相对所有 SM 的并发块容量。即使 $N$ 恰好整除 512、块内没有一个尾线程，只有 4 个块的网格在拥有许多 SM 的设备上仍可能从一开始就填不满；反过来，长网格有数百波时，最后一个不满块的固定成本很容易摊薄。因此不能用 $48/2048$ 推导设备利用率，也不能把小网格尾延迟归因于这 48 个线程。前者可通过块宽改善，后者需要更多独立块或任务融合，修复手段并不相同。
\discussionpoint{块大小应在真实长度分布上选择。}
把块从 512 改成 256，本题仍启动 2048 个线程，却会形成 8 个块，给调度器更细粒度的分配机会，同时改变每块寄存器与共享内存总量。更小并不总更好：块级协作空间会缩小，固定开销增加，某些归约还需要更多阶段。可以先用资源模型排除无法获得合理驻留数的配置，再在业务实际出现的短、中、长向量上测量中位和高分位延迟。这样选出的块宽反映长度分布与目标架构，而不是只对 $N=2000$ 这个单点偶然合适。若网格本来很小，还应单独记录 SM 覆盖率，防止较大块造成并行块数量不足。
\variantheading 若长度为 2049、块大小仍为 512，则需 5 个块、2560 个线程，其中 511 个线程越界。
\practiceheading 在生产配置中用 \code{cudaOccupancyMaxPotentialBlockSize} 给出候选块大小，再通过 Nsight Compute 比较占用率与尾块浪费。
\sourcecheck{https://docs.nvidia.com/cuda/cuda-c-programming-guide/}{NVIDIA CUDA C++ Programming Guide 的线程层次与网格维度定义；算式独立复核，置信度高}
\end{solutionexercise}

\begin{solutionexercise}{3}
\problemheading 对上一题，考虑长度为 2000 的边界检查，预计有多少个 warp 会产生分歧？

\approachheading 按 32 个连续索引一组检查最后的有效/无效线程分界。

\answerheading $2000=62\times32+16$。前 62 个 warp 全部有效；下一个 warp 中
16 个线程有效、16 个线程越界，因而发生分歧；最后一个 warp 的 32 个线程全部越界，
整个 warp 走同一分支，不发生分歧。

\conclusionheading 只有 1 个 warp 因边界检查发生分歧。
\discussionheading
\discussionpoint{没有分歧的 warp 也可能完全没有有效工作。}
$2000=62\times32+16$，第 63 个 warp 有 16 个有效 lane 和 16 个越界 lane，因而在边界分支处分歧；再后的 warp 若全部越界，会一致跳过主体，分支效率反而可能是百分之百。后者仍被创建、调度并执行条件，只是没有 lane 写输出，所以“零分歧”不能推出“利用率高”。同理，一个全真分支可以无分歧却执行昂贵路径，一个半真分支也可能因主体极短而成本很小。性能语言必须把控制一致性与有用工作比例分开，否则优化会追逐错误目标。分析器中的活动线程指标与输出计数应同时使用，才能识别这种全空 warp。
\discussionpoint{余数决定部分 warp，块取整决定全空 warp。}
若按连续索引映射，$N\bmod32$ 给出最后一个部分 warp 的有效 lane 数；当余数为零时，边界恰落在 warp 之间，部分 warp 数为零。块大小再向上取整到完整线程块，可能在其后留下若干完全空 warp，本题 512 线程块的最后一块便同时含有部分和空闲区域。分别统计有效线程、部分 warp 和全空 warp，能精确解释边界判断的动态行为。这个计数还可预测改变块宽后哪些成本会消失，而一项总体“分支效率”无法提供同样信息。若采用二维映射，余数应沿每一维分别计算，角块不能用一维公式直接代替。
\discussionpoint{为了消除一个尾分支而增加第二条路径常得不偿失。}
可以让主体只覆盖 $\lfloor N/32\rfloor\times32$ 个元素，再由一个 warp、另一个内核甚至 CPU 处理余数，从而使主体完全一致。代价是多一套调度与边界代码，若尾部只有一次加法，额外启动远大于 16 个被屏蔽 lane 的损失。带掩码的库原语可保持一条调用路径，但底层仍必须确保无效 lane 不读写越界，并没有消除物理空闲。只有尾部分支很长、被大量短数组反复触发，或能把许多尾部合并处理时，专门化才具有系统层面的收益。决定拆分前可建立阈值实验，找出尾部成本何时超过额外调度与同步成本。
\discussionpoint{变长数据把一个尾部扩散到每个样本。}
ragged batch、稀疏矩阵行和自然语言序列都具有不同长度，尾部不再只出现于整个数组末端，而可能在每个 warp 负责的样本中反复出现。按长度分桶可让相近样本共用 warp，padding 则用额外存储换规则控制流，压缩索引还能只调度有效项；三者分别支付排序、填充流量和间接寻址成本。最佳选择取决于长度分布是否稳定、后续算子能否复用分桶以及填充值是否参与昂贵计算。这个更广场景说明，本题的一次余数计算其实是负载整形问题的最小模型。线上长度分布若漂移，固定桶边界也应重新测量，而不是永久沿用离线最优值。
\variantheading 若向量长度为 1984，即 $62\times32$，边界落在 warp 边界上，分歧 warp 为 0，最后两个 warp 全部越界。
\practiceheading Nsight Compute 的 Branch Efficiency 应结合 Active Threads Per Warp 和 Executed Ipc 分析，避免把“无分歧但全被屏蔽”误判为高效。
\sourcecheck{https://docs.nvidia.com/cuda/cuda-c-programming-guide/}{NVIDIA CUDA C++ Programming Guide 的 SIMT 分支行为；边界余数独立复核，置信度高}
\end{solutionexercise}

\begin{solutionexercise}{4}
\problemheading 假设一个有 8 个线程的线程块在到达屏障前执行一段代码。8 个线程执行
这段代码所需时间（微秒）分别为 2.0、2.3、3.0、2.8、2.4、1.9、2.6 和
2.9，其余时间都在等待屏障。线程总执行时间中有百分之多少用于等待屏障？

\approachheading 屏障释放时间由最慢线程的 $3.0\ \mu\mathrm{s}$ 决定；把每个线程
相对该时刻的空闲时间相加，再除以 8 个线程的总驻留时间。

\answerheading 有效执行时间之和为
\[
2.0+2.3+3.0+2.8+2.4+1.9+2.6+2.9=19.9\ \mu\mathrm{s}.
\]
到屏障释放为止的总线程时间为 $8\times3.0=24.0\ \mu\mathrm{s}$，故等待时间为
$24.0-19.9=4.1\ \mu\mathrm{s}$，比例为
\[
 \frac{4.1}{24.0}\times100\%=17.0833\%\approx17.1\%.
\]

\conclusionheading 约 $17.1\%$ 的线程总执行时间用于等待屏障。
\discussionheading
\discussionpoint{把等待量看成关键路径周围的“松弛”。}
为什么八个线程都做了有用工作，仍能累计出 $4.1\ \mu\mathrm{s}$ 的等待？在题设的同时开始模型中，屏障释放时刻是 $t_{\max}=3.0\ \mu\mathrm{s}$，线程 $i$ 的松弛量为 $t_{\max}-t_i$，所以等待比例一般写成 $\sum_i(t_{\max}-t_i)/(n t_{\max})$。这个比值既取决于最慢者，也取决于其余时间分布；仅把最快线程再加速，甚至可能让总等待比例上升而不缩短关键路径。它与并行任务图中的关键路径和工作量下界相呼应，但结论依赖“同时开始、持续占用且时间可相加”的抽象，不能直接当作真实 GPU 的能耗或利用率。

\discussionpoint{单线程计时怎样映射到 warp 调度。}
真实 GPU 不会给八个 lane 各配一个完全独立的指令调度器，因此题中的八个时间不能机械解释成八条硬件时间线。同一 warp 内的差异通常来自控制流分歧、不同长度的依赖链或不同访存完成时刻，而调度器发射的是 warp 指令；若这些线程分属多个 warp，其他就绪 warp 又可能在等待期间发射，从而隐藏一部分延迟。屏障自身还有到达速率和同步指令吞吐量，最后一个 warp 何时到达才决定释放。因此 $17.1\%$ 是给定教学模型下的线程时间占比，不是分析器中 SM 利用率、占用率或“stall barrier”百分比的同义词。

\discussionpoint{减少等待不能等同于删除屏障。}
若慢线程处理高次数循环而快线程只处理短任务，首先应问工作能否重新分桶、用动态队列分配，或让每个 warp 承担更均匀的批次；这些做法缩小到达时间分布，才真正缩短屏障前关键路径。某些算法还能让装载 warp 与计算 warp 专门化，或把相邻批次流水化，使等待与独立工作重叠。可是只要下一阶段会读取所有生产者写入的数据，直接删掉屏障就破坏到达条件与内存可见性，性能问题会转化成数据竞争。更一般的设计原则是先画出生产者--消费者依赖，再缩小同步作用域、减少同步次数或均衡作用域内工作，而不是凭一次计时猜测屏障“多余”。

\discussionpoint{从症状反推真正的慢到达者。}
分析器报告某 warp 因 barrier 停顿，只说明它已经到达而同组参与者尚未满足条件，并不告诉我们迟到者是在算术、分支还是内存依赖上耗时。较可靠的实验会分别关闭或替换候选阶段，记录阶段级 GPU 时间戳，并比较线程块之间与块内 warp 之间的到达分布；若短阶段被时间戳写入本身明显扰动，还应改用多个独立内核和事件做差分。随后把 barrier stall 与分支效率、内存延迟、eligible warps 和调度器发射情况一起看，才能判断均衡工作是否有效。这样的诊断也适用于 CPU 线程池和分布式栅栏：等待者容易观测，真正需要优化的却通常是尚未到达的路径。
\variantheading 若最慢线程由 $3.0$ 降至 $2.9\ \mu\mathrm{s}$，有效执行时间之和变为 $19.8\ \mu\mathrm{s}$，总线程时间为 $23.2\ \mu\mathrm{s}$，等待为 $3.4\ \mu\mathrm{s}$，比例约 $14.7\%$。
\practiceheading 用 Nsight Compute 的 Barrier、Warp Stall Barrier 和 Scheduler Statistics 指标区分屏障等待与内存依赖停顿。
\sourcecheck{https://docs.nvidia.com/cuda/cuda-c-programming-guide/}{NVIDIA CUDA C++ Programming Guide 对块级屏障的语义；时间总量独立求和，置信度高}
\end{solutionexercise}

\begin{solutionexercise}{5}
\problemheading CUDA 程序员认为，如果每个线程块只启动 32 个线程，那么在需要屏障同步
的地方可以省略 \code{__syncthreads()}。你认为这是好主意吗？请解释原因。

\approachheading 区分“线程属于同一 warp”和“程序获得了内存可见性及 happens-before
保证”。前者不能自动替代后者。

\answerheading 这不是可靠做法。自 Volta 起，独立线程调度允许同一 warp 的线程以
更细粒度暂停和推进；依赖隐式 lockstep 的 warp-synchronous 代码可能读到尚未写入的
共享数据。即使在较早架构上，编译器重排和内存可见性也不应靠未写入程序的调度假设
保证。若通信确实只涉及一个 warp，可用带正确参与掩码的 \code{__syncwarp(mask)}；
若语义要求整个线程块到齐并让先前共享/全局内存访问对块内线程可见，就必须使用
\code{__syncthreads()}。所有参与线程还必须以一致控制流到达块级屏障。

\conclusionheading 不能仅因块大小为 32 就删去同步。应根据所需作用域选用 warp
同步或块级同步。
\discussionheading
\discussionpoint{锁步直觉缺少哪一半语义。}
假设 lane 0 写共享变量而 lane 1 随后读取，即使人眼觉得一个 warp 会“整齐向前”，程序仍需要同时回答两个问题：消费者是否等到生产者到达，以及写入是否按规定对消费者可见。屏障把到达同步与相应作用域内的内存顺序组合起来，未写入程序的调度习惯并不能替代这份契约。自 Volta 起的独立线程调度使同一 warp 中不同 lane 可以更细粒度地暂停和推进，让旧式 warp-synchronous 技巧更容易暴露；但即使面向更早设备，也不应依赖编译器重排和硬件时序的偶然结果。把同步写明还能使代码审查围绕 happens-before 关系进行，而不是围绕某一代 GPU 的经验猜测。

\discussionpoint{同步原语应覆盖恰好需要通信的参与者。}
若数据只在掩码指定的一个 warp 内交换，\code{__syncwarp(mask)} 可以等待这些 lane 并提供规定的内存排序；若生产者和消费者跨越多个 warp，块级 \code{__syncthreads()} 才覆盖完整通信域。再扩大到线程块集群、整个网格或不同内核时，需要 cooperative groups、受支持的集群同步、协作启动，或以内核边界和事件建立顺序，不能把块级屏障想象成全设备栅栏。作用域越大，必须同时到达的工作越多，慢参与者带来的等待也越广。因而选择“最小但充分”的作用域既是正确性问题，也是可扩展性问题，这与 CPU 原子操作中的内存序和分布式系统中的集合通信具有相同结构。

\discussionpoint{边界线程提前退出为何会破坏整个块。}
设一个二维块覆盖图像右下角，只有部分线程对应有效像素；若无效线程在共享装载前直接 \code{return}，剩余线程随后执行块级屏障，参与集合就不再一致，可能挂起或产生未定义行为。安全结构是让所有线程用谓词装载有效输入或零，所有线程无条件到达屏障，然后只让有效线程计算和写回。类似地，把 \code{__syncthreads()} 放进仅部分线程成立的普通分支也不安全，除非能证明条件对整个块一致。这个反例说明边界检查的位置属于同步证明的一部分：谓词可以屏蔽数据操作，却不能随意改变集体操作的成员集合。

\discussionpoint{动态检查只能辅助依赖证明。}
Compute Sanitizer 的 racecheck 能暴露许多共享内存冲突，synccheck 能发现一部分非法屏障或 warp 掩码使用，但一次测试通过并不证明所有输入和调度下都存在正确的 happens-before。更稳妥的审查会为每个共享数组列出写者、读者、索引域和跨线程边，再在每条边上标注负责排序的同步原语；如果某条边跨块，块级屏障立即显得不充分。还应构造尾块、分支两侧和不同 warp 到达次序等对抗用例，而不是只运行整齐尺寸。这种“依赖图加动态检测”的组合也能迁移到 CPU 并发代码：工具寻找实际竞态，结构化证明覆盖尚未发生的调度。
\variantheading 若通信只发生在 lane 0--15，可用参与掩码 \code{0x0000ffff} 调用 \code{__syncwarp}，前提是这 16 个 lane 都会到达该调用。
\practiceheading 采用 Compute Sanitizer racecheck/synccheck 检查共享内存竞态和非法屏障参与，并为 Volta 及更新架构编译测试独立线程调度路径。
\sourcecheck{https://docs.nvidia.com/cuda/cuda-c-programming-guide/}{NVIDIA CUDA C++ Programming Guide 的独立线程调度与同步函数语义；对抗审查确认隐式 warp 同步不是可移植保证}
\end{solutionexercise}

\begin{solutionexercise}{6}
\problemheading 如果 CUDA 设备的 SM（流式多处理器）最多容纳 1536 个线程和 4 个线程块，
下列哪一种线程块配置会使 SM 中的线程数最多？
\begin{enumerate}[label=\alph*.,leftmargin=2.2em]
  \item 每块 128 个线程；
  \item 每块 256 个线程；
  \item 每块 512 个线程；
  \item 每块 1024 个线程。
\end{enumerate}

\approachheading 对块大小 $T$，驻留块数为 $\min(4,\lfloor1536/T\rfloor)$。

\answerheading
\begin{center}
\begin{tabular}{c@{\qquad}c@{\qquad}c}
\toprule
每块线程数 & 驻留块数 & 驻留线程数\\
\midrule
128  & 4 & 512\\
256  & 4 & 1024\\
512  & 3 & 1536\\
1024 & 1 & 1024\\
\bottomrule
\end{tabular}
\end{center}

\conclusionheading 配置 (c)，每块 512 个线程，使 SM 中驻留线程数最多，为 1536。
\discussionheading
\discussionpoint{占用率首先是一个整数装箱问题。}
为什么 512 线程块能装满 1536 个线程，而 1024 线程块反而只能驻留 1024 个？对题中只考虑线程数与块数的模型，块大小为 $T$ 时驻留块数是 $B=\min(4,\lfloor1536/T\rfloor)$，驻留线程数才是 $BT$；资源不能按半个块分配，所以每次向下取整都会留下空隙。由此得到的是阶梯函数而非连续曲线，块大小略变就可能跨过一个整数边界。这个现象与内存分页、箱装和批处理中的内部碎片相似，也解释了为何“块越大越接近上限”的直觉会在 1024 线程处失效。

\discussionpoint{题目省略的资源会怎样改变答案。}
手算中 512 线程块能驻留三块，是因为我们按题意假定寄存器、共享内存和 warp 数都不限制；真实内核却要同时满足每 SM 寄存器总量、每块静态与动态共享内存、驻留 warp 数以及架构规定的分配粒度。若每块共享内存用量使 SM 只能放两块，同一配置便从 1536 线程降到 1024；若寄存器按 warp 粒度向上取整，简单的“每线程乘线程数”也可能偏乐观。编译器版本、目标 SM、循环展开和模板参数都会改变这些资源。因此这个答案是题给抽象下的确定结论，而不是同一源代码跨设备都能获得的常数。

\discussionpoint{更多驻留线程为何不保证更快。}
高占用率的主要价值是提供更多就绪 warp，使一个 warp 等待内存或依赖时调度器仍有工作可发射；它不是计算单元利用率或程序吞吐量本身。一个算术密集、指令级并行充分的内核可能用较少 warp 就填满管线，而满占用的内核若访问未合并、分支严重或每条指令都有长依赖，仍会很慢。块过小还可能削弱共享瓦片复用、归约树效率或特定矩阵指令的组织方式。因而“512 线程使驻留线程最多”只回答资源容量问题，性能选择还要同时看有效带宽、发射槽利用和每个输出的工作量。

\discussionpoint{把手算结果变成可解释的调优起点。}
工程上可以先用 Occupancy API 结合真实函数属性，排除因寄存器或动态共享内存而不可行的块大小，再在 128、256、512 等候选间进行基准测试。测试不应只选一个整齐大数组，还要覆盖生产中的尺寸分布，因为末波块数、缓存工作集和短内核启动成本都会改变排名；记录时间时同时保存寄存器、共享内存、实际驻留块和网格波次数，才能解释结果。成熟库常把最佳配置按架构、数据类型和形状缓存，而不是宣称一个全局最佳数字。这样保留了题目整数模型的预测力，也允许测量指出另一个资源或算法瓶颈已经成为主导。
\variantheading 若块上限提高到 8 而线程上限仍为 1536，则 256 线程/块可驻留 6 块并达到 1536 线程，与 512 线程/块并列最优。
\practiceheading 用 CUDA Occupancy API 或 Nsight Compute Occupancy 面板输入真实寄存器和动态共享内存用量，避免仅按线程数选择块大小。
\sourcecheck{https://docs.nvidia.com/nsight-compute/NsightCompute/index.html}{NVIDIA Nsight Compute Occupancy Calculator 对线程与块资源上限的联合约束；逐项整数除法复核，置信度高}
\end{solutionexercise}

\begin{solutionexercise}{7}
\problemheading 假设设备每个 SM 最多允许 64 个线程块和 2048 个线程。指出下列每种
每个 SM 的分配是否可行；若可行，给出占用率：
\begin{enumerate}[label=\alph*.,leftmargin=2.2em]
  \item 8 个块，每块 128 个线程；
  \item 16 个块，每块 64 个线程；
  \item 32 个块，每块 32 个线程；
  \item 64 个块，每块 32 个线程；
  \item 32 个块，每块 64 个线程。
\end{enumerate}

\approachheading 题目只给出块数和线程数限制，故检查两项上限，并以驻留线程数除以
2048。这里不额外假设寄存器、共享内存或架构分配粒度。

\answerheading
\begin{center}
\begin{tabular}{c@{\quad}c@{\quad}c@{\quad}c}
\toprule
小题 & 分配 & 驻留线程 & 可行性与占用率\\
\midrule
(a) & $8\times128$  & 1024 & 可行，$50\%$\\
(b) & $16\times64$  & 1024 & 可行，$50\%$\\
(c) & $32\times32$  & 1024 & 可行，$50\%$\\
(d) & $64\times32$  & 2048 & 可行，$100\%$\\
(e) & $32\times64$  & 2048 & 可行，$100\%$\\
\bottomrule
\end{tabular}
\end{center}

\conclusionheading 五种分配都可行；占用率依次为 $50\%,50\%,50\%,100\%,100\%$。
\discussionheading
\discussionpoint{同样的驻留线程数可以代表不同的调度结构。}
配置 (d) 与 (e) 都放入 2048 个线程并得到题设意义下的 $100\%$ 占用率，但前者是 64 个 32 线程块，后者是 32 个 64 线程块；相同总数并没有抹去块边界。较多小块给调度器更多独立完成与替换的单位，却让每块只能在较小范围内共享数据和同步；较大块能组织跨 warp 的归约或瓦片，也可能让一个慢块占据资源更久。若网格总块数很少，两种配置的末波利用率还会不同。因此占用率是一个容量投影，不是对块级协作、负载均衡和调度尾部的完整描述。比较两者时应保持算法工作相同，避免把块形状改变引起的复用差异误算成调度收益。

\discussionpoint{资源分配粒度会制造看不见的碎片。}
题目允许把 64 个单 warp 块精确装入 SM，是因为只给了线程数和块数上限；实际设备往往按 warp 或块的粒度分配寄存器、共享内存和调度状态。假如每块哪怕保留一份固定共享结构，64 个小块的固定开销就可能先触及块或内存限制，而较大块能摊薄它；反过来，单个大块的资源向上取整也可能留下无法再容纳另一块的空洞。动态共享内存和模板特化还会在运行或编译配置间改变这个结果。故比较等占用配置时必须重新读取内核属性并做多资源的整数装箱，不能只把总线程数相除。实际分配粒度应查目标架构资料或工具输出，不能从一个占用率结果反向猜定。

\discussionpoint{理论占用率与执行利用率回答不同问题。}
理论占用率表示可驻留 warp 数相对硬件上限的比例，却没有保证这些 warp 在某周期已就绪，更没有保证 ALU、加载存储或特殊函数管线都在做有用工作。长依赖链会使许多驻留 warp 同时等待，bank 冲突和缓存未命中会延长数据到达，分支分歧则让一次发射只有部分 lane 有效；于是 $100\%$ 占用率仍可能伴随很低吞吐。反之，只要较少 warp 已足以覆盖延迟，$50\%$ 也可能接近峰值。分析时应把 achieved occupancy、eligible warps、发射槽和各管线吞吐联系起来，而不是把一个百分数解释为“GPU 忙碌程度”。

\discussionpoint{块形状还编码了算法的数据几何。}
在矩阵转置中，同样 256 个线程可以排成 $32\times8$ 或 $16\times16$，两者对连续内存方向、共享存储布局和边界浪费的影响不同；在模板计算中，块边长还决定 halo 与有效输出的比例。归约需要足够线程形成树，而每线程多元素又会改变寄存器和访存数量，所以只比较线程总数会遗漏真正的算法代价。生产调优通常把二维或三维形状、每线程工作量、共享瓦片和资源上限作为联合参数，并用真实的非方形尺寸分布评价。题中五项计算教会我们先筛掉非法点，后续性能模型则要把任务几何重新放回来。
\variantheading 若设备块上限改为 32，则配置 (d) 的 64 块不可行，其余四项仍可行，(e) 仍以 2048 线程达到 $100\%$。
\practiceheading 在 Nsight Compute 中同时观察 Theoretical/ achieved Occupancy、Eligible Warps 和 SM Throughput，确认提高驻留数是否真正改善延迟隐藏。
\sourcecheck{https://docs.nvidia.com/nsight-compute/NsightCompute/index.html}{NVIDIA Nsight Compute Occupancy Calculator 的占用率定义；只按题给资源进行对抗复算，置信度高}
\end{solutionexercise}

\begin{solutionexercise}{8}
\problemheading 某 GPU 的硬件限制为：每个 SM 2048 个线程、32 个线程块和 64 K（65,536）
个寄存器。对下列内核特征，判断内核能否达到满占用率；如果不能，指出限制因素：
\begin{enumerate}[label=\alph*.,leftmargin=2.2em]
  \item 每块 128 个线程，每线程 30 个寄存器；
  \item 每块 32 个线程，每线程 29 个寄存器；
  \item 每块 256 个线程，每线程 34 个寄存器。
\end{enumerate}

\approachheading 先求容纳 2048 个线程所需块数，再核对块数和寄存器总量。按题目
抽象模型使用“每线程寄存器数乘线程数”，不引入未给出的架构分配粒度。

\answerheading
\begin{enumerate}[label=\alph*.,leftmargin=2.2em]
\item 需要 $2048/128=16$ 个块，寄存器数为
$2048\times30=61{,}440<65{,}536$，且 $16<32$，可以达到满占用率。
\item 需要 $2048/32=64$ 个块，超过 32 块上限。最多驻留
$32\times32=1024$ 个线程，即 $50\%$；限制因素是每 SM 块数。
\item 需要 $2048/256=8$ 个块，但寄存器需求
$2048\times34=69{,}632>65{,}536$。每块用 $256\times34=8704$ 个寄存器，
最多驻留 $\lfloor65536/8704\rfloor=7$ 个块，即 1792 个线程、$87.5\%$；限制因素
是寄存器。
\end{enumerate}

实际 GPU 还按架构规定的粒度分配寄存器，实测占用率可能更低，应以目标设备属性和
占用率计算器为准。

\conclusionheading (a) 可以；(b) 受块数限制；(c) 受寄存器限制。
\discussionheading
\discussionpoint{简化的寄存器方程为何有用又为何不够。}
在题设模型中，只要把每线程寄存器数 $r$ 与目标驻留线程数 $T$ 相乘，就能看出 (c) 所需 $2048\times34=69{,}632$ 已超过 65,536；这个容量反例说明线程上限允许并不代表资源组合可行。真实 GPU 却常按 warp 或块的寄存器分配单元向上取整，所以每块实际消耗可能大于 $r$ 与线程数的朴素乘积，临界的 (a) 也应交给目标架构的 occupancy 计算器复核。教学方程适合揭示限制因素和给出上界，分配粒度则解释为何实测驻留块数有时低于手算值。两层模型并不矛盾，而是分别服务于推理与部署。

\discussionpoint{寄存器数是编译产物而非源码常量。}
同一段 CUDA C++ 在循环展开、函数内联、模板实例化或目标 SM 改变后，变量的活跃区间和指令选择都会变化，因而 \code{ptxas -v} 报告的寄存器数只能归属于某个具体构建。程序员有时用 \code{--maxrregcount} 强迫编译器降低用量，希望多驻留一个块；若被挤出的值发生 spill，它们会进入 local address space，并可能增加缓存和显存流量，最终比低占用更慢。还要区分编译器报告的 local memory、stack 和 spill load/store，不能看到“local”就断言都是寄存器溢出。正确做法是保存构建参数与资源报告，再用运行时间验证这一交换是否真的获利。

\discussionpoint{算法复用与寄存器压力常来自同一个选择。}
让一个线程计算多个输出，可以把已加载的数据或系数重复使用，并暴露相互独立的乘加来提高指令级并行；代价是多个累加器和预取值同时活跃，寄存器数随粗化因子增长。软件流水线也会在当前数据计算时保留下一批数据，具有同样矛盾。若寄存器增加使驻留 warp 跨过一个装箱台阶，延迟隐藏可能突然下降；若复用减少了更多显存访问，较低占用仍可能更快。因此优化目标不是把每线程寄存器数压到最小，而是在数据复用、独立指令和可驻留并行度之间寻找可测量的平衡。扫描粗化因子时应同时记录 spill，防止把寄存器压力转移成隐藏的设备内存流量。

\discussionpoint{如何把资源跃迁纳入性能回归。}
一次编译器升级可能让模板实例多用两个寄存器，恰好使驻留块数从八个降为七个，即使源码 diff 看不出性能风险。可靠的回归应按目标架构保存寄存器数、静态与动态共享内存、stack/local 用量、spill 指令、理论和实际占用率以及内核时间，并用代表性尺寸识别是资源跃迁还是缓存变化。对靠近硬件上限的内核，还可在启动前查询函数属性并检查动态共享内存设置，避免换到资源较小设备才得到配置错误。这样的记录把题中的一次手算扩展为可追踪的构建契约，也让“更高占用却更慢”的结果有证据可解释。
\variantheading 若 (c) 每线程降为 32 个寄存器，则 8 块共需 65,536 个寄存器，可在题给简化模型下达到满占用率。
\practiceheading 用 \code{nvcc --ptxas-options=-v} 查看静态寄存器数，再调用 CUDA
Occupancy API，并用 Nsight Compute 核对实际驻留块数。
\sourcecheck{https://docs.nvidia.com/nsight-compute/NsightCompute/index.html}{NVIDIA Nsight Compute Occupancy Calculator 对寄存器、块和线程限制的说明；按题给简化模型复核，置信度高}
\end{solutionexercise}

\begin{solutionexercise}{9}
\problemheading 一名学生说，他用每块 $32\times32$ 个线程的矩阵乘法内核成功计算了两个
$1024\times1024$ 矩阵的乘法。该学生使用的 CUDA 设备每块最多允许 512 个线程、
每个 SM 最多允许 8 个线程块；并且每个线程计算结果矩阵的一个元素。你对此有
什么反应？为什么？

\approachheading 首先检查启动配置是否合法，而不是从矩阵维度推断程序已经正确执行。

\answerheading $32\times32=1024$ 个线程，超过设备每块 512 个线程的硬上限。因此
该内核启动应失败，不能完成所声称的计算；若程序没有紧接启动调用
\code{cudaGetLastError()} 并在适当位置同步检查异步错误，旧输出或未初始化输出可能
被误判为成功。每 SM 最多 8 个块并不能放宽单块线程上限。

可改用 $32\times16$ 或 $16\times16$ 的块，并把网格两维分别设为覆盖 1024 行、
1024 列的向上取整值；每个线程仍可计算一个结果元素。

\conclusionheading 该说法与硬件限制矛盾；原配置不可能合法启动，必须缩小线程块并
检查 CUDA 错误。
\discussionheading
\discussionpoint{先区分“不能启动”与“启动后可能较慢”。}
$32\times32=1024$ 超过题设每块 512 线程，是启动配置的硬约束：运行时不能把这个块拆成两个合法块，也不会以较低占用率勉强执行。每 SM 最多八块描述的是若干已经合法的块能否并发驻留，它与每块线程数、单维尺寸和每块共享内存上限处在不同层次，不能互相抵消。相反，寄存器较多而仍能驻留一块通常属于性能或资源装箱问题，不必然使启动非法。审查 CUDA 程序时先按“设备级、SM 级、块级”分层列约束，可以避免把软性能建议误当错误，也避免把真正的非法配置包装成占用率问题。

\discussionpoint{异步执行怎样制造“我算成功了”的假象。}
内核启动通常只把工作排入队列，主机若不立即检查 \code{cudaGetLastError()}，又不在读取结果前检查同步 API 的返回值，非法配置可能直到很晚才被发现。此时程序可能打印先前运行残留的缓冲区、初始化常量或根本未更新的数据；若测试矩阵全零、单位阵或输出未清空，错误结果甚至会碰巧像正确答案。可靠证据必须同时包含启动状态、完成状态和与独立 CPU 参考的逐元素比较，并让负测试确认非法配置确实被测试框架捕获。这个原则也适用于其他异步系统：提交成功不等于任务成功，更不等于结果正确。

\discussionpoint{合法块形状只是矩阵乘法重构的第一步。}
把块改为 $32\times16$ 会恰好满足 512 线程上限，改为 $16\times16$ 则给资源和调度更多余量；两者都必须重新计算网格覆盖和行列边界，才能保持每线程一个输出的语义。若追求性能，成熟 GEMM 往往让线程协作加载共享瓦片并让每线程累加多个输出，这会改变合并访问方向、数据复用、寄存器累加器和同步次数，而非简单缩小二维尺寸。还要检查矩阵维度不是块尺寸整数倍的尾块，以及 leading dimension、转置和批处理布局。由非法配置出发，可以看到“合法性、正确性、性能”是三道连续但不同的设计关卡。

\discussionpoint{跨设备配置应是一组经验证的契约。}
即使某台设备允许 1024 线程块，也不能据此否定题设设备的 512 上限，更不能把一次运行经验推广到全部部署节点。生产代码可以查询 \code{cudaDeviceProp} 和函数属性，按计算能力、数据类型与问题形状选择经过验证的配置族，或把选择交给具有调优记录的库；查询结果仍须与内核的动态共享内存和寄存器需求联合检查。持续集成应至少编译主要目标 SM，运行合法边界和故意超限的负测试，并确认错误从异步调用链传到测试结果。这样“曾经在一台 GPU 上运行”才会升级为可审计、可移植的部署证据。
\variantheading 改为 $32\times16$ 线程块后每块正好 512 线程，覆盖 $1024^2$ 输出需 $32\times64=2048$ 个块。
\practiceheading 每次启动后立即检查 \code{cudaGetLastError}，并在测试同步点检查执行错误；部署时查询 \code{cudaDeviceProp.maxThreadsPerBlock} 而非硬编码设备能力。
\sourcecheck{https://docs.nvidia.com/cuda/cuda-c-programming-guide/}{NVIDIA CUDA C++ Programming Guide 的每块最大线程数与内核启动错误规则；配置乘积独立复核，置信度高}
\end{solutionexercise}

\chapter{内存架构与数据局部性}

\begin{solutionexercise}{1}
\problemheading 考虑矩阵加法。能否使用共享内存减少全局内存带宽消耗？提示：分析每个
线程访问的元素，看看不同线程之间是否存在共同访问。

\approachheading 统计每个输入元素在一个内核调用中被多少线程使用。

\answerheading 对 $C_{ij}=A_{ij}+B_{ij}$，负责 $(i,j)$ 的线程各读取一个 $A_{ij}$、
一个 $B_{ij}$ 并写一个 $C_{ij}$；其他线程不会复用这两个输入。把它们先搬到共享内存
仍需同样的两次全局加载，反而增加共享内存读写和同步。因此，单纯分块不能降低必要的
全局内存字节数。只有当后续算子融合后会多次使用同一元素，或数据布局转换等额外目标
存在时，共享内存才可能有收益。

\complexity{$O(N^2)$ 工作}{$O(1)$ 每线程额外空间}{$N^2$ 个独立元素任务}
\conclusionheading 对纯矩阵加法，使用共享内存不能减少全局内存带宽消耗。
\discussionheading
\discussionpoint{先问共享副本能否被第二个消费者使用。}
对 $C_{ij}=A_{ij}+B_{ij}$，负责 $C_{ij}$ 的线程各读取两个输入一次并写一个输出，别的线程并不需要这两个输入；因此每个元素的必要流量已经是两读一写。若先把 $A_{ij}$ 和 $B_{ij}$ 写进共享内存，同一线程随后再读回来，只增加共享指令和可能的同步，却没有消掉任何首次全局读取。分块的收益来自跨线程或跨阶段复用，而不是“共享内存比全局内存快”这一孤立事实。这个反例提供了通用判断法：先画生产者--消费者关系并计算复用次数，复用为一时通常没有搬运到显式缓存的经济基础。

\discussionpoint{算子链会改变单个加法的流量下界。}
若矩阵加法之后马上执行缩放与激活，把三个算子分别启动会写出并重新读取两个中间矩阵；融合后，同一线程可让中间标量留在寄存器中，通常比把它绕经共享内存更直接。可是后续若是转置、邻域滤波或同一输入产生多个输出，消费者可能换成其他线程，此时共享内存就成为块内交换和复用的介质。由此可见“矩阵加法不适合共享分块”只针对孤立逐元素表达式，并不禁止围绕更大计算图设计瓦片。编译器与框架中的算子融合，正是把优化单位从一条公式提升到完整数据流。融合边界还受数值语义和中间结果是否被其他消费者复用约束，不能只按内核数量决定。

\discussionpoint{硬件缓存也不能消除必经的首次流量。}
连续的矩阵访问通常能被 warp 合并，L1 或 L2 也可能保留某些 cache line，但只读一次的输入至少要从更低层进入芯片一次；缓存命中不会凭空删除这次 compulsory traffic。主动复制到共享内存还占用片上容量，可能减少驻留块，并增加地址、存储、读取甚至屏障指令，所以“小矩阵能放下”并不构成优化证据。若数据已因另一个算子留在缓存，实测 DRAM 流量可能低于孤立模型，但这属于跨内核时间局部性而非共享分块的贡献。分析时应明确所讨论的是源码请求、L2 流量还是 DRAM 字节，避免把不同层次的“带宽减少”混为一谈。

\discussionpoint{用端到端字节而非单个内核外观评价融合。}
图像预处理或深度学习推理常把加法、偏置和激活合为一个 elementwise 内核，目标是省掉中间数组的写回、再次读取和额外启动，而不是让每一步都占一块共享瓦片。实验可分别运行三内核版本与融合版本，记录完整流水线时间、DRAM 读写字节、启动次数和输出误差；若只看融合内核的占用率，寄存器增加可能让数字变低却仍取得更高吞吐。再把实测 FLOP/s 与算术强度放到 Roofline 上，可判断节省字节是否让瓶颈向计算侧移动。这样的测量把本题的“无复用”结论扩展成可验证的编译和框架决策。
\variantheading 若每个 $A_{ij}$ 同时参与两次后续计算，协作加载一次可把该输入的两次全局读取降为一次。
\practiceheading 用 Nsight Compute 的 DRAM Bytes、L1/L2 命中率和 Barrier Stall 比较直接内核与融合内核。
\sourcecheck{https://docs.nvidia.com/cuda/cuda-c-best-practices-guide/}{NVIDIA CUDA Best Practices Guide 关于共享内存消除重复全局读取的条件；复用计数为 1，置信度高}
\end{solutionexercise}

\begin{solutionexercise}{2}
\problemheading 分别为使用 $2\times2$ 瓦片和 $4\times4$ 瓦片的 $8\times8$ 矩阵乘法，
画出图 5.7 的等价图。验证全局内存带宽的减少量确实与瓦片的维度成正比。

为使题意独立可读，图 5.7 描述的是一个 $4\times4$ 矩阵乘法的 $2\times2$ 输出瓦片：
线程 $(r,c)$ 在 $k=0,1,2,3$ 四个时刻依次读取 $M_{r,k}$ 与 $N_{k,c}$；同一输出行的
两个线程重复读取相同的 $M$ 行元素，同一输出列的两个线程重复读取相同的 $N$ 列元素。
本题要求对 $8\times8$ 情形分别画出同类“线程--时间--输入元素”访问关系，而不是复制
原书插图。

\approachheading 用“输出块数 $\times$ 阶段数 $\times$ 每阶段输入瓦片元素数”计数，
并把访问过程画成阶段表。

\editorialnote 本题所称英文底本图 5.7，在中文译稿中因新增示意图而重编号为图 5.8；
译稿练习中的“图 5.7”没有随正文图号更新。下文按底本图号叙述，并同时给出这一对应关系。

\answerheading 下表是原书图 5.7 的等价阶段图；$M_{p,q}$、$N_{q,r}$ 表示相应输入
瓦片，每个方格只从全局内存协作加载一次。

\begin{center}
\begin{tabular}{c@{\quad}c@{\quad}c@{\quad}c}
\toprule
输出瓦片 & 阶段 & $M$ 输入瓦片 & $N$ 输入瓦片\\
\midrule
$P_{p,r}$，$T=2$ & $q=0,1,2,3$ & $M_{p,q}$（$2\times2$） & $N_{q,r}$（$2\times2$）\\
$P_{p,r}$，$T=4$ & $q=0,1$ & $M_{p,q}$（$4\times4$） & $N_{q,r}$（$4\times4$）\\
\bottomrule
\end{tabular}
\end{center}

不分块时，64 个输出各读取 8 个 $M$ 和 8 个 $N$ 元素，总请求数为
$64\times16=1024$（每个输入矩阵 512 次）。$T=2$ 时有 $(8/2)^2=16$ 个输出块、
4 个阶段，每阶段每块加载 $2\times2\times2=8$ 个输入元素，总请求数
$16\times4\times8=512$（每个矩阵 256 次），正好减少 2 倍。$T=4$ 时有 4 个块、
2 个阶段，每阶段加载 $2\times4^2=32$ 个元素，总请求数 $4\times2\times32=256$
（每个矩阵 128 次），正好减少 4 倍。

\complexity{$O(8^3)$ 工作；一般为 $O(N^3)$}{$2T^2$ 个共享元素/块}{$(N/T)^2$ 个输出瓦片可并行}
\conclusionheading $2\times2$ 与 $4\times4$ 瓦片分别把输入全局请求减少 2 倍和
4 倍，减少因子等于瓦片边长 $T$。
\discussionheading
\discussionpoint{理想的 $T$ 倍节省来自消费者分组。}
在某个 $k$ 阶段，$T\times T$ 输出瓦片的同一行需要相同的 $T$ 个 $M$ 元素，同一列需要相同的 $T$ 个 $N$ 元素；若每个输出线程自行加载，每个输入会被同块的 $T$ 个消费者重复请求。协作装载只需把两个 $T\times T$ 输入瓦片各取一次，于是每阶段输入请求从 $2T^3$ 降为 $2T^2$，比值正好是 $T$。对本题，$T=2$ 与 $T=4$ 因而分别给出理想的 2 倍和 4 倍。这个推导的核心不是矩阵形状本身，而是把共享同一数据的消费者放进一个通信域，因此也能迁移到卷积 halo、模板更新和块式注意力。

\discussionpoint{扩大瓦片同时扩大了哪些片上承诺。}
从 $2\times2$ 换成 $4\times4$，朴素实现的输出线程数由 4 增到 16，两块共享输入的元素数也由 8 增到 32；在更实际的尺寸上，这些量按面积增长，而每线程累加器与循环展开还会增加寄存器压力。大瓦片能提高复用，却可能跨过每块线程、每 SM 共享内存或寄存器的装箱边界，并让边界块出现更多无效线程。屏障等待也取决于所有装载者到达，而不是只取决于字节数。因此 $T$ 不是越大越好的连续旋钮，它是一项把复用、片上容量、并行度与尾部浪费绑在一起的离散设计选择。

\discussionpoint{请求数降低不等于 DRAM 事务严格降低同倍数。}
教学图统计的是线程发出的逻辑输入请求，真实硬件会把一个 warp 的相邻地址合并成内存事务，并可能让重复请求在 L1 或 L2 命中；所以未分块版本从 DRAM 取得的数据未必真比瓦片版多 $T$ 倍。反过来，地址未对齐、边界部分瓦片和 cache line 过取也会使物理字节偏离元素计数，输出写回又不随输入复用一起减少。应分别报告 load instruction、全局扇区、L2 字节和 DRAM 字节，才能知道收益来自共享复用、合并访问还是缓存。理论 $T$ 倍仍有价值，但它是理想输入请求模型和算术强度上界，不是端到端加速承诺。

\discussionpoint{高性能 GEMM 把一个瓦片继续分成多级数据运动。}
教学内核让一个线程产生一个输出，真实 GEMM 常先从全局内存搬大瓦片到共享内存，再由 warp 取子瓦片到寄存器或矩阵指令 fragment，并让每线程累加多个结果。这样同一全局元素在共享层复用，同一共享元素又服务多个寄存器累加器；双缓冲或异步复制还能让下一瓦片装载与当前计算重叠。每一级都引入自己的形状、对齐、容量和同步约束，因此不能用单一 $T$ 描述整个内核。可复现的性能报告应列出块/warp/线程瓦片、每线程输出数、共享字节、寄存器数、数据类型和目标架构，才能解释不同设备为何选择不同形状。
\variantheading 对 $8\times8$ 乘法使用 $T=8$，仅 1 块、1 阶段，输入请求为 $2\times64=128$，比 1024 减少 8 倍。
\practiceheading 以 cuBLAS SGEMM 为基线，并用 Nsight Compute 的 Global Load Efficiency 与 Occupancy 选择瓦片。
\sourcecheck{https://docs.nvidia.com/cuda/cuda-c-best-practices-guide/}{NVIDIA CUDA Best Practices Guide 的共享内存矩阵乘法复用分析；访问总数独立枚举，置信度高}
\end{solutionexercise}

\begin{solutionexercise}{3}
\problemheading 如果忘记在图 5.9 的内核中使用一个或两个
\code{__syncthreads()}，可能发生什么类型的错误执行行为？

图 5.9 的每个阶段先由线程协作把当前 $M$、$N$ 输入瓦片写入共享数组
\code{Mds}/\code{Nds}，经第一个块级屏障后用这些共享值计算；第二个块级屏障
位于本阶段计算结束与下一阶段覆盖共享数组之间。

\approachheading 分别检查“加载当前瓦片后”和“覆盖为下一瓦片前”的跨线程数据依赖。

\editorialnote 本题所称英文底本图 5.9，对应中文译稿图 5.10；译稿练习中的旧图号
没有随正文新增插图而更新。这里分析的是含两个块级屏障的分块矩阵乘法内核。

\answerheading 若遗漏第一个屏障，较快线程可能在其他线程尚未完成加载时读取
\code{Mds} 和 \code{Nds}，得到未初始化值或上一阶段残留值，这是读后写竞态。若遗漏
第二个屏障，较快线程可能开始加载下一阶段并覆盖共享数组，而较慢线程仍在使用当前
阶段数据，这是写后读竞态。两个都遗漏时两类竞态同时存在，结果随 warp/块调度变化，
可能偶尔看似正确但没有正确性保证。屏障还必须由块内所有线程在一致控制流中到达。

\conclusionheading 两个屏障分别保护“加载完成再读取”和“读取完成再覆盖”，都不能
随意删除。
\discussionheading
\discussionpoint{第一道屏障关闭当前瓦片的生产窗口。}
设线程 $(0,0)$ 已把一个 $M$ 元素写入 \code{Mds}，便开始沿共享行做乘加，而负责同一行另一元素的线程尚在等待全局内存；没有装载后的屏障，前者就可能读到旧阶段数据或未初始化值。每个线程只生产瓦片的一小部分，却会消费其他线程生产的整行或整列，所以依赖跨越线程边界。第一道 \code{__syncthreads()} 要求整个块到达，并使此前相关共享写入对块内线程可见，恰好闭合这组生产--消费边。错误取决于调度与内存延迟，小尺寸偶然正确只说明竞态未被触发，不能把结果稳定性当成同步证明。

\discussionpoint{第二道屏障保护的是旧瓦片不被过早复用。}
当前阶段计算结束的线程可能立刻进入下一轮，把新 $M$、$N$ 元素写到相同 \code{Mds}/\code{Nds} 槽位；若另一个线程仍在读取旧值完成最后几次乘加，新写会破坏尚未结束的消费。这是一条从“旧读”到“下一轮写”的反依赖，位置跨越循环迭代，所以比第一道屏障更容易被遗漏。第二道屏障把共享缓冲区的生命周期划成装载、完整消费、允许覆盖三个阶段。它与环形缓冲区或双缓冲中的 ownership 转移是同一原则：重用存储前必须证明所有旧消费者已经释放，而不是只确认新生产者准备好了。

\discussionpoint{边界谓词不能改变屏障参与集合。}
当矩阵尺寸不是瓦片整数倍时，尾块中的部分线程没有有效输入或输出；若这些线程在第一道屏障前直接返回，块内其余线程执行集体屏障时便缺少参与者，可能挂起或出现未定义行为。可靠写法让每个线程都计算地址，以谓词把有效值或零写入共享槽，随后全块无条件同步，再仅在最终写回阶段屏蔽越界输出。这样既保证每个共享位置已定义，也保持所有循环阶段的参与集合一致。审查分块内核时应把“数据是否有效”和“线程是否参加集体操作”分开，前者可用谓词处理，后者必须满足原语的收敛要求。

\discussionpoint{流水化改变同步形式，却不会消除依赖。}
异步全局到共享复制、双缓冲或 producer/consumer warp 能让下一瓦片搬运与当前乘加重叠，看起来不再在每一步执行相同的两个全块屏障；实际上它们用阶段计数、pipeline 或更细粒度 barrier 表达相同的“数据已到达”和“缓冲可复用”条件。若等待了错误阶段，或生产者覆盖一个仍被消费者使用的槽，竞态只会变得更难复现。验证可先为每个共享槽画出生产、提交、等待、消费和再次写入的状态机，再用 racecheck、synccheck、非整除尺寸与延迟扰动测试。正确的流水化是把等待与独立工作重叠，而不是把 happens-before 从程序中省略。
\variantheading 若只有一个阶段，第二个屏障可删除，因为不存在下一瓦片覆盖；第一个屏障仍不可缺。
\practiceheading 用 Compute Sanitizer racecheck 与 synccheck 检测共享内存竞态和不一致屏障。
\sourcecheck{https://docs.nvidia.com/cuda/cuda-c-programming-guide/}{NVIDIA CUDA C++ Programming Guide 的块级同步与内存可见性说明；高置信度}
\end{solutionexercise}

\begin{solutionexercise}{4}
\problemheading 假设寄存器和共享内存容量都不是问题，使用共享内存而不是寄存器保存
从全局内存取出的值，一个重要原因是什么？请解释。

\approachheading 比较两种存储空间的线程可见范围。

\answerheading 寄存器是线程私有的：线程 A 载入寄存器的值不能由线程 B 直接读取。
若多个线程需要同一个输入，全部使用寄存器会迫使它们各自重复执行全局加载。共享内存
由整个线程块可见，线程可协作把值加载一次，再让块内多个消费者复用。代价是需要正确
同步并避免 bank 冲突。

\conclusionheading 共享内存的关键价值是块内通信和跨线程复用，而不只是容量。
\discussionheading
\discussionpoint{首要差别是“谁能命名这个值”。}
线程从全局内存取得一个值后，若只有本线程反复使用，寄存器最自然：它属于线程私有状态，适合累加器、地址和局部中间量。可是转置中线程 A 装载的元素要由线程 B 读取，寄存器变量没有让块内任意线程按共享索引访问的地址语义；共享内存则提供块级可见的命名空间，并能配合屏障建立交换顺序。即使假定两种容量都充足，这个通信需求仍足以选择共享内存。因而“用共享而不用寄存器”的重要原因不是一张延迟表，而是算法的生产者和消费者是否跨越线程边界。若消费者只在同一 warp，shuffle 才提供另一种显式且受作用域限制的交换路径。

\discussionpoint{warp shuffle 是一个作用域受限的第三条路。}
若所有生产者和消费者都在同一 warp，shuffle 指令可以让 lane 直接读取另一 lane 的寄存器值，常用于小归约、扫描和矩阵微瓦片，省去共享数组与全块屏障。它要求正确的活动掩码，来源 lane 必须是参与者，数据宽度和交换模式也受指令形式约束；分支后盲用完整掩码会把正确性问题重新带回来。一旦通信跨越 warp，普通 shuffle 就无法触达对方寄存器，仍需共享内存、块级同步或更高层机制。它不是证明“寄存器可共享”，而是硬件为一个明确通信域提供的特殊寄存器交换原语。

\discussionpoint{两种片上资源会以不同方式限制并发。}
增加每线程寄存器可能使 SM 少驻留一个 warp 或线程块，压力更高时编译器还可能把值放入 local address space；增加每块共享内存则直接减少可同时驻留的块数。共享访问还可能因 bank 映射发生冲突，并需要同步来保护跨线程依赖，而寄存器交换一般不承担这些成本。高性能内核因此常混合使用：共享内存保存能被多个线程交换的块瓦片，寄存器保存线程微瓦片和部分和。选择的关键是用哪种资源表达复用关系后，总数据移动、同步和驻留并行度的组合最好，而不是单独最小化某一个资源数字。

\discussionpoint{应用结构决定共享内存扮演的角色。}
GEMM 用共享内存重排并复用输入瓦片，FFT 用它交换蝶形阶段的数据，直方图可用块内私有桶减少全局原子竞争，模板计算则缓存重叠邻域；这些场景都含明确的块内共享或布局变换。编译器通常自动决定标量寄存器分配，程序员则显式设计共享数组、bank 布局和屏障，所以两者的调优证据也不同。实验应把 \code{ptxas} 的寄存器、stack/local 和共享报告，与占用率、bank conflict、同步停顿及真实时间一起检查。只有这样才能判断瓶颈来自通信结构、资源装箱还是编译器放置，而不是凭“片上都很快”得出笼统结论。
\variantheading 若值只由加载它的线程使用，则寄存器更合适，既不需要共享写读也不需要屏障。
\practiceheading 结合 ptxas 寄存器报告与 Nsight Compute 的 Shared Bank Conflicts、Occupancy 决定存储位置。
\sourcecheck{https://docs.nvidia.com/cuda/cuda-c-best-practices-guide/}{NVIDIA CUDA Best Practices Guide 关于共享内存支持块内协作及消除重复传输的说明；置信度高}
\end{solutionexercise}

\begin{solutionexercise}{5}
\problemheading 对分块矩阵乘法内核，如果使用 $32\times32$ 瓦片，输入矩阵 $M$ 和 $N$
的内存带宽使用量减少多少？

\approachheading 一个输入瓦片元素由同一输出瓦片中的 32 个线程复用。

\answerheading 与不分块实现相比，每个 $M$ 或 $N$ 元素的全局请求次数降为原来的
$1/32$，即带宽减少 32 倍。若用百分比表示，减少量为
\[
 1-\frac1{32}=\frac{31}{32}=96.875\%.
\]
这是假设维度整除、忽略边界和缓存效应的理论请求流量。

\conclusionheading $M$ 和 $N$ 的输入全局内存带宽各理论减少 32 倍（减少
$96.875\%$）。
\discussionheading
\discussionpoint{32 倍是完整输出瓦片中的消费者数量。}
对一个 $32\times32$ 输出瓦片，装入共享内存的某个 $M_{ik}$ 会被同一输出行的 32 个线程使用，某个 $N_{kj}$ 会被同一输出列的 32 个线程使用；不分块时，这 32 个消费者各自发出请求。协作装载把这组请求合成一次块级装载，所以理想输入请求降为原来的 $1/32$，即减少 $31/32=96.875\%$。这个比例只针对输入流量，输出仍须写回，地址计算、乘加与屏障也没有按 32 倍消失。因而它提高的是理想算术强度，而不是宣称整个内核时间必然缩短到三十二分之一。

\discussionpoint{尾瓦片与缓存会把理想比例向两边修正。}
若矩阵维度不是 32 的倍数，最后一个输出瓦片只有部分有效消费者，但实现仍可能启动 1024 个线程并为越界位置装零；有效复用低于完整瓦片，谓词和 cache line 过取也要付出成本。另一方面，不分块版本的重复读取可能在 L1 或 L2 命中，物理 DRAM 字节本来就少于源码请求数，因此分块后的 DRAM 降幅可能小于 96.875\%。对齐和合并访问改善又可能提供额外收益。最稳妥的表述是：32 倍是完整瓦片、无缓存基线下的输入请求上界，再用全局扇区、L2 与 DRAM 字节逐层验证实际发生了什么。

\discussionpoint{逻辑 $32\times32$ 瓦片不必绑定 1024 个线程。}
朴素教材实现让每个输出对应一个线程，于是 $32^2=1024$ 已达到许多设备的每块线程上限；再叠加寄存器与分配粒度，SM 可能只能驻留一个这样的块，延迟隐藏和末波调度都受影响。高性能实现可让 256 个线程各计算 $2\times2$ 个输出，仍覆盖同一逻辑瓦片并保留输入复用，只是每线程需要更多累加器和更复杂的共享读取。块规模、逻辑瓦片与每线程微瓦片因此是三个独立参数。把它们分开后，设计者才能在复用、寄存器压力和可驻留块数之间调节，而不是被“瓦片边长等于线程块边长”限制。

\discussionpoint{32 倍复用只是多级 GEMM 数据运动的第一站。}
共享瓦片把全局元素带入块内后，如果每个线程只读一次共享值，片上带宽仍可能成为瓶颈；寄存器微瓦片让一次共享读取服务多个累加器，warp 级矩阵指令又按规定的 fragment 组织更细粒度复用。生产 GEMM 还会用双缓冲或异步复制，使下一块数据传输与当前乘加重叠，并针对 FP32、FP16 或整数选择不同瓦片形状。于是完整模型要分别计算 DRAM--共享、共享--寄存器和寄存器--执行单元的字节与 FLOP。教材中的 32 倍推导提供了第一层直觉，多级模型则解释为何库内核不会只靠扩大一个方形块达到峰值。
\variantheading 使用 $16\times16$ 瓦片时理论减少 16 倍，即减少 $93.75\%$。
\practiceheading 用 Nsight Compute 的 Requested/Actual Global Throughput 与 Occupancy 验证理论带宽收益。
\sourcecheck{https://docs.nvidia.com/cuda/cuda-c-best-practices-guide/}{NVIDIA CUDA Best Practices Guide 的分块矩阵乘法全局读取复用原则；比例独立复核，置信度高}
\end{solutionexercise}

\begin{solutionexercise}{6}
\problemheading 假设启动一个 CUDA 内核，网格有 1000 个线程块，每块有 512 个线程。
如果在内核中声明一个局部变量，在该内核的整个执行生命周期中会创建多少个
该变量的版本？

\approachheading 自动局部变量按线程私有，每个逻辑线程在其生命周期中拥有一份。

\answerheading
\[
1000\ \text{块}\times512\ \text{线程/块}=512{,}000\ \text{个版本}.
\]
即使硬件分多批调度这些块，逻辑变量实例总数也不变。

\conclusionheading 整个内核执行生命周期中创建 512,000 个局部变量版本。
\discussionheading
\discussionpoint{“一份变量”首先是语言层的线程私有状态。}
网格含 $1000\times512=512{,}000$ 个逻辑线程，自动局部变量在每个线程执行该作用域时都有独立的语义值，因此答案按整个内核生命周期计为 512,000 个版本。某个值若从未影响可观察结果，编译器可以把它完全删掉；若仍需要，它可能留在寄存器，也可能进入 local address space，但程序行为都必须像每个线程拥有自己的版本。这里数的是抽象机器中的实例，不是询问芯片同时焊接了多少存储单元。区分语义实例与物理实现，也是理解编译优化不改变并发程序可观察行为的基础。

\discussionpoint{逻辑线程总量并不等于峰值物理并发。}
设备不可能一次让 1000 个 512 线程块全部驻留，SM 只接纳线程、寄存器、共享内存和块数约束允许的若干块；一个块结束后，其寄存器和调度状态便可复用于后续块。于是生命周期内共有 512,000 个局部版本，不代表必须为 512,000 个值同时预留物理位置，峰值需求只由同时驻留线程及各自活跃状态决定。这个区别类似线程池处理百万任务：每个任务有自己的局部上下文，但工作线程和栈槽可按批次复用。容量规划若把逻辑总数直接乘每线程寄存器数，会严重高估某一时刻的 SM 资源，却可能仍低估 spill 的总传输次数。

\discussionpoint{变量名数量与寄存器压力之间没有一一对应。}
编译器按值的活跃区间分配寄存器：若变量 $u$ 在 $v$ 产生前已经最后一次使用，两者可以占用同一物理寄存器，即使源码写了两个名字。相反，循环展开让多轮中间值同时活跃，线程粗化保留多个累加器，或函数内联延长数据依赖，都可能在名字不变时增加压力。压力跨过资源装箱边界会降低驻留 warp，再高时还可能产生 spill。因此优化寄存器不能靠机械删除局部声明，而应检查编译器资源报告和生成代码，缩短真正的活跃区间或重新安排计算，并用性能测量确认没有以额外重算换取更差吞吐。

\discussionpoint{线程私有数组的实际放置需要多层证据。}
像 \code{float x[4]} 这样的局部数组若索引可静态解析，编译器常把元素标量化到寄存器；动态索引、取地址或较高压力则可能让它位于 local address space，该空间语义上仍线程私有，物理上经过缓存并由设备内存支持。不能只看到 profiler 中的 spill 指标就涵盖所有 local 用量，因为显式局部数组、栈帧与寄存器溢出是不同成因。应联合查看 \code{ptxas -v} 的 \code{lmem}/stack 与 spill 报告、PTX 的 \code{.local} 声明以及最终 cubin/SASS 访问。这样才能把“有多少逻辑版本”与“同时分配多少字节、产生多少流量”分别回答清楚。
\variantheading 若改为 750 块、每块 256 线程，则生命周期内共有 192,000 个局部变量版本。
\practiceheading 用 ptxas 和 Nsight Compute 检查局部变量是否发生 local-memory spill，而非只数源码变量。
\sourcecheck{https://docs.nvidia.com/cuda/cuda-c-programming-guide/}{NVIDIA CUDA C++ Programming Guide 的线程私有自动变量语义；乘法独立复核，置信度高}
\end{solutionexercise}

\begin{solutionexercise}{7}
\problemheading 对上一题，如果声明一个共享内存变量，在该内核的整个执行生命周期中会
创建多少个该变量的版本？

\approachheading 静态共享变量按线程块实例化，同一块线程共享一份，不同块互不共享。

\answerheading 1000 个逻辑线程块各有一份，因此整个执行生命周期中创建 1000 个版本。
块分批驻留或复用物理 SM 不会把不同逻辑块的共享变量合并成同一实例。

\conclusionheading 共享内存变量共有 1000 个版本。
\discussionheading
\discussionpoint{共享变量按块实例化，但现代集群带来受限例外。}
普通执行模型中，每个逻辑线程块都有一份静态共享变量，块内 512 个线程共同看到它，而 1000 个块在整个网格生命周期中对应 1000 份逻辑实例。一个块不能用普通共享地址直接访问任意另一块的实例。计算能力 9.0 起，显式线程块集群可通过 Distributed Shared Memory 映射并访问同一集群内其他块的共享段，但每块仍拥有自己的段，访问范围也没有扩大到整个网格。把这项例外写清，既保留教材经典模型，也避免把“共享内存永远严格块内不可访问”误当成所有新架构的绝对命题。

\discussionpoint{生命周期实例数与同时分配份数必须分开。}
1000 个块不会同时驻留：每个 SM 只为当前活动块分配共享内存，块完成后同一物理容量可交给后续块。若每块共享数组从 4 KB 增到 48 KB，可同时驻留的块数可能阶梯式下降，但网格语义上仍会依次创建 1000 个独立版本；减少并发不等于让不同块共用同一逻辑变量。峰值片上容量由每块用量乘活动块数决定，总生命周期中的装载和计算次数则由网格规模决定。这个区分说明共享内存为何既是程序通信对象，又是决定 occupancy 的可复用硬件资源。容量复用由运行时管理，程序不能据此让先后执行的块通过旧共享内容传递状态。

\discussionpoint{块级通信边界决定跨块算法的结构。}
在非 cluster 模型中，同一块可用共享数组与 \code{__syncthreads()} 做归约，两个普通块却不能凭各自共享变量形成生产--消费；通常要写全局内存，再用原子、第二个内核边界或满足条件的协作网格同步建立顺序。DSM 允许显式 cluster 内的块访问彼此共享段，但要求目标块仍处于集群生命周期中，并使用相应映射与同步，超出 cluster 仍需更高层通信。于是共享内存的“快”始终附带一个作用域契约。算法若先确定通信跨度，再选择块、集群或网格机制，会比先选存储类型再勉强拼接同步更可靠。

\discussionpoint{不要用共享指针数值证明块间隔离。}
共享地址处在每块的地址语义中，不同块打印或比较出的指针数值即使相同，也不能推出它们指向同一物理对象；数值不同也不是隔离性的正式证明。隔离性来自 CUDA 语言与体系结构规定的作用域契约：普通线程块不能构造指向另一普通块共享段的可解引用地址，因而也没有一个合法的跨块 shared 读取实验可用来“证明”该契约。让块 $b$ 在自己的共享段写入由 \code{blockIdx.x} 派生的哨兵、同步后再写入全局数组，只能演示当前实现确实得到各自的值，适合作为回归样例，不能从有限观察推出所有执行都隔离。cluster DSM 应另用规定的地址映射和集群同步构造跨块访问正例；静态与动态共享内存的资源后果则用函数属性和 Occupancy API 独立核对。
\variantheading 若网格为 240 块，则共有 240 个逻辑版本，但任一时刻只驻留设备资源允许的子集。
\practiceheading 通过 \code{cudaFuncGetAttributes} 和 Occupancy API 核对每块共享内存及最大活动块数。
\sourcecheck{https://docs.nvidia.com/cuda/cuda-c-programming-guide/\#distributed-shared-memory}{NVIDIA CUDA C++ Programming Guide 的块级共享实例及计算能力 9.0 thread-block cluster 的 Distributed Shared Memory；对抗检查：DSM 扩大可访问范围但不把每块实例合并为网格全局变量}
\end{solutionexercise}

\begin{solutionexercise}{8}
\problemheading 考虑把两个 $N\times N$ 输入矩阵相乘。输入矩阵中的每个元素从全局内存
请求多少次？
\begin{enumerate}[label=\alph*.,leftmargin=2.2em]
  \item 不使用分块时；
  \item 使用大小为 $T\times T$ 的瓦片时。
\end{enumerate}

\approachheading 固定一个 $M$ 元素，数有多少输出列使用它；分块后，一次加载可供
同一输出瓦片的 $T$ 个线程复用。$N$ 矩阵同理。

\answerheading
\begin{enumerate}[label=\alph*.,leftmargin=2.2em]
\item 不分块时，一个 $M_{ik}$ 被第 $i$ 行的 $N$ 个输出使用，一个 $N_{kj}$ 被第
$j$ 列的 $N$ 个输出使用，故每个输入元素请求 $N$ 次。
\item 使用 $T\times T$ 瓦片且 $T\mid N$ 时，一次加载服务一个含 $T$ 个消费者的
输出瓦片；每个元素只在另一输出维的 $N/T$ 个瓦片中各加载一次，故请求 $N/T$ 次。
若不能整除，边界瓦片使上界写为 $\lceil N/T\rceil$ 次。
\end{enumerate}

\complexity{总算术工作仍为 $O(N^3)$}{每块 $O(T^2)$ 共享空间}{$O((N/T)^2)$ 个输出瓦片并行}
\conclusionheading (a) $N$ 次；(b) 整除时 $N/T$ 次，理论减少 $T$ 倍。
\discussionheading
\discussionpoint{请求次数由一个输入元素服务的输出集合决定。}
固定 $M_{ik}$ 会参与输出行 $i$ 的全部 $N$ 个 $C_{ij}$，不分块时每个输出线程独立读取它，所以逻辑请求为 $N$ 次；$N_{kj}$ 对输出列也完全对称。$T\times T$ 分块把相邻 $T$ 个消费者放在同一块，一次协作装载服务这一组，因此沿另一输出维只需经过 $N/T$ 个完整瓦片。请求次数由此从 $N$ 降为 $N/T$，复用因子为 $T$。这个推导把单个元素的扇出与块划分联系起来，也揭示了瓦片增大时矩阵乘法算术强度提高的来源。若两个输入采用不同瓦片形状，应分别按各自在输出方向覆盖的消费者数计算。

\discussionpoint{非整除尺寸要把除法改成瓦片覆盖数。}
若 $N=70,T=32$，输出行或列会被三个瓦片覆盖，而不是 $70/32$ 个；一个有效输入元素因此至多由对应方向的 $\lceil70/32\rceil=3$ 个块各装载一次。实现可能让第三块对越界位置装零或用谓词跳过请求，但有效元素仍需要服务该块中剩余消费者。对规则方阵和标准分块，统一请求计数可写成 $\lceil N/T\rceil$；若算法裁剪部分瓦片或采用不规则分区，精确总量则应按每块有效范围求和。这个例子说明性能公式中的整除假设必须显式写出，尾部不是把连续商数简单四舍五入。

\discussionpoint{逻辑 load、内存事务与 DRAM 字节是三种计量。}
一个 warp 的 32 条源码 load 可以被合并为若干扇区事务，同一 cache line 的后续请求又可能在 L1 或 L2 命中，所以“某元素被请求 $N$ 次”不等于 DRAM 真传输 $N$ 份。相反，未对齐地址、扇区过取和写策略会让物理字节大于有效元素字节，边界谓词也可能使事务利用率下降。比较模型与 profiler 时应先声明分母在源码、L2 还是 DRAM 层，再选相应指标。只有口径一致，$N/T$ 才能解释复用，而不会被误用成精确的总线访问次数或运行时间比例。缓存预热状态也要固定，否则两次基准可能比较了不同的首次访问成本。

\discussionpoint{共享瓦片之上还存在寄存器与批次复用。}
若每线程只算一个 $C_{ij}$，从共享内存读取的 $M$、$N$ 值仍会被多个线程分别取走；寄存器 blocking 让一个线程持有多个累加器，一次共享读取便可更新若干输出。warp 级矩阵指令把这种微瓦片组织成固定 fragment，而批量 GEMM 在某些工作负载中还能跨多个输入矩阵重用权重或元数据。于是现代内核需要分别核算 DRAM 到共享、共享到寄存器以及寄存器到执行单元的数据运动，每层都有自己的 $T$ 或形状。题目的 $N/T$ 是全局请求层的清晰结论，但不能替代对后续片上带宽和寄存器容量的分析。
\variantheading 若 $N=1000,T=32$，每个输入元素最多由 32 个输出瓦片请求，而不是 $31.25$ 次。
\practiceheading 用 Nsight Compute 的 DRAM sectors 和 L2 hit rate 区分算法请求复用与缓存复用。
\sourcecheck{https://docs.nvidia.com/cuda/cuda-c-best-practices-guide/}{NVIDIA CUDA Best Practices Guide 的分块矩阵乘法数据复用；逐元素消费者计数复核，置信度高}
\end{solutionexercise}

\begin{solutionexercise}{9}
\problemheading 一个内核每线程执行 36 次浮点运算，并进行 7 次 32 位全局内存访问。
对下面每组设备属性，判断该内核是计算受限还是内存受限：
\begin{enumerate}[label=\alph*.,leftmargin=2.2em]
  \item 峰值 FLOPS = 200 GFLOPS，峰值内存带宽 = 100 GB/s；
  \item 峰值 FLOPS = 300 GFLOPS，峰值内存带宽 = 250 GB/s。
\end{enumerate}

\approachheading 算术强度为 FLOP/传输字节；与设备 ridge point
$P_{\max}/B_{\max}$ 比较。

\answerheading 每线程传输 $7\times4=28$ B，算术强度为
\[
I=\frac{36}{28}=\frac97\approx1.286\ \mathrm{FLOP/B}.
\]
(a) 设备阈值为 $200/100=2\ \mathrm{FLOP/B}$。因为 $I<2$，带宽屋顶只有
$100\times9/7\approx128.6$ GFLOP/s，小于 200 GFLOP/s，故内存受限。

(b) 设备阈值为 $300/250=1.2\ \mathrm{FLOP/B}$。因为 $I>1.2$，带宽屋顶为
$250\times9/7\approx321.4$ GFLOP/s，高于 300 GFLOP/s，故计算受限。

\conclusionheading (a) 内存受限；(b) 计算受限。
\discussionheading
\discussionpoint{ridge point 把算法强度与设备比例放在同一尺度。}
本题每线程做 36 FLOP、传输 $7\times4=28$ B，所以算术强度为 $I=9/7\approx1.286$ FLOP/B；设备的 ridge point 是峰值计算率除以峰值带宽。对 (a)，交点为 2 FLOP/B，$I$ 位于其左侧，带宽屋顶仅约 128.6 GFLOP/s；对 (b)，交点为 1.2，算法落到右侧，计算屋顶 300 GFLOP/s 先被触及。这个比较不是凭两个峰值大小猜瓶颈，而是判断“每传一个字节可供多少计算”是否足以喂满计算单元。它也解释了同一内核为何换设备后能从内存受限变成计算受限。

\discussionpoint{算术强度的分母必须绑定一个内存层次。}
题解把七次 32 位全局访问直接计为 28 B，这是源码层的教学口径；若其中若干读取在 L2 命中，DRAM 实测字节会更少，而 cache line 过取、写回策略和 ECC 等因素也会让总线字节不同。若改画 L2 Roofline 或共享内存 Roofline，分母和相应带宽屋顶都必须一起更换，不能拿 DRAM 峰值配 L2 字节。分子同样要统一：一次 FMA 通常按乘和加计 2 FLOP，整数地址指令与控制指令不能塞进 FLOP 峰值比较。只有层次和计数规则一致，图上的位置才具有可比意义。

\discussionpoint{落在某条屋顶下方不证明已经碰到那条屋顶。}
Roofline 给出 $\min(P_{\max},I B_{\max})$ 的理想上界，真实内核可能因长依赖链、低并行度、分支分歧、未合并访问或指令发射限制远低于该线。比如 (a) 虽按模型归为内存侧，若只有很小网格，GPU 可能主要花在启动和延迟上，实际带宽根本未饱和；(b) 位于计算侧也不代表已达到峰值浮点管线。分析器若显示低 DRAM 吞吐且大量依赖停顿，更准确的诊断是延迟或执行效率不足。Roofline 先界定理论限制方向，scheduler、memory 和 instruction 指标再解释为何没有接近该界。

\discussionpoint{优化应区分移动算法点与提高屋顶利用率。}
对带宽侧内核，融合中间结果、增加数据复用、压缩表示或在误差允许时降低存储精度，能减少每 FLOP 对应字节，从而把点向右移动；单纯提高占用率若没有增加有效带宽，算术强度并未改变。对计算侧内核，可改善指令级并行、采用适合数据类型的矩阵指令或减少非浮点开销，让实际点向计算屋顶靠近。层次 Roofline 同时绘制 DRAM 和 L2 口径，可以判断复用发生在哪一级。工程报告若保留优化前后的 FLOP/s、各层字节和强度，就能分辨是算法数据运动变了，还是同一强度下硬件执行效率提高。
\variantheading 若设备为 400 GFLOPS、200 GB/s，ridge point 为 2 FLOP/B，故本内核仍是内存受限。
\practiceheading 用 Nsight Compute Roofline Chart 的实测 FLOP/s 与 DRAM 带宽定位瓶颈。
\sourcecheck{https://docs.nvidia.com/cuda/cuda-c-best-practices-guide/}{NVIDIA CUDA Best Practices Guide 的有效带宽与算术工作量分析方法；单位和两条屋顶独立复核，置信度高}
\end{solutionexercise}

\begin{solutionexercise}{10}
\problemheading 为操作瓦片，一名 CUDA 新手编写了下面的设备内核，用于转置矩阵中的
每个瓦片。瓦片大小为 \code{BLOCK_WIDTH}，矩阵各维度都是
\code{BLOCK_WIDTH} 的倍数。内核调用和代码如下；\code{BLOCK_WIDTH} 在编译
时已知，可以取 1 到 20 之间的任意值。
\begin{lstlisting}[language=C++,firstnumber=1]
dim3 blockDim(BLOCK_WIDTH, BLOCK_WIDTH);
dim3 gridDim(A_width/blockDim.x, A_height/blockDim.y);
BlockTranspose<<<gridDim, blockDim>>>(A, A_width, A_height);
__global__ void
BlockTranspose(float* A_elements, int A_width, int A_height)
{
    __shared__ float blockA[BLOCK_WIDTH][BLOCK_WIDTH];
    int baseIdx = blockIdx.x * BLOCK_SIZE + threadIdx.x;
    baseIdx += (blockIdx.y * BLOCK_SIZE + threadIdx.y) * A_width;
    blockA[threadIdx.y][threadIdx.x] = A_elements[baseIdx];
    A_elements[baseIdx] = blockA[threadIdx.x][threadIdx.y];
}
\end{lstlisting}
\begin{enumerate}[label=\alph*.,leftmargin=2.2em]
  \item 在 \code{BLOCK_SIZE} 的可能取值范围内，哪些值能让该内核在设备上
  正确执行？
  \item 如果不是所有值都能正确执行，错误执行行为的根本原因是什么？如何
  修改代码，使其对所有 \code{BLOCK_SIZE} 值都能工作？
\end{enumerate}
\approachheading 关键依赖是每个线程先写共享瓦片的 $(y,x)$ 元素，再由另一个线程
读 $(x,y)$ 元素；必须在读前建立块级同步。另需统一底本混用的宏名。

\answerheading 按教材传统 lockstep 模型，$B^2\le32$ 时整个块不超过一个 warp，
所以常见预期答案是 $B=1,2,3,4,5$；$B\ge6$ 跨越多个 warp，某个 warp 可能在另一个
warp 完成加载前读取共享数组。严格按现代 CUDA C++ 内存模型，未同步的跨线程读写
没有可移植保证，因此除了不发生跨线程通信的 $B=1$ 外，不应宣称原程序对任何
$B>1$ 必然正确。Volta 及更新架构的独立线程调度进一步否定了“单 warp 自然同步”
这一假设。

统一使用 \code{BLOCK_WIDTH} 并在加载之后加入块级屏障即可覆盖全部 $1\ldots20$：
\begin{lstlisting}[language=C++]
__global__ void BlockTranspose(float* A, int width, int height) {
    __shared__ float tile[BLOCK_WIDTH][BLOCK_WIDTH];
    int x = blockIdx.x * BLOCK_WIDTH + threadIdx.x;
    int y = blockIdx.y * BLOCK_WIDTH + threadIdx.y;
    int idx = y * width + x;
    tile[threadIdx.y][threadIdx.x] = A[idx];
    __syncthreads();
    A[idx] = tile[threadIdx.x][threadIdx.y];
}
\end{lstlisting}
题设保证矩阵维度整除块宽，所以无需额外边界判断。所有线程无条件到达屏障，写回位置
互异，不存在全局写写竞态。

\editorialnote 底本声明用 \code{BLOCK_WIDTH}、索引用 \code{BLOCK_SIZE}；上面的
修复把二者统一。小题 (a) 的 $1\ldots5$ 是底本教学模型下的答案，现代内存模型的
可移植答案更严格，故同时列出。

\complexity{$O(B^2)$ 工作/块}{$4B^2$ B 共享内存}{$B^2$ 个加载与写回线程}
\conclusionheading 教材预期为 $B=1\ldots5$；标准兼容代码必须加入
\code{__syncthreads()}，之后 $B=1\ldots20$ 均可工作。
\discussionheading
\discussionpoint{转置把一次块内读写变成了跨线程依赖。}
线程 $(y,x)$ 写入 \code{tile[y][x]}，随后读取 \code{tile[x][y]}；除 $x=y$ 的对角线外，读者与写者是两个不同线程。共享数组只让双方能够命名同一位置，并不会自动保证生产者已经完成写入或消费者能看见该值，因此加载与读取之间需要明确的 happens-before。若缺少同步，某些线程会读到未初始化或上一块物理存储留下的内容，表现随调度变化而变化。把每个共享位置的写者和读者列出来，能立即看出 $B=1$ 是唯一没有跨线程边的特例，也解释了问题根源不是矩阵整除性。

\discussionpoint{“一个 warp 内”是旧教学模型，不是现代可移植契约。}
当 $B\le5$ 时块至多 25 个线程，传统答案假定它们在一个 warp 中隐式锁步，所以列出 $B=1\ldots5$；这能解释底本预期，却不能作为现代 CUDA C++ 的保证。自 Volta 起独立线程调度允许同一 warp 的 lane 更细粒度地暂停和推进，即使早期目标偶然按锁步执行，编译器和内存模型也没有为未声明的跨线程通信承诺顺序。因而严格可移植答案只有不交换数据的 $B=1$ 原样正确，所有 $B>1$ 都应写出同步。区分“历史机器上常见行为”与“语言保证”，是对抗审查并发代码时尤其重要的一步。

\discussionpoint{先修正确性，再单独优化共享布局。}
统一混用的 \code{BLOCK_WIDTH}/\code{BLOCK_SIZE}，并在全块写入共享瓦片后执行 \code{__syncthreads()}，就为 $B=1\ldots20$ 建立块级到达与内存可见性；题设最大 $20^2=400$ 个线程也在给定设备能力允许时合法。若实现明确限制为单 warp，可以改用带正确参与掩码的 \code{__syncwarp}，但跨 warp 不能这样替代。转置的按列共享访问还可能产生 bank conflict，常见性能布局是在第二维增加 padding；这不会替代同步，只改变 bank 映射。把屏障修复与 padding 优化分开验证，可避免用更快的错误程序掩盖竞态。

\discussionpoint{测试数据应主动破坏“偶然正确”的对称性。}
若输入是全一矩阵或对称矩阵，读错 \code{tile[x][y]} 与读对 \code{tile[y][x]} 可能给出同值，未同步内核便会伪装成正确；单次运行也可能恰好采用有利调度。应让每个位置含唯一的非对称编码，例如 $1000y+x$，覆盖 $B=2,5,6,20$ 等单 warp 边界和跨 warp 情形，并重复运行多个目标 SM。Compute Sanitizer racecheck 可提供动态竞态证据，但未报告并不等价于所有调度已证明安全。源码层仍要为每条跨线程读写找到屏障或 warp 同步，这与数值参考比较共同构成完整验证。
\variantheading 若 $B=6$，块有 36 个线程并跨两个 warp，原代码可能跨 warp 读到未写值；加屏障后正确。
\practiceheading 针对 \code{sm_70} 及更新目标运行 Compute Sanitizer racecheck，并用非对称瓦片数据验证转置。
\sourcecheck{https://docs.nvidia.com/cuda/cuda-c-programming-guide/}{NVIDIA CUDA C++ Programming Guide 的独立线程调度、块级屏障及内存可见性；对抗审查发现单 warp 假设不可移植}
\end{solutionexercise}

\begin{solutionexercise}{11}
\problemheading 考虑下面的 CUDA 内核及调用它的主机函数：
\begin{lstlisting}[language=C++,firstnumber=1]
__global__ void foo_kernel(float* a, float* b) {
    unsigned int i = blockIdx.x*blockDim.x + threadIdx.x;
    float x[4];
    __shared__ float y_s;
    __shared__ float b_s[128];
    for (unsigned int j = 0; j < 4; ++j) {
        x[j] = a[j*blockDim.x*gridDim.x + i];
    }
    if (threadIdx.x == 0) {
        y_s = 7.4f;
    }
    b_s[threadIdx.x] = b[i];
    __syncthreads();
    b[i] = 2.5f*x[0] + 3.7f*x[1] + 6.3f*x[2] + 8.5f*x[3]
         + y_s*b_s[threadIdx.x] + b_s[(threadIdx.x + 3)%128];
}
void foo(int* a_d, int* b_d) {
    unsigned int N = 1024;
    foo_kernel<<<(N + 128 - 1)/128, 128>>>(a_d, b_d);
}
\end{lstlisting}
\begin{enumerate}[label=\alph*.,leftmargin=2.2em]
  \item 变量 \code{i} 有多少个版本？
  \item 数组 \code{x[]} 有多少个版本？
  \item 变量 \code{y_s} 有多少个版本？
  \item 数组 \code{b_s[]} 有多少个版本？
  \item 每个线程块使用多少字节共享内存？
  \item 该内核的浮点运算量与全局内存访问量之比是多少（OP/B）？
\end{enumerate}
\approachheading 自动变量按线程计数，共享变量按块计数；OP/B 只统计浮点算术和全局
内存传输，不把共享访问计入分母。

\answerheading 网格有 $1024/128=8$ 个块和 1024 个线程。
\begin{enumerate}[label=\alph*.,leftmargin=2.2em]
\item \code{i}：1024 个版本。
\item \code{x[]}：1024 份数组（合计 $1024\times4=4096$ 个 \code{float} 元素）。
\item \code{y_s}：每块一份，共 8 个版本。
\item \code{b_s[]}：每块一份，共 8 份数组。
\item 每块声明的共享内存有效载荷为
$4+128\times4=516$ B；实现可能按硬件对齐粒度保留更多物理空间。
\item 每线程有 5 次乘法和 5 次加法，共 10 FLOP；全局内存有 4 次
\code{a} 加载、1 次 \code{b} 加载和 1 次 \code{b} 存储，共 6 个 32 位字，
即 24 B。因此
\[
\frac{10}{24}=\frac5{12}\approx0.417\ \mathrm{OP/B}.
\]
\end{enumerate}
若编译器生成 FMA，指令条数会变化，但 FLOP 计数仍把一次乘和一次加计为 2 FLOP。

\editorialnote 底本主机函数把参数写成 \code{int*}，与内核的 \code{float*}
不一致；实际可编译程序必须统一为 \code{float*}。

\conclusionheading (a)--(f) 为 1024、1024 份、8、8 份、516 B、约
$0.417$ OP/B。
\discussionheading
\discussionpoint{变量版本数来自执行作用域，而不是声明出现次数。}
网格有 1024 个线程和 8 个块，所以线程私有的 \code{i} 与 \code{x[4]} 各有 1024 个逻辑版本，块共享的 \code{y_s} 与 \code{b_s[128]} 各有 8 份；数组版本数与其中元素总数又是两个不同问题。指针 \code{a}/\code{b} 虽作为每线程参数存在，其所指数据处于全局地址空间，不能因参数寄存器私有就误判数组也私有。按“线程、块、网格”逐级标注每个对象，既能复核本题计数，也能直接提示哪些访问需要跨线程同步。这个方法比看变量名中的后缀更可靠，因为名称没有内存空间语义。

\discussionpoint{局部数组可能进入 local address space，但原因不止 spill。}
循环若被展开且 \code{x[j]} 的索引静态可知，四个元素可能标量化为寄存器；动态索引、取地址或栈布局要求则可能让数组直接拥有 local 存储，即使并非寄存器压力导致的溢出。local address space 语义上每线程私有，物理访问经过缓存并由设备内存层次支持，会增加题目朴素模型未计的指令和流量。因而只查看 spill load/store 不能断言 \code{x} 的放置，还应结合 \code{ptxas} 的 \code{lmem}/stack frame、PTX 中的 \code{.local} 声明以及最终 cubin/SASS。多层证据能区分显式局部对象、栈和真正的寄存器溢出。

\discussionpoint{$0.417$ OP/B 是一个口径明确的静态下层模型。}
题解按每线程 5 次乘、5 次加计 10 FLOP，并按四次 \code{a} 读、一次 \code{b} 读和一次 \code{b} 写计 24 B，故得 $5/12\approx0.417$ OP/B。编译器若生成 FMA，机器指令数变少，但常用 FLOP 口径仍把一次乘加算作 2 FLOP；共享内存访问不进入“全局字节”分母。真实 DRAM 字节还会受缓存、事务过取和潜在 local 访问影响，所以静态强度不是实际吞吐，也可能不等于 profiler 的 DRAM Roofline 强度。比较工具前必须同时统一 FLOP、读写与内存层次定义。

\discussionpoint{从静态计数到优化结论需要闭合测量链。}
四个 \code{x} 输入沿全网格连续分段读取，首先应检查各轮 warp 地址是否合并；\code{y_s} 是块内广播常量，而 \code{b_s} 的邻居访问需要屏障但可能具有规则 bank 映射。算子融合可省掉后续中间数组，数据重用可能降低 DRAM 流量，然而更多局部状态也可能提高寄存器或 local 压力。工程上应把源码的 10 FLOP/24 B 与 SASS 指令、\code{ptxas} 资源、Nsight Compute 的 FLOP、L2/DRAM 字节和实际时间并列。模型与测量不一致时，再追查缓存、spill 或编译优化，而不是直接用理论 OP/B 宣称瓶颈。
\variantheading 若删去最后一个原始 \code{b_s} 加项，则每线程变为 5 乘、4 加，即 $9/24=0.375$ OP/B。
\practiceheading 同时检查 ptxas 的寄存器、spill、\code{lmem} 与 stack frame，搜索 PTX \code{.local} 并核对 cubin/SASS，再以 Nsight Compute 实测 DRAM Bytes 与 FLOP 计数。
\sourcecheck{https://docs.nvidia.com/cuda/cuda-c-programming-guide/}{NVIDIA CUDA C++ Programming Guide 的变量内存空间与线程/块作用域；字节数和 FLOP 数独立复核，置信度高}
\end{solutionexercise}

\begin{solutionexercise}{12}
\problemheading 某 GPU 的硬件限制为：每个 SM 2048 个线程、32 个线程块、65,536 个
寄存器和 96 KB 共享内存。对下列内核特征，判断内核能否达到满占用率；如果
不能，指出限制因素：
\begin{enumerate}[label=\alph*.,leftmargin=2.2em]
  \item 每块 64 个线程，每线程 27 个寄存器，每个 SM 4 KB 共享内存；
  \item 每块 256 个线程，每线程 31 个寄存器，每个 SM 8 KB 共享内存。
\end{enumerate}

\approachheading 按底本原文把共享内存数解释为“每 SM 使用量”；分别核对达到 2048
线程所需的块数、寄存器数和共享内存。

\answerheading
\begin{enumerate}[label=\alph*.,leftmargin=2.2em]
\item 64 线程/块需要 32 个块，恰好等于块上限；寄存器需求
$2048\times27=55{,}296<65{,}536$，共享内存 $4$ KB $<96$ KB。因此可以满占用。
\item 256 线程/块需要 8 个块；寄存器需求
$2048\times31=63{,}488<65{,}536$，共享内存 $8$ KB $<96$ KB。因此也可以满占用。
\end{enumerate}

\editorialnote 底本明确写的是 shared memory/SM，而非 shared memory/block。若原意
其实是每块用量，则 (a) 会因 $32\times4=128$ KB 超过 96 KB 而受共享内存限制，
最多 24 块、1536 线程（$75\%$）；(b) 的 $8\times8=64$ KB 仍可满占用。题解不静默
替换底本单位。

\conclusionheading 按题面字面含义，两种配置都可达到满占用率。
\discussionheading
\discussionpoint{“每 SM 使用量”与“每块使用量”会导出不同答案。}
硬件规格通常说 SM 总共有 96 KB，共享资源报告则通常说某内核每块用多少；题面却把 4 KB、8 KB 写成每 SM 使用量，因此按字面只需分别与 96 KB 比较，两项都不限制满占用。若原意其实是每块，(a) 为达到 2048 线程需要 32 块，总需求 $32\times4=128$ KB，只能驻留 24 块、1536 线程，即 $75\%$；(b) 的 8 块共 64 KB 仍可满占用。单位歧义不是文字小节，它改变了限制资源和数值结论。题解同时给出两条解释分支，比静默选择常见含义更可审计。

\discussionpoint{活动块数由所有资源商的最小整数决定。}
对每块线程数 $T$、寄存器用量 $R_b$ 和共享字节 $S_b$，简化模型的驻留块数取线程商、块上限、寄存器商与共享内存商向下取整后的最小值；满占用还要求该块数乘 $T$ 达到 2048。任何一项跨过整数边界都能成为新的限制因素，因此只验证寄存器总量小于 65,536 并不充分。真实设备还会按架构粒度分配寄存器和共享区，并把静态与启动时动态共享内存合并计算。这个“多维整数装箱”模型能统一前几章的占用率题，也说明为何临界手算必须用目标设备工具复核。若两项同时达到相同最小值，应把它们都报告为共同限制，而非任意挑一个名称。

\discussionpoint{满占用只是延迟隐藏容量，不是性能终点。}
按题面字面，两种配置都能达到理论 $100\%$，但 (a) 需要 32 个小块，(b) 只需 8 个大块，块级同步、调度自由度和数据复用方式并不相同。若为了维持满占用而强压寄存器，spill 可能增加 local 流量；若缩小共享瓦片，输入复用下降也可能让 DRAM 成为瓶颈。相反，一个计算密集内核在较低占用时仍可能有足够就绪 warp 填满管线。配置选择应同时观察实际吞吐、eligible warps、资源使用和数据移动，把占用率当作必要的容量诊断而非单一优化目标。

\discussionpoint{资源结论要附带可复现的构建与设备证据。}
同一源码因目标 SM、编译器版本、内联和模板参数不同，寄存器与静态共享内存都可能变化；动态共享内存又来自具体启动参数，所以文档中的一行数字不足以复现 occupancy。应记录设备属性、\code{cudaFuncAttributes}、编译命令和动态共享字节，再用 Occupancy API 或分析器计算活动块，并以实测驻留情况核对。遇到本题这样的单位冲突，还应在状态记录中保留原文、解释假设和两套数值，而不是只留下最终百分数。这样后来者能判断差异来自题面、编译产物还是硬件，而不会把假设变化误认为计算错误。
\variantheading 按每块 4 KB 解释 (a)，96 KB 只能容纳 24 块，即 1536 线程、$75\%$ 占用率。
\practiceheading 读取 \code{cudaFuncAttributes.sharedSizeBytes} 并用 Occupancy API 计算目标 GPU 的真实驻留数。
\sourcecheck{https://docs.nvidia.com/nsight-compute/NsightCompute/index.html}{NVIDIA Nsight Compute Occupancy Calculator 的资源约束模型；并对底本 shared memory/SM 单位作对抗审查}
\end{solutionexercise}

\chapter{性能考量}

\begin{solutionexercise}{1}
\problemheading 考虑下面的 CUDA 内核：
\begin{lstlisting}[language=C++,firstnumber=1]
__global__ void foo_kernel(float* a, float* b, float* c, float* d, float* e) {
    unsigned int i = blockIdx.x*blockDim.x + threadIdx.x;
    __shared__ float a_s[256];
    __shared__ float bc_s[4*256];
    a_s[threadIdx.x] = a[i];
    for (unsigned int j = 0; j < 4; ++j) {
        bc_s[j*256 + threadIdx.x] = b[j*blockDim.x*gridDim.x + i] + c[i*4 + j];
    }
    __syncthreads();
    d[i + 8] = a_s[threadIdx.x];
    e[i*8] = bc_s[threadIdx.x*4];
}
\end{lstlisting}
对下面每个内存访问，说明它是合并的、未合并的，还是不适用合并概念：
\begin{enumerate}[label=\alph*.,leftmargin=2.2em]
  \item 第 05 行对数组 \code{a} 的访问；
  \item 第 05 行对数组 \code{a_s} 的访问；
  \item 第 07 行对数组 \code{b} 的访问；
  \item 第 07 行对数组 \code{c} 的访问；
  \item 第 07 行对数组 \code{bc_s} 的访问；
  \item 第 10 行对数组 \code{a_s} 的访问；
  \item 第 10 行对数组 \code{d} 的访问；
  \item 第 11 行对数组 \code{bc_s} 的访问；
  \item 第 11 行对数组 \code{e} 的访问。
\end{enumerate}

\approachheading 合并只适用于同一 warp 的全局内存指令；共享内存应改用 bank 冲突
分析。对循环体，固定同一次迭代的 $j$ 再比较相邻线程地址。

以下分类假设 \code{blockDim.x <= 256}（教材常用配置为 256），并且所有全局数组都覆盖
题面产生的最大索引。若块宽超过 256，固定大小的共享数组会先发生越界，此时不能只讨论
合并访问或 bank 冲突。

\answerheading
\begin{enumerate}[label=\alph*.,leftmargin=2.2em]
\item \code{a[i]}：合并，相邻线程读取连续 \code{float}。
\item \code{a_s[threadIdx.x]}：不适用全局合并；共享 bank 连续且无冲突。
\item \code{b[j*blockDim.x*gridDim.x+i]}：合并；固定 $j$ 时仅 $i$ 随线程连续变化。
\item \code{c[i*4+j]}：未合并；相邻线程相差 4 个 \code{float}（16 B）。
\item \code{bc_s[j*256+threadIdx.x]}：不适用；共享 bank 连续且无冲突。
\item \code{a_s[threadIdx.x]}：不适用；共享 bank 连续且无冲突。
\item \code{d[i+8]}：地址连续，属于合并访问；固定偏移可能使首事务未对齐并增加
一个内存事务，但不把单位步长模式变成跨距访问。
\item \code{bc_s[threadIdx.x*4]}：不适用合并；共享内存 bank 步长为 4，形成
4 路 bank 冲突。
\item \code{e[i*8]}：未合并；相邻线程相差 8 个 \code{float}（32 B）。
\end{enumerate}

\conclusionheading 全局访问 (a)、(c)、(g) 合并，(d)、(i) 未合并；其余为共享
内存，合并概念不适用。
\discussionheading
\discussionpoint{合并访问描述的是一条 warp 指令的地址集合。}
单个线程读取 \code{a[i]} 时，不能脱离其他 lane 称它“合并”或“不合并”；硬件处理的是同一 warp 在同一条动态指令上提出的整组地址。若 lane $t$ 的字节地址是 $A+4t$，32 个 lane 覆盖连续 128 B，可由少量内存扇区服务；若地址是 $A+32t$，请求字节散落在更大区域，需要更多事务。循环中的不同 $j$ 属于不同动态指令实例，应固定 $j$ 后再横向比较 lane，而不能把一个线程四次加载混成一组。这个观察法把“源码看起来连续”的直觉转化为可计算的步长、对齐和事务覆盖。
\discussionpoint{共享内存使用 bank，而不是全局合并器。}
\code{a_s} 与 \code{bc_s} 位于片上共享内存，访问不会经过全局内存的合并路径；它们的主要问题是同一 warp 的不同地址是否映射到同一 bank。连续 32 位字通常轮流落到 32 个 bank，所以 \code{a_s[t]} 可并行服务，而 \code{bc_s[4t]} 只使用部分 bank并形成多路冲突。若多个 lane 读取完全相同的共享地址，硬件还能广播该字，不应把它误判为普通冲突；多个 lane 写同址则另有数据竞争语义。将两套模型分开，才能为全局跨距选择布局重排，为共享冲突选择 padding 或 shuffle，而不是用错优化工具。
\discussionpoint{“合并”是随架构与对齐连续变化的效率。}
本题按 32 位元素和有效数组范围分类，\code{d[i+8]} 虽然首地址可能相对事务边界错位，lane 间仍保持单位步长，因此与 \code{e[8i]} 的大跨距本质不同。现代架构常把请求分解为固定字节扇区，错位可能让连续 128 B 覆盖多一个扇区，但不会突然变成 32 个完全独立事务；元素宽度、子数组偏移和缓存路径又会改变具体数量。因而“合并/未合并”适合作为教学分类，工程分析更应计算请求字节与实际传输字节的比值。先声明目标架构和地址假设，才能避免把一次设备上的事务结果写成永久定律。
\discussionpoint{数据布局是在替整个 warp安排邻接关系。}
结构数组中，相邻线程若各取一个结构的同一字段，地址间距会等于结构大小；改成字段数组后，同一字段连续排列，lane 便可消费相邻值，这就是 AoS 到 SoA 转换常见的动机。矩阵转置和批量张量重排也在做相似事情：把算法最常同时访问的逻辑对象放到物理相邻位置。重排本身需要额外内存流量，只有多个后续阶段复用新布局或原跨距十分严重时才值得。诊断应在对齐、偏移和边界块上比较实际传输量，并同时排除共享 bank 冲突，才能知道瓶颈来自哪一条存储路径。接口若长期采用新布局，还应把字段对齐和版本写入数据格式契约，防止生产者与消费者再次失配。
\variantheading 若 \code{d[i+8]} 改为 \code{d[i*2]}，相邻线程跨 8 B，访问不再连续并增加事务数。
\practiceheading 用 Nsight Compute 的 Global Memory Efficiency 与 Shared Bank Conflicts 分别验证两类访问。
\sourcecheck{https://docs.nvidia.com/cuda/cuda-c-best-practices-guide/}{NVIDIA CUDA Best Practices Guide 的全局内存合并与共享内存 bank 规则；地址差逐项复核，置信度高}
\end{solutionexercise}

\begin{solutionexercise}{2}
\problemheading 编写一个使用转角的矩阵乘法内核函数，设计应对应图 6.4。

为使题意独立可读，图 6.4 的“转角”设计把第二个输入矩阵按列主序存储：线程以
连续地址协作读取其列片段，再把数据按转置方向写入带填充的共享内存瓦片；计算阶段
仍按行主序输出 $P=MN$。内核应给出边界检查和阶段间同步。

\approachheading 假设 $M$ 行主序、$N$ 列主序、$P$ 行主序。让相邻 \code{threadIdx.x}
沿 $N$ 的一列读取连续 $k$，写入带一列填充的共享瓦片后转置索引。

\answerheading 以下核心实现支持任意 $m\times k$ 与 $k\times n$ 尺寸；$N$ 的列主序
索引为 \code{row + col*k}。
\begin{lstlisting}[language=C++]
constexpr int TILE = 32;
__global__ void matmul_corner(const float* M, const float* Ncol,
                              float* P, int m, int kdim, int n) {
    __shared__ float Ms[TILE][TILE];
    __shared__ float Ns[TILE][TILE + 1];
    int tx = threadIdx.x, ty = threadIdx.y;
    int row = blockIdx.y * TILE + ty;
    int col = blockIdx.x * TILE + tx;
    float sum = 0.0f;

    for (int ph = 0; ph < (kdim + TILE - 1) / TILE; ++ph) {
        int mk = ph * TILE + tx;
        Ms[ty][tx] = (row < m && mk < kdim)
                   ? M[row * kdim + mk] : 0.0f;

        // tx walks down one column of column-major N: coalesced.
        int nk = ph * TILE + tx;
        int ncol = blockIdx.x * TILE + ty;
        Ns[tx][ty] = (nk < kdim && ncol < n)
                   ? Ncol[nk + ncol * kdim] : 0.0f;
        __syncthreads();

        #pragma unroll
        for (int q = 0; q < TILE; ++q)
            sum += Ms[ty][q] * Ns[q][tx];
        __syncthreads();
    }
    if (row < m && col < n) P[row * n + col] = sum;
}
// dim3 block(TILE, TILE);
// dim3 grid((n + TILE - 1) / TILE, (m + TILE - 1) / TILE);
\end{lstlisting}
两个无条件块级屏障分别防止读未完成和下一阶段覆盖；所有全局访问均有边界检查。
\code{Ns} 的额外一列把列写入的 bank 步长从 32 改成 33，避免 32 路冲突。

\complexity{$O(mkn)$ 总工作}{$O(TILE^2)$ 共享空间/块}{$mn$ 个输出点，按瓦片并行}
\conclusionheading 该实现以共享内存“转角”，同时保持 $M$ 行主序加载与 $N$ 列主序
加载合并，并处理边界、同步和 bank 冲突。
\discussionheading
\discussionpoint{“转角”是在片上改变数据观看方向。}
行主序矩阵乘中，$M$ 的一行适合连续读取，而点积又需要沿 $N$ 的列移动；若线程直接按计算方向读全局内存，其中一个输入容易形成大步长。转角方案先让线程按全局内存喜欢的连续方向协作加载，再把值写进共享瓦片的转置坐标，计算阶段便能从另一个方向读取同一批数据。额外写读共享内存看似绕路，却用低延迟片上重排换回规则的外存事务，并让一个输入值服务多个输出。FFT 蝶形重排、图像转置和某些卷积数据变换都体现同一原则：外存布局与计算布局不一致时，可在片上建立适配层。
\discussionpoint{加一列填充为何能打散 bank 冲突。}
对 32 个 bank 和 32 位字，\code{tile[32][32]} 相邻两行同一列相隔 32 个字，bank 号增加 32 后模 32 不变，因此一个 warp 沿列访问时可能全部落到同一 bank。把第二维改为 33 后，行间步长模 32 等于 1，连续 lane 的列访问便轮流经过所有 bank；只付出每行一个额外元素的容量，就把序列化访问变成可并行访问。这个数论解释依赖 bank 数、bank 字宽和元素类型，若元素是 64 位或 tile 形状改变，必须重新按字节地址计算。padding 是针对特定映射的构造，而不是看到二维共享数组就机械加一列的规则。
\discussionpoint{两个屏障保护共享瓦片的完整生命周期。}
第一次屏障位于协作加载之后，确保所有生产者完成写入，任何线程才开始读取整行或整列；第二次屏障位于本阶段计算之后，防止快线程用下一阶段数据覆盖慢线程仍在消费的槽位。边界线程即使没有有效全局元素，也必须写入零并到达屏障，若提前返回，整个块可能违反屏障参与条件。tile 增大能提高复用，却同步增加线程、共享容量与每线程累加状态，最终还要满足合法块维和驻留资源。正确性分析应先画出共享槽的“写入--读取--再次覆盖”时间线，再谈用更大瓦片换取性能。
\discussionpoint{异步流水化并没有取消依赖关系。}
异步复制和双缓冲可在计算当前 $k$ 瓦片时预取下一瓦片，但它们只是把两个共享缓冲区的阶段显式化，仍需等待“数据已到达”和“旧数据已消费”两个事件。Tensor Core 路径还要求按照 fragment 形状安排 warp 与数据布局，转角原则依然存在，只是适配层更细。评估时可先用成熟矩阵乘库作为性能参照，再检查外存流量、共享冲突、同步等待和寄存器压力是否随流水化改善。若减少了加载停顿却因双缓冲占用翻倍而降低并发，优化只是把瓶颈移动到另一资源，并不必然缩短总时间。
\variantheading 若 TILE=16，则块为 256 线程，共享数组约 2112 B，网格各维仍向上取整。
\practiceheading 与 cuBLAS SGEMM 对比，并查看 Nsight Compute 的 bank conflict、DRAM sectors 和 occupancy。
\sourcecheck{https://docs.nvidia.com/cuda/cuda-c-best-practices-guide/}{NVIDIA CUDA Best Practices Guide 的转置瓦片、合并读取及共享数组加一列消除 bank 冲突示例；索引独立复核}
\end{solutionexercise}

\begin{solutionexercise}{3}
\problemheading 对第五章的分块矩阵乘法，在 \code{BLOCK_SIZE} 的可能取值范围内，
哪些 \code{BLOCK_SIZE} 值能让内核完全避免对全局内存的未合并访问？（只需
考虑正方形线程块。）

\approachheading CUDA 将二维块按 x 维最快线性化。按本章把“一个 warp 跨越矩阵行”
统一归为未完全合并的简化模型，要让 $M,N$ 瓦片加载都落在同一矩阵行的连续位置，
块宽至少为 32，且 warp 边界必须与行边界重合。

\answerheading 正方形块有 $B^2$ 个线程，硬件每块最多 1024 个线程，所以
$B\le32$。当 $B=32$ 时，每个 warp 恰好对应线程块的一整行，加载 $M$ 和 $N$ 时
32 个线程都访问同一全局行中的连续元素。任何 $B<32$ 都会使至少一个 warp 跨越两条
或更多线程块行；在教材模型中，这种访问不算完全合并。因此题目预期的唯一答案是 32。

这不是跨架构的硬件定理。现代 GPU 会根据一个 warp 的实际地址集合生成一个或多个内存
事务；$B=8,16,24$ 的跨行片段在特定 leading dimension、基址对齐与事务粒度下也可能
充分利用事务。旧架构或不同缓存路径还可能得到另一结果。实际代码应以目标计算能力和
Nsight Compute 的事务/字节指标验证，而不能只看“是否跨行”。

\conclusionheading 在题设可能范围内，\code{BLOCK_SIZE = 32}。
\discussionheading
\discussionpoint{教材答案 32 来自“一个 warp 对齐一行”的模型。}
正方形块以 x 维最快线性化，若边长 $B=32$，每 32 个连续线程恰好构成一整行，warp 加载矩阵行时地址既连续又不会跨到下一行。又因为单块最多 1024 个线程，正方形边长不能超过 32，于是教材模型得到唯一候选 32。这个推理很适合建立“warp 地址集合”的直觉，却把跨行片段一概归入未完全合并，属于有意简化。现代硬件按实际地址覆盖生成内存扇区，两个对齐的半行片段也可能没有浪费，因此应把 32 称为题设预期，而不是跨架构必要条件。实际判断必须列出每个 lane 的字节地址集合，而不能只观察二维块的外形。
\discussionpoint{行首对齐与主导维会改变同一块宽的事务数。}
即使 $B=32$，若矩阵 leading dimension 让下一行从非事务边界开始，或子矩阵首地址偏移了一个元素，连续 128 B 也可能跨越额外扇区。反过来，$B=16$ 时一个 warp 横跨两条半行，若两段各自对齐并恰好覆盖完整扇区，实际传输字节仍可能全部有用。元素由 32 位改成 64 位后，每个 lane 的字节宽度翻倍，事务覆盖又会变化。块宽、物理行跨距、基址与元素类型因而共同决定效率，任何只保留 $B$ 的规则都不足以描述真实内存系统。分配函数给出的基址对齐也不能自动保证任意子视图仍保持相同对齐。
\discussionpoint{块形状还同时决定资源与协作结构。}
$32\times32$ 虽让每个 warp 对齐一行，却使用 1024 个线程，可能使每个 SM 只能驻留很少的块，并为每块分配较多共享内存和寄存器。改成 $32\times8$ 可保持 x 方向整 warp 连续，让每线程循环计算多行，以较少线程维持相同逻辑瓦片；代价是线程粗化和更长寄存器生存期。归约、卷积或 Tensor Core 内核还有固定的 warp 协作形状，不能只为全局读取选择线程块。合理设计应把访存方向、共享瓦片、同步域和驻留资源作为联合问题，而不是先找到“合并块宽”再忽略其余限制。
\discussionpoint{目标设备测量负责修正教学模型。}
可以把 $B=8,16,24,32$ 与若干非方形块作为候选，在不同 leading dimension 和起始偏移下运行相同数学任务。先比较请求字节与实际传输字节，确认跨行片段是否真的浪费事务，再结合内核总时间和资源报告判断这项浪费是否主导性能。若 $B=16$ 与 32 的传输效率接近而前者允许更多独立块驻留，最优配置完全可能不是教材答案。用模型提出可解释的候选、用硬件证据选择部署配置，正是架构知识随设备演进时仍保持有效的方法。基准还应固定输出正确性和工作量，避免不同边界处理使字节指标失去可比性。
\variantheading 在计算能力 6.x+、32 B 事务且行宽按 32 B 对齐时，$B=16$ 的两个半行片段也可完整利用四个事务。
\practiceheading 用 Nsight Compute 的 sectors/request 和 Global Load Efficiency 实测候选 $B$，不要把 32 当跨架构定律。
\sourcecheck{https://docs.nvidia.com/cuda/cuda-c-best-practices-guide/}{NVIDIA CUDA Best Practices Guide 对 warp 地址集合、对齐与内存事务效率的官方说明；结论 32 是本章简化定义下的教材预期答案，不外推为所有架构的必要条件}
\end{solutionexercise}

\begin{solutionexercise}{4}
\problemheading 实现一个使用向量加载的向量加法内核，并正确处理边界条件。

\approachheading 每个常规线程处理一个 16 B 的 \code{float4}，另用一个线程以标量
方式处理不足 4 个元素的尾部；网格按“完整向量数加尾线程”启动。

\answerheading
\begin{lstlisting}[language=C++]
__global__ void vecadd4(const float* x, const float* y, float* z, int n) {
    int tid = blockIdx.x * blockDim.x + threadIdx.x;
    int n4 = n / 4;
    if (tid < n4) {
        float4 a = reinterpret_cast<const float4*>(x)[tid];
        float4 b = reinterpret_cast<const float4*>(y)[tid];
        float4 c = make_float4(a.x + b.x, a.y + b.y,
                               a.z + b.z, a.w + b.w);
        reinterpret_cast<float4*>(z)[tid] = c;
    } else if (tid == n4) {
        for (int i = 4 * n4; i < n; ++i) z[i] = x[i] + y[i];
    }
}

void launch_vecadd(const float* x, const float* y, float* z, int n) {
    if (n <= 0) return;
    int work = n / 4 + ((n % 4) != 0);
    vecadd4<<<(work + 255) / 256, 256>>>(x, y, z, n);
    // Production code checks cudaGetLastError() and later sync/copy errors.
}
\end{lstlisting}
该实现要求 $x,y,z$ 的基地址满足 \code{float4} 的 16 B 对齐；\code{cudaMalloc}
返回的基地址满足该条件。若传入的是带任意元素偏移的子数组，调用方应保证对齐，或在
内核中改走全标量路径。线程只写互异元素，无竞态，也不需要块级同步。

\complexity{$O(n)$ 总工作}{$O(1)$ 每线程额外空间}{$\lceil n/4\rceil$ 个粗化任务}
\conclusionheading 完整四元组使用向量加载/存储，0--3 个尾元素由唯一线程标量处理，
无越界访问。
\discussionheading
\discussionpoint{向量类型减少的是指令开销，不是数据总量。}
\code{float4} 把四个相邻浮点数视为一个 16 B 对齐对象，编译器因而可能用一次宽加载和一次宽存储代替四组地址计算与标量指令。一个 warp 仍然请求相同数量的有效字节，设备 DRAM 的峰值也不会因为类型名字改变；若原标量访问已完美合并，物理事务数甚至可能相同。收益通常来自更少的指令、更清晰的对齐信息和每线程一次处理多项的粗化，而不是凭空提高带宽。查看最终机器指令很重要，因为源码中的向量类型并不保证编译器在所有上下文都保留单条宽访问。尾部不足四项时仍需标量或掩码路径，并验证它不会越过已分配对象边界。
\discussionpoint{16 字节对齐是访问契约而非转换效果。}
\code{cudaMalloc} 返回的基础地址通常满足充分对齐，但若调用方传入 \code{x+1}，地址只移动 4 B，强制转换为 \code{float4*} 不会把它重新排列到 16 B 边界。通用实现可先用少量标量元素推进到对齐地址，再执行向量主体，最后处理标量尾部；也可在接口层明确要求对齐并为不满足者选择另一内核。长度与三个指针的偏移都要检查，因为输入对齐而输出未对齐同样非法。将对齐写成可验证的前置条件，比依赖某个测试分配恰好从基础地址开始更可靠。
\discussionpoint{尾部不能通过越界读取后屏蔽写回来补救。}
若 $n\bmod4=3$，最后一个完整向量之后只剩三个浮点数，读取一个 \code{float4} 已经越过对象边界，即使随后不写第四项也无法撤销非法读取。稳妥方案让唯一线程按标量处理 0--3 个余数，或在分配层提供明确的可安全过读 padding，并保证该约定对所有地址空间成立。C++ 还对对象类型、别名和对齐有规则，任意把来源不明的字节区解释成向量对象可能产生额外未定义行为。余数类与指针偏移的笛卡尔积应进入单元测试，才能证明向量主体和尾路径之间既无重叠也无遗漏。
\discussionpoint{更宽的向量可能以寄存器压力换指令减少。}
一次处理四项意味着线程要同时持有更多输入和输出，简单加法中这通常便宜，复杂表达式却可能延长许多值的活跃区间并增加寄存器用量。可驻留 warp 因此下降，极端时还会出现 local-memory spill；与此同时，原本已经合并的标量访问并不会因 128 位指令而减少 DRAM 字节。选择 \code{float2}、\code{float4} 或标量路径时，应确认生成的宽指令，并比较资源用量、实际传输与总时间。向量宽度是编译器和微架构的联合参数，而不是越宽越接近硬件峰值的单调旋钮。
\variantheading 若 $n=1027$，有 256 个完整 \code{float4} 和 3 个尾元素，工作任务数为 257。
\practiceheading 用 Compute Sanitizer memcheck 覆盖四种余数，并以 SASS 确认生成 128 位访问。
\sourcecheck{https://docs.nvidia.com/cuda/cuda-c-programming-guide/}{NVIDIA CUDA C++ Programming Guide 的内建向量类型及对齐要求；边界 $n\bmod4$ 全枚举复核}
\end{solutionexercise}

\begin{solutionexercise}{5}
\problemheading 考虑下面的共享内存数组，以及一个包含一个 warp（32 个线程）的线程块
对该数组的访问：
\begin{lstlisting}[language=C++,numbers=none]
__shared__ float a[1024];
a[threadIdx.x*stride] = ...;
\end{lstlisting}
假设 \code{a[0]} 位于 bank 0。对于下面每个 \code{stride} 值，指出该 warp
访问了哪些 bank，以及每个 bank 上的 bank 冲突数量（如果有）：
\begin{enumerate}[label=\alph*.,leftmargin=2.2em]
  \item \code{stride = 32}；
  \item \code{stride = 31}；
  \item \code{stride = 24}；
  \item \code{stride = 16}；
  \item \code{stride = 12}；
  \item \code{stride = 8}；
  \item \code{stride = 7}。
\end{enumerate}

\approachheading 对 32 位 \code{float}，线程 $t$ 访问
$b_t=(t\,s)\bmod32$ 号 bank。访问 bank 数为 $32/\gcd(s,32)$，每个所用 bank 上有
$\gcd(s,32)$ 个线程。

\answerheading
\begin{center}
\begin{tabular}{c@{\quad}l@{\quad}l}
\toprule
stride & 访问的 bank & 冲突\\
\midrule
32 & $\{0\}$ & bank 0 上 32 路（31 个额外请求）\\
31 & $\{0,1,\ldots,31\}$（逆序置换） & 无冲突\\
24 & $\{0,8,16,24\}$ & 每个 8 路（7 个额外请求）\\
16 & $\{0,16\}$ & 每个 16 路（15 个额外请求）\\
12 & $\{0,4,8,12,16,20,24,28\}$ & 每个 4 路（3 个额外请求）\\
8  & $\{0,8,16,24\}$ & 每个 8 路（7 个额外请求）\\
7  & $\{0,1,\ldots,31\}$（置换） & 无冲突\\
\bottomrule
\end{tabular}
\end{center}

\conclusionheading 冲突路数依次为 32、1、8、16、4、8、1；路数 1 表示无冲突。
\discussionheading
\discussionpoint{最大公因数揭示了步长访问的周期。}
对连续 32 位字和 32 个 bank，lane $t$ 访问 \code{a[t*s]} 时落入的 bank 是 $(ts+c)\bmod32$。若 $s$ 与 32 互素，序列在 32 步内遍历所有 bank，像 $s=7$ 或 31 都只是对 bank 的置换；若最大公因数为 8，就只访问 $32/8=4$ 个 bank，每个 bank 接收 8 个不同地址。于是冲突路数恰由 $\gcd(s,32)$ 给出，不必为每个 stride 手画 32 个 lane。这个数论视角还能解释为何二维数组加一列常有效：它把行步长从 32 改成与 32 互素的 33。
\discussionpoint{同一 bank 不一定意味着同一种硬件行为。}
当多个 lane 读取完全相同的共享字时，硬件可以广播该值，虽然它们都落在一个 bank，也不会像访问 32 个不同地址那样串行。若这些 lane 写同一地址，问题首先是多个写者产生的数据竞争或未定义的最终值，不能因为硬件可能合并某些请求就宣称程序正确。本题的 stride 访问让 lane 指向不同数组元素，所以 $s=32$ 确实是多路冲突；把它与同址读取对照，能看出 bank 号只是一半信息，字内偏移和读写方向同样重要。性能规则永远不能替代生产者、消费者和竞态的语义分析。
\discussionpoint{先换算成字节地址，才能推广到其他类型。}
简单公式假设每个元素恰好占一个 32 位 bank 单元；若改成 64 位值，一个元素可能跨越两个这样的单元，同一逻辑 stride 便对应不同 bank 序列。结构体含 padding、向量类型一次覆盖多个字，二维数组又把第二维长度变成物理行步长，都使“元素步长”不再等同于“bank 步长”。不同计算能力的共享内存组织也应以目标文档为准，不应把历史上可配置的模式泛化到所有现代设备。可靠做法是从基础地址加字节偏移开始，按目标 bank 宽度映射，再判断同一 warp 的实际请求。
\discussionpoint{消除冲突也要核算所付出的资源代价。}
二维数组加一列可改变模映射，warp shuffle 可直接在线程寄存器间交换值，重新排列索引或分阶段访问也能避免多个不同地址同时落入一 bank。每种方法都有成本：padding 占共享容量，shuffle 受 warp 作用域与参与掩码限制，分阶段则增加指令和同步。轻微冲突有时被其他延迟隐藏，而多占的共享内存可能少驻留一个完整块，所以理论冲突路数下降不保证总时间缩短。应先构造无冲突的功能等价版本，再比较共享吞吐、资源占用和内核时间，确认优化消除的是主导瓶颈而非漂亮的单一计数。
\variantheading \code{stride=6} 时 $\gcd(6,32)=2$，访问 16 个 bank，每个 bank 2 路冲突。
\practiceheading 在目标 GPU 上查看 Nsight Compute 的 Shared Bank Conflicts，并注意广播读与不同地址写的差别。
\sourcecheck{https://docs.nvidia.com/cuda/cuda-c-best-practices-guide/}{NVIDIA CUDA Best Practices Guide 对 32 个 bank、连续 32 位字映射及冲突串行化的说明；模 32 序列独立枚举}
\end{solutionexercise}


\part{并行模式}
\chapter{卷积}

\begin{solutionexercise}{1}
\problemheading 计算图 7.3 中的 $P[0]$ 值。

为使题意独立可读，图 7.3 给出输入
$N=\{8,2,5,4,1,7,3\}$ 和五点滤波器 $F=\{1,3,5,3,1\}$；滤波器中心与当前
输出位置对齐，数组外的 ghost cell 取 0，并沿用本章不反转滤波器的相关运算约定。
底本使用 $N,F,P$，译稿图中对应写作 $x,f,y$。

\approachheading 原书示例采用零填充，输入为
$\{8,2,5,4,1,7,3\}$，滤波器为 $\{1,3,5,3,1\}$；$P[0]$ 左侧缺两个输入。

\editorialnote 练习沿用底本符号 $N,F,P$；中文译稿图 7.3 改用 $x,f,y$。二者逐项对应，
所以本题的 $P[0]$ 就是译稿图中的 $y[0]$。

\answerheading
\[
P[0]=1\times0+3\times0+5\times8+3\times2+1\times5=51.
\]
这里沿用本章代码的相关运算约定（滤波器不反转）。

\conclusionheading $P[0]=51$。
\discussionheading
\discussionpoint{卷积与互相关的名称必须落实到索引公式。}
看到“卷积”二字，为什么还不能立即决定核的方向？数学卷积通常写成 $P[i]=\sum_j F[j]N[i-j]$，滑动前会相对相关运算反转核；本章代码却按 $N[i-r+j]F[j]$ 读取，因此严格说是离散互相关，对非对称核会给出不同答案。深度学习库常把名为 convolution 的前向算子实现为相关，因为可学习权重能在训练中吸收方向，但预训练权重、手工滤波器和信号处理公式并不能自动互换。接口与测试应保存完整索引定义，而非只记算子名称；本题的 51 正是确认方向契约的一个小型见证。

\discussionpoint{边界规则是算子定义的一部分。}
为什么内部公式完全相同，首个输出仍可能不是 51？零填充把左侧两个缺失样本解释为 0；边缘复制会使用 8，镜像会取域内对称位置，循环边界则从数组末端取值，valid 模式甚至不产生这个位置的输出。每种规则都隐含了对观测域外信号的假设：周期谱分析适合循环延拓，图像边缘可能选镜像，PDE 则由物理边界条件决定。因而端点不是需要“修补”的性能例外，而是模型语义；跨库复现必须一起记录 padding、输出长度和偶数核 anchor，不能只比较内部样本。

\discussionpoint{核的矩把局部内积连接到物理作用。}
一串系数除了产生 51，还说明滤波器在做什么？核系数和 $\sum_jF[j]$ 是对常量输入的直流增益；一阶矩 $\sum_j(j-r)F[j]$ 描述对斜坡的响应，正负相消的核常用来近似导数或突出边缘。对称平滑核通常没有方向偏置，非对称核则可能表达时间因果性、运动方向或空间偏移，因此核反转会改变可解释含义，而不仅是数组顺序。工程上可在部署前检查增益、对称性和矩条件，再把这些性质同归一化、单位和允许的边界伪影对应起来，这比只对一个随机输入验数更能发现模型层错误。

\discussionpoint{结构化测试比随机抽查更能定位卷积错误。}
若 GPU 输出只在首尾错几个元素，怎样最快区分核方向、边界和索引问题？单位脉冲会把核的方向与 anchor 直接显露出来；全常量输入可检查内部直流增益，线性斜坡能暴露一阶矩和一格偏移，非对称短数组则专门区分相关与卷积。随后再用独立 CPU 循环逐元素比较，并把首尾 $r$ 个位置单独列出，便能把边界错误同主体循环错误隔离。随机测试适合扩大覆盖，但不应替代这些具有诊断性的基例；CUDA 版本还需配合越界检查工具，确认“数值碰巧正确”不是由非法读取造成。
\variantheading 若左侧采用边缘复制，则 $P[0]=1(8)+3(8)+5(8)+3(2)+1(5)=83$。
\practiceheading 用独立 CPU 参考实现固定边界与核方向，再逐元素比较 CUDA 输出，并以 Compute Sanitizer memcheck 排除端点越界。
\sourcecheck{https://github.com/NVIDIA/cuda-samples/tree/master/cpp/2_Concepts_and_Techniques/convolutionSeparable}{NVIDIA convolutionSeparable 示例直接展示 halo 越界位置写零；该示例会反转卷积核，和本章相关运算约定不同，因此 51 只由底本数据按本章公式独立代入核验}
\end{solutionexercise}

\begin{solutionexercise}{2}
\problemheading 对数组
\[
  N=\{4,1,3,2,3\}
\]
执行一维卷积，滤波器为
\[
  F=\{2,1,4\}.
\]
求得到的输出数组。沿用本章约定：输出与输入等长，数组外 ghost cell 取 0，
滤波器中心与当前输出对齐且系数不反转。

\approachheading 半径为 1，数组两端按 0 延拓，并沿用本章不翻转滤波器的定义
$P[i]=2N[i-1]+N[i]+4N[i+1]$。

\answerheading
\[
\begin{aligned}
P[0]&=2(0)+1(4)+4(1)=8,\\
P[1]&=2(4)+1(1)+4(3)=21,\\
P[2]&=2(1)+1(3)+4(2)=13,\\
P[3]&=2(3)+1(2)+4(3)=20,\\
P[4]&=2(2)+1(3)+4(0)=7.
\end{aligned}
\]

\complexity{$O(NM)$，本题 15 个候选乘法}{$O(N)$ 输出空间}{$N$ 个输出彼此独立}
\conclusionheading 输出数组为 $\{8,21,13,20,7\}$。
\discussionheading
\discussionpoint{滑动窗口可以看成带结构的矩阵乘法。}
五个输出为何能由同一小段代码计算？把输入写成列向量后，每个输出对应矩阵的一行，核系数沿对角线重复，内部呈 Toeplitz 结构；零填充只是在首尾行删去越界列，因此线性和局部性都清楚可见。这个表示还能直接导出伴随算子：反向传播对输入的梯度相当于用转置矩阵聚合相邻输出梯度，边界行同样不能忽略。实际实现不会真的构造这个大矩阵，但矩阵视角有助于证明线性、检查维度，并把卷积同稀疏矩阵、有限冲激响应系统和自动微分连接起来。对非对称核，转置还会再次改变方向。

\discussionpoint{same、valid、full 与 circular 描述不同的问题。}
输出数组究竟应该有多长，并非由核系数单独决定。valid 只保留核完全落在输入内的位置，长度通常为 $N-M+1$；full 保留一切非空重叠，长度为 $N+M-1$；same 通过选择填充维持输入长度，而 circular 则按 $N$ 取模并假定信号周期延拓。奇数核有自然中心，本题因而较少歧义；偶数核的 anchor 偏左或偏右会带来一格相位差，不同框架还可能采用不对称 padding。跨系统交换模型时应记录左右填充量、stride、dilation 与 anchor，这些参数共同定义输出坐标系。

\discussionpoint{直接卷积与 FFT 卷积服务于不同尺度。}
既然卷积可在频域相乘，为什么小核实现仍普遍使用直接循环？直接法做 $O(NM)$ 工作，控制简单、启动延迟低且边界易处理；线性 FFT 卷积需先零填充到至少 $N+M-1$，进行正逆变换和逐点乘法，复杂度为 $O((N+M)\log(N+M))$。只有在 $M=O(N)$ 等可把总变换长度按 $N$ 表示的条件下，才可简写成 $O(N\log N)$；短核时常数、额外缓冲和变换成本可能更差。流式长信号可用 overlap-add 或 overlap-save 分块，把有限工作区、延迟预算与 FFT 长度一并纳入选择。

\discussionpoint{累加类型和测试域共同决定工程准确度。}
本题整数很小，为什么生产接口仍要说明累加精度？若输入与系数为 8 位或 16 位整数，$M$ 项乘积之和可能超出原类型；浮点长核虽然不发生整数回绕，却会因加法次序、FMA 和混合精度累加得到不同末位。实现应根据最坏幅值推导累加范围，规定饱和或回绕语义，并用更高精度参考评估误差。测试除了固定样例，还要覆盖核长于输入、全零、单位脉冲、极值和随机数组，并确保 CPU 与 GPU 使用完全相同的核方向和边界；否则差异可能来自契约而非并行代码。参考值也应避免使用同一低精度累加路径。
\variantheading 滤波器改为 $\{1,1,1\}$ 时，零填充 same 输出为 $\{5,8,6,8,5\}$。
\practiceheading 与 cuDNN 或 CPU reference 对照时显式设置 padding 和核方向。
\sourcecheck{https://github.com/NVIDIA/cuda-samples/tree/master/cpp/2_Concepts_and_Techniques/convolutionSeparable}{NVIDIA convolutionSeparable 示例核验显式零 halo 的边界策略；系数方向按本章定义，五个输出由独立手算核验}
\end{solutionexercise}

\begin{solutionexercise}{3}
\problemheading 你认为下面的一维卷积滤波器分别在做什么？
\begin{enumerate}[label=\alph*.,leftmargin=2.2em]
  \item $[0,1,0]$
  \item $[0,0,1]$
  \item $[1,0,0]$
  \item $\left[-\frac12,0,\frac12\right]$
  \item $\left[\frac13,\frac13,\frac13\right]$
\end{enumerate}
沿用本章的 same、零填充和不反转滤波器约定。

\approachheading 把滤波器直接代入
$P[i]=F[0]N[i-1]+F[1]N[i]+F[2]N[i+1]$。

\answerheading
\begin{enumerate}[label=\alph*.,leftmargin=2.2em]
\item $[0,1,0]$：恒等滤波，输出当前样本。
\item $[0,0,1]$：输出 $N[i+1]$，把信号向较小下标方向平移一格（提前一格）。
\item $[1,0,0]$：输出 $N[i-1]$，把信号向较大下标方向平移一格（延迟一格）。
\item $[-\tfrac12,0,\tfrac12]$：单位采样间隔下的中心一阶差分
$(N[i+1]-N[i-1])/2$，可检测斜率或边缘。
\item $[\tfrac13,\tfrac13,\tfrac13]$：三点移动平均，抑制高频变化并平滑信号。
\end{enumerate}
零填充会使平移和平均滤波在端点出现边界效应。

\complexity{$O(N)$（滤波器长度固定为 3）}{$O(N)$ 输出空间}{$N$ 个输出可并行}
\conclusionheading 五者依次为恒等、左移/提前、右移/延迟、中心差分和三点平均。
\discussionheading
\discussionpoint{核的矩与频率响应给出互补解释。}
怎样把“平滑”或“检测变化”的直觉变成可计算性质？系数和等于零频增益，因此恒等核和三点平均保留常量，而中心差分的系数和为零，会消除直流分量；系数相对中心的一阶矩则决定对线性斜坡的响应。进一步令 $H(\omega)=\sum_jF[j]e^{\mathrm{i}\omega j}$，其幅值描述各频率被放大或衰减多少，相位描述波形如何平移。矩适合检查低频局部行为，频率响应展示全谱，两者共同把手工核连接到数字信号处理中的低通、高通和相位设计。若响应零点或相位方向不符，往往意味着系数顺序、anchor 或采样间隔有误。

\discussionpoint{平移结论依赖索引方向与相关约定。}
核 $[0,0,1]$ 为什么在本题中表示提前，而在另一份资料里可能写成延迟？这里输出位置 $i$ 读取 $N[i+1]$，所以原来位于较大下标的样本出现在较小输出下标；若采用数学卷积，核先反转，结论便互换。对实时信号，读取未来样本的提前算子不可因果实现，除非系统允许缓存或离线处理；对数组图像，它可能只是一次坐标重排。非对称滤波器还会引入频率相关相位，因此文档应同时给出索引正方向、时间含义和 anchor，不能用“左移”一个词承担全部语义。单位脉冲可直接检验这一方向。

\discussionpoint{差分核还需要网格尺度才能成为导数。}
为什么 $[-\tfrac12,0,\tfrac12]$ 只在单位采样间隔下近似一阶导数？泰勒展开给出 $f(x+h)-f(x-h)=2hf'(x)+O(h^3)$，所以一般公式要除以 $2h$，截断误差为 $O(h^2)$；二阶核 $[1,-2,1]$ 则还要除以 $h^2$。若忘记尺度，网格加密后数值会改变物理单位，无法讨论收敛性，即使边缘图看起来仍然合理。PDE 离散、图像梯度和传感器采样都依赖这一区别，实际代码应把 spacing 作为接口参数或把它明确吸收到核系数中。收敛测试还应让 $h$ 系统缩小并核对误差阶。

\discussionpoint{平滑与求导之间存在噪声和分辨率取舍。}
中心差分能突出边缘，为何常不直接作用于原始测量？求导的频率响应随频率增大，会同时放大高频噪声；三点平均能抑制噪声，却会削弱窄峰并模糊真实边缘，所以二者的先后和尺度决定最终分辨率。Sobel 或 Gaussian derivative 把局部平滑与方向导数组合起来，代价是更宽支持域、更多边界影响及位置偏差。验证时可依次使用脉冲、常量、斜坡和不同频率正弦波，分别检查方向、直流增益、导数尺度与截止特性，并按任务的信噪比和空间分辨率选择核宽度。
\variantheading $[1,-2,1]$ 给出二阶中心差分 $N[i-1]-2N[i]+N[i+1]$。
\practiceheading 用脉冲和线性斜坡作为 CUDA 单元测试输入，验证方向、增益和边界，并用 Compute Sanitizer memcheck 覆盖两个端点。
\sourcecheck{https://github.com/NVIDIA/cuda-samples/tree/master/cpp/2_Concepts_and_Techniques/convolutionSeparable}{NVIDIA convolutionSeparable 示例用于核验卷积的邻域访问结构；五个滤波器的方向与作用由本章相关运算公式逐项代入，不借用示例的核反转约定}
\end{solutionexercise}

\begin{solutionexercise}{4}
\problemheading 对大小为 $N$ 的数组执行一维卷积，滤波器大小为 $M$：
\begin{enumerate}[label=\alph*.,leftmargin=2.2em]
  \item 总共有多少个 ghost cell？
  \item 如果把 ghost cell 视为乘法（乘以 0），会执行多少次乘法？
  \item 如果不把 ghost cell 视为乘法，会执行多少次乘法？
\end{enumerate}
沿用本章居中、零填充、输出与输入同尺寸的约定，因此 $M$ 为奇数；闭式计数还假设
$N\ge (M-1)/2$。

\approachheading 令半径 $r=(M-1)/2$，假设 $N\ge r$ 且两端零填充。区分填充数组中
“唯一 ghost cell 数”和所有边界输出累计引用 ghost 的次数。

\answerheading
\begin{enumerate}[label=\alph*.,leftmargin=2.2em]
\item 左右各需 $r$ 个唯一 ghost cell，总数 $2r=M-1$。
\item 若乘零也执行，每个 $N$ 个输出均做 $M$ 次乘法，共 $NM$ 次。
\item 左边各输出跳过的 ghost 引用数为 $r+(r-1)+\cdots+1=r(r+1)/2$，右边相同，
总共跳过 $r(r+1)$ 次。因此实际乘法数为
\[
 NM-r(r+1)=NM-\frac{M^2-1}{4}.
\]
\end{enumerate}
若题目把“ghost cell”理解为计算过程中的累计 ghost 引用，则该数是 $r(r+1)$，
而非 (a) 的唯一填充位置数。

\complexity{$O(NM)$}{$O(N)$ 输出；显式填充时另需 $M-1$ 个元素}{$N$ 个输出可并行}
\conclusionheading (a) $M-1$；(b) $NM$；(c) $NM-(M^2-1)/4$。
\discussionheading
\discussionpoint{边界缺口形成两个可以精确求和的三角形。}
为什么不执行 ghost 乘法时，要从 $NM$ 中减去 $r(r+1)$？最左输出缺 $r$ 项，下一个缺 $r-1$ 项，直到第 $r$ 个边界输出只缺 1 项，单侧缺口是 $1+\cdots+r=r(r+1)/2$；右侧对称，故总缺口恰为 $r(r+1)$。这里 $M-1$ 统计的是填充数组中不同的 ghost 位置，而三角数统计它们被各输出累计引用的次数，两者回答不同问题。把完整工作矩形减去边界缺口的思路也用于 stencil 有效点、带状稀疏矩阵非零元和 halo 流量分析，关键是先说明统计对象是否去重。

\discussionpoint{闭式公式必须连同 anchor 与尺寸假设使用。}
若核比输入还宽，原公式能否照抄？推导假定 $M$ 为奇数、anchor 位于中心且 $N\ge r=(M-1)/2$，这样左右边界的缺口序列具有上述三角形结构；若 $N<r$，同一输出可能同时越过两侧，应逐位置求输入区间与核支持的交集长度。偶数核不存在唯一中心，anchor 偏左或偏右会让两侧缺口不对称；valid、full 或 dilation 又会改变输出坐标集合。可靠实现可以把通用交集公式作为参考，再把闭式仅用于满足前提的快速路径，这也是防止边界优化悄悄扩大适用域的常见做法。

\discussionpoint{跳过乘零不必然缩短运行时间。}
少执行 $r(r+1)$ 次乘法，为什么 GPU 仍可能更慢？逐项边界判断会增加整数索引、谓词或分支；同一 warp 中内部线程和端点线程走不同路径时，节省的算术可能被控制流和低利用率抵消。显式预填充可让主体内核无分支且访问规则，却要支付额外存储、初始化和读流量；另一个选择是拆成大规模内部内核与小型边界内核。对固定小核且 $N$ 很大，边界比例约按 $M/N$ 消失，专门优化端点价值有限；对大量很短序列则相反，所以应测端到端时间而非只比较源代码乘法数。

\discussionpoint{改变核结构能比省去边界乘法带来更大收益。}
若性能瓶颈在全部 $NM$ 项，继续雕琢端点是否抓住主要矛盾？稀疏核可存储非零系数和偏移，只遍历有效支持；可分离核把二维外积拆成两次一维运算；长而稠密的一维核可转向 FFT，其工作结构与直接卷积完全不同。每种变化同时影响数据布局、中间缓冲、舍入路径和适用尺寸，例如稀疏访问可能损害合并，FFT 对短核有较大固定成本。性能报告应区分数学有效乘加、实际发射指令、DRAM 字节和预处理时间，并以同一数值契约验证，这样才不会把“少算了零”误当成整体算法已经最优。
\variantheading $N=5,M=3$ 时 ghost 为 2、计零乘法 15 次、不计零乘法 13 次。
\practiceheading 枚举不同 $N,M$ 的边界索引并与 CPU 计数对照，再用 Compute Sanitizer memcheck 覆盖 $N<M$ 的小数组边界。
\sourcecheck{https://github.com/NVIDIA/cuda-samples/tree/master/cpp/2_Concepts_and_Techniques/convolutionSeparable}{NVIDIA convolutionSeparable 示例核验零边界与 halo 加载；$M-1$ 和左右三角数完全由独立代数推导核验}
\end{solutionexercise}

\begin{solutionexercise}{5}
\problemheading 对大小为 $N\times N$ 的方阵执行二维卷积，方形滤波器大小为 $M\times M$：
\begin{enumerate}[label=\alph*.,leftmargin=2.2em]
  \item 总共有多少个 ghost cell？
  \item 如果把 ghost cell 视为乘法（乘以 0），会执行多少次乘法？
  \item 如果不把 ghost cell 视为乘法，会执行多少次乘法？
\end{enumerate}
沿用本章居中、零填充、输出与输入同尺寸的约定，因此 $M$ 为奇数；闭式计数还假设
$N\ge (M-1)/2$。

\approachheading 令 $r=(M-1)/2$。唯一 ghost cell 是外扩后方阵与原方阵的面积差；
不乘 ghost 时，二维有效配对数可分解成两个相同的一维有效配对数之积。

\answerheading
\begin{enumerate}[label=\alph*.,leftmargin=2.2em]
\item 填充后边长为 $N+2r=N+M-1$，唯一 ghost cell 数为
\[
 (N+M-1)^2-N^2.
\]
\item 若 ghost 也乘以 0，共有 $N^2$ 个输出，每个做 $M^2$ 次乘法，总数
$N^2M^2$。
\item 一维所有输出的有效输入配对数为
$S=NM-r(r+1)=NM-(M^2-1)/4$。二维行列选择独立，故有效乘法总数为
\[
 S^2=\left(NM-\frac{M^2-1}{4}\right)^2.
\]
\end{enumerate}
上述闭式假设 $N\ge r$；极小输入应直接用截断区间求和。

\complexity{$O(N^2M^2)$}{$O(N^2)$ 输出空间}{$N^2$ 个输出可并行}
\conclusionheading 三项依次为 $(N+M-1)^2-N^2$、$N^2M^2$ 和
$[NM-(M^2-1)/4]^2$。
\discussionheading
\discussionpoint{支持域计数可分解并不意味着滤波器可分离。}
二维有效乘法数为什么能写成一维数量 $S^2$？矩形输入与方形支持使行方向和列方向的合法索引选择相互独立，每一个合法行配对都能同每一个合法列配对组合，因此笛卡尔积的大小等于两者乘积。这个论证只依赖“哪些偏移存在”，并未约束每个位置的系数；一般 $M\times M$ 核仍有 $M^2$ 个独立权重。只有系数矩阵满足低秩条件、特别是能写成列向量与行向量外积时，计算本身才能拆为两次一维卷积。区分组合可分解与代数可分离，可避免从计数公式错误推出算法复杂度已经降低。秩检验才能回答后一个问题。

\discussionpoint{ghost 面积需要用容斥避免重复计算角点。}
若沿四条边各加一条宽 $r$ 的带，角落应算几次？直接相加边带会重复覆盖四个角，最简洁的办法是用扩展方阵面积 $(N+2r)^2$ 减去原面积 $N^2$，等价展开后再用容斥校正交集。三维 halo 更明显：六个面相交成十二条棱，棱又相交于八个角，若按消息或存储区域逐块统计，很容易重复或漏项。多 GPU stencil 中还要区分“存储的唯一 halo 单元”和“按邻居分别发送的元素”，后者可能有意重复角数据；因此几何公式必须与通信策略和所有权定义配套。

\discussionpoint{非矩形支持域需要回到偏移集合。}
圆盘、十字或膨胀卷积为何不能直接套用 $S^2$？这些支持不再是行偏移集合与列偏移集合的完整笛卡尔积，某个行偏移是否合法还取决于列偏移，应对每个核偏移 $(\delta_y,\delta_x)$ 统计它与输入域重叠的输出位置。稀疏支持可只存非零系数及偏移，显著减少算术和系数流量，但不规则地址可能破坏连续访问，并增加索引开销。规则十字核、结构元素和图邻域各有更专用的遍历方式，工程选择应同时比较支持稀疏度、地址规律和分支一致性，而不是只看非零项个数。偏移排序还会改变缓存局部性。

\discussionpoint{全 1 测试只有输入与支持核都为 1 才能计数邻点。}
把输入全设为 1，输出是否自动等于有效邻居数量？若核系数任意，边界位置得到的是仍落在域内的那些系数之和，而不是系数个数；只有输入为全 1，且核或支持掩码也在每个有效位置取 1，输出才恰好是邻点计数图。内部应为 $M^2$，边上减去一条缺失带，角上同时缺两维，这个模式能快速暴露填充、anchor 与二维索引错误。随后再换回真实系数，可用单位脉冲检查方向；把结构测试与随机 CPU 对照结合，诊断能力远强于只核对总乘法数。这一区分防止把加权和误报成拓扑度数。
\variantheading $N=5,M=3$ 时 ghost 为 24，含零乘法 225 次，不含零乘法 169 次。
\practiceheading 用全 1 输入和全 1 支持核生成有效邻点计数图，并以 Compute Sanitizer 检查 halo 边界。
\sourcecheck{https://github.com/NVIDIA/cuda-samples/tree/master/cpp/2_Concepts_and_Techniques/convolutionSeparable}{NVIDIA convolutionSeparable 示例核验二维可分离实现的零 halo；面积差和有效配对平方由独立代数推导核验}
\end{solutionexercise}

\begin{solutionexercise}{6}
\problemheading 对大小为 $N_1\times N_2$ 的矩形矩阵执行二维卷积，矩形滤波器大小为
$M_1\times M_2$：
\begin{enumerate}[label=\alph*.,leftmargin=2.2em]
  \item 总共有多少个 ghost cell？
  \item 如果把 ghost cell 视为乘法（乘以 0），会执行多少次乘法？
  \item 如果不把 ghost cell 视为乘法，会执行多少次乘法？
\end{enumerate}
沿用本章居中、零填充、输出与输入同尺寸的约定，$M_1,M_2$ 为奇数；闭式计数还假设
$N_i\ge (M_i-1)/2$。

\approachheading 令 $r_i=(M_i-1)/2$，分别沿两维计数，再利用笛卡尔积相乘。

\answerheading
\begin{enumerate}[label=\alph*.,leftmargin=2.2em]
\item 唯一 ghost cell 数为
\[
(N_1+M_1-1)(N_2+M_2-1)-N_1N_2.
\]
\item 把 ghost 乘零也计入时，共执行 $N_1N_2M_1M_2$ 次乘法。
\item 令
\[
S_i=N_iM_i-r_i(r_i+1)
   =N_iM_i-\frac{M_i^2-1}{4},\qquad i=1,2.
\]
不执行 ghost 乘法时总数为
\[
S_1S_2=\left(N_1M_1-\frac{M_1^2-1}{4}\right)
        \left(N_2M_2-\frac{M_2^2-1}{4}\right).
\]
\end{enumerate}
闭式同样假设 $N_i\ge r_i$；否则用每个输出处的有效交集长度求和。

\complexity{$O(N_1N_2M_1M_2)$}{$O(N_1N_2)$ 输出空间}{$N_1N_2$ 个输出可并行}
\conclusionheading 结果由两维的一维计数相乘得到，见上述三个公式。
\discussionheading
\discussionpoint{各向异性几何必须逐维保留。}
为什么不能先把矩形问题近似成一个边长为 $N$ 的方形再计数？高度方向的半径 $r_1$ 与宽度方向的半径 $r_2$ 独立决定边界缺口，一维有效配对分别为 $S_1$ 和 $S_2$，二维总数才是 $S_1S_2$；长条输入中某一方向可能几乎全是边界，另一方向却以内部点为主。时频谱、扫描线和批量短序列都具有这种不对称。保留逐维公式不仅使乘法数正确，还能预测哪一侧 halo、分支和缓存浪费更严重，从而指导矩形 tile、分区方向与专用边界内核的选择。这一预测还可用于选择分区方向。

\discussionpoint{各向异性核把数组尺寸连接到物理尺度。}
矩形核较长的一边是否必然表示更强平滑？只有在两个数组轴的采样间隔相同且系数按同一单位设计时才成立；运动模糊、时空卷积和医学体数据常在不同轴上拥有不同物理 spacing，$M_1$ 与 $M_2$ 不能脱离单位解释。若系数矩阵可写成 $uv^{\mathsf T}$，两次一维卷积可把每输出工作从 $M_1M_2$ 降到 $M_1+M_2$，但会引入中间舍入和缓冲。近似低秩分解还可用少数可分离项换取误差，因此需要把秩、任务误差与内存流量一起选择。标定数据应覆盖各轴的典型频率与极端幅值。

\discussionpoint{数学对称并不保证内存访问对称。}
把长和宽互换后公式形式不变，GPU 时间为何可能明显不同？行主序数组中，连续线程沿宽度读取可形成合并访问，沿高度读取则跨越 leading dimension；即使有效乘法数量相同，缓存行和事务数也不同。两阶段可分离卷积常在竖向阶段使用共享内存转置，或选取宽而矮的 tile 让全局加载保持连续，但这又影响 halo 比例与占用率。长宽互换是发现索引写反的好测试，却不能作为性能等价证明；分析时应分别报告逻辑工作、地址步长和实际内存事务。跨步大小必须纳入性能记录。

\discussionpoint{高维卷积还需要说明哪些维度参与求和。}
把二维公式推广到视频或深度学习张量，只增加一个空间下标就够了吗？矩形支持域在独立空间维上仍可逐维计数，但 batch 通常只是并行样本，channel 往往参与求和并产生多个输出通道，group 又限制通道连接；布局可能是 NCHW、NHWC 或带阻塞的物理格式。分布式体数据还会让某些空间维跨设备，halo 通信量由分区轴和半径共同决定。因而接口必须区分逻辑维度、卷积维度、存储顺序和分区方式，单给 $N_1,N_2,M_1,M_2$ 只能定义本题的单通道二维情形。
\variantheading $N_1=5,N_2=4,M_1=M_2=3$ 时不计 ghost 的乘法数为 $13\times10=130$。
\practiceheading 将长宽互换各运行一次，并与 cuDNN 非方形张量结果比较，可发现 width/height 或 leading dimension 写反。
\sourcecheck{https://github.com/NVIDIA/cuda-samples/tree/master/cpp/2_Concepts_and_Techniques/convolutionSeparable}{NVIDIA convolutionSeparable 示例核验二维零边界加载；矩形闭式以两维相等和一维退化为 1 的回代独立审查}
\end{solutionexercise}

\begin{solutionexercise}{7}
\problemheading 对大小为 $N\times N$ 的数组执行二维分块卷积，使用图 7.12 所示内核，
滤波器大小为 $M\times M$，输出瓦片大小为 $T\times T$：
\begin{enumerate}[label=\alph*.,leftmargin=2.2em]
  \item 需要多少个线程块？
  \item 每个块需要多少个线程？
  \item 每个块需要多少共享内存？
  \item 如果使用图 7.15 中的内核，重复回答上述问题。
\end{enumerate}
图 7.12 的块与完整输入瓦片同尺寸，输入瓦片含半径 halo，边长为 $T+M-1$；每个线程
加载一个输入元素，只有中央 $T\times T$ 个线程计算输出。图 7.15 的块、共享瓦片与
输出瓦片均为 $T\times T$，halo 从原输入经缓存路径读取。元素为单精度浮点数，
滤波器位于常量内存，不计入每块共享内存。

\approachheading 图 7.12 把完整 halo 放入共享内存，输入瓦片边长为 $T+M-1$；
图 7.15 依赖缓存读取 halo，共享瓦片与输出瓦片同为 $T$。

\answerheading 两种内核都需要
\[
B=\left\lceil\frac NT\right\rceil^2
\]
个线程块。

图 7.12 每块线程数为 $(T+M-1)^2$，共享内存为
$4(T+M-1)^2$ B（单精度）。图 7.15 每块线程数为 $T^2$，共享内存为 $4T^2$ B。
常量内存中的 $M^2$ 个滤波器系数是设备级资源，不重复计入每块共享内存。

必须另行验证每块线程数不超过设备上限；图 7.12 的输入瓦片更大，较容易超过上限。

\complexity{$O(N^2M^2)$ 总工作}{分别为 $O((T+M)^2)$ 与 $O(T^2)$ 共享空间/块}{$\lceil N/T\rceil^2$ 个块}
\conclusionheading 图 7.12：$B$ 块、$(T+M-1)^2$ 线程/块、
$4(T+M-1)^2$ B；图 7.15：$B$ 块、$T^2$ 线程/块、$4T^2$ B。
\discussionheading
\discussionpoint{完整输入瓦片与缓存 halo 选择不同的复用责任。}
两个内核为何生成相同输出，却需要不同线程数和共享内存？图 7.12 让块协作装入中央区及全部 halo，随后每个邻域读取都来自片上瓦片，复用位置和同步点明确；图 7.15 只把 $T\times T$ 中央区放入共享内存，halo 仍从原数组读取，并希望相邻块访问曾被缓存。前者用 $(T+M-1)^2$ 个装载位置换取可预测复用，后者用 $T^2$ 个线程和较小共享状态换取更高驻留可能性。正确性都不应依赖缓存命中，但性能责任落在不同存储层，因此选型必须结合资源限制和实测命中行为。

\discussionpoint{halo 的表面积成本随瓦片尺寸和半径变化。}
为什么增大 $T$ 往往提高理想复用，却不能无限增大？完整输入边长为 $T+M-1$，相对输出的装载比为 $((T+M-1)/T)^2$；$T$ 增大时该比值趋近 1，邻块重复的边界占比下降。与此同时，线程数、共享内存、寄存器活跃状态和边界块浪费都会增加，甚至先碰到每块线程上限。大半径或小图像中 halo 可能占瓦片大部分，此时每线程循环载入多个元素、分阶段搬运或使用矩形 tile 更合适；这与域分解中的表面积/体积权衡属于同一几何规律。最优点必须用合法配置实测。

\discussionpoint{缓存 halo 是性能假设而非语义前提。}
若 L1 或 L2 完全没有命中，图 7.15 是否还应算对？每次 halo 读取仍访问合法的全局地址，所以数值必须正确；变化的是同一边界数据会不会被多个块反复从 DRAM 取回。命中率受缓存容量、替换策略、只读路径、块调度顺序和其他流量干扰，源代码中“刚被邻块读取过”并不等于当前块一定命中。评估时要区分加载请求、L1/L2 命中、DRAM sector 和实际字节，并在目标问题尺寸与并发负载上测量。这样才能避免把某一架构或空闲机器上的偶然局部性写成普遍硬件结论。

\discussionpoint{生产卷积选择跨越多个算法层级。}
把两种 tile 中较快的一种固定下来，是否就得到通用卷积库？它们都属于直接卷积的局部实现；可分离核能拆成一维阶段，某些小核与通道形状适合 Winograd，长核可能适合 FFT，隐式 GEMM 又能复用成熟矩阵乘路径。不同方案需要不同 workspace、预处理、布局变换和数值路径，首次调用延迟也可能与稳态吞吐相反。生产库通常按形状、数据类型、确定性和内存预算调优并缓存选择；评价本题方案时也应同时观察 occupancy、共享吞吐、缓存命中和端到端时间，而不是用单一计数器决定胜负。
\variantheading $N=1024,M=5,T=16$ 时有 4096 块；两方案分别为 400/256 线程和 1600/1024 B 共享内存。
\practiceheading 用 Nsight Compute 比较 L2 hit rate、shared throughput 和 occupancy 后选型。
\sourcecheck{https://docs.nvidia.com/cuda/cuda-c-programming-guide/}{NVIDIA CUDA C++ Programming Guide 的二维网格和共享内存资源规则；halo 尺寸按原书两内核独立复核}
\end{solutionexercise}

\begin{solutionexercise}{8}
\problemheading 修改图 7.7 中的二维内核，使其执行三维卷积。图 7.7 的完整基线为：
\begin{lstlisting}[language=C++,firstnumber=1]
__global__ void convolution_2D_basic_kernel(float *N, float *F, float *P,
                                            int r, int width, int height) {
    int outCol = blockIdx.x*blockDim.x + threadIdx.x;
    int outRow = blockIdx.y*blockDim.y + threadIdx.y;
    float Pvalue = 0.0f;
    for (int fRow = 0; fRow < 2*r+1; fRow++) {
        for (int fCol = 0; fCol < 2*r+1; fCol++) {
            int inRow = outRow - r + fRow;
            int inCol = outCol - r + fCol;
            if (inRow >= 0 && inRow < height &&
                inCol >= 0 && inCol < width) {
                Pvalue += F[fRow*(2*r+1) + fCol]*
                    N[inRow*width + inCol];
            }
        }
    }
    if (outRow < height && outCol < width) {
        P[outRow*width + outCol] = Pvalue;
    }
}
\end{lstlisting}
二维数组均按行主序线性化，ghost cell 取 0；三维版本也应明确三维线性索引、
输出边界和输入邻域边界。

\approachheading 将线程映射扩展到 $(x,y,z)$，把滤波器和输入按行主序三维线性化，
并在每次乘加前同时检查三维边界。

\editorialnote 英文底本图 7.7 把一维指针误写成二维下标且缺少输出边界保护；中文译稿
已经改为线性索引并补上保护。以下代码以译稿的显式修正版为基线，再扩展到三维。

\answerheading
\begin{lstlisting}[language=C++]
__global__ void convolution_3D_basic(const float* N, const float* F,
                                     float* P, int r,
                                     int width, int height, int depth) {
    int x = blockIdx.x * blockDim.x + threadIdx.x;
    int y = blockIdx.y * blockDim.y + threadIdx.y;
    int z = blockIdx.z * blockDim.z + threadIdx.z;
    if (x >= width || y >= height || z >= depth) return;
    int K = 2 * r + 1;
    float sum = 0.0f;
    for (int fz = 0; fz < K; ++fz)
      for (int fy = 0; fy < K; ++fy)
        for (int fx = 0; fx < K; ++fx) {
          int ix = x - r + fx, iy = y - r + fy, iz = z - r + fz;
          if (ix >= 0 && ix < width && iy >= 0 && iy < height &&
              iz >= 0 && iz < depth) {
            int ni = (iz * height + iy) * width + ix;
            int fi = (fz * K + fy) * K + fx;
            sum += F[fi] * N[ni];
          }
        }
    P[(z * height + y) * width + x] = sum;
}
\end{lstlisting}
可用如 \code{dim3 block(8,8,4)}，网格三维分别向上取整。无跨线程共享数据，故无需
屏障且没有写写竞态；主机负责在内核结束前保持三个设备数组有效，并检查启动/同步错误。

\complexity{$O(WHD(2r+1)^3)$}{$O(1)$ 每线程额外空间}{$WHD$ 个输出体素可并行}
\conclusionheading 三重滤波循环、三维边界检查和线性索引共同构成完整三维基本内核。
\discussionheading
\discussionpoint{三维推广同时放大工作量和索引风险。}
二维循环多加一层是否只是机械改写？半径为 $r$ 时，每个体素要访问 $(2r+1)^3$ 个邻点，核边长翻倍会使算术和系数容量近似增长八倍；输出并行度则由体积 $WHD$ 决定，两种增长并不相同。行主序线性地址应为 $(zH+y)W+x$，若把 $H$、$W$ 或 slice stride 写反，访问仍可能落在已分配范围内，产生形状平滑却轴向错置的隐蔽结果。推广高维代码时应先写明每轴含义、跨度与边界，再用非立方尺寸测试，避免方形或立方样例掩盖置换错误。非对称脉冲核最适合发现此类轴置换。

\discussionpoint{三维复用受面、棱和角共同约束。}
相邻体素的邻域高度重叠，为什么基本内核仍可能消耗大量带宽？每个线程独立加载整个立方邻域时，只能寄希望于缓存捕获跨线程重复；共享瓦片可显式只装一次，但输入边长比输出多 $2r$，halo 不仅有六个面，还包含棱和角，其相对体积在小 tile 中很大。立方块的线程数又按边长三次方增长，常在共享内存耗尽前就碰到线程上限。实际实现常用非立方 tile、每线程多载或分平面滑动窗口，以连续 x 访问为核心控制工作集；选择应由字节模型和目标设备资源共同验证。

\discussionpoint{核结构决定直接法之外的算法空间。}
面对 $M^3$ 个系数，是否只能继续扩大三重循环？若三维核是三个一维核的外积，可按 x、y、z 做三次一维卷积，把每输出算术从 $O(M^3)$ 降为 $O(3M)$，代价是中间缓冲和额外读写；低秩近似还能用若干可分离项交换误差。大而稠密的周期或零填充问题可考虑 3D FFT，稀疏十字或模板核则更适合偏移列表。算法选择会改变边界处理、舍入次序、启动次数和工作区，因此应在相同输出语义与误差容限下比较，而不能只用渐近 FLOP 判断。转换成本也要计入小体积问题。

\discussionpoint{领域验证必须把体素间距和轴方向纳入参考。}
医学体数据或流体网格中的三个数组下标，是否代表相同物理距离？CT、MRI 和地震体常具有各向异性 spacing，同样的三格半径在不同轴上覆盖不同长度，导数或平滑核必须按物理单位缩放；方向元数据若丢失，数值可能正确却对应错误解剖或坐标轴。验证可先令 $D=1$ 退化为已验证的二维算子，再用单位脉冲检查三轴核方向，用非对称尺寸和 CPU 双精度参考逐体素比较。性能上再结合 Roofline 判断主要受邻域字节还是乘加限制，把物理正确性与硬件调优分成两条证据链。
\variantheading 半径 $r=2$ 时每个内部输出访问 $5^3=125$ 个系数和输入，执行 125 次乘加。
\practiceheading 以 cuDNN 三维卷积或 CPU reference 验证，并用 Nsight Compute Roofline 定位瓶颈。
\sourcecheck{https://docs.nvidia.com/cuda/cuda-c-programming-guide/}{NVIDIA CUDA C++ Programming Guide 的三维线程索引与全局内存规则；以 $D=1$ 退化回二维内核作对抗审查}
\end{solutionexercise}

\begin{solutionexercise}{9}
\problemheading 修改图 7.9 中的二维内核，使其执行三维卷积。图 7.9 假设半径为编译期
常量 \code{FILTER_RADIUS}，并把滤波器放在设备常量内存；完整二维基线为：
\begin{lstlisting}[language=C++,firstnumber=1]
__constant__ float F[2*FILTER_RADIUS+1][2*FILTER_RADIUS+1];
__global__ void convolution_2D_const_mem_kernel(float *N, float *P, int r,
                                                int width, int height) {
    int outCol = blockIdx.x*blockDim.x + threadIdx.x;
    int outRow = blockIdx.y*blockDim.y + threadIdx.y;
    float Pvalue = 0.0f;
    for (int fRow = 0; fRow < 2*r+1; fRow++) {
        for (int fCol = 0; fCol < 2*r+1; fCol++) {
            int inRow = outRow - r + fRow;
            int inCol = outCol - r + fCol;
            if (inRow >= 0 && inRow < height &&
                inCol >= 0 && inCol < width) {
                Pvalue += F[fRow][fCol]*N[inRow*width + inCol];
            }
        }
    }
    if (outRow < height && outCol < width) {
        P[outRow*width + outCol] = Pvalue;
    }
}
\end{lstlisting}
三维版本应相应扩展常量滤波器、网格映射和行主序线性索引，并保持零 ghost cell 约定。

\approachheading 半径在编译期固定，把 $K^3$ 个系数存入设备常量符号；线程和边界逻辑
沿用上一题，但不再传入滤波器指针。

\editorialnote 英文底本图 7.9 同样含一维指针的二维下标与缺失输出边界问题；中文译稿
已经修正。以下三维版本保留这些必要修正，不把底本错误静默带入答案。

\answerheading 以下示例取半径 1；修改 \code{R} 后重新编译即可改变滤波器尺寸。
\begin{lstlisting}[language=C++]
constexpr int R = 1;
constexpr int K = 2 * R + 1;
__constant__ float F3[K * K * K];

__global__ void convolution_3D_const(const float* N, float* P,
                                     int width, int height, int depth) {
    int x = blockIdx.x * blockDim.x + threadIdx.x;
    int y = blockIdx.y * blockDim.y + threadIdx.y;
    int z = blockIdx.z * blockDim.z + threadIdx.z;
    if (x >= width || y >= height || z >= depth) return;
    float sum = 0.0f;
    for (int fz = 0; fz < K; ++fz)
      for (int fy = 0; fy < K; ++fy)
        for (int fx = 0; fx < K; ++fx) {
          int ix = x - R + fx, iy = y - R + fy, iz = z - R + fz;
          if (ix >= 0 && ix < width && iy >= 0 && iy < height &&
              iz >= 0 && iz < depth)
            sum += F3[(fz * K + fy) * K + fx] *
                   N[(iz * height + iy) * width + ix];
        }
    P[(z * height + y) * width + x] = sum;
}
// Before launch:
// cudaMemcpyToSymbol(F3, host_filter, K*K*K*sizeof(float));
\end{lstlisting}
同一 warp 在固定循环迭代读取同一系数，适合常量缓存广播。覆盖 \code{F3} 前必须确保
仍使用旧系数的内核已经完成，尤其是跨 stream 使用同一符号时。

\complexity{$O(WHDK^3)$}{$K^3$ 个设备级常量元素，$O(1)$ 每线程私有空间}{$WHD$ 个输出可并行}
\conclusionheading 三维常量内核减少滤波器参数流量，并完整处理输出/输入边界和符号
生命周期。
\discussionheading
\discussionpoint{常量缓存擅长 warp 内同地址广播。}
把滤波器放进常量内存，为何卷积比随机查表更容易受益？循环的某次迭代中，同一 warp 的线程通常都读取同一个 $F[f_z,f_y,f_x]$，常量路径可把一个缓存值广播给所有 lane，系数又被许多块反复复用。若线程使用不同半径、控制流失步，或每个样本读取不同核，同一 warp 会请求多个地址，这些不同地址可能需要分批服务，广播优势随之减弱。优化前应确认实际访问的一致性，并结合常量请求和序列化指标测量；“只读且很小”是必要背景，却不自动等于常量内存最佳。

\discussionpoint{常量符号的容量和版本具有设备级生命周期。}
两个 stream 能否同时把不同滤波器写入同一个 \code{F3} 再各自启动内核？常量符号是设备上的共享状态，不是每个 stream 或每次调用的私有副本；若更新与读取没有事件或 stream 顺序，内核可能看到错误版本，问题属于数据竞争而非缓存失效。符号尺寸通常在编译期确定，复制量和核大小还必须受容量约束，旧版本在最后一个使用者完成前不能覆盖。服务多个并发请求时，可为版本分配不同符号，或改用每次启动传入的只读缓冲，从而以指针和缓存行为换取清晰所有权。

\discussionpoint{系数存储与输入复用是两个独立优化维度。}
常量内存让核系数很便宜后，内核是否就不再受访存限制？每个输出仍需读取大量重叠的输入邻域，系数流量下降并不会自动消除这些输入字节；共享瓦片、只读缓存和滑动窗口解决的是另一条数据流。小核还可作为内核参数传递，或由块协作载入共享内存，再由编译器把热点值保存在寄存器中，选择取决于大小、地址一致性、更新频率和资源占用。建立模型时应分别计算系数与输入的请求、复用范围和生命周期，再用计数器判断真正瓶颈，避免一个成功优化掩盖另一侧流量。两类流量的最优存储层往往不同。

\discussionpoint{复用周期决定常量上传是否值得。}
同一个小核连续处理上千帧时，一次符号上传可被大量内核摊薄，常量广播通常很合适；若每个样本都有动态权重，频繁 \code{cudaMemcpyToSymbol} 可能串行化流水，并使共享符号难以支持并发请求。生产系统还要处理模型热更新：新权重必须在旧请求完成后发布，或者采用双缓冲和事件建立版本切换。测量应把上传、内核和同步时间一并纳入，而不只看稳态 kernel duration，并用多 stream 压力测试核对结果版本。这个问题与配置缓存和只读快照相似，性能建立在明确的发布与回收协议之上。
\variantheading 半径改为 2 后常量数组含 125 个 float，即 500 B。
\practiceheading 跨 stream 更新 \code{F3} 时用事件建立依赖，并观察 constant-memory 请求效率。
\sourcecheck{https://docs.nvidia.com/cuda/cuda-c-programming-guide/}{NVIDIA CUDA C++ Programming Guide 的常量内存、常量符号复制 API 与三维网格语义；置信度高}
\end{solutionexercise}

\begin{solutionexercise}{10}
\problemheading 修改图 7.12 中的二维分块内核，使其执行三维卷积。图 7.12 使用与完整
输入瓦片同尺寸的线程块，完整二维基线为：
\begin{lstlisting}[language=C++,firstnumber=1]
#define IN_TILE_DIM 32
#define OUT_TILE_DIM ((IN_TILE_DIM) - 2*(FILTER_RADIUS))
__constant__ float F[2*FILTER_RADIUS+1][2*FILTER_RADIUS+1];
__global__ void convolution_tiled_2D_const_mem_kernel(float *N, float *P,
                                                       int width, int height) {
    int col = blockIdx.x*OUT_TILE_DIM + threadIdx.x - FILTER_RADIUS;
    int row = blockIdx.y*OUT_TILE_DIM + threadIdx.y - FILTER_RADIUS;
    __shared__ float N_s[IN_TILE_DIM][IN_TILE_DIM];
    if(row>=0 && row<height && col>=0 && col<width)
        N_s[threadIdx.y][threadIdx.x] = N[row*width + col];
    else
        N_s[threadIdx.y][threadIdx.x] = 0.0f;
    __syncthreads();
    int tileCol = threadIdx.x - FILTER_RADIUS;
    int tileRow = threadIdx.y - FILTER_RADIUS;
    if (col >= 0 && col < width && row >= 0 && row < height &&
        tileCol >= 0 && tileCol < OUT_TILE_DIM &&
        tileRow >= 0 && tileRow < OUT_TILE_DIM) {
        float Pvalue = 0.0f;
        for (int fRow = 0; fRow < 2*FILTER_RADIUS+1; fRow++)
          for (int fCol = 0; fCol < 2*FILTER_RADIUS+1; fCol++)
            Pvalue += F[fRow][fCol]*N_s[tileRow+fRow][tileCol+fCol];
        P[row*width+col] = Pvalue;
    }
}
\end{lstlisting}
三维版本应扩展输入/输出瓦片、常量滤波器、行主序索引、边界和同步。

\approachheading 使用输入立方瓦片大小 \code{IN3} 和输出边长
\code{OUT3=IN3-2*R3}；每线程加载一个输入/ghost 元素，块内同步后仅内部线程计算。

\editorialnote 英文底本图 7.12 的共享数组声明缺少元素类型，中文译稿已补为
\code{float}。以下实现沿用该最小修正，并同时显式处理三维边界。

\answerheading 下面给出半径 1、$8^3=512$ 线程块的关键完整实现。
\begin{lstlisting}[language=C++]
constexpr int R3 = 1, IN3 = 8, OUT3 = IN3 - 2 * R3;
__constant__ float F3_t[(2*R3+1)*(2*R3+1)*(2*R3+1)];

__global__ void convolution_3D_tiled(const float* N, float* P,
                                     int width, int height, int depth) {
    int tx=threadIdx.x, ty=threadIdx.y, tz=threadIdx.z;
    int x=blockIdx.x*OUT3 + tx-R3;
    int y=blockIdx.y*OUT3 + ty-R3;
    int z=blockIdx.z*OUT3 + tz-R3;
    __shared__ float tile[IN3][IN3][IN3];
    tile[tz][ty][tx] = (x>=0 && x<width && y>=0 && y<height &&
                         z>=0 && z<depth)
                       ? N[(z*height+y)*width+x] : 0.0f;
    __syncthreads();

    int lx=tx-R3, ly=ty-R3, lz=tz-R3;
    if (lx>=0 && lx<OUT3 && ly>=0 && ly<OUT3 && lz>=0 && lz<OUT3 &&
        x>=0 && x<width && y>=0 && y<height && z>=0 && z<depth) {
      constexpr int K3 = 2*R3+1;
      float sum=0.0f;
      for (int fz=0; fz<K3; ++fz)
        for (int fy=0; fy<K3; ++fy)
          for (int fx=0; fx<K3; ++fx)
            sum += F3_t[(fz*K3+fy)*K3+fx] *
                   tile[lz+fz][ly+fy][lx+fx];
      P[(z*height+y)*width+x]=sum;
    }
}
// block = dim3(IN3,IN3,IN3); grid dimensions use ceil(size/OUT3).
\end{lstlisting}
所有线程先加载（越界者写 0）并无条件到达屏障，避免死锁；每个活跃线程写唯一输出。
一般配置必须满足 $\code{IN3}^3$ 不超过每块线程上限，并核对共享内存容量。

\complexity{$O(WHD(2R3+1)^3)$}{$4\,IN3^3$ B 共享内存/块}{$OUT3^3$ 个输出/块}
\conclusionheading 这是三维的“输入瓦片大小等于块大小”实现，边界、同步和资源限制均
已显式处理。
\discussionheading
\discussionpoint{共享输入瓦片只有内部体积能够产生输出。}
为什么 $IN3^3$ 个线程中只有 $OUT3^3$ 个负责写结果？半径为 $R3$ 的输出需要左右各一层厚度 $R3$ 的邻域，故共享输入边长 $IN3$ 只能支撑 $OUT3=IN3-2R3$ 的无 halo 内部，活跃计算比例为 $(OUT3/IN3)^3$。三次方会放大薄边界的代价，例如 $IN3=8,R3=1$ 时只有 $6^3/8^3\approx42.2\%$ 的线程产生输出，其余线程仍对装载有贡献。这个比例不是单纯“浪费率”，而是显式复用的成本指标，应同每个输入被多少输出复用、块驻留和全局字节一起评价。

\discussionpoint{每块线程上限在三维中比共享容量更早成为约束。}
把瓦片边长从 8 增至 10 看似只增加两格，为什么配置已接近极限？一个线程对应一个输入位置时，线程数按 $IN3^3$ 增长，$10^3=1000$ 已逼近常见的每块 1024 线程上限；此时单精度共享数组只有 4000 B，可能远未耗尽共享容量。继续增大复用范围不能靠扩大立方块，而应让较少线程用循环载入多个元素，或让每线程计算多个输出。所有具体上限仍需查询目标设备属性；这一例子说明多维 tile 的首要资源可能随映射方式改变，不能只按共享字节选尺寸。

\discussionpoint{边界线程也必须参加屏障前的协作装载。}
体数据边缘的线程最终不写输出，能否一进入内核就返回？它们仍对应共享瓦片中的输入或零填充位置，且块级 \code{__syncthreads()} 要求参与块的线程以一致控制流到达；提前返回会让其余线程等待，或者让共享位置未初始化。正确顺序是所有线程先写入域内值或 0，共同同步，再仅让内部且对应合法全局坐标的线程计算。若使用双缓冲或多阶段流水，还需在覆盖共享区前确认上一批消费者完成。三维边界有面、棱和角，使用非立方尺寸及极小深度测试能更容易暴露遗漏条件。

\discussionpoint{非立方 tile 可以匹配布局、核形状与子域几何。}
物理问题是三维的，线程块是否也必须是立方体？行主序数据让 x 方向连续，扩大 x、适度保留 y 并缩短 z，常能控制 slice 工作集和线程总数；各向异性核或扁平子域也可能在某一轴需要更宽 halo。形状变化会同时影响合并访问、共享 bank 行为、halo 比、寄存器和边界块利用率，不存在只按体积排序的通用最优解。可先用设备属性与 Occupancy API 排除非法配置，再以有效输出/启动线程、DRAM 字节和实际吞吐搜索；医学体卷积与模板 PDE 即使尺寸相同，也可能因核和精度目标选择不同形状。
\variantheading 对 $R3=1,IN3=10$，OUT3=8、线程数 1000、输出数 512、共享内存 4000 B。
\practiceheading 用 Occupancy API 筛除非法配置，并以 Compute Sanitizer synccheck 验证屏障。
\sourcecheck{https://docs.nvidia.com/cuda/cuda-c-programming-guide/}{NVIDIA CUDA C++ Programming Guide 的三维线程块、共享内存与块级屏障规则；边界块全线程到达屏障复核}
\end{solutionexercise}

\begin{solutionexercise}{11}
\problemheading 修改图 7.12 中的分块二维内核，使每个线程块使用与输出瓦片
大小相同数量的线程，并让每个线程在把输入瓦片加载到共享内存时加载多个元素。

基线中输入瓦片边长为 \code{IN_TILE_DIM}，输出瓦片边长满足
\code{OUT_TILE_DIM} $=$ \code{IN_TILE_DIM} $-2\times$ \code{FILTER_RADIUS}；原实现启动
\code{IN_TILE_DIM}$^2$ 个线程、每线程加载一个输入元素，再只让中央
\code{OUT_TILE_DIM}$^2$ 个线程计算。本题要求改为只启动
\code{OUT_TILE_DIM}$^2$ 个线程，以循环协作加载全部 \code{IN_TILE_DIM}$^2$ 个
共享元素；必须覆盖零填充边界、装载后的块级屏障及输出写回边界。

\approachheading 用 $T^2$ 个线程线性遍历 $(T+2R)^2$ 个共享元素，步长为块内线程
总数；同步后，每线程直接计算一个输出。

\editorialnote 本题继续以译稿已修正的图 7.12 为基线；底本共享数组缺少元素类型，
不能按原样通过 CUDA C++ 编译。

\answerheading
\begin{lstlisting}[language=C++]
constexpr int R2 = 2, OUT2 = 16, IN2 = OUT2 + 2*R2;
__constant__ float F2[(2*R2+1)*(2*R2+1)];

__global__ void convolution_2D_output_threads(const float* N, float* P,
                                               int width, int height) {
    __shared__ float tile[IN2][IN2];
    int tx=threadIdx.x, ty=threadIdx.y;
    int lane=ty*OUT2+tx, workers=OUT2*OUT2;
    for (int q=lane; q<IN2*IN2; q+=workers) {
        int sy=q/IN2, sx=q%IN2;
        int x=blockIdx.x*OUT2 + sx-R2;
        int y=blockIdx.y*OUT2 + sy-R2;
        tile[sy][sx]=(x>=0 && x<width && y>=0 && y<height)
                    ? N[y*width+x] : 0.0f;
    }
    __syncthreads();

    int x=blockIdx.x*OUT2+tx, y=blockIdx.y*OUT2+ty;
    if (x<width && y<height) {
      constexpr int K2=2*R2+1;
      float sum=0.0f;
      for (int fy=0; fy<K2; ++fy)
        for (int fx=0; fx<K2; ++fx)
          sum += F2[fy*K2+fx]*tile[ty+fy][tx+fx];
      P[y*width+x]=sum;
    }
}
// block = dim3(OUT2,OUT2); grid = ceil(width/OUT2) x ceil(height/OUT2).
\end{lstlisting}
循环覆盖每个共享位置恰好一次；越界输入显式写 0。边界输出线程也先参与加载和屏障，
之后才被条件屏蔽，因此不会造成条件屏障死锁或未初始化共享元素。

\complexity{$O(W H(2R2+1)^2)$}{$4\,IN2^2$ B 共享内存/块}{$OUT2^2$ 个输出线程/块}
\conclusionheading 线程块恰有 $T^2$ 个线程，并通过步长循环完整加载
$(T+2R)^2$ 个输入瓦片元素。
\discussionheading
\discussionpoint{协作载入把线程所有权与数据数量解耦。}
共享瓦片比线程块大时，是否必须再创建一圈只负责 halo 的线程？把块内二维坐标压成线性编号 $lane$，令每个线程访问 $q=lane+kT^2$，便能用恰好 $T^2$ 个输出线程覆盖 $(T+2R)^2$ 个输入位置，域外位置照常写 0。线程数由输出并行度决定，待载数据量则由循环次数承担，这种解耦使较大 halo 不必直接推高块线程数。GEMM、注意力和 stencil 也常用同一模式协作搬运 tile；真正的共同前提是所有生产者完成后再同步，消费者才可读取共享状态。

\discussionpoint{模与整除关系给出无遗漏、无重复的覆盖证明。}
循环步长取 $T^2$，为何每个共享位置恰好由一个线程写？任意 $0\le q<(T+2R)^2$ 都有唯一分解 $q=(q\bmod T^2)+\lfloor q/T^2\rfloor T^2$，前一项指定线程，后一项指定迭代轮次，因此两个线程不可能得到同一 $q$，也没有索引被跳过。再用 $sy=\lfloor q/IN2\rfloor$、$sx=q\bmod IN2$ 还原坐标并做全局边界检查。这个证明只保证覆盖，不保证访问连续或高效；屏障仍须放在整个循环之后并由全块到达，才能把所有写入发布给输出计算。覆盖证明与同步可见性契约缺一不可，二者应分别测试。

\discussionpoint{尾轮负载与向量载入需要单独处理。}
当 $(T+2R)^2$ 不能整除 $T^2$ 时，一部分线程会不会长期落后？静态步长分配使任意两线程的载入数相差至多一项，通常负载很均衡，但最后一轮只有部分 lane 活动，且线性 tile 跨行处可能破坏理想向量对齐。可以按 warp 分配整行、使用自然对齐的宽加载，或在支持的架构上采用异步复制和双缓冲；每种方案都必须处理行尾、域边界与等待组完成。宽指令减少指令数不等于减少逻辑字节，优化应以实际事务和流水重叠证明，而不能从类型名推断收益。对齐契约应贯穿基址与每一行。

\discussionpoint{正确覆盖之后还要验证同步成本与真实复用。}
输出线程比例提高，为什么新内核仍可能输给“一个线程装一个输入”的版本？每线程多轮载入增加地址计算，尾轮部分活动造成 lane 空闲，块级屏障和共享访问也可能成为瓶颈；另一方面，更少线程可能改善驻留或减少无输出线程，收益取决于 $T$、$R$ 和缓存层次。调试时可先把共享数组填入哨兵，或在测试构建中为每格维护写入计数，证明恰写一次且消费前已同步。性能阶段再检查 barrier stall、global sector、shared throughput 与 occupancy，并把边界块纳入数据集，避免内部样例掩盖不规则成本。
\variantheading $T=8,R=1$ 时 64 个线程加载 100 个元素，其中前 36 个线程各多加载一次。
\practiceheading 用共享内存哨兵或 Compute Sanitizer 验证每个 tile 位置恰写一次，并测量 barrier stall。
\sourcecheck{https://docs.nvidia.com/cuda/cuda-c-programming-guide/}{NVIDIA CUDA C++ Programming Guide 的协作式共享内存加载和屏障规则；线性循环覆盖性与边界控制流独立复核}
\end{solutionexercise}

\chapter{模板计算与线程粗化}

\begin{solutionexercise}{1}
\problemheading 考虑在包含边界单元的 $120\times120\times120$ 网格上执行三维模板计算。
\begin{enumerate}[label=\alph*.,leftmargin=2.2em]
  \item 每次模板扫掠会计算多少个输出网格点？
  \item 对图 8.6 的基本内核，假设线程块大小为 $8\times8\times8$，
  需要多少个线程块？
  \item 对图 8.8 的共享内存分块内核，假设线程块大小为
  $8\times8\times8$，需要多少个线程块？
  \item 对图 8.10 同时采用共享内存分块和线程粗化的内核，假设线程块
  大小为 $32\times32$，需要多少个线程块？
\end{enumerate}
沿用本章一阶三维七点模板：只更新每维下标 $1\ldots N-2$ 的内部点。图 8.6 的
$8^3$ 块直接覆盖全网格；图 8.8 的 $8^3$ 输入瓦片产生 $6^3$ 输出瓦片；图 8.10
的 $32\times32$ 输入平面产生 $30\times30$ 输出平面，并在 $z$ 方向连续粗化
处理 30 个输出平面。

\approachheading 一阶七点模板只更新每维去掉两侧各一层后的内部点。基本内核按 8
覆盖全网格；共享内存内核的 $8^3$ 输入块产生 $6^3$ 输出；$32^2$ 粗化块每维有效
输出边长为 30，并沿 $z$ 连续处理 30 个平面。

\answerheading
\begin{enumerate}[label=\alph*.,leftmargin=2.2em]
\item 输出点数为
\[
(120-2)^3=118^3=1{,}643{,}032\ \text{个网格点}.
\]
\item 基本内核每维需 $\lceil120/8\rceil=15$ 个块，共
$15^3=3375$ 个块。
\item 原书图 8.8 中输入瓦片边长 8、输出瓦片边长 $8-2=6$，每维需
$\lceil118/6\rceil=20$ 个块，共 $20^3=8000$ 个块。
\item 原书图 8.10 中 $32\times32$ 输入平面对应 $30\times30$ 输出平面，并沿
$z$ 粗化 30 层。每维需 $\lceil118/30\rceil=4$ 个块，共 $4^3=64$ 个块。
\end{enumerate}
边界块中的越界或边界线程由内核条件屏蔽，不改变启动块数。

\complexity{$O(118^3)$ 有效工作}{基本内核 $O(1)$；分块内核使用片上瓦片}{输出点在三维块网格中并行}
\conclusionheading (a) 1,643,032；(b) 3375；(c) 8000；(d) 64。
\discussionheading
\discussionpoint{七点模板把离散物理写成局部数据依赖。}
中心点与六个轴向邻点的加权和看似一种小卷积，但这些系数通常来自偏微分方程的有限差分，例如三维扩散或 Poisson 方程，而不是任意图像滤波器。固定边界单元不更新，等价于把边界值作为 Dirichlet 条件；若改成周期边界，越过左面会从右面取值，Neumann 条件则可能通过镜像或 ghost 值表达导数。边界规则改变后，有效输出点数、邻域索引和数值解的物理含义都会改变。模板优化因此不能只验证 GPU 与某段代码一致，还应检查离散方程的守恒、收敛或边界不变量。
\discussionpoint{启动线程、加载线程和输出线程是三种角色。}
基本内核可为整个 $120^3$ 体积启动线程，再让边界线程跳过更新；分块内核还可能让外圈线程只负责把 halo 载入共享内存，而不生成输出。于是“启动了多少线程”不能直接换算成“计算了多少网格点”，更不能单独判断效率。一个 tile 的额外加载如果让许多相邻输出复用同一全局值，可能比每线程都计算却重复读取更划算；反之，小 tile 的 halo 面积占比很高，协作开销可能吞掉收益。应分别统计有效输出、边界屏蔽、协作加载和相邻块重复 halo，才能把线程冗余与数据复用放在同一账本中比较。
\discussionpoint{时间分块把复用从空间推进到迭代维。}
共享内存空间分块只复用当前时间步的邻点，而 Jacobi 等迭代会连续读取下一时刻刚生成的网格；若在片上推进多个时间步，就能避免每一步都把整个体积写回并重新读取。代价是依赖锥随时间扩张：计算 $k$ 步后的内部点需要初始 tile 周围厚度为 $k$ 的 halo，每推进一步可安全输出的内部区域都会收缩。还要保存多个时间层或采用交替缓冲，并确保同步不会混用新旧状态。时间分块因此提升算术强度的同时加剧容量、寄存器和边界复杂度，它是依赖图变换而不只是多循环几次。
\discussionpoint{多 GPU 模板受表面积与体积比例支配。}
把三维域切成子域后，每个 GPU 主要计算自己的体积，却在每轮交换六个面上的 halo；计算量随边长立方增长，通信量近似随边长平方增长，所以较大的子域更容易摊薄通信。子域过小或切分方向不合适时，网络延迟和边界复制会压过单 GPU 内核优化，即使局部 tile 的带宽表现很好。工程实现常把内部点计算与 halo 传输重叠，再在数据到达后处理边界层，但这要求明确 stream、通信库和设备拓扑依赖。评价结果应同时报告每秒网格更新、通信隐藏比例与数值收敛一致性，才能连接局部内核吞吐和整体求解时间。
\variantheading 若网格为 $62^3$，内部点为 $60^3=216000$；$8^3$ 基本块需 $8^3=512$ 块。
\practiceheading 用 Nsight Compute 比较三内核的 DRAM Bytes、active threads 和 achieved occupancy。
\sourcecheck{https://docs.nvidia.com/cuda/cuda-c-programming-guide/}{NVIDIA CUDA C++ Programming Guide 的三维网格向上取整规则；并按原书图 8.6、8.8、8.10 的索引独立复核}
\end{solutionexercise}

\begin{solutionexercise}{2}
\problemheading 考虑一个同时采用共享内存分块和线程粗化的三维七点模板实现。该实现与图
8.10 和图 8.12 类似，但分块不是规则立方体；它采用
$32\times32$ 的线程块，并把粗化因子设为 16，也就是说，每个线程块沿 $z$ 维
处理连续的 16 个输出平面。
\begin{enumerate}[label=\alph*.,leftmargin=2.2em]
  \item 在线程块的整个生命周期中，它所加载的输入分块大小是多少（以元素个数
  计）？
  \item 在线程块的整个生命周期中，它所处理的输出分块大小是多少（以元素个数
  计）？
  \item 该内核的算术强度是多少（以 FLOP/B 计）？
  \item 如果不使用寄存器分块，如图 8.10 所示，每个线程块需要多少
  字节共享内存？
  \item 如果使用寄存器分块，如图 8.12 所示，每个线程块需要多少
  字节共享内存？
\end{enumerate}
每个网格元素为 4 B 单精度浮点数；七点模板每个输出执行 7 次乘法和 6 次加法。
输入在 $x,y$ 两维各含一层 halo，在 $z$ 向的 16 个输出平面前后也各含一个输入平面。
算术强度按每个输入元素从全局内存加载一次、每个输出元素写回一次计，忽略系数流量
与边界块浪费。图 8.10 同时在共享内存保存 prev/curr/next 三个平面，图 8.12 只在
共享内存保存 curr 平面而把另外两个平面保存在寄存器。

\approachheading 一阶模板在 $x,y$ 两维各保留一层 halo，所以每个输入平面为
$32\times32$、输出平面为 $30\times30$；16 个输出平面在 $z$ 向还需前后各一个
输入平面。

\answerheading
\begin{enumerate}[label=\alph*.,leftmargin=2.2em]
\item 生命周期中加载的输入瓦片为 $32\times32\times(16+2)$，共
\[
32^2\times18=18{,}432\ \text{个元素}.
\]
\item 输出瓦片为 $30\times30\times16$，共
$14{,}400$ 个元素。
\item 每个输出执行 13 FLOP（7 乘、6 加）。全局流量按每个输入加载一次、每个输出
存储一次计，因此
\[
I=\frac{13\times14{,}400}
        {4(18{,}432+14{,}400)}
 =\frac{325}{228}\approx1.42544\ \mathrm{FLOP/B}.
\]
\item 不使用寄存器分块时同时保存 prev/curr/next 三个输入平面，共
$3\times32^2\times4=12{,}288$ B，即 12 KiB 共享内存。
\item 使用寄存器分块后，prev/next 留在线程寄存器，只需共享当前平面，共
$32^2\times4=4096$ B，即 4 KiB。
\end{enumerate}
算术强度忽略系数流量、缓存重复和边界块浪费，正是题目采用的理想化模型。

\complexity{$13\times14{,}400$ FLOP/块}{12 KiB 或 4 KiB 共享内存/块}{1024 个线程协作处理 16 个输出平面}
\conclusionheading (a) 18,432；(b) 14,400；(c) 约 1.42544 FLOP/B；
(d) 12,288 B；(e) 4096 B。
\discussionheading
\discussionpoint{沿 z 粗化形成一个寄存器滑动窗口。}
若同一线程连续计算多个 z 位置，当前输出所需的“前、当前、后”三个中心值会在下一输出中平移角色：旧当前成为新前，旧后成为新当前，只需再载入一个新后值。把这三个值留在寄存器中，就不必为每个 z 输出重新经过共享内存，类似 CPU 模板代码中的滑动窗口或流式管线。复用成立的前提是线程始终沿同一 x、y 坐标前进，并严格区分尚未进入、正在使用和已经离开的平面。这个变换把数据局部性从块内线程协作提升为线程私有时间序列，也解释了粗化为何能减少某些片上访问。z 边界的 ghost 值还必须在窗口推进前定义，不能让首尾迭代读取未初始化寄存器。
\discussionpoint{不同方向的邻点适合不同通信范围。}
x、y 邻点属于同一输出平面，一个线程始终可以按自己的邻点下标直接从全局内存重载它们；缓存和合并访问可能使这条朴素路径已经足够有效。只有当设计意图是复用相邻线程已经装入寄存器的值、避免再次访问全局内存时，才需要共享内存，或在活动掩码与线程映射允许时使用 warp shuffle 等线程交换机制。z 邻点来自同一线程即将连续处理的位置，因此更适合保存在寄存器滑动窗口中。优化方案可把当前 x/y 平面放进共享内存、把前后中心列留在线程私有状态，但这是一种复用选择而非正确性的强制条件；共享容量下降还会与滚动变量延长寄存器活跃区间相互制约。应依据重载流量、线程作用域和实际缓存行为选择存储层次。
\discussionpoint{理想算术强度会被粗化尾部和物理流量修正。}
增加 z 粗化因子可以让上下 halo 平面由更多内部输出摊薄，所以按“每个输入只从 DRAM 载入一次、每个输出写一次”的模型，单位输出字节数逐渐下降。然而线程数量也随每线程输出数增加而减少，最后一个 z 分段若不足粗化长度，会出现闲置循环与边界谓词。缓存行过取、未合并的跨平面地址、寄存器 spill 以及相邻块重复 halo 都会让实测 DRAM 字节偏离公式。Roofline 上的理论点只是上界假设，只有把实际传输量代回算术强度，才能判断优化真正移动了数据屋顶还是仅改变了源码计数。
\discussionpoint{调优应搜索形状组合并守住数值更新语义。}
x/y tile、z 粗化长度、共享平面数量和是否双缓冲会共同决定线程数、共享容量、寄存器及同步次数，单独扫描一个参数容易错过资源间的阶梯变化。可以先用编译器资源报告排除发生 spill 的组合，再在目标尺寸上比较实际流量和时间，并用 CPU 参考覆盖非整除边界。若进一步跨多个迭代做时间分块，依赖锥会扩大，内部有效区域收缩，而且 Jacobi 的新旧数组语义不能无意中变成 Gauss--Seidel 的就地更新。性能搜索必须与守恒量、残差或收敛曲线一起验证，才能确保更快的内核仍在求解同一个离散问题。
\variantheading 粗化因子改为 8 时，输入为 10240、输出为 7200 个元素，强度约 1.3417 FLOP/B。
\practiceheading 用 ptxas 和 Nsight Compute 同时检查寄存器数、shared bytes、spill 与 Roofline 位置。
\sourcecheck{https://docs.nvidia.com/cuda/cuda-c-best-practices-guide/}{NVIDIA CUDA Best Practices Guide 的共享内存与算术/带宽计量原则；元素数、FLOP 和字节三路独立复核}
\end{solutionexercise}

\chapter{直方图：原子操作与私有化}

\begin{solutionexercise}{1}
\problemheading 假设 DRAM 系统中每次原子操作的总延迟为 100 ns。针对同一个全局内存变量
执行原子操作时，能够获得的最大吞吐量是多少？

\approachheading 同一地址的读--改--写必须串行化，因此理想吞吐量是单次总延迟的倒数。

\answerheading
\[
R=\frac{1}{100\times10^{-9}\ \mathrm{s}}
 =10^7\ \mathrm{atomic\ ops/s}
 =10\ \mathrm{Mops/s}.
\]

\conclusionheading 最大约为每秒 1000 万次原子操作。
\discussionheading
\discussionpoint{从倒数公式走向排队模型。}把同一地址想成只有一个服务窗口，读、修改、写回构成不可重叠的状态转移；在题设的固定服务时间模型里，服务率就是 $\mu=1/L=10^7$ 次/秒。只有当到达过程带有随机性并且长期到达率 $\lambda$ 接近 $\mu$ 时，排队论才进一步预言等待时间急剧增加；若请求严格等间隔，不能无条件宣称存在同样的排队增长。若暂用指数服务时间的 M/M/1 近似，$1/(1-\rho)$ 是平均逗留时间相对服务时间 $L$ 的放大倍数，而纯排队等待为 $W_q=\rho L/(1-\rho)$，两者不能混称。题设把服务时间固定为 $L$，在泊松到达假设下更接近 M/D/1，此时 $W_q=\rho L/[2(1-\rho)]$；真实 GPU 请求的成批到达和调度相关性还会偏离这两个模型。这个区分说明吞吐上界与响应时间不是同一概念，也解释了为何热点计数器常在平均吞吐尚可时先出现长尾延迟。

\discussionpoint{地址级并行如何打破单热点上界。}十兆次每秒只约束同一状态变量，因为不同地址没有“必须看到上一次结果”的共同依赖链。若更新均匀落到 $H$ 个真正独立且由足够多内存分区服务的地址，粗略上界可以增长到 $H/L$；但缓存行映射、分区冲突和有限的原子执行单元会使增长早于理想值饱和。算法因此常先在寄存器、warp、线程块和全局分片中逐级聚合，再提交少量原子更新，这与树形归约、数据库分区聚合以及分布式计数器的分片思想属于同一种“先拆热点、后合并”的设计。例如八个分片的理想值是 80 Mops/s，若实测仅 45 Mops/s，就应继续检查这些地址是否落到少数内存分区。

\discussionpoint{原子性不等于完整的并发契约。}一次原子读--改--写保证该位置不会出现撕裂更新，却不自动把周围普通内存访问组成事务，也不必然提供应用所需的先后可见性；CUDA 的作用域和内存序仍需按生产者--消费者关系选择。即使竞争已经消除，32 位整数计数器仍可能回绕，浮点加法也会因线程到达次序改变舍入路径，例如 $(10^{20}+(-10^{20}))+1$ 与 $10^{20}+((-10^{20})+1)$ 的结果不同。由此可见，“没有数据竞争”“数学结果正确”和“逐位可复现”是三层不同要求，金融累计、科学复现实验或确定性训练必须分别设计验证策略。

\discussionpoint{怎样把微基准变成可解释证据。}只测一个线程规模会把原子管线、普通带宽和调度不足混在一起，无法判断十兆次上界是否真正生效。可信实验应固定更新总数，再系统改变热点地址数、地址到缓存行的分布、输入偏斜和并发块数，并把吞吐曲线与原子请求、缓存扇区和 warp 停顿关联起来；若增加独立地址后近似线性提升，才支持“同址序列化”判断。这样的实验方法也适用于锁、队列尾指针和图遍历 frontier 等共享状态：先构造可控竞争，再用不变量核对结果，最后才解释性能计数器。每组运行还应检查最终计数等于请求总数，否则一个丢更新的“更快”内核会制造完全错误的性能结论。
\variantheading 若总延迟降为 50 ns，同址上限为 $2\times10^7$ ops/s，即 20 Mops/s。
\practiceheading 用 Nsight Compute 的 atomic throughput 与 L2 sectors 在目标 GPU 上验证。
\sourcecheck{https://docs.nvidia.com/cuda/cuda-c-programming-guide/}{NVIDIA CUDA C++ Programming Guide 对原子读--改--写和同址序列化的定义；量纲独立复核，置信度高}
\end{solutionexercise}

\begin{solutionexercise}{2}
\problemheading 某处理器支持在 L2 缓存中执行原子操作。假设每次原子操作在 L2 缓存中需要
4 ns，在 DRAM 中需要 100 ns，并且 90\% 的原子操作命中 L2 缓存。针对同一个
全局内存变量执行原子操作时，近似吞吐量是多少？

\approachheading 对同一地址的串行操作，长期平均服务时间是两种延迟按命中概率的
加权平均，再取倒数；不能把两条独立吞吐率直接相加。

\answerheading
\[
\bar L=0.9(4)+0.1(100)=13.6\ \mathrm{ns},
\qquad
R\approx\frac1{13.6\times10^{-9}}
=7.35294\times10^7\ \mathrm{ops/s}.
\]
即约 $73.5$ Mops/s。

\conclusionheading 近似吞吐量为 $7.35\times10^7$ 次/秒（约 73.5 Mops/s）。
\discussionheading
\discussionpoint{期望延迟为何不要求独立命中。}设第 $i$ 次原子的服务时间 $L_i$ 只可能为 4 ns 或 100 ns，只要长期命中比例是 $0.9$，总时间除以请求数就趋近 $0.9\times4+0.1\times100=13.6$ ns。这里使用的是期望的线性性，并不要求相邻命中事件独立；即使九十次命中集中在一起、十次未命中也集中在一起，固定服务时间的串行模型仍给出相同平均值。真正错误的是平均两条吞吐率，因为那等价于把两类请求交给并行服务台，而题目明确要求同一地址按一个序列完成。相关性可以改变短时间窗口中的方差和尾延迟，却不会在这个固定成本模型中改变长期总服务时间，这两种统计量应分别讨论。

\discussionpoint{什么时候命中顺序才会改变现实吞吐。}真实缓存中的 4 ns 和 100 ns 很少是永远不变的常数：连续未命中可能填满队列，写回会占用链路，其他流量也会改变缓存行驻留和内存分区等待。只有加入这些状态相关服务时间、有限队列或可重叠请求后，命中成簇与否才会影响总耗时，而不仅是影响单次延迟的排列。例如两次 DRAM 请求若能并行在不同分区飞行，交错序列可能比严格同址模型更快；若都争用同一分区，则可能更慢。因而题目答案是明确抽象下的精确期望，工程外推必须先说明新增机制，不能把“相关性”本身误当成加权平均失效的原因。

\discussionpoint{缓存命中与拆分热点是两个优化维度。}让原子停留在 L2 可以缩短每次状态转移，却没有消除同一变量的顺序依赖，所以命中率从低到高最终仍会碰到单地址服务率。若把一个计数器按 key、warp 或线程块拆成多个副本，既可能提高这些缓存行的驻留概率，又创造真正的地址级并行；代价是初始化副本、占用容量并在末尾再做一次合并。这个权衡与数据库分区计数、分布式 sharded counter 十分相似：缓存优化降低一次访问的成本，分片优化改变依赖图，两者必须在模型中分开记账。若 32 个副本最终仍同时提交到一个地址，合并阶段会重新形成短时热点，所以还要把阶段峰值纳入容量设计。

\discussionpoint{由两点估计走向可校准性能模型。}可以先用常驻单热点和主动驱逐两组实验估计快、慢路径，再在不同命中比例下检查 $1/E[L]$ 是否预测实测饱和吞吐。若误差随并发请求数增加而扩大，通常说明服务可以重叠或队列开始饱和；若误差主要随工作集扩大而变化，则应调查容量失效、缓存行迁移和分区映射。进一步可把模型写成测得的条件延迟 $E[L\mid\text{state}]$ 与状态占比之和，并用独立数据集验证，而不是不断为同一条曲线调参数；这种校准--验证流程也是分析缓存、TLB 和远端内存系统的通用方法。
\variantheading 命中率为 $80\%$ 时平均延迟为 23.2 ns，吞吐约 43.1 Mops/s。
\practiceheading 用 Nsight Compute 的 L2 hit rate 和 atomic 指标代入实测值，并按热点分组分析。
\sourcecheck{https://docs.nvidia.com/cuda/cuda-c-programming-guide/}{NVIDIA CUDA C++ Programming Guide 的原子操作语义与缓存层次说明；以平均服务时间而非平均速率作对抗复核}
\end{solutionexercise}

\begin{solutionexercise}{3}
\problemheading 在第 1 题的条件下，假设内核每执行一次原子操作就执行 5 次浮点运算。受
原子操作吞吐量限制时，内核执行能达到的最大浮点吞吐量是多少？

\approachheading 每秒原子操作数乘每次原子操作关联的浮点运算数。

\answerheading
\[
(10^7\ \mathrm{atomic\ ops/s})(5\ \mathrm{FLOP/atomic\ op})
=5\times10^7\ \mathrm{FLOP/s}
=0.05\ \mathrm{GFLOP/s}.
\]

\conclusionheading 最大浮点吞吐量为 50 MFLOP/s（0.05 GFLOP/s）。
\discussionheading
\discussionpoint{把屋顶模型的分母换成原子操作。}传统 roofline 用 FLOP/byte 描述每搬运一字节能做多少计算，本题则可定义 $I_a=5$ FLOP/atomic，原子屋顶为 $P_a=I_aR_a=0.05$ GFLOP/s。完整设备模型还应取 $P=\min(P_{\mathrm{peak}},I_bB,I_aR_a)$，分别表示计算峰值、普通内存带宽和原子服务率；题目只给出最后一项，因此答案是条件上界而不是 GPU 的总体能力。这个写法的价值在于把“很慢”变成可比较的资源预算，也可推广到每次锁获取、消息发送或稀疏索引查找所支撑的有效工作量。

\discussionpoint{指令能否重叠取决于数据依赖。}若五次浮点运算只使用早已准备好的寄存器，调度器可以在某个 warp 等待原子时发射其他 warp 的算术，因而单次原子延迟未必全部暴露在时间线上。反之，若运算读取原子返回值，关键路径会形成“原子、五次运算、下一原子”的串行链；此时增加独立线程只能在资源允许时交错多条链，不能消掉链内依赖。用依赖图而非简单指令计数看问题，可以解释为何两段拥有相同 FLOP 和原子次数的代码性能仍可能相差很大，也连接到指令级并行、软件流水和异步访存的基本方法。可用两个微基准分别让算术依赖或不依赖原子返回值，二者的差距就是关键路径作用的直接证据。

\discussionpoint{聚合会重新定义每次原子的有效工作量。}假设一个 warp 中有 $u$ 个线程更新同一 key，先用 ballot 和 shuffle 求和后只提交一次全局原子，理想情况下 $I_a$ 可从 5 增到 $5u$。但聚合本身需要投票、交换和分组，若 32 个线程的 key 全不相同，原子数没有减少，额外指令反而成为纯开销；线程块私有化也还要支付共享空间、同步和最终提交。因而收益由 key 的重复度和偏斜决定，这与图顶点度分布、稀疏梯度重复索引以及哈希聚合中的局部性直接相关，不能只用平均线程数推断。若聚合固定成本相当于四次普通操作，就至少需要足够重复 key 省下超过该成本的原子请求才可能获利。

\discussionpoint{有效 FLOP 才能回答业务生产率。}分析器可能统计地址计算、重试路径或编译器生成的浮点指令，但这些并不都代表算法向答案推进的五次运算。以图 frontier 为例，重复边造成的原子竞争会增加硬件工作，却不增加新发现顶点；以梯度累加为例，数值精度和确定性又可能比峰值操作数更重要。可靠报告应从算法定义给出有效 FLOP，再同时记录原子请求、普通字节流量和端到端时间，并用输入偏斜分层展示结果。这样的口径能避免把无效忙碌误写成高性能，也便于比较私有化、排序归约和直接原子三类完全不同的实现。
\variantheading 每原子改为 12 FLOP 时上限为 $1.2\times10^8$ FLOP/s，即 0.12 GFLOP/s。
\practiceheading 在 Nsight Compute Roofline 中把 atomic throughput 与 FLOP 计数联合核对。
\sourcecheck{https://docs.nvidia.com/cuda/cuda-c-programming-guide/}{NVIDIA CUDA C++ Programming Guide 的原子串行化语义；单位约分与十进制前缀独立复核，置信度高}
\end{solutionexercise}

\begin{solutionexercise}{4}
\problemheading 在第 1 题的条件下，假设内核把全局内存变量私有化为共享内存变量，共享
内存访问延迟为 1 ns。原来的全局内存原子操作全部改成共享内存原子操作。为简化
分析，假设把私有变量累加到全局变量所增加的全局内存原子操作，会使总执行时间
增加 10\%。又假设内核每执行一次原子操作就执行 5 次浮点运算。受原子操作吞吐量
限制时，内核执行能达到的最大浮点吞吐量是多少？

\approachheading 本章模型把一次原子操作近似为一次加载加一次存储，因此 1 ns 的
共享内存访问延迟对应 2 ns 原子服务时间；再把合并开销造成的总时间乘 1.1。

\answerheading 不计合并时，共享原子上限为
$1/(2\ \mathrm{ns})=5\times10^8$ ops/s，对应 $2.5\times10^9$ FLOP/s。加入
$10\%$ 时间开销后：
\[
P=\frac{2.5\times10^9}{1.10}
=2.2727\times10^9\ \mathrm{FLOP/s}
\approx2.27\ \mathrm{GFLOP/s}.
\]

\conclusionheading 最大约为 2.27 GFLOP/s。
\discussionheading
\discussionpoint{私有化是一棵分层归约树。}把全局计数器复制到每个线程块，相当于先让局部更新在近处汇合，再把每块结果送回根节点；原来每个输入一次的远端竞争因此变成每个非空分箱至多一次提交。还可继续设置每 warp 副本，甚至让线程先在寄存器中累计短游程，形成寄存器、warp、块、全局四级结构。层次越深，热点越容易被拆散，但初始化、同步、额外存储和合并轮次也越多；这与多核 CPU 的 per-core counter 和分布式 map--reduce 都体现相同原则，即通过增加中间状态换取更短的共享依赖链。

\discussionpoint{延迟模型不能代替吞吐模型。}题目把共享原子抽象成一次 1 ns 读加一次 1 ns 写，所以得到 2 ns；真实硬件可能由专门原子管线完成，并能让访问不同 bank 的请求流水重叠。于是“某次操作从发出到完成要多久”和“饱和后每秒能完成多少次”不再互为简单倒数，尤其当多个 warp 同时访问不同地址时更明显。相反，若所有 lane 命中同一地址或同一 bank，序列化会重新显露依赖链。理解这一区别有助于避免用指针追逐测得的延迟预测流式吞吐，也解释了缓存、Tensor Core 和网络交换机中普遍存在的 latency--throughput 分离。

\discussionpoint{用成本方程寻找私有化盈亏点。}设直接方案需要 $N$ 次全局原子，私有化后有 $N$ 次共享更新和 $S$ 次全局提交，则可比较 $T_d\approx NL_g$ 与 $T_p\approx NL_s+SL_g+T_{init}+T_{sync}$。题目把后几项合并成额外 10\% 时间，便于算出 2.27 GFLOP/s；现实中的 $S$ 取决于块数和每块非空分箱数，绝不是固定比例。低竞争、小输入或分箱极多时，清零与提交可能吞掉收益；高偏斜、大输入时，节省的全局序列化则通常占主导。类似的盈亏方程也可指导缓存分块、局部哈希表和批量消息合并。

\discussionpoint{共享容量会反过来限制并行度。}每块保存 $K$ 个 32 位计数需要 $4K$ 字节，若再为多个 warp 建副本，容量按副本数成倍增长；共享内存过大时，每个 SM 可驻留块数下降，隐藏延迟的能力也随之减弱。使用 16 位计数虽能节省空间，却必须证明单块命中数不会溢出，或设计定期冲刷；分批处理大直方图则会重复读取输入。生产实现因此要在均匀、Zipf 偏斜和单热点数据上同时观察局部冲突、全局提交与驻留量，并以端到端时间决定设计，而不能孤立追求最高共享原子吞吐。例如 $K=4096$ 的单副本已需 16 KiB，八个 warp 各建副本就达到 128 KiB，许多配置根本无法同时驻留。
\variantheading 若合并开销为 $25\%$，则吞吐为 $2.5/1.25=2.0$ GFLOP/s。
\practiceheading 用共享原子 microbenchmark 和 Nsight Compute 的 shared atomic/bank 指标测量目标 SM。
\sourcecheck{https://docs.nvidia.com/cuda/cuda-c-programming-guide/}{NVIDIA CUDA C++ Programming Guide 对共享/全局原子读改写的说明；对抗审查同时计入读写两次访问，并确认时间增加 10\% 应使吞吐除以 1.1}
\end{solutionexercise}

\begin{solutionexercise}{5}
\problemheading 某直方图内核处理包含 524,288 个元素的输入，生成包含 128 个分箱的直方图。
内核配置为每块 1024 个线程。
\begin{enumerate}[label=\alph*.,leftmargin=2.2em]
  \item 图 9.6 的内核不使用私有化、共享内存和线程粗化，它会在全局
  内存上执行多少次原子操作？
  \item 图 9.10 的内核使用私有化和共享内存，但不使用线程粗化。它在
  全局内存上最多可能执行多少次原子操作？
  \item 图 9.14 的内核使用私有化、共享内存和线程粗化，粗化因子为 4。
  它在全局内存上最多可能执行多少次原子操作？
\end{enumerate}
图 9.6 每个输入元素直接对全局分箱做一次原子更新；图 9.10 每块先在共享内存形成
128 个私有分箱，再把每个非零分箱至多原子累加到全局一次；图 9.14 的提交方式相同，
但每块由线程粗化处理 $4\times1024$ 个输入元素。

\approachheading 基本内核每输入一次全局原子；私有化内核每块最终对每个非零私有
分箱至多提交一次全局原子，因此最大值是块数乘 128。粗化会把块数缩小 4 倍。

\answerheading
\begin{enumerate}[label=\alph*.,leftmargin=2.2em]
\item 原书图 9.6 每个输入元素执行一次全局原子，共
$524{,}288$ 次。
\item 无粗化时块数为 $524{,}288/1024=512$。每个块的 128 个私有分箱都可能非零，
故全局提交至多 $512\times128=65{,}536$ 次。
\item 粗化因子 4 时，每块处理 $4\times1024=4096$ 个输入，块数为
$524{,}288/4096=128$；至多执行 $128\times128=16{,}384$ 次全局原子。
\end{enumerate}
私有直方图内部仍对输入执行共享内存原子；上面的 (b)、(c) 只统计题目所问的全局
内存原子。若某块没有命中某分箱，实际提交次数会低于上界。

\complexity{$O(524{,}288)$ 输入工作}{$128$ 个共享分箱/块}{块数从 512 降至 128；块内更新有分箱争用}
\conclusionheading (a) 524,288；(b) 最多 65,536；(c) 最多 16,384 次全局原子。
\discussionheading
\discussionpoint{上界和典型提交量回答不同问题。}每块 128 个分箱都非空时，私有化方案确实提交 $128$ 次全局原子，但随机输入未必达到这个上界。若块内 $m$ 个样本独立且均匀落入 $K$ 个箱，某箱为空的概率是 $(1-1/K)^m$，所以非空箱期望为 $K[1-(1-1/K)^m]$；例如 $m$ 很小时，它更接近 $m$ 而不是 $K$。这个占位模型还能用非均匀概率 $p_k$ 推广成 $\sum_k[1-(1-p_k)^m]$，从而把真实偏斜映射到提交量，连接到缓存工作集估计、布隆过滤器占用分析和数据库 distinct 计数。

\discussionpoint{私有化只是把竞争搬到更近的位置。}每个输入仍要更新一个共享分箱，因此单热点数据会让同一块的线程在共享地址上排队，只是远端全局竞争被隔离成多个局部队列。若 warp 中许多 lane 命中同一箱，可先用 ballot 找出同组 lane，再由一个代表执行原子；也可为每个 warp 建独立直方图，最后在块内归约。另一条路线是先按分箱排序再计算游程长度，它几乎消除原子却增加排序流量。三种方法分别适合高重复、容量充足和可复用有序结果的场景，说明优化必须匹配数据分布而非只匹配算子名称。可以先抽样估计每 warp 的唯一 key 数，再据此选择直接原子或聚合路径，避免为均匀数据支付无效分组。

\discussionpoint{粗化同时改变块数和块内覆盖率。}因子 4 让每块处理 4096 项，块数由 512 降到 128，这是全局提交上界缩小四倍的直接原因；但每块样本更多，也提高了其触及全部 128 个箱的概率。均匀模型下某箱仍为空的概率约为 $(127/128)^{4096}$，已经极小，因此本题的上界几乎就是期望值；若分箱数扩大到数万，结论便可能反转。这个双重效应与批处理大小选择相似：更大批次摊薄固定开销，却扩大局部工作集并减少并行任务数，需要把两条变化同时放进模型。在偏斜分布下，应把统一的 $1/128$ 换成各箱概率 $p_k$，稀有箱的空概率会远高于热门箱。

\discussionpoint{不同业务会把同一计数问题变成不同算法。}图像颜色直方图的箱数小且固定，共享数组通常合适；数据库 group-by 的键空间可能巨大，需要局部哈希表或排序；词频统计又常呈长尾分布，少数热词适合专门缓存、冷词则走普通路径。验证这些实现时，首先要用总和等于输入数等不变量保证没有丢计数，再观察每块非空箱分布和端到端吞吐，判断时间花在更新、初始化还是合并。这样从题目的三个确定计数出发，就能建立一个按键空间、偏斜度和结果复用需求选择算法的完整决策框架。若结果随后需要按 key 有序输出，排序归约的前期成本还能被下游直接复用，这会改变原本只看直方图阶段的选择。
\variantheading 若分箱数为 64，其余不变，则 (b) 最多 32768 次，(c) 最多 8192 次全局原子。
\practiceheading 用 Nsight Compute 统计 global/shared atomic transactions，并以均匀、单热点和真实输入基准测试。
\sourcecheck{https://docs.nvidia.com/cuda/cuda-c-programming-guide/}{NVIDIA CUDA C++ Programming Guide 的共享/全局原子作用域；按原书图 9.6、9.10、9.14 的提交循环逐块计数复核}
\end{solutionexercise}

\chapter{归约}

\begin{solutionexercise}{1}
\problemheading 对图 10.5 的简单归约内核，若元素数为 1024、warp 大小为 32，
第 5 次迭代时块内有多少个 warp 会发生分歧？图 10.5 的基线内核为：
\begin{lstlisting}[language=C++,firstnumber=1]
__global__ void SimpleSumReductionKernel(float* input, float* output) {
    unsigned int i = 2*threadIdx.x;
    for(unsigned int stride = 1; stride <= blockDim.x; stride *= 2) {
        if(threadIdx.x % stride == 0) {
            input[i] += input[i + stride];
        }
        __syncthreads();
    }
    if(threadIdx.x == 0) {
        *output = input[0];
    }
}
\end{lstlisting}
对 1024 个输入元素，单块启动 512 个线程；“第 5 次迭代”从步幅 1 的第一次开始计数。

\approachheading 一个线程先负责两个元素，故块中有 $1024/2=512$ 个线程，即 16 个
warp。把“第 1 次”对应到步幅 1，逐轮列出活跃线程。

\answerheading 第 5 次迭代的步幅为 $1,2,4,8,16$ 中的 16。条件
\code{threadIdx.x \% 16 == 0} 使每个 warp 的 lane 0 和 lane 16 执行加法，其余 lane
不执行。每个 warp 内都有真假两条控制路径，因此 16 个 warp 全部分歧。

\conclusionheading 发生分歧的 warp 数为 $\boxed{16}$。
\discussionheading
\discussionpoint{分歧描述的是执行组内部而不是线程总数。}SIMT 让一个 warp 的 lane 共享取指，但每条指令带有活动掩码；当同一 warp 中部分 lane 满足条件、部分不满足时，硬件才需要分别覆盖不同路径。第 5 轮的条件是线程号能被 16 整除，因此每个完整 warp 都有 lane 0 和 16 为真、其余为假，恰好形成 16 个混合掩码。若整个 warp 都为假，它虽然没有做加法，却不算分歧；这个反例能避免把“闲置线程多”与“控制流分歧严重”混成一个概念，也适用于稀疏分支和边界谓词分析。反过来，一个仅有 33 个线程的块即使只有最后一线程有效，尾 warp 也可在活动掩码范围内保持一致，不能只按块中真假线程总数判断。

\discussionpoint{改变索引映射可以改变控制流几何。}交错归约把活跃线程按步幅散布在所有 warp 中，因此即使活跃线程数已经很少，仍可能让许多 warp 各自执行一小部分工作。收敛归约则让活跃线程始终组成从 0 开始的连续前缀，使完整 warp 要么全做、要么全不做；最后缩到 32 项时还可用 shuffle 在寄存器间交换值。这个演进没有改变求和的数学树，却改变了线程到树节点的映射，说明并行算法优化常是“重新安排所有权”而非减少算术工作，类似思想也见于稀疏矩阵行分组和图遍历 frontier 压缩。若重新映射同时把相邻线程对应到连续共享下标，还可能连带减少 bank 冲突，因此控制流与访存收益往往来自同一所有权变换。

\discussionpoint{分歧计数不能单独预测内核时间。}基线代码每轮还执行取模、地址计算、内存读写和块级屏障，后半段又因活跃线程减少而失去并行度；这些成本可能比真假路径本身更大。编译器还可能把很短的条件体变成谓词化指令，此时分支事件看似消失，假 lane 仍占据发射机会而不产生有效加法。因而性能模型至少要同时考虑每轮活跃 lane、指令数、共享事务和同步等待，并用时间验证各项权重；这也是为什么“分支效率提高”并不自动等于应用加速。例如把取模改成位运算后，分歧数完全不变却可能少发多条整数指令；反之消除一个分支却保留十轮屏障，时间也可能几乎不动。应以逐轮耗时和 stall 分类验证谁在关键路径上。

\discussionpoint{执行组宽度是算法参数而不是宇宙常数。}本题明确给出 warp 为 32，所以 lane 0 和 16 的结论成立；在采用 64-lane wavefront 或可变 subgroup 的平台上，同一条件会形成不同数量的混合执行组。即使仍在 CUDA 上，最后一个不足 32 个有效元素的 warp 也必须用实际活动掩码参与 shuffle，不能假定满组。可移植实现应从运行环境或编程模型取得组宽，并用 31、32、33 等边界规模检查掩码。把这种假设显式化后，归约代码才可能安全迁移到合作组、HIP 或着色器子组，而不是只在某个固定块大小下偶然正确。
\variantheading 若仍有 512 个线程，但问第 6 次迭代（步幅 32），每个 warp 只有 lane 0
满足条件，16 个 warp 仍全部分歧；活跃加法则从 32 次降为 16 次。
\practiceheading 可在 Nsight Compute 中比较 Branch Efficiency 与 Warp State Statistics，
并以不同块大小复测；分析时假定设备 warp 宽度为运行时属性报告的 32，而非把它隐含推广到
其他并行体系结构。
\sourcecheck{https://docs.nvidia.com/cuda/cuda-c-programming-guide/}{NVIDIA CUDA C++ Programming Guide 对 warp 与分支分歧的定义；高置信度}
假设审查：若把迭代从 0 编号，第 5 号迭代会是步幅 32；题目用“第 5 次”而非“编号 5”，
应按步幅 16 解释。即使取步幅 32，各 warp 也仍只有 lane 0 活跃，答案仍为 16。
\end{solutionexercise}

\begin{solutionexercise}{2}
\problemheading 对图 10.8 的改进归约内核，若元素数为 1024、warp 大小为 32，
第 5 次迭代时有多少个 warp 会发生分歧？图 10.8 的基线内核为：
\begin{lstlisting}[language=C++,firstnumber=1]
__global__ void ConvergentSumReductionKernel(float* input, float* output) {
    unsigned int i = threadIdx.x;
    for(unsigned int stride = blockDim.x; stride >= 1; stride /= 2) {
        if(threadIdx.x < stride) {
            input[i] += input[i + stride];
        }
        __syncthreads();
    }
    if(threadIdx.x == 0) {
        *output = input[0];
    }
}
\end{lstlisting}
对 1024 个输入元素，单块启动 512 个线程；“第 5 次迭代”从步幅 512 的第一次开始计数。

\approachheading 该内核让活跃线程始终构成从线程 0 开始的连续前缀。第 5 次迭代的
步幅是 $512,256,128,64,32$ 中的 32。

\answerheading 条件为 \code{threadIdx.x < 32}。warp 0 的 32 个 lane 全部为真，
warp 1--15 的 lane 全部为假；没有任何一个 warp 同时执行真假两条路径。因此分歧数为 0。
这里“非活跃 warp”不等同于“发生分歧”：分歧只在同一 warp 的活动 lane 对分支条件
取值不一致时产生。

\conclusionheading 发生分歧的 warp 数为 $\boxed{0}$。
\discussionheading
\discussionpoint{没有分歧仍可能只有很少有效工作。}第 5 轮的 32 个活跃线程恰好填满 warp 0，其余 15 个 warp 对条件全为假，所以控制流是收敛的；然而有效加法只剩初始 512 个线程容量的 $1/16$。若只观察 branch efficiency，会把这个阶段误判为高效，而算术利用率和屏障等待会揭示大部分线程资源已经闲置。这个例子建立了三个独立维度：路径是否一致、lane 是否执行有效指令、硬件是否仍有其他 warp 可调度；分析稀疏计算、条件卷积和尾块时也必须保持这种区分。若一个 SM 同时驻留多个独立归约块，其他块的活跃 warp 尚可填补空槽；只有单块或低占用时，后期闲置才会完全暴露为时间。

\discussionpoint{归约天然需要跨越多个并行层次。}当活动范围大于一个 warp 时，共享内存和块级屏障提供任意线程间的交换；缩到一个 warp 后，shuffle 能直接传递寄存器值，避免继续支付全块同步。数组再大时，单块先产生部分和，第二层再归约这些部分和，必要时递归多层。这个“局部聚合、上层合并”的结构从 GPU 块归约一直延伸到 CPU SIMD、NUMA 节点和分布式 all-reduce，区别只在每层通信代价和同步范围，因此理解本题的边界正是理解大型归约系统的入口。例如多 GPU 先在卡内块归约，再经 NVLink 或网络合并卡级结果，正是把同一树继续向外扩展。

\discussionpoint{删除屏障必须同时维护可见性和参与集合。}连续前缀减少了混合分支，却没有自动消除每轮的 \code{__syncthreads()}；后几轮只有几十次加法，数百线程仍等待屏障，固定成本会占据主要比例。改用 \code{__shfl_down_sync} 时，调用者要传递包含所有参与 lane 的正确掩码，并保证这些 lane 以一致方式执行原语；在独立线程调度架构上，依赖过去“warp 天然锁步”的隐含假设可能读到旧值。类似规则也适用于 warp 级投票和矩阵片段操作：局部原语降低同步范围，却把活动集合变成了显式正确性契约。

\discussionpoint{规模决定应该比较哪一种实现。}几十个元素的归约往往被 kernel launch 和主机往返主导，此时再优化一轮共享访问几乎不可见；数百万元素则更接近内存带宽上限，首层加载与部分和流量成为关键。手写版本应分别与块级库原语、两阶段设备归约和可能的融合消费者比较，并在相同数值语义下核对误差。若只报告一个大规模峰值，会错过小批量推理和在线查询的延迟需求；若只看一个小数组，又无法评价树形映射对带宽的影响。分段绘制规模--时间曲线比单一加速比更能指导生产选择。曲线的转折点还能直接告诉调度器何时走 CPU、小型融合 kernel 或完整设备归约，而不必把一个实现强用于所有输入。
\variantheading 若第 5 次的活跃前缀改为 20 个线程，则 warp 0 有 20 个真 lane 和 12 个
假 lane，恰有 1 个 warp 分歧，其余 15 个 warp 全假。
\practiceheading 用 Nsight Compute 同时检查 Branch Efficiency 和每轮 Executed
Instructions，才能区分“没有分歧”与“没有工作”；生产代码通常在最后一个 warp 改用
shuffle 归约以减少块级屏障。
\sourcecheck{https://docs.nvidia.com/cuda/cuda-c-best-practices-guide/}{NVIDIA CUDA C++ Best Practices Guide 的分支分歧说明；高置信度}
反例审查：若步幅不是 warp 大小的整数倍，例如 20，则第一个 warp 会有 20 个真 lane 与
12 个假 lane；本题步幅恰为 32，不能把“只有一个活跃 warp”误判为分歧。
\end{solutionexercise}

\begin{solutionexercise}{3}
\problemheading 修改图 10.8 的内核，采用下面所示的访问模式：
\begin{center}
\scriptsize
\setlength{\tabcolsep}{2.2pt}
\begin{tabular}{c|*{16}{c}}
  &0&1&2&3&4&5&6&7&8&9&10&11&12&13&14&15\\\hline
  线程所有者&&&&&&&&&$T_0$&$T_1$&$T_2$&$T_3$&$T_4$&$T_5$&$T_6$&$T_7$\\
  步 0&&&&&&&&&\colorbox{blue!25}{\strut$\Sigma$}&\colorbox{blue!25}{\strut$\Sigma$}&
    \colorbox{blue!25}{\strut$\Sigma$}&\colorbox{blue!25}{\strut$\Sigma$}&
    \colorbox{blue!25}{\strut$\Sigma$}&\colorbox{blue!25}{\strut$\Sigma$}&
    \colorbox{blue!25}{\strut$\Sigma$}&\colorbox{blue!25}{\strut$\Sigma$}\\
  步 1&&&&&&&&&&&&&\colorbox{blue!25}{\strut$\Sigma$}&
    \colorbox{blue!25}{\strut$\Sigma$}&\colorbox{blue!25}{\strut$\Sigma$}&
    \colorbox{blue!25}{\strut$\Sigma$}\\
  步 2&&&&&&&&&&&&&&&\colorbox{blue!25}{\strut$\Sigma$}&
    \colorbox{blue!25}{\strut$\Sigma$}\\
  步 3&&&&&&&&&&&&&&&&\colorbox{blue!25}{\strut$\Sigma$}
\end{tabular}
\\[0.35em]
\begin{tabular}{c@{\quad}l}
  步 0 & $\code{input[8+k]}\leftarrow\code{input[8+k]}+\code{input[k]}$，$0\le k<8$\\
  步 1 & $\code{input[12+k]}\leftarrow\code{input[12+k]}+\code{input[8+k]}$，$0\le k<4$\\
  步 2 & $\code{input[14+k]}\leftarrow\code{input[14+k]}+\code{input[12+k]}$，$0\le k<2$\\
  步 3 & $\code{input[15]}\leftarrow\code{input[15]}+\code{input[14]}$
\end{tabular}
\end{center}
图表以 16 个输入、8 个线程说明一般模式：线程拥有右半区，活跃所有者逐轮向区段
末端收敛，最终和位于最后一个输入位置。

\approachheading 对 $2B$ 个输入启动 $B$ 个线程。步幅依次为 $B,B/2,\ldots,1$；步幅
$s$ 时只让 $t\ge B-s$ 的线程把位置 $B+t-s$ 加到所有者位置 $B+t$。

\answerheading 以下内核适用于 $B$ 为 2 的幂且输入恰有 $2B$ 个元素：
\begin{lstlisting}[language=C++]
__global__ void right_owned_sum(float* input, float* output) {
    const unsigned int B = blockDim.x;
    const unsigned int t = threadIdx.x;
    const unsigned int owner = B + t;
    for (unsigned int stride = B; stride >= 1; stride >>= 1) {
        if (t >= B - stride)
            input[owner] += input[owner - stride];
        __syncthreads();
    }
    if (t == B - 1) *output = input[2 * B - 1];
}
// Example: right_owned_sum<<<1, 512>>>(d_input, d_output);
\end{lstlisting}
步幅 $B$ 更新位置 $B\ldots2B-1$；步幅 $B/2$ 只更新最后 $B/2$ 个位置。归纳地，
每轮把两个不相交的索引子集聚合到右侧所有者，子集大小逐轮倍增，最后位置 $2B-1$
包含全部 $2B$ 项的和。这里的子集一般并不连续：例如 $B=2$ 时，若把操作数写成“左源在
前”，拓扑顺序仍是 $a,c,b,d$。对本题加法，交换律使该置换不影响总和；每轮源、目的位置
不重叠，块级屏障又隔开相邻轮，所以没有读写竞态。

\complexity{$O(B)$ 总工作、$O(\log B)$ 并行深度}{$O(1)$ 额外设备空间}{$B$ 个线程，活跃数逐轮减半}
\conclusionheading 对加法归约，结果位于 \code{input[2*blockDim.x-1]}，访问模式与题图完全一致。
\discussionheading
\discussionpoint{所有者在右端不代表它拥有连续区间。}首轮直接把位置 $t$ 与 $B+t$ 配对，所以位置 $B+t$ 保存的是两个相隔 $B$ 的元素之和，而不是以它结尾的长度 2 连续片段。后续轮次继续合并这种交错集合，最终结果确实落在下标 $2B-1$，加法数和树深也与常见归约相同，但中间状态的集合含义完全不同。因而该布局可以省去最终结果搬到数组首端的一步，却不能直接充当游程结束值或分段扫描摘要；算法设计中“结果放在哪里”和“结果代表哪一段输入”必须分别证明。可在调试表中为每个槽保存叶子下标集合而非只保存和，这会立即显示第一轮的集合是 $\{t,B+t\}$。

\discussionpoint{非交换运算给出最直接的反例。}加法满足交换律，所以子树顺序被置换后结果不变；字符串连接和矩阵乘法只满足结合律，不能承受这种置换。取 $B=2$、输入为 $a,b,c,d$，即使每轮都让源操作数写在左边，首轮得到 $ac$ 与 $bd$，末轮仍是 $acbd$ 而不是 $abcd$。要保持原序，应改用先合并相邻区间的树，或先按此拓扑的逆置换重排输入，并明确单位元与尾部补齐规则。这个小反例也说明测试非交换结合算子是发现并行扫描、归约中隐式重排的有效方法。矩阵乘法可用不对易的两个 $2\times2$ 矩阵构造同类数值反例，避免字符串被误认为只是展示技巧。

\discussionpoint{设备范围扩展需要新的同步边界。}单块只能覆盖 $2B$ 项，更长输入通常让每块输出一个部分结果，再启动下一层内核归约这些结果；每次 launch 天然形成设备范围的阶段边界。合作启动允许用 \code{grid.sync()} 把若干层放入一个内核，但要求整个合作网格满足驻留约束，普通内核不能用共享标志自旋来伪造屏障。对交换求和，部分结果的排列通常无碍；对需要固定次序的算子，块分区和上层树还必须保持原始区段顺序。相同问题会出现在多 GPU all-reduce 中，只是同步边界从 kernel launch 变成通信轮次。

\discussionpoint{数值误差使“数学等价”仍不等于逐位相同。}浮点加法在实数上可交换结合，在有限精度中却会因每棵树的配对和舍入点不同而得到不同末位，量级悬殊或强抵消数据尤其明显。固定树只能提供确定的执行顺序，不必然提高准确度；双精度部分和、成对求和或补偿算法则用额外资源换取误差控制。调试时应保存小输入每轮的所有者集合和值，与按同一树序运行的 CPU 模拟比较，同时断言未拥有位置是否允许被覆盖。这样可以把索引错误、缺屏障和纯舍入差异分开，而不是看到最终和不同就笼统归咎于竞态。误差测试可同时给出高精度参考、绝对误差和 ULP 距离，使“可接受差异”成为量化契约。
\variantheading 若 $B=256$，内核读取 512 项并在 9 轮后把总和留在
输入数组下标 511；若 $B$ 不是 2 的幂，当前倍减循环会遗漏项，必须先填充到 2 的幂或改写树。
\practiceheading 可用 CUB 块归约作数值基线，并在 NVIDIA 竞争与同步检查工具下覆盖
$B=32,256,512$；启动前查询每块线程上限并验证输入至少有
$2B$ 项。
\sourcecheck{https://docs.nvidia.com/cuda/cuda-c-programming-guide/}{NVIDIA 对 \code{__syncthreads()} 的块级内存可见性保证；另以 $B=2$ 的非交换字符串反例对抗核验索引顺序，确认该拓扑不能仅凭结合律推广}
边界审查：内核不能跨块同步，且没有长度参数；调用方必须保证单块、$2B$ 个有效元素以及
$B$ 为 2 的幂。任一线程在屏障前提前返回都会令屏障行为未定义。
\end{solutionexercise}

\begin{solutionexercise}{4}
\problemheading 修改图 10.20 的内核，使它执行最大值归约而不是求和归约。图 10.20
给出的粗化关键代码为：
\begin{lstlisting}[language=C++,firstnumber=1]
__global__ void CoarsenedSumReductionKernel(float* input, float* output) {
    ...
    unsigned int segment = COARSE_FACTOR*2*blockDim.x*blockIdx.x;
    unsigned int i = segment + threadIdx.x;
    float partialSum = input[i];
    for(unsigned int c = 1; c < COARSE_FACTOR*2; ++c) {
        partialSum += input[i + c*BLOCK_DIM];
    }
    ...
}
\end{lstlisting}
省略号表示与图 10.18 相同的共享内存块内归约和块结果提交部分；改写时还必须选择
适用于任意（包括全负）输入的最大值单位元。

\approachheading 把二元运算从加法换成 \code{fmaxf}，并把单位元从 0 换成负无穷。
各块先产生一个最大值；对块结果递归调用同一内核，避免依赖浮点 \code{atomicMax} 的
架构与 NaN 语义差异。

\answerheading 下面是完整的单层内核；主机反复令输入指向上一层 \code{partial}，直到
只剩一个元素。
\begin{lstlisting}[language=C++]
#include <cuda_runtime.h>
#include <math_constants.h>

template<int B, int C>
__global__ void coarsened_max(const float* in, float* partial, size_t n) {
    __shared__ float s[B];
    size_t segment = size_t(blockIdx.x) * (2 * C * B);
    size_t i = segment + threadIdx.x;
    float v = -CUDART_INF_F;
#pragma unroll
    for (int c = 0; c < 2 * C; ++c) {
        size_t p = i + size_t(c) * B;
        if (p < n) v = fmaxf(v, in[p]);
    }
    s[threadIdx.x] = v;
    __syncthreads();
    for (int stride = B / 2; stride > 0; stride >>= 1) {
        if (threadIdx.x < stride)
            s[threadIdx.x] = fmaxf(s[threadIdx.x],
                                  s[threadIdx.x + stride]);
        __syncthreads();
    }
    if (threadIdx.x == 0) partial[blockIdx.x] = s[0];
}
// blocks = (n + 2*C*B - 1) / (2*C*B), B must be a power of two.
\end{lstlisting}
不变量是 \code{v} 始终等于该线程已访问元素的最大值；共享内存归约每轮合并两个不相交
集合的最大值，因此块结果正确。不同块写不同的 \code{partial} 元素，不存在竞态。

\complexity{$O(n)$ 总工作、每层 $O(C+\log B)$}{$O(n/(2CB))$ 首层部分结果空间}{$\lceil n/(2CB)\rceil$ 个块，块内 $B$ 线程}
\conclusionheading 用负无穷作为单位元后，负数输入也能正确处理；递归归约得到全局最大值。
\discussionheading
\discussionpoint{最大值首先是一份接口语义。}对普通有限浮点数，最大运算满足结合律，负无穷又是单位元，因此树形归约可以从任意尾部补入负无穷而不改变答案。NaN 和有符号零却迫使接口作选择：\code{fmaxf} 在只有一个操作数为 NaN 时通常返回数值操作数，而数据质量系统可能要求任何 NaN 都传播；$+0$ 与 $-0$ 的选择也可能影响后续倒数。因而“求最大值”并非天然唯一的位级函数，库实现、CPU 参考和测试必须共享同一契约。类似问题也出现在最小值、argmax 和缺失数据聚合中。例如输入只有 NaN 时，是返回 NaN、负无穷还是错误状态，必须由空值政策而非归约树偶然决定。

\discussionpoint{线程粗化是在减少层次与隐藏延迟之间交易。}每线程读取 $2C$ 项时，块数约缩为原来的 $1/C$，部分结果数组和递归 launch 随之减少，地址计算也被多个输入摊薄。另一方面，每个线程的串行加载链更长，独立线程总数下降；即使只保留一个最大值累加器而没有明显寄存器膨胀，设备也可能缺少足够 warp 来遮蔽内存延迟。例如短数组取很大 $C$ 会只启动少数块，SM 大量空闲。最优粗化因子因此取决于数组规模、带宽延迟和驻留资源，这也是向量化、批处理与循环展开普遍面对的平衡。可把首层块数至少覆盖若干倍 SM 数作为下限，再在满足该并行度的候选 $C$ 中测量。

\discussionpoint{设备级归约有递归与原子两条主路线。}块内缩到一个 warp 后可用 shuffle 完成，块结果则可以写入数组并递归归约，或直接原子更新一个全局最大值。递归方案多读写一轮部分结果，却让每层访问规则、同步边界和吞吐比较稳定；原子方案省空间和 launch，但所有块争用一个地址，且浮点类型、NaN 语义与架构支持需要逐项确认。成熟的 DeviceReduce 还会按规模选择多种策略，所以手写内核应把库结果当作语义与性能基线，而不是假定某一路线普遍最快。当块数只有几十时原子热点可能可接受，达到数万时同址提交尾部往往主导，这给出了自然的分派阈值。

\discussionpoint{很多应用真正需要的是带证据的 argmax。}最大池化只要值，异常检测、数据库 top-1 和注意力解码常还要知道最大值来自哪个位置，因此归约元素应是值--索引对。相等值时选最小索引可以获得稳定结果，若不同层采用不同 tie-break，最终位置就可能随块大小改变；含 NaN 时还要决定它是否胜出。测试集应覆盖全负数、正负无穷、重复最大值、NaN 和不完整尾块，并同时检查值和索引。由单纯 max 扩展到带伴随数据的单子，是理解 top-k、分段最优和动态规划回溯的自然一步。组合运算可定义为先比较值、再比较索引，从而仍满足结合律并可安全放入任意归约树。
\variantheading 若输入类型改为 \code{double}，应使用 \code{-CUDART_INF} 与
\code{fmax}；对空输入则没有数组最大值，主机接口应返回错误或显式约定的空集合结果。
\practiceheading 以 CUB \code{DeviceReduce::Max} 为性能与结果基线，使用全负、含正负无穷、
含 NaN 和非整除尾部的测试集；用 Nsight Compute 检查增大 $C$ 后的占用率与内存吞吐。
\sourcecheck{https://github.com/NVIDIA/cccl/tree/main/cub}{CUB DeviceReduce 对归约运算和初值的官方接口说明；高置信度}
反例审查：以 0 作单位元会把全负数组错误地归约为 0。NaN 的处理取决于 \code{fmaxf}
语义；若应用要求“遇到 NaN 即返回 NaN”，必须另行累计 NaN 标志。
\end{solutionexercise}

\begin{solutionexercise}{5}
\problemheading 修改图 10.20 的内核，使它能处理不一定为
\code{COARSE_FACTOR*2*blockDim.x} 整数倍的任意长度输入。为内核增加一个表示输入
长度的参数 \code{N}。图 10.20 的粗化加载关键代码为：
\begin{lstlisting}[language=C++,firstnumber=1]
unsigned int segment = COARSE_FACTOR*2*blockDim.x*blockIdx.x;
unsigned int i = segment + threadIdx.x;
float partialSum = input[i];
for(unsigned int c = 1; c < COARSE_FACTOR*2; ++c) {
    partialSum += input[i + c*BLOCK_DIM];
}
\end{lstlisting}
其后仍执行共享内存块内求和归约并提交块结果；所有线程必须继续到达块级屏障，
越界输入应贡献加法单位元而不能提前返回。

\approachheading 每次全局读取前检查下标；越界 lane 贡献加法单位元 0。不要在检查失败时
让线程返回，因为所有线程仍必须到达块级屏障。

\answerheading
\begin{lstlisting}[language=C++]
template<int B, int C>
__global__ void coarsened_sum_any(const float* input, float* partial,
                                  size_t N) {
    __shared__ float s[B];
    const size_t base = size_t(blockIdx.x) * (2 * C * B);
    const size_t owner = base + threadIdx.x;
    float sum = 0.0f;
#pragma unroll
    for (int c = 0; c < 2 * C; ++c) {
        const size_t i = owner + size_t(c) * B;
        if (i < N) sum += input[i];
    }
    s[threadIdx.x] = sum;
    __syncthreads();
    for (int stride = B / 2; stride > 0; stride >>= 1) {
        if (threadIdx.x < stride) s[threadIdx.x] += s[threadIdx.x + stride];
        __syncthreads();
    }
    if (threadIdx.x == 0) partial[blockIdx.x] = s[0];
}
// If N==0, return 0.0f on the host; a zero-block CUDA launch is invalid.
// Otherwise grid.x = (N + 2*C*B - 1) / (2*C*B); reduce partial recursively.
\end{lstlisting}
每个合法下标恰落入一个块、一个线程和一个粗化循环迭代；越界下标不读取。故各块部分和
覆盖输入的不相交划分，递归相加得到总和。使用 \code{size_t} 还避免大数组的 32 位乘法
溢出。

\complexity{$O(N)$ 总工作}{$O(\lceil N/(2CB)\rceil)$ 首层部分和空间}{$O(N/(CB))$ 个线程，块间独立}
\conclusionheading 边界检查与单位元填充使尾部不完整区段正确且无越界访问。
\discussionheading
\discussionpoint{单位元把不规则尾部嵌入规则树。}最后一个块不足 $2CB$ 项时，为每个越界位置贡献 0，等价于把原数组扩展为更长数组，因为 $x+0=x$；归约树的控制流因此无需随有效长度变化。换成最大值时应填负无穷，逻辑与应填真，前提都是该运算存在左右单位元。若操作只有结合律却没有单位元，就不能随意构造虚拟元素，而应让节点携带“是否有效”标志或使用分支合并。这个代数视角把边界检查从 CUDA 技巧提升为通用方法，也适用于扫描、分块卷积和树形数据库聚合。对平均值则不能直接填 0 后除以填充长度，必须同时归约总和与有效计数，这就是错误单位元选择的具体反例。

\discussionpoint{屏障要求无效线程继续参与协议。}若越界线程在加载后直接 \code{return}，同一块的有效线程还会到达 \code{__syncthreads()}，于是参与集合不一致，行为未定义且可能永久等待。正确模式是让所有线程执行相同的同步骨架，只在内存访问和最终写回处用谓词决定是否真正读写；无效线程把单位元写入共享槽后仍参加每一轮。这个原则也解释了为何 tile 边界常用掩码填零，而不是让线程提前退出。更一般地，任何块级集体操作都要求调用集合一致，条件路径必须在进入集体之前重新收敛。

\discussionpoint{一次边界判断不能覆盖整个规模契约。}即使每次读取都检查 $p<N$，表达式 $2CB\cdot\code{blockIdx.x}$ 若先在 32 位中溢出，得到的 $p$ 仍可能落回合法范围并错误读取旧位置。地址、网格数和部分结果容量应使用 \code{size_t} 或经证明足够宽的类型，主机还需特判 $N=0$，因为零块 launch 与递归终止各有接口约束。转换到 CUDA 的有限网格维度前也要验证上界，而非直接截断。这样的全链路整数审查同样适用于张量步幅、文件偏移和分布式计数，许多大规模错误并不发生在最终解引用处。

\discussionpoint{尾部成本应按余数而非单一规模观察。}常见情形只有最后一个块包含无效循环，所以当数组很大时，浪费比例至多约为一个块容量除以 $N$；但短输入或很大的 $C$ 会让大多数 lane 都只执行谓词和同步。可按 $N\bmod(2CB)$ 分组测量，使满块、差一个元素和只含一个元素三种极端清晰可见，再决定是否增加小规模专用 kernel 或 CPU 路径。网格步进加载能让固定块数反复处理长输入，却会改变缓存和归约层次。这样建立的分段成本模型比一句“尾块可忽略”更可靠，也能指导批处理系统如何把多个小数组拼接起来。
\variantheading 若 $B=256,C=2,N=1025$，每块覆盖 1024 项，网格为 2 块；第二块仅一个
合法输入，其余 lane 都贡献 0，两个部分和再归约得到总和。
\practiceheading 对 $N=0,1,2CB-1,2CB,2CB+1$ 做参数化测试，并在 Compute Sanitizer
memcheck 下运行；生产封装还应检查 \code{size_t} 网格计算是否超过设备 \code{gridDim.x}
上限以及每次启动的返回错误。
\sourcecheck{https://docs.nvidia.com/cuda/cuda-c-best-practices-guide/}{NVIDIA Best Practices Guide 的合并访问与边界处理原则；高置信度}
同步审查：把 \code{if (i>=N) return} 放在屏障之前会导致部分块线程缺席屏障；这里仅跳过
加载而不退出。浮点加法不满足实数意义上的严格结合律，结果可能随分组产生末位差异。
\end{solutionexercise}

\begin{solutionexercise}{6}
\problemheading 假设要对下面的输入数组执行并行归约：
\[
  \{6,2,7,4,5,8,3,1\}.
\]
分别采用以下内核时，画出每次迭代后数组内容的变化：
\begin{enumerate}[label=\alph*.,leftmargin=2.2em]
  \item 图 10.5 中未经优化的内核；
  \item 图 10.8 中针对访问合并与分歧优化的内核。
\end{enumerate}
图 10.5 令线程 $t$ 拥有位置 $2t$，步幅从 1 开始逐轮加倍，满足
\code{threadIdx.x \% stride == 0} 的线程执行
\code{input[2*t] += input[2*t + stride]}。图 10.8 令线程 $t$ 拥有位置 $t$，
步幅从 4 开始逐轮减半，满足 $t<\code{stride}$ 的线程执行
\code{input[t] += input[t + stride]}；每轮之后均有块级屏障。

\approachheading 两种内核都以 4 个线程处理 8 个元素。严格按每轮屏障前的旧值同时计算，
未被写入的位置保持不变。

\answerheading 交错活跃版本的步幅依次为 1、2、4：
\[
\begin{array}{c|rrrrrrrr}
\text{初始}&6&2&7&4&5&8&3&1\\
s=1&8&2&11&4&13&8&4&1\\
s=2&19&2&11&4&17&8&4&1\\
s=4&36&2&11&4&17&8&4&1
\end{array}
\]
收敛式版本的步幅依次为 4、2、1：
\[
\begin{array}{c|rrrrrrrr}
\text{初始}&6&2&7&4&5&8&3&1\\
s=4&11&10&10&5&5&8&3&1\\
s=2&21&15&10&5&5&8&3&1\\
s=1&36&15&10&5&5&8&3&1
\end{array}
\]
两条路径的最终和都可独立复算为 $6+2+7+4+5+8+3+1=36$。

\complexity{两者均为 $7$ 次加法、$O(N)$ 工作}{$O(1)$ 额外空间}{$N/2$ 个线程，$O(\log N)$ 轮}
\conclusionheading 两种状态演化不同，但最终都在位置 0 得到 $\boxed{36}$。
\discussionheading
\discussionpoint{树形拓扑编码了允许的代数变换。}两种求和过程都把八个叶子逐层合并，在没有整数溢出的数学模型中，结合律和交换律保证最终都是 36；不同之处是中间节点落在哪些下标、代表哪些叶子集合。若操作只结合而不交换，树可以重新加括号，却不能调换叶子顺序，右侧所有权的交错配对便可能失效。用字符串连接作为测试，正确结果必须保持 $1,2,\ldots,8$ 的次序，任何重排都会立刻暴露。这种“用更弱代数结构审查实现”的方法也适合扫描、矩阵链和函数复合。测试框架还可随机生成短字符串并与顺序左折叠比较，使隐藏置换不再依赖一个人工样例。

\discussionpoint{浮点的确定性与准确度是两项目标。}IEEE 加法每次都舍入，因此 $(a+b)+c$ 与 $a+(b+c)$ 可能不同；例如大数与小数混合、正负抵消时，差异甚至会累积到多个 ULP。固定块大小、树形和编译选项可以让同一平台重复得到相同位型，却不保证它最接近实数真值；使用双精度、成对求和或补偿求和可能提高准确度，又会改变吞吐和结果位型。科学计算应先决定需要误差界、统计稳定还是逐位复现，再选择算法。把这三者混称为“数值正确”会使测试标准无法落实。还应固定是否允许 FMA 收缩和快速数学优化，因为编译器变换本身也能改变舍入位置。

\discussionpoint{控制流和地址轨迹由同一映射共同决定。}交错树用取模选出稀疏线程，使活动 lane 分散在多个 warp，同时访问下标也跨越越来越大的间隔；收敛树把活动线程聚为前缀，通常减少混合分支和地址计算。可是两者都要在每轮隔离读写并同步，所以共享 bank 映射、屏障次数和后期并行度仍会影响性能。中间数组表不仅是算术演示，也是逐轮内存轨迹：给每个单元标注读者、写者和所属 bank，就能把抽象树转成可检查的体系结构预测。这种方法同样适合扫描和动态规划波前。例如连续前缀若让 lane $t$ 访问 $t+s$，可直接按地址字模 32 推导当前轮是否存在 bank 冲突。

\discussionpoint{中间状态是比最终答案更强的测试预言。}若只断言最终值为 36，两个相反的索引错误可能偶然抵消，缺少屏障也可能在某次调度中侥幸通过。把每轮数组作为黄金结果，可以精确指出从哪个步幅开始偏离；CPU 参考还应按完全相同的树序更新，而不是只调用顺序 \code{accumulate}。调试构建可让一个小块写出快照并配合竞争检查，发布构建则关闭记录以免改变时序和带宽。属性测试还能随机生成输入，验证每个中间节点等于其声明叶子集合的和，把一个固定手算例子扩展成持续回归工具。对浮点输入，参考模拟还应在每个节点用相同精度舍入，否则会把参考模型差异误报为内核错误。
\variantheading 若把最后一项从 1 改为 $-1$，两条树的中间状态相应改变，但精确总和均为
$34$；按各轮旧值重算是评分时的核心要求。
\practiceheading 可把两张状态表编码为 CPU 单元测试的黄金结果，再用 CUB
\code{BlockReduce} 验证最终值；浮点生产测试应采用误差界或 ULP 容差，而非默认要求与顺序
累加逐位一致。
\sourcecheck{https://github.com/NVIDIA/cccl/tree/main/cub}{CUB BlockReduce 官方文档对块级归约的运算语义；高置信度}
对抗审查：若在同一轮使用刚写回的新值，会得到错误状态；CUDA 屏障把轮次分开，表中每一
行必须只依赖上一行。对非结合运算，两棵树不保证给出相同结果。
\end{solutionexercise}

\chapter{前缀扫描}

\begin{solutionexercise}{1}
\problemheading 对数组 $[4,6,7,1,2,8,5,2]$ 使用 Kogge--Stone 算法执行并行包含式前缀扫描，
写出每一步后的数组中间状态。第 $r$ 轮步幅为 $2^r$；下标 $i\ge2^r$ 的元素把本轮开始时
下标 $i-2^r$ 的值加到自身，所有线程完成该轮后才进入下一轮。

\approachheading 长度为 8，故步幅依次为 1、2、4。步幅 $s$ 时，位置 $i\ge s$ 使用
上一轮的值更新为 $x_i+x_{i-s}$。

\answerheading
\[
\begin{array}{c|rrrrrrrr}
\text{初始}&4&6&7&1&2&8&5&2\\
s=1&4&10&13&8&3&10&13&7\\
s=2&4&10&17&18&16&18&16&17\\
s=4&4&10&17&18&20&28&33&35
\end{array}
\]
最终一行与顺序前缀和逐项一致。每轮必须读取上一轮快照；若原地读取本轮较早 lane 的新值，
例如步幅 1 时位置 2 读取已经更新的位置 1，就会重复计入 4。

\complexity{$N\log_2N-N+1$ 次加法，深度 $O(\log N)$}{$O(N)$ 原地缓冲}{$N$ 个位置并行演化}
\conclusionheading 包含式扫描结果为 $[4,10,17,18,20,28,33,35]$。
\discussionheading
\discussionpoint{扫描处理的是前缀代数而不只是整数加法。}每个输出都概括从序列开头到当前位置的全部输入，所以算法真正要求的是二元运算满足结合律，使括号可以从顺序链改写为并行树。它不要求交换律：字符串连接或矩阵乘法也能扫描，但合并时必须始终让左区段位于右区段之前；若把两个操作数颠倒，整数加法测试仍会通过，字符串结果却立刻乱序。单位元还决定排他扫描的首项以及空输入语义。由此可把同一框架用于最大前缀、逻辑状态传播、仿射函数组合和有限状态机摘要，并把正确性证明统一为“每个节点表示一个连续输入区间”。例如区间摘要若记录仿射映射 $x\mapsto ax+b$，组合次序一错，常数项就会给出立即可见的反例。

\discussionpoint{前缀网络展示了工作量与依赖深度的交换。}Kogge--Stone 每轮把距离 $2^r$ 的前缀传播到更远节点，只需 $\log_2N$ 层，而且经典拓扑把单节点扇出控制在较小范围；代价是大量前缀节点、运算和密集跨距连线。Brent--Kung 或 Blelloch 用更深的上扫、下扫树把工作降到线性数量，布线也更稀疏。不能简单把前者称为“高扇出网络”，更准确的区别是低深度与节点/连线成本。这个谱系最早与并行加法器密切相关，在 GPU 上则对应屏障数、共享事务和能耗之间的同一组权衡。芯片上长连线增加电容与布线拥塞，软件中同一跨距则表现为更多共享读写，这使抽象网络成本能落到可测指标上。

\discussionpoint{原地更新必须隔离相邻轮次。}第 $r$ 轮的每个线程都应读取第 $r-1$ 轮形成的完整状态；若线程先把新值写回共享数组，另一个线程随后可能把这个本轮值误当成旧值，导致一次轮次传播超过规定距离。安全实现可先读入寄存器、在所有读者完成后写回，并用块级屏障隔开轮次，或使用两个缓冲交替。warp 内 shuffle 以同步掩码交换寄存器，跨 warp 仍需共享状态；跨块更不能依靠 \code{__syncthreads()}，而要扫描块和并回加或使用严格的发布--获取协议。这个分层同步模型也是 stencil 时间步和动态规划波前正确性的共同基础。

\discussionpoint{扫描是动态输出空间的地址生成器。}过滤先把谓词变成 0/1，再用排他前缀给每个保留元素分配唯一地址；基数排序用桶标志前缀确定稳定位置，CSR 构建则扫描每行非零数得到行指针。积分图、变长编码和并行内存分配看似不同，也都在询问“我之前一共需要多少资源”。测试因此不能只覆盖一组正整数：非 2 的幂会暴露尾部处理，空输入检验单位元，计数接近类型上限检验溢出，非交换结合运算能发现隐藏换序。把这些性质写成参数化测试后，一个扫描原语才真正可供多种上层算法复用。例如长度 $[3,0,2]$ 的三条记录经排他扫描得到偏移 $[0,3,3]$，同时展示零长度对象共享边界而不发生写冲突。
\variantheading 若改求排他扫描，答案为 $[0,4,10,17,18,20,28,33]$；它可由本题包含式
结果右移一位并在首项填加法单位元 0 得到。
\practiceheading 以 CUB \code{DeviceScan::InclusiveSum} 对照随机与边界规模，并用 Nsight
Compute 比较 Kogge--Stone 教学内核的全局流量和同步开销；浮点测试应允许由树形顺序导致的
末位差异。
\sourcecheck{https://doi.org/10.1109/TC.1973.5009159}{Kogge 与 Stone 的原始并行前缀网络论文；高置信度}
假设审查：这里的运算是加法且要求包含当前位置；排他扫描应右移一位并在位置 0 填单位元。
\end{solutionexercise}

\begin{solutionexercise}{2}
\problemheading 分析图 11.3 的 Kogge--Stone 内核，证明步幅不超过 warp 大小一半时，
控制分歧只发生在每个块的第一个 warp 中。也就是说，warp 大小为 32 时，步幅为
1、2、4、8、16 的各轮会出现控制分歧。图 11.3 的基线内核为：
\begin{lstlisting}[language=C++,firstnumber=1]
__global__ void scan_kernel(float *input, float *output, unsigned int N) {
    __shared__ float buffer_s[SEG_SIZE];
    unsigned int i = blockIdx.x * blockDim.x + threadIdx.x;
    if (i < N) buffer_s[threadIdx.x] = input[i];
    else       buffer_s[threadIdx.x] = 0.0f;
    for (unsigned int stride = 1; stride < blockDim.x; stride *= 2) {
        __syncthreads();
        float temp;
        if (threadIdx.x >= stride)
            temp = buffer_s[threadIdx.x] + buffer_s[threadIdx.x-stride];
        __syncthreads();
        if (threadIdx.x >= stride)
            buffer_s[threadIdx.x] = temp;
    }
    if (i < N) output[i] = buffer_s[threadIdx.x];
}
\end{lstlisting}
二维或多块的全局合并不在本题范围；只分析单个块内条件
\code{threadIdx.x >= stride} 的 warp 内真假分布。

\approachheading 内核条件是 \code{threadIdx.x >= stride}。步幅为 2 的幂，逐 warp
分析真假 lane 的集合。

\answerheading 设 warp 宽 $W=32$，且 $1\le s\le W/2$。第一个 warp 的 lane
$0,\ldots,s-1$ 条件为假，lane $s,\ldots,W-1$ 为真，因此恰有这个 warp 分歧。对任意
后续 warp，其最小线程下标至少为 $W>s$，全部 lane 条件为真，不分歧。

当 $s\ge W$ 且 $s$ 是 $W$ 的整数倍时，前 $s/W$ 个 warp 全假，其余 warp 全真，边界
落在 warp 之间，也不分歧。因此 $s=1,2,4,8,16$ 五轮各只有第一个 warp 分歧。

\conclusionheading 结论成立；总计五个“warp--轮次”分歧事件，每轮都是 warp 0。
\discussionheading
\discussionpoint{一个阈值如何切开 warp。}条件 $t\ge s$ 把线程号轴分成假前缀和真后缀，分歧只会出现在阈值穿过的那个执行组。设 warp 宽为 $W$，当 $0<s<W$ 时，warp 0 同时含有 $s$ 个假 lane 和 $W-s$ 个真 lane，其后完整 warp 全真；当 $s=W$ 时，warp 0 全假而 warp 1 起全真，反而没有混合路径。这个几何视角比逐线程枚举更一般，还能处理 $t<k$ 的尾部谓词和按区间划分的稀疏任务。只要把阈值除以执行组宽度并观察余数，就能预测分歧发生在哪一组以及真假 lane 数量。

\discussionpoint{少量分歧不代表扫描工作高效。}条件 $t\ge s$ 的真线程始终构成后缀 $[s,N)$，数量为 $N-s$；步幅从 1 翻倍到 $N/2$ 时，活跃数从 $N-1$ 降到 $N/2$，最后一轮仍有一半线程做加法，并非只剩很小前缀。因此本算法的问题不是后期几乎没有工作，而是同一批元素跨 $\log_2N$ 轮反复参与，总加法为 $N\log_2N-N+1$，并且每轮都要支付共享通信和块级屏障。branch efficiency 只能说明阈值是否切开 warp，不能据此推出 work efficiency 或同步延迟；与只做 $O(N)$ 运算的 Blelloch 树比较，才看得见低深度网络用额外工作换关键路径的真实代价。类似地，规则树内核应同时报告活跃 lane、总动态指令和 barrier stalls，而不能把“分支很整齐”当成整体高效。

\discussionpoint{谓词化只改变控制实现而不创造有效工作。}编译器可能把短条件体翻译为带谓词的指令，使硬件不执行显式跳转；假 lane 虽不会写结果，仍可能占据指令发射槽，所以源代码没有分支事件不等于没有掩码浪费。warp shuffle 可在执行组内交换前缀并减少共享往返，但要求正确的活动掩码和一致参与；跨 warp 的总和仍要通过共享内存或块级集体原语传播。理解这两层后，就能选择让编译器谓词化短路径、用 shuffle 完成局部网络、再用共享数组连接各 warp，而不是把“无 branch 指令”误当作优化终点。

\discussionpoint{测量必须把硬件事件对应回算法轮次。}一个内核汇总出的分支效率无法告诉我们问题发生在距离 1 还是距离 512 的传播，因此可在实验版本中拆分轮次、插入范围标记或对小规模版本逐轮计时。分析时应同时看活动 lane、已执行指令、屏障等待和共享事务，并以源代码条件推导的阈值位置作预测。证明还依赖块从线程 0 对齐分组、完整 warp 与给定宽度 $W$；尾 warp 或其他平台 subgroup 会改变掩码。这样先预测、再采样、最后解释偏差的过程，比看到一个低百分比便归因于“分歧”更有可证伪性，也可迁移到任何规则阈值内核。
\variantheading 若假设 warp 宽为 64 且算法步幅不超过 32，则仍恰有第一个 warp 分歧；
步幅 32 时它由 32 个假 lane 和 32 个真 lane 构成。
\practiceheading 可在 Nsight Compute 的 Source Counters 中逐轮核对分支与活动线程掩码；
生产扫描通常用 warp shuffle，使 $s<W$ 的数据交换不必经过共享内存，但仍须传入正确掩码。
\sourcecheck{https://docs.nvidia.com/cuda/cuda-c-programming-guide/}{NVIDIA 对 warp 分支路径串行化的官方说明；高置信度}
反例审查：若步幅不是 2 的幂或不与 32 对齐，例如 $s=40$，边界会落在第二个 warp 内；
证明依赖算法的倍增步幅，不能推广到任意整数步幅。
\end{solutionexercise}

\begin{solutionexercise}{3}
\problemheading 图 11.3 的 Kogge--Stone 内核扫描 2048 个元素时，总共执行多少次
加法？在抽象算法计数中令步幅依次为 $1,2,4,\ldots,1024$，每轮由满足
$i\ge\code{stride}$ 的元素执行一次加法。另请注意，实际 CUDA 设备通常不允许单块
启动 2048 个线程，因此实现时还需说明如何分块以及分层扫描会引入的额外加法。

\approachheading 第 $r$ 轮步幅 $2^r$，有 $N-2^r$ 个线程执行加法；$r=0,\ldots,10$。

\answerheading
\[
\sum_{r=0}^{10}(2048-2^r)
=11\cdot2048-(2^{11}-1)
=22528-2047
=\boxed{20481}.
\]
这里统计的是源代码层面的有效加法，不把分歧造成的空闲 lane 计为加法，也不计地址运算。

\editorialnote 20,481 是一个含 2048 个逻辑节点的 Kogge--Stone 网络计数，不是合法的
“单个 CUDA 块含 2048 个线程”启动配置；当前 CUDA 每块最多 1024 个线程。实际处理
2048 项必须采用每线程多项或多块层次化扫描。例如两个 1024 项局部 Kogge--Stone 扫描
先做 $2\times9217=18434$ 次加法，再扫描两个块和并把首块总和加到第二块的 1024 项，
还需 $1+1024$ 次加法，总计 19,459 次；这已是另一种层次化算法，不能与抽象网络计数混同。

\complexity{$O(N\log N)$ 工作、$O(\log N)$ 深度}{$O(N)$ 缓冲}{抽象网络有 $N$ 个逻辑节点；CUDA 实现须受每块线程上限约束并分层}
\conclusionheading 总共执行 20,481 次加法。
\discussionheading
\discussionpoint{从逐轮计数可以看见非工作高效性。}当 $N=2^m$ 时，第 $r$ 轮只有前 $2^r$ 个位置缺少足够远的左前缀，因此执行 $N-2^r$ 次加法。求和得到 $\sum_{r=0}^{m-1}(N-2^r)=N\log_2N-N+1$，它比顺序扫描的 $N-1$ 次多一个对数因子，却把依赖深度降到 $m$。这个公式只计算数学加法，不包含条件、地址、共享读写和屏障，所以不能直接乘某个指令周期预测时间。区分 work 与 depth 后，才能讨论资源无限时的延迟、资源有限时的吞吐以及硬件实际执行成本三种不同评价。

\discussionpoint{抽象网络规模不等于合法线程块规模。}2048 个前缀节点适合纸面推导 20,481 次加法，但常见 CUDA 设备的单块线程上限低于 2048，不能把每个节点直接映射为同一块线程。实际程序会把输入切成至多设备允许的块，先扫描各块，再扫描块总和并把前缀加回；若用两个 1024 项块，额外阶段和加回操作改变了工作量、内存流量与同步位置。因而算法网络计数和某个具体 kernel 的动态指令数必须分别报告。这个边界也提醒我们，PRAM 式无限处理器分析需要经过层次化映射才能成为可执行并行程序。实现还应在运行时查询线程数和共享内存上限，因为同一源码换到资源更小的设备后，合法分块本身就可能变化。

\discussionpoint{低深度网络适合延迟目标，线性工作树适合带宽目标。}Kogge--Stone 对 2048 项只需 11 轮，在运算昂贵、数据已在片上且并行资源充分时，较短关键路径很有吸引力；代价是每项参与多轮传播。Brent--Kung 或 Blelloch 增加轮数，却显著减少总加法和共享事务，对长数组的流式扫描往往更接近带宽上限。若元素是大结构或二元运算本身昂贵，多出的操作代价会进一步放大；若只是极小整数且 launch 主导，网络差异又可能不可见。最佳算法由目标是低延迟还是高吞吐、元素成本和硬件容量共同决定。尤其当操作数是几十字节的状态摘要时，多做一轮不仅多一次算术，还会同步放大片上搬运，不能再用整数加法基准外推。

\discussionpoint{端到端流水线可能奖励融合而不是最快的独立扫描。}基数排序、过滤和动态分配都把扫描夹在生产标志与消费偏移之间；若融合能够让中间值留在寄存器或共享内存，少一次全局往返可能比扫描原语本身快几个百分点更重要。基准应同时列出每元素时间、有效字节、数学加法与 launch 数，并特别测试 1024、1025 和 2048 这类跨越块层次的规模。还要把临时空间分配和同步计入端到端时间，而不是只截取一个 kernel。由此可建立组件微基准与应用基准的对应关系，避免为了孤立排名选择一个使整个流水线更慢的实现。若中间标志和偏移均为 32 位，单独物化就会新增至少两次写读，按每项十余字节估算可先判断融合的收益上限。
\variantheading 若规模改为 $N=1024$，共有 10 轮，加法数为
$10\cdot1024-(2^{10}-1)=\boxed{9217}$。
\practiceheading 用 Nsight Compute 分别验证抽象单块可实现规模和 2048 项层次化版本的
指令计数，并与 CUB \code{DeviceScan} 的吞吐量比较；测试规模应覆盖非 2 的幂，避免只验证
恰好填满扫描树的输入。
\sourcecheck{https://github.com/NVIDIA/cccl/tree/main/cub}{CUB BlockScan 官方文档列出的块级扫描语义与复杂度背景；高置信度}
假设审查：若问的是 Brent--Kung 或工作高效扫描，答案会是线性量级；本题明确指定
Kogge--Stone，不能用 $2N-2-\log_2N$ 代替。
\end{solutionexercise}

\begin{solutionexercise}{4}
\problemheading 修改图 11.13 的内核，使用向量加载与存储，并消除共享内存 bank 冲突。
图 11.13 的完整基线为：
\begin{lstlisting}[language=C++,firstnumber=1]
__global__ void scan_kernel(float *input, float *output, unsigned int N) {
    unsigned int blockSegment = blockIdx.x * COARSE_FACTOR * BLOCK_DIM;
    __shared__ float buffer_s[COARSE_FACTOR * BLOCK_DIM];
    for (unsigned int c = 0; c < COARSE_FACTOR; ++c)
        buffer_s[c * BLOCK_DIM + threadIdx.x] =
            input[blockSegment + c * BLOCK_DIM + threadIdx.x];
    __syncthreads();

    unsigned int threadSegment = threadIdx.x * COARSE_FACTOR;
    float buffer_r[COARSE_FACTOR];
    buffer_r[0] = buffer_s[threadSegment];
#pragma unroll
    for (unsigned int c = 1; c < COARSE_FACTOR; ++c)
        buffer_r[c] = buffer_s[threadSegment + c] + buffer_r[c - 1];

    float threadSum = buffer_r[COARSE_FACTOR - 1];
    threadSum = blockScan(threadSum);
    __shared__ float threadSums[BLOCK_DIM];
    threadSums[threadIdx.x] = threadSum;
    __syncthreads();
    float prevPartialSum = (threadIdx.x > 0)
                           ? threadSums[threadIdx.x - 1] : 0.0f;
#pragma unroll
    for (unsigned int c = 0; c < COARSE_FACTOR; ++c)
        buffer_s[threadSegment + c] = buffer_r[c] + prevPartialSum;
    __syncthreads();
    for (unsigned int c = 0; c < COARSE_FACTOR; ++c)
        output[blockSegment + c * BLOCK_DIM + threadIdx.x] =
            buffer_s[c * BLOCK_DIM + threadIdx.x];
}
\end{lstlisting}
\code{blockScan} 表示块级包含式扫描。题目要求保持全局装载/存储合并，使用向量访问，
并通过共享数组布局或填充避免线程私有连续子段访问造成的 bank 冲突；还应补齐
\code{N} 的尾部边界处理。

\approachheading 令每线程粗化因子 $C$ 为 4 的倍数，$V=C/4$。全局阶段让 warp 以
striped 次序读取相邻的 \code{float4}，再用块交换把它转为每线程持有 $V$ 个连续向量的
blocked 次序；写回时执行逆交换。这样向量事务在全局地址上连续，线程内仍能顺序扫描自己的
$C$ 项，也不会出现朴素 \code{s[thread][item]} 在固定 item 上按 $C$ 跨步的 bank 模式。

\answerheading 下列核心内核用 CUB 的 warp-transpose BlockLoad/BlockStore 实现真实的
striped--blocked 交换；\code{blockInclusive()} 返回每线程总和的块级包含式扫描。完整且
16 字节对齐的分段走 \code{float4}，尾段由同一算法的标量有界路径处理。
\begin{lstlisting}[language=C++]
template<int B, int C>
__global__ void scan_vec_exchange(const float* in,float* out,size_t n) {
    static_assert((B & (B-1)) == 0 && C % 4 == 0, "shape");
    constexpr int V=C/4;
    using VL=cub::BlockLoad<float4,B,V,cub::BLOCK_LOAD_WARP_TRANSPOSE>;
    using VS=cub::BlockStore<float4,B,V,cub::BLOCK_STORE_WARP_TRANSPOSE>;
    using SL=cub::BlockLoad<float,B,C,cub::BLOCK_LOAD_WARP_TRANSPOSE>;
    using SS=cub::BlockStore<float,B,C,cub::BLOCK_STORE_WARP_TRANSPOSE>;
    union X { typename VL::TempStorage vl; typename VS::TempStorage vs;
              typename SL::TempStorage sl; typename SS::TempStorage ss; };
    __shared__ X exchange;
    __shared__ float totals[B];
    const size_t seg = size_t(blockIdx.x) * B * C;
    const int valid = seg<n ? int(min(size_t(B*C),n-seg)) : 0;
    const bool vec = valid==B*C &&
        ((reinterpret_cast<size_t>(in+seg)|
          reinterpret_cast<size_t>(out+seg))&15)==0;
    float r[C];
    if (vec) {
        float4 p[V];
        VL(exchange.vl).Load(reinterpret_cast<const float4*>(in+seg),p);
        for(int q=0;q<V;++q) {
            r[4*q]=p[q].x; r[4*q+1]=p[q].y;
            r[4*q+2]=p[q].z; r[4*q+3]=p[q].w;
        }
    } else {
        SL(exchange.sl).Load(in+seg,r,valid,0.0f);
    }
    __syncthreads();
    for(int c=1;c<C;++c) r[c]+=r[c-1];
    float inclusive=blockInclusive<B>(r[C-1]);
    totals[threadIdx.x] = inclusive;
    __syncthreads();
    const float before=threadIdx.x ? totals[threadIdx.x-1] : 0.0f;
    for(int c=0;c<C;++c) r[c]+=before;
    __syncthreads();
    if (vec) {
        float4 p[V];
        for(int q=0;q<V;++q) {
            p[q]=make_float4(r[4*q],r[4*q+1],r[4*q+2],r[4*q+3]);
        }
        VS(exchange.vs).Store(reinterpret_cast<float4*>(out+seg),p);
    } else {
        SS(exchange.ss).Store(out+seg,r,valid);
    }
}
\end{lstlisting}
BlockLoad 先以 warp-striped 方式取得连续全局事务，再由内部 BlockExchange 返回 blocked
线程私有数组；BlockStore 执行逆过程。\code{vec} 是块内一致分支，union 临时区每次换用途前
都有屏障。标量尾路径以 0 填充无效项并只写前 \code{valid} 项，所以不会越界。

\complexity{$O(BC)$ 工作、块内深度 $O(C+\log B)$}{$O(BC)$ 交换临时区和 $B$ 个块扫描总和}{$B$ 线程，每线程 $C$ 个元素}
\conclusionheading 完整分段使用合并的 \code{float4} warp-striped 事务，经真实布局交换后按线程连续扫描；尾段安全退化为合并标量路径。
\discussionheading
\discussionpoint{布局转换的核心是把数据所有权交给另一个线程。}全局访问希望 lane $t$ 读取与 lane $t+1$ 相邻的位置，这叫 striped 布局；线程内顺序扫描却希望同一线程最终拥有连续 $C$ 项，这叫 blocked 布局。若线程始终只写、只读 \code{s[t][c]}，共享数组只是暂存，所有权没有变化，自然不构成转置。真正的 BlockExchange 先按 striped 次序装入，让不同线程把值放到可交换的共享槽，再按 blocked 下标取回；写回时执行逆过程。这个区别连接到矩阵转置、结构数组变换和 Tensor Core fragment 排列，凡声称“转置”都应能指出元素从哪个线程移交给哪个线程。

\discussionpoint{bank 冲突必须从地址映射公式推导。}共享内存把地址字映射到有限个 bank，同一 warp 的不同地址若落到同一 bank，访问会拆成多批；为二维行增加一列常使行跨度与 bank 数互质，从而把固定列访问置换到不同 bank。以 32 个四字节 bank 和 \code{float} 为例，跨度 $C+1$ 需要结合 $\gcd(C+1,32)$ 分析，不能把“加一永远无冲突”推广到 \code{double}、复数或不同 bank 配置。向量指令还可能分成多个共享事务，每个子事务都要单独检查。参数化推导比记忆 padding 常数更可靠，也适用于 tiled GEMM 和卷积共享布局。

\discussionpoint{向量指令和 warp 合并是两层不同条件。}一次 \code{float4} 要求起始地址 16 字节对齐并且四项都在有效对象内，它可以减少单线程加载指令；但若线程 $t$ 的起点为 $tC$ 且 $C>4$，相邻 lane 的向量之间仍留有间隔，warp 会覆盖更多内存扇区。warp-striped BlockLoad 让 lane 先读取相邻向量，从而获得高事务利用率，再经共享交换恢复每线程连续所有权。尾部、未对齐子视图和行首偏移必须用标量或受保护路径，不能先形成非法向量指针再避免解引用。这个双层分析同样适用于 SIMD gather 与网络批量包。

\discussionpoint{粗化因子会同时占用寄存器、共享容量和并行机会。}增大 $C$ 能让一次块扫描服务更多元素，摊薄同步和块总和处理；每线程却需要保存更长的局部前缀，BlockExchange 临时空间也随项目数增长。寄存器不足时数组可能落入 local memory，较大共享空间又减少每 SM 驻留块数，结果是内存延迟更难隐藏。现代实现可比较直接 blocked、warp-striped 转置和异步搬运三类策略，并在相同边界语义下观察每元素时间、global sectors、bank 冲突、寄存器和 occupancy。只有这些量随 $C$ 的曲线一致，才能判断收益来自少指令、好合并还是阶段融合。
\variantheading 若 $C=12$，则 $V=3$：每个 lane 在全局阶段参与读取连续的 warp-striped
\code{float4}，交换后得到属于自己的 3 个连续向量；若省略交换而从 \code{t*C} 直接加载，
相邻 lane 起点相隔 48 字节，无法得到与 $C=4$ 相同的事务利用率。
\practiceheading 用 Nsight Compute 的 Shared Load/Store Bank Conflicts 分别测
$C=4,8,12,16$，并用 CUB \code{BlockScan} 校验结果；启动配置还须满足 $B$ 为 2 的幂、
$BC$ 的交换临时区不超过设备上限；随附可执行代码覆盖 $C=12$ 和 997 项尾块。
\sourcecheck{https://github.com/NVIDIA/cccl/tree/main/cub}{NVIDIA CCCL/CUB 的 BlockLoad、BlockStore 与 BlockExchange 官方实现；对抗检查：只有 striped 装载后发生真实所有权交换，才能同时得到高事务利用率与 blocked 线程布局}
边界审查：\code{float4} 要求 16 字节对齐且整段有效；否则统一退回有 \code{valid} 计数的
标量 warp-transpose 路径。共享交换的 bank 行为还应在目标架构实测，不能由二维数组外形臆断。
\end{solutionexercise}

\begin{solutionexercise}{5}
\problemheading 修改图 11.17 的 \code{interBlockScan} 设备函数，用解耦回看替代
单次回看。图 11.17 的完整单次回看基线为：
\begin{lstlisting}[language=C++,firstnumber=1]
__device__ inline float interBlockScan(float val, unsigned int bid,
                                       float *partialSums,
                                       unsigned int *flags) {
    __shared__ float prevBlockPartialSum_s;
    if (threadIdx.x == blockDim.x - 1) {
        if (bid > 0) {
            cuda::atomic_ref<unsigned int, cuda::thread_scope_device>
                prevBlockFlag_ref(flags[bid - 1]);
            while (prevBlockFlag_ref.fetch_add(
                       0, cuda::memory_order_acquire) == 0) {}
            prevBlockPartialSum_s = partialSums[bid - 1];
        } else {
            prevBlockPartialSum_s = 0.0f;
        }
        partialSums[bid] = prevBlockPartialSum_s + val;
        cuda::atomic_ref<unsigned int, cuda::thread_scope_device>
            myBlockFlag_ref(flags[bid]);
        myBlockFlag_ref.fetch_add(1, cuda::memory_order_release);
    }
    __syncthreads();
    return prevBlockPartialSum_s;
}
\end{lstlisting}
其中 \code{bid} 必须在线程块开始执行后动态分配，以避免后继块占满设备并等待尚未调度
的前驱块。解耦回看版本应区分“块聚合值已就绪”和“完整前缀已就绪”两类状态：向前
跨过只有聚合值的块时累加其块和，遇到已有完整前缀的块时终止回看，并以设备作用域的
release/acquire 发布和读取状态。

\approachheading 每块先发布自己的块和；向前回看时，遇到“块和就绪”便累加并继续，
遇到“前缀就绪”便累加该前缀并停止。块和与完整前缀分开存放，避免同一非原子位置在读写
间发生数据竞态。

\answerheading 状态 0、1、2 分别表示 EMPTY、AGGREGATE、PREFIX。动态块号必须按块开始
执行的先后原子分配。
\begin{lstlisting}[language=C++]
#include <cuda/atomic>
enum : unsigned { EMPTY=0, AGGREGATE=1, PREFIX=2 };

__device__ float decoupledLookback(float blockSum, unsigned bid,
    float* aggregates, float* prefixes, unsigned* states) {
  __shared__ float exclusive_s;
  if (threadIdx.x == blockDim.x-1) {
    aggregates[bid] = blockSum;
    cuda::atomic_ref<unsigned, cuda::thread_scope_device> mine(states[bid]);
    mine.store(AGGREGATE, cuda::memory_order_release);

    float before = 0.0f;
    int p = int(bid)-1;
    while (p >= 0) {
      cuda::atomic_ref<unsigned, cuda::thread_scope_device> st(states[p]);
      unsigned kind;
      do { kind=st.load(cuda::memory_order_acquire); } while(kind==EMPTY);
      if (kind == PREFIX) { before += prefixes[p]; break; }
      before += aggregates[p--];
    }
    exclusive_s = before;
    prefixes[bid] = before + blockSum;
    mine.store(PREFIX, cuda::memory_order_release);
  }
  __syncthreads();
  return exclusive_s;
}
\end{lstlisting}
不变量是 \code{before} 等于已越过连续块的块和；遇到完整前缀 $P_p$ 后，
$P_p+\sum_{q=p+1}^{bid-1}A_q$ 正好是当前块的排他基址。release/acquire 使状态发布前的
数据写入对等待块可见。

\complexity{最坏 $O(G^2)$ 回看工作，通常远低于单次链式等待}{$O(G)$ 状态、块和与前缀空间}{$G$ 块可重叠回看}
\conclusionheading 解耦回看允许跳过尚未形成前缀的前驱块，以冗余加法换取更短等待路径。
\discussionheading
\discussionpoint{状态值其实是一份跨块可见性契约。}AGG 表示该块的局部总和已经写好，PREFIX 表示从开头到该块的完整前缀已经写好；消费者看到不同状态，允许读取的载荷也不同。生产者必须先写普通数组，再用 device-scope release 存储发布状态；消费者以 acquire 读到状态后，前面的载荷写才与其建立 happens-before。若先写状态再写数值，或只把标志声明为 \code{volatile}，另一个 SM 可能观察到“就绪”却读到旧载荷。这个消息传递模式与无锁队列、GPU 任务图和 CPU 原子发布完全同构，原子性与内存排序缺一不可。

\discussionpoint{跨块等待首先受调度前向进度约束。}若逻辑编号直接等于静态 \code{blockIdx.x}，调度器可能先让高编号块占满所有驻留槽，而它们又自旋等待尚未获准运行的低编号前驱，于是形成资源死锁。让块在真正开始执行后原子领取递增逻辑号，可以保证较小编号至少已经启动；不过它不承诺固定时间内获得调度，自旋次数仍可能很大。生产级持久 kernel 通常限制驻留块数量、在轮询中降低资源压力并保留诊断。这个案例说明无锁并不等于无等待，算法证明必须包含执行平台的进度保证。例如设备只能驻留 8 块而先调度逻辑 8--15 时，八个等待者就能共同阻止逻辑 0 获得任何执行槽。

\discussionpoint{解耦回看用冗余相加换取不必串行等待。}若前驱只有 AGG，当前块无需等它算出完整 PREFIX，可以把该块总和加入临时累计并继续向左；一旦遇到某个 PREFIX，就用它一次覆盖更长历史并停止。多个后继可能重复累加同一组 AGG，所以总工作最坏可达平方量级，状态轮询也产生缓存一致性流量；换来的好处是多数块能与前缀传播重叠，不形成逐块接力。它体现了并发算法常见的空间--工作--等待三方交换，与路径压缩、帮助完成和乐观并发中的冗余工作具有相似思想。以四块全都只发布 AGG 的瞬间为例，末块可能依次累加前三块，而相邻块也重复走过其中一段，冗余来源由此清晰可见。

\discussionpoint{单遍扫描是许多融合算法的基础设施。}过滤需要前缀来分配输出地址，基数排序需要桶偏移，动态编码需要为变长记录预留空间；解耦回看让这些步骤有机会在一次输入遍历中完成。协议复杂且依赖微妙内存序，所以工程上优先使用成熟实现，并用低 occupancy、随机延迟和竞争检查主动扩大进度与可见性缺陷。若元素是浮点数，不同块观察 PREFIX 的时机还可能改变合并括号，从而破坏逐位确定性；需要固定结果时应采用固定层次树或明确误差契约。这里性能优化最终会变成 API 语义选择，而不只是少一次 launch。
\variantheading 若逻辑块号 $bid=0$，回看循环立即结束，排他基址为 0，并发布
\code{prefixes[0]=blockSum}；这是协议的启动边界。
\practiceheading 生产中优先使用 CUB \code{DeviceScan} 已验证的解耦回看实现；自定义版本
应在不同驻留块限制下做压力测试，并用 Compute Sanitizer racecheck 检查状态与载荷的
release/acquire 配对。
\sourcecheck{https://research.nvidia.com/publication/2016-03_single-pass-parallel-prefix-scan-decoupled-look-back}{Merrill 与 Garland 的 NVIDIA 原始技术报告；高置信度}
对抗审查：若让 \code{aggregates} 与 \code{prefixes} 共用一个普通 \code{float} 数组，
读块和与覆写前缀可能并发，属于数据竞态。静态 \code{blockIdx.x} 还可能使后继块占满设备
并等待尚未调度的前驱块；必须动态分配逻辑块号。
\end{solutionexercise}

\begin{solutionexercise}{6}
\problemheading 对数组 $[4,6,7,1,2,8,5,2]$ 使用 Brent--Kung 算法执行并行包含式前缀扫描，
写出每一步后的数组中间状态。先按距离 1、2、4 构造归约树，再按距离 2、1 执行反向树；
每轮更新必须读取该轮开始时的值。

\approachheading 先做归约树（距离 1、2、4），再做反向树（距离 2、1）。所有更新都使用
本轮开始时的值。

\answerheading
\[
\begin{array}{c|rrrrrrrr}
\text{初始}&4&6&7&1&2&8&5&2\\
\text{归约 }1&4&10&7&8&2&10&5&7\\
\text{归约 }2&4&10&7&18&2&10&5&17\\
\text{归约 }4&4&10&7&18&2&10&5&35\\
\text{反向 }2&4&10&7&18&2&28&5&35\\
\text{反向 }1&4&10&17&18&20&28&33&35
\end{array}
\]
归约树产生区间和；反向距离 2 更新位置 5，距离 1 更新位置 2、4、6。总加法数为
$2N-2-\log_2N=16-2-3=11$。

\complexity{$O(N)$ 工作，$2\log_2N-1$ 轮}{$O(N)$ 原地缓冲}{每轮最多 $N/2$ 个并行加法}
\conclusionheading 最终扫描结果同样是 $[4,10,17,18,20,28,33,35]$。
\discussionheading
\discussionpoint{Brent--Kung 通过稀疏节点换取更长路径。}上扫阶段只让每个长度 $2^r$ 区段的末端保存总和，因此许多位置暂时没有完整前缀；反向树再把左侧已知区段传播到这些缺口。对 $N$ 项，它把工作控制在线性数量附近，却需要约 $2\log_2N-1$ 层。Kogge--Stone 只有 $\log_2N$ 层且保持受限扇出，但创建更多前缀节点和密集连线；Brent--Kung 的主要优势是节点和布线较少，不能简化为“最大扇出更低”。这种网络取舍既影响芯片加法器布线，也映射为 GPU 的加法数、共享事务和屏障深度。

\discussionpoint{渐近工作相同仍可能对应不同硬件瓶颈。}Brent--Kung 少做加法，通常也少读写共享内存，适合带宽或能耗受限的大块；反向阶段增加轮次，每轮又需要同步，会拉长单个块的关键路径。若二元运算本身昂贵，省一次运算可能远胜多一个屏障；若输入很小且延迟敏感，浅网络反而更快。还要考虑块数量是否足以让其他 warp 遮蔽屏障等待，不能只从 $O(N)$ 与 $O(N\log N)$ 下结论。用 roofline 加同步成本的简化模型，可以把算法网络参数转成设备上的可测假设。实验中还应固定块数与数据位置，否则缓存热度和可并发块数的变化会被误认成网络拓扑收益。

\discussionpoint{包含式和排他式共享骨架但边界语义不同。}包含式第 $i$ 项含 $x_i$，排他式只含其左侧输入，所以可把包含式结果右移一位并在开头填单位元；也可在树的根和下扫规则中直接构造排他结果，省去额外搬移。对加法，两种实现很难暴露操作数次序；对字符串连接，反向阶段必须把左区段放在当前区段之前，否则前缀内部会倒序。空数组、单元素和没有单位元的运算还需要 API 明确定义。理解这些差别后，同一树才能安全泛化到函数组合、区间变换和分段扫描，而不只是数值数组。例如输入字符串 $[a,b,c]$ 的排他结果应为 $[\epsilon,a,ab]$，它能同时检查单位元、方向与移位边界。

\discussionpoint{逐轮状态可以成为可执行规范。}题目列出的五轮数组不仅用于手算，它准确规定了每一屏障前后哪些位置允许变化、每个新值由哪些旧值组成。CPU 模拟器可复制同一轮次并把每轮状态作为黄金断言，GPU 调试版本则只在小输入时保存快照；一旦偏离，就能定位错误下扫下标、缺失同步或误用本轮新值。进一步比较 Kogge--Stone、Brent--Kung 与库实现时，应同时统计数学加法、屏障、共享事务、时间和浮点误差。这样建立的选择图谱比单一吞吐排名更完整，也让网络图真正成为工程验证工具。对随机大输入不必保存全部快照，可逐轮计算哈希并在首个不一致轮次缩小复现规模，降低调试数据量。
\variantheading 若要求同一输入的排他结果，可把包含式结果右移并填 0，得到
$[0,4,10,17,18,20,28,33]$；上扫与反向树的中间状态本身不必改写。
\practiceheading 将每轮状态作为单元测试断言，可定位遗漏屏障或错误下扫下标；部署时用
CUB \code{BlockScan} 作基线，并在 Nsight Compute 中比较同步停顿和共享内存流量。
\sourcecheck{https://doi.org/10.1109/TC.1982.1675982}{Brent 与 Kung 的原始并行加法器论文；高置信度}
假设审查：反向树不能从距离 4 再更新最终元素；位置 7 已是完整前缀。把经典排他扫描的
下扫交换步骤直接套到这里，会得到不同的中间状态。
\end{solutionexercise}

\begin{solutionexercise}{7}
\problemheading 实现一个在块级扫描数组的 Brent--Kung CUDA 内核。内核应产生包含式
前缀结果，处理不足一个块的尾段，并明确块大小为 2 的幂这一前提；还应说明单个块级
内核为何不能单独完成跨多个线程块的全数组扫描。

\approachheading 用共享内存保存一个长度为 2 的幂的块区段。上扫阶段更新每个分组末端，
下扫阶段补齐其余位置；尾块以单位元 0 填充。

\answerheading
\begin{lstlisting}[language=C++]
template<int B>
__global__ void brent_kung_scan(const float* in, float* out, size_t n) {
  static_assert((B & (B-1)) == 0, "B must be a power of two");
  __shared__ float x[B];
  size_t g = size_t(blockIdx.x)*B + threadIdx.x;
  x[threadIdx.x] = (g<n) ? in[g] : 0.0f;
  __syncthreads();

  // Reduction tree.
  for (unsigned stride=1; stride<B; stride<<=1) {
    unsigned i=(threadIdx.x+1)*(stride<<1)-1;
    if (i<B) x[i] += x[i-stride];
    __syncthreads();
  }
  // Reverse tree for inclusive prefixes.
  for (unsigned stride=B>>2; stride>0; stride>>=1) {
    unsigned i=(threadIdx.x+1)*(stride<<1)+stride-1;
    if (i<B) x[i] += x[i-stride];
    __syncthreads();
  }
  if (g<n) out[g]=x[threadIdx.x];
}
// brent_kung_scan<1024><<<(n+1023)/1024,1024>>>(in,out,n);
\end{lstlisting}
归纳可证上扫后位置 $2^k-1$ 保存相应前缀区间和；下扫从大距离到小距离，把紧邻左侧的
完整区间加到缺失前缀的位置。每个位置每轮至多一个写者，屏障隔开读写轮次。

\complexity{每满块 $2B-2-\log_2B$ 次加法，$O(B)$ 工作}{$B$ 个共享元素}{$B$ 线程，深度 $2\log_2B-1$}
\conclusionheading 内核正确产生每块包含式扫描；全数组扫描还需扫描块和并把块前缀加回。
\discussionheading
\discussionpoint{全数组前缀可以分解为块外与块内两部分。}单块内核得到的只是相对于块首的局部前缀；全局第 $i$ 项还需要组合所有前块的总和。标准三步是保存每块总和、扫描这些总和，再把相应排他块前缀组合回块内每项；块数仍太多时，中间扫描可以递归。若两块局部结果分别为 $L_0,L_1$，第二块每项应变成 $s_0\mathbin\circ L_1[i]$，其中 $s_0$ 是第一块总和，而第一块保持原值。正确性不变量正是“全局前缀等于前块摘要前缀与当前块局部前缀的有序组合”，因此非交换运算也必须让前块结果在左。这个分解与 MPI Scan、分层时钟和分布式日志偏移同构，展示了局部知识如何通过小型摘要扩展到全局。

\discussionpoint{尾部填充值由代数单位元决定。}最后一个块不足 $B$ 项时，把无效共享槽填 0 能保持所有线程执行相同上扫、下扫流程，因为 0 不改变加法结果；最大值应填负无穷，逻辑与则填真。写回仍必须屏蔽虚拟位置，避免越界。若 $B$ 不是 2 的幂，当前索引公式描述的规则树不再覆盖所有叶子，可以把容量提升到下一幂并填单位元，也可换任意长度算法。由单位元统一边界比在每轮插入不同分支更容易证明，但无单位元的运算需要显式有效标志，不能捏造中性值。例如全为负数的最大前缀若误填 0，虚拟叶就会胜过真实数据，使边界错误污染整个尾块。

\discussionpoint{块大小同时控制上层规模和片上占用。}增大 $B$ 会把上层块和数组规模降为 $\lceil N/B\rceil$，并可能减少递归层；每块却需要更多线程、共享槽和屏障参与者，寄存器总量也可能限制同时驻留块数。若 $B=1024$ 只允许一个块驻留，较小的 256 线程块反而可能用更多并发隐藏内存延迟；若二元运算昂贵，大块减少上层工作又可能获益。常见的 256 或 512 只是初始候选，应针对元素宽度、操作成本和输入规模测量。这个多级权衡与 GEMM tile、排序批次和网络消息大小调优具有相同结构，而且尾块利用率也要纳入比较，不能只测整除规模。

\discussionpoint{结果类型和数值契约决定生产可用性。}浮点树的加法顺序不同于 CPU 顺序循环，块边界变化还会改变高层配对，所以默认只能承诺某个误差容限而非逐位相同；需要确定性时应固定分块与编译模式，需要准确度时可提高累加精度。整数扫描常用于过滤位置、基数桶偏移和稀疏行指针，其风险不是舍入而是总数超过 32 位后静默回绕，必须从最大输出容量选择类型。块级结果可与 BlockScan 对照，全局结果可与 DeviceScan 对照，再用应用不变量验证消费端地址。算法正确、数值语义和容量安全三者都通过后，扫描才是可靠基础设施。
\variantheading 若 $B=512,n=700$，需两个块：首块扫描 512 项，次块扫描其余 188 项并以
0 填充；在扫描两个块和并回加前，第二块输出仍缺少首块总和。
\practiceheading 用 CUB \code{BlockScan} 和 \code{DeviceScan} 分别验证块内与全局两层结果，
并用 Compute Sanitizer synccheck 覆盖尾块；启动前查询设备是否允许所选 $B$ 个线程。
\sourcecheck{https://github.com/NVIDIA/cccl/tree/main/cub}{CUB BlockScan 官方块级接口与同步约束；高置信度}
边界审查：该代码只完成块级扫描，不能把不同块误当作已全局同步。输入运算必须有单位元；
浮点加法的树形顺序可能与 CPU 顺序累加有末位差异。
\end{solutionexercise}

\chapter{过滤}

\begin{solutionexercise}{1}
\problemheading 使用三个独立内核实现带排他扫描的稳定过滤：条件求值、排他扫描和输出生成。

输入为 $x_0,\ldots,x_{N-1}$，谓词产生 $f_i\in\{0,1\}$；排他扫描应得到
$p_i=\sum_{j<i}f_j$，输出生成阶段在 $f_i=1$ 时把 $x_i$ 写到 $y_{p_i}$，并返回
保留元素总数。实现应保持输入中被保留元素的相对次序，并处理空输入与尾块边界。

\approachheading 对输入 $x_i$ 产生标志 $f_i\in\{0,1\}$，令
$p_i=\sum_{j<i}f_j$。若 $f_i=1$，唯一目的位置就是 $p_i$。中间扫描采用生产级 CUB
实现，避免在本题中重复实现任意长度全局扫描。

\answerheading 以下程序片段以“保留非零值”为具体谓词；替换 \code{keep()} 即可使用
其他条件。
\begin{lstlisting}[language=C++]
#include <cuda_runtime.h>
#include <cub/device/device_scan.cuh>
#include <climits>
#include <stdexcept>

__device__ __forceinline__ int keep(int x) { return x != 0; }
__global__ void eval(const int* x, int* flags, size_t n) {
  size_t i=size_t(blockIdx.x)*blockDim.x+threadIdx.x;
  if(i<n) flags[i]=keep(x[i]);
}
__global__ void scatter(const int* x,const int* flags,const int* pos,
                        int* y,size_t n,int* outSize) {
  size_t i=size_t(blockIdx.x)*blockDim.x+threadIdx.x;
  if(i<n) {
    if(flags[i]) y[pos[i]]=x[i];
    if(i==n-1) *outSize=pos[i]+flags[i];
  }
}

// Device arrays flags[n], pos[n], outSize[1] are already allocated.
void stable_filter(const int* x,int* y,size_t n,int* flags,int* pos,
                   int* outSize,void* temp,size_t tempBytes) {
  constexpr unsigned B=256;
  if(n==0) { cudaMemset(outSize,0,sizeof(*outSize)); return; }
  if(n>size_t(INT_MAX)) throw std::length_error("int positions overflow");
  size_t G64=(n+B-1)/B;
  int device; cudaDeviceProp prop;
  cudaGetDevice(&device); cudaGetDeviceProperties(&prop,device);
  if(G64>size_t(UINT_MAX) || G64>size_t(prop.maxGridSize[0]))
    throw std::length_error("grid.x overflow");
  unsigned G=static_cast<unsigned>(G64);
  eval<<<G,B>>>(x,flags,n);                         // stage 1
  cub::DeviceScan::ExclusiveSum(temp,tempBytes,flags,pos,n); // stage 2
  scatter<<<G,B>>>(x,flags,pos,y,n,outSize);        // stage 3, including size
}
// Query tempBytes once with temp=nullptr using the same CUB call.
\end{lstlisting}
若 $i<j$ 且二者都保留，则 $p_j-p_i=\sum_{k=i}^{j-1}f_k\ge1$，所以目的下标严格递增，
既不冲突又保持输入相对次序。主机在复用或释放临时存储前必须等待同一 stream 完成。

\complexity{$O(N)$ 工作、扫描深度 $O(\log N)$}{$2N+O(1)$ 个整数及 CUB 临时空间}{$N$ 个条件与散射任务}
\conclusionheading 排他扫描给出唯一且有序的输出位置，因而实现稳定过滤。
\discussionheading
\discussionpoint{稳定过滤把真假判断转化为地址分配。}谓词为真可看成当前元素申请一个输出槽，为假则申请零个；对这些 0/1 做排他扫描，$p_i=\sum_{j<i}f_j$ 就是第 $i$ 项之前已经批准的槽数。若 $i<j$ 且二者保留，则 $p_i<p_j$，所以写地址唯一并保持原相对次序，这正是稳定性的核心证明。把“一个槽”改成每条记录的可变字节数，同一方法就成为压缩编码和序列化的偏移生成器；换成每行非零数，又得到稀疏矩阵行指针。扫描因此不仅求累计量，更是在并行任务间无冲突地切分输出空间。以 flags $[1,0,1,1]$ 为例，排他秩 $[0,1,1,2]$ 让三个真项依次写入 0、1、2，唯一性和稳定性可同时观察。

\discussionpoint{三阶段设计以额外流量换取可组合性。}条件内核只负责 flags，扫描只把 flags 变成 positions，散射再消费两者；每个阶段都有清楚输入输出，可以单独与 CPU 参考比较，也能在多个消费者间复用位置。代价是两个长度 $N$ 的数组被写入又读回，当谓词只是一次整数比较时，这些字节往往比算术更昂贵。若谓词本身代价很高或 positions 还供后续索引使用，分阶段成本则更容易被摊薄。这个取舍与编译器循环融合、数据库物化中间表相同：融合减少通信，物化增加可观察性与复用，不能只按 launch 数决定。

\discussionpoint{容量和异步生命周期是算法接口的一部分。}最坏情况下全部元素保留，因此调用方若一次完成散射，输出缓冲必须容纳 $N$ 项；位置类型也要表示 $N-1$ 和总长度，使用 32 位 \code{int} 时应拒绝超出 \code{INT_MAX} 的规模。设备端 \code{outSize} 只有在所属 stream 完成后主机才能安全读取，提前据它分配精确缓冲会形成时序错误。生产流水线通常预留最大容量、使用内存池，或先只计数再进行第二阶段精确分配。由此可见过滤返回的不只是数据，还包括容量、类型和完成事件三项契约。

\discussionpoint{选择率把同一算法变成不同的性能问题。}数据库谓词下推可能只保留百分之一，图 frontier 有时几乎全保留，粒子淘汰又可能随时间步剧烈变化；相同代码的散射事务和输出字节因此完全不同。若选择率为 $r$、记录宽度为 $s$，理想输出流量约为 $rNs$，但谓词与扫描仍要触及全部 $N$ 项，所以低 $r$ 并不会让总时间同比缩小。基准应从 0\% 扫到 100\% 选择率，并把谓词成本、扫描流量与输出写入分开观察，再用携带原始下标的记录检查稳定次序。重复值不足以验证稳定性，因为交换两个相等值在数值上不可见。这样的数据集设计能把过滤原语扩展为数据库、图计算和流处理系统共同可比较的基础操作。
\variantheading 若谓词改为“保留偶数”，输入 $[3,4,2,5]$ 的 flags 为 $[0,1,1,0]$、排他
位置为 $[0,0,1,2]$，输出为 $[4,2]$，长度为 2。
\practiceheading 生产代码可直接采用 CUB \code{DeviceSelect}，并用 $N=0,1,257$ 及全保留、
全丢弃测试校验；当 \code{pos} 为 \code{int} 时拒绝 $N>\code{INT_MAX}$，网格转为
\code{unsigned} 前还要查询 \code{maxGridSize[0]}。
\sourcecheck{https://github.com/NVIDIA/cccl/tree/main/cub}{CUB DeviceScan 官方排他扫描接口；高置信度}
假设审查：接口层严格只有条件、扫描、生成三个阶段，输出长度已并入生成阶段；CUB 仍可能在
扫描阶段内部启动多个内核。若要求恰好三次底层 launch，应换成单次扫描实现并证明其前向进度。
\end{solutionexercise}

\begin{solutionexercise}{2}
\problemheading 把过滤代码加入单遍扫描内核，以单个内核实现稳定过滤。单内核实现有哪些主要
优点？

仍采用第 1 题的稳定过滤定义：谓词标志的排他前缀给出保留元素目的位置。单内核版本
应在块内完成局部扫描，通过正确的跨块前缀协议取得此前逻辑区段的保留数，然后直接
散射输出并给出总长度；必须说明如何避免仅按静态 \code{blockIdx.x} 等待前驱造成的
前向进度问题。

\approachheading 每块先做局部排他扫描，再用解耦回看取得此前逻辑块的保留总数，最后
直接散射。逻辑块号在块开始执行后动态分配，使较小编号一定已有执行机会。

\answerheading 下列内核使用 CUB 块扫描和三状态回看；\code{aggregates}、
\code{prefixes}、\code{states} 长度均为 $\lceil N/B\rceil$，启动前清零
\code{states} 与 \code{counter}。
\begin{lstlisting}[language=C++]
#include <cub/block/block_scan.cuh>
#include <cuda/atomic>
enum : unsigned { EMPTY=0, AGG=1, PREFIX=2 };

__device__ int publishLookback(int total,unsigned bid,int* aggregates,
    int* prefixes,unsigned* states) {
  __shared__ int base;
  if(threadIdx.x==0) {
    aggregates[bid]=total;
    cuda::atomic_ref<unsigned,cuda::thread_scope_device> me(states[bid]);
    me.store(AGG,cuda::memory_order_release);
    int sum=0;
    for(int p=int(bid)-1;p>=0;) {
      cuda::atomic_ref<unsigned,cuda::thread_scope_device> st(states[p]);
      unsigned k;
      do{k=st.load(cuda::memory_order_acquire);}while(k==EMPTY);
      if(k==PREFIX){sum+=prefixes[p];break;}
      sum+=aggregates[p--];
    }
    base=sum; prefixes[bid]=sum+total;
    me.store(PREFIX,cuda::memory_order_release);
  }
  __syncthreads();
  return base;
}

template<int B>
__global__ void stable_filter_onepass(const int* x,int* y,size_t n,
    unsigned* counter,int* aggregates,int* prefixes,unsigned* states,
    int* outSize) {
  using Scan=cub::BlockScan<int,B>;
  __shared__ typename Scan::TempStorage scan;
  __shared__ unsigned bid_s;
  if(threadIdx.x==0) bid_s=atomicAdd(counter,1u);
  __syncthreads();
  unsigned bid=bid_s; size_t i=size_t(bid)*B+threadIdx.x;
  int value=(i<n)?x[i]:0, flag=(i<n && value!=0), rank,total;
  Scan(scan).ExclusiveSum(flag,rank,total);
  int base=int(publishLookback(total,bid,aggregates,prefixes,states));
  if(flag) y[base+rank]=value;
  unsigned segments=unsigned((n+B-1)/B);
  if(threadIdx.x==0 && bid==segments-1) *outSize=base+total;
}
\end{lstlisting}
局部 rank 保持块内次序；动态 \code{bid} 同时决定输入区段与全局基址，因此块间也保持
输入次序。每个保留元素的目的位置唯一。主要优点是只读一次输入、只写一次保留输出，省去
两个 $N$ 长中间数组和阶段间全局流量，也减少启动次数。

\complexity{$O(N)$ 常规工作；回看最坏 $O(G^2)$}{$O(G)$ 跨块状态与 CUB 块临时空间}{$G$ 块并行并沿前缀依赖推进}
\conclusionheading 单内核融合降低流量和 launch 开销，但增加跨块协议、状态空间及调试复杂度。
\discussionheading
\discussionpoint{单遍融合没有消灭前缀依赖，只是搬了同步边界。}三阶段方案把 flags 和 positions 写到全局，并依靠 kernel 边界保证全设备完成；融合方案让标志与局部秩留在寄存器或共享内存，只输出保留项和每块少量状态。省掉的是 $O(N)$ 中间流量和若干 launch，仍需知道前面逻辑块一共保留多少项，才能形成全局基址。因此全局同步被替换为运行中的状态发布和回看协议，复杂度从显式阶段进入并发控制。即使某块本地保留 17 项，它也不能仅凭这个数写回，因为前块可能已占据 0 个或数千个槽；全局偏移仍是不可消除的数据依赖。这个变化与流式数据库算子融合相似：减少物化会提升局部性，却把背压和顺序保证变成在线协议。

\discussionpoint{动态逻辑编号解决的是一种资源死锁。}若静态高编号块先占满设备并自旋等待未调度的低编号前驱，前驱永远拿不到驻留资源，算法就无法前进。块开始执行后再原子领取递增逻辑号，可以保证被等待的较小逻辑块至少已经启动；逻辑号同时决定输入区段，因而稳定顺序仍与编号一致。不过 CUDA 不承诺每个已启动块在固定时间内公平运行，自旋次数不能写成常数，资源压力仍需控制。若设备只能驻留 8 块而恰由静态 8--15 号先占满，这八块等待 0--7 号的状态就形成具体反例。生产级持久化任务通常限制网格为可驻留规模，并加入超时诊断，这与线程池中的有序任务调度有相同风险。

\discussionpoint{发布状态必须携带载荷的可见性。}块先写 aggregate 或完整 prefix，再用 device-scope release 原子把状态从 EMPTY 改为 AGG/PREFIX；回看块用 acquire 读到新状态后，才能合法消费对应普通数组。若把顺序颠倒，状态可能先跨缓存可见，而数值仍是旧内容；\code{volatile} 只能限制部分编译优化，并不建立跨线程块的 happens-before。压力测试偶尔成功只说明某次时序有利，不能替代内存模型证明。相同的发布--获取模式广泛用于 GPU 工作队列、无锁环形缓冲和 CPU 消息传递，因此本题也是理解现代并发正确性的具体入口。

\discussionpoint{单遍方案最适合通信成本显著的长流水。}大数组配廉价谓词时，中间 flags 和 positions 的两次全局往返占比较高，融合通常有明确收益；昂贵谓词可能掩盖这部分流量，小输入又可能主要由 launch 延迟决定。成熟库还优化回看窗口、缓存布局和架构差异，自定义实现应在相同选择率与稳定语义下和 DeviceSelect 比较，而不是只与朴素三阶段代码比较。低 occupancy、随机延时和全保留/全丢弃数据能主动放大死锁与边界错误。只有吞吐、稳定性和前向进度都验证后，少一次遍历才是可接受的工程收益。
\variantheading 若 $B=256,N=513$，逻辑区段数为 3；最后区段只有一个有效元素，但仍由
256 个线程共同执行块扫描，最后逻辑块写出的 \code{base+total} 才是总长度。
\practiceheading 与 CUB \code{DeviceSelect::Flagged} 对照输出稳定性和吞吐量，并用受限
占用率压力测试动态编号协议；调用方对 $N=0$ 应直接清零 \code{outSize}，且在使用
\code{int} 前缀时限制 $N\le\code{INT_MAX}$。
\sourcecheck{https://research.nvidia.com/publication/2016-03_single-pass-parallel-prefix-scan-decoupled-look-back}{NVIDIA 解耦回看原始报告；高置信度}
对抗审查：直接按 \code{blockIdx.x} 等待前块可能死锁；把块和与前缀存入同一普通位置也
会发生竞态。输出大小可在内核完成后读取，不能把最后逻辑块发布大小误当作整个内核已结束。
\end{solutionexercise}

\begin{solutionexercise}{3}
\problemheading 对练习 2 的内核应用私有化。

练习 2 的基线是单内核稳定过滤：每块对谓词标志做局部排他扫描，取得跨块前缀后把
保留元素直接散射到全局输出。本题要求先在每块的私有共享缓冲中按局部秩稳定打包，
再由块内线程协作把连续区段写到全局输出；说明所需同步、共享内存容量和稳定性依据。

\approachheading 用共享内存建立每块私有输出。直接使用排他 rank 时，同一 warp 的有效
线程已经写向单调且紧凑的地址；共享 staging 的额外价值是把整个块的保留项重新交给从 0
开始的稠密线程前缀冲刷，从而合并跨稀疏 warp 边界的短区段，而不是把真正随机地址变连续。

\answerheading 下列内核沿用上一题的 \code{publishLookback()} 和状态数组协议：
\begin{lstlisting}[language=C++]
template<int B>
__global__ void stable_filter_private(const int* x,int* y,size_t n,
    unsigned* counter,int* aggregates,int* prefixes,unsigned* states,
    int* outSize) {
  using Scan=cub::BlockScan<int,B>;
  __shared__ typename Scan::TempStorage scan;
  __shared__ int packed[B], base_s, total_s;
  __shared__ unsigned bid_s;
  if(threadIdx.x==0) bid_s=atomicAdd(counter,1u);
  __syncthreads();
  unsigned bid=bid_s; size_t i=size_t(bid)*B+threadIdx.x;
  int v=(i<n)?x[i]:0, f=(i<n && v!=0), rank,total;
  Scan(scan).ExclusiveSum(f,rank,total);
  if(f) packed[rank]=v;
  __syncthreads();
  int blockBase=int(publishLookback(total,bid,aggregates,prefixes,states));
  if(threadIdx.x==0){base_s=blockBase; total_s=total;}
  __syncthreads();
  for(int j=threadIdx.x;j<total_s;j+=B) y[base_s+j]=packed[j];
  if(threadIdx.x==0 && bid==(n+B-1)/B-1) *outSize=base_s+total_s;
}
\end{lstlisting}
\code{packed[0..total-1]} 按原输入下标递增排列，块基址也按逻辑块递增，因此私有化不
破坏稳定性。屏障保证共享数组写完后才冲刷；冲刷时每个 $j$ 只有一个写者。

\complexity{$O(N)$ 常规工作加回看开销}{$B$ 个私有整数/块及 $O(G)$ 全局状态}{$G$ 个块，写出阶段为连续协作访问}
\conclusionheading 私有化把每块若干局部紧凑的秩写重新打包为一次稠密连续冲刷；收益取决于保留率和 warp 分布，代价是共享内存和一次块级屏障。
\discussionheading
\discussionpoint{直接按秩写回并不等于完全随机散射。}排他 rank 随输入位置单调增加，所以同一 warp 中所有保留 lane 的目标编号仍是连续序列；只是这些 lane 在执行掩码中可能稀疏，不同 warp 又分别发出很短的连续区段。共享 staging 把整块保留项先压紧到 \code{packed[0..total)}，再让稠密线程前缀统一写出，主要改善跨 warp 的事务打包，而不改变总输出字节。若原来每个 warp 已有足够多连续有效 lane，DRAM 扇区数可能并不下降。准确理解地址序列可以避免把任何 scatter 都笼统标成“随机访问”。

\discussionpoint{选择率决定共享打包究竟是帮助还是额外绕路。}接近 100\% 保留时，直接秩写已是连续存储，共享写、屏障和再读只增加成本；中等选择率且有效 lane 分散在多个 warp 时，块级压紧更可能把多个短事务合成饱满区段。选择率极低时，原方案本来只有几次存储，执行打包和同步又可能得不偿失。本内核并没有清零整个共享数组，成本模型应只计实际 packed 写、连续冲刷与必要屏障。例如四个 warp 各只保留两个 lane 时，直接路径会形成四段很短的写事务，块级打包则可由八个连续线程一次冲刷。由这种分段曲线可设计按选择率自适应路径，但路径判定本身也需低成本且块内一致。

\discussionpoint{记录宽度会把共享容量问题放大。}若元素是一个整数，缓冲需要 $B\times4$ 字节；若是包含键、值和时间戳的 32 字节记录，同样的线程块要占八倍共享空间，可能使驻留块数急降。以 $B=512$ 为例，两者分别约为 2 KiB 和 16 KiB，后者可能直接减少同一 SM 可驻留的块数。结构体布局还会引入对齐、无用填充和向量搬运问题，可以只缓存索引再从源数组聚集，也可采用结构分离布局分别压紧字段。无论选择哪条路线，各字段必须使用同一 rank，不能让键和值因独立压紧而错配。这个所有权约束连接到列式数据库、粒子结构和图边列表的内存布局设计。

\discussionpoint{诊断应同时证明稳定性和解释事务变化。}图 frontier、碰撞候选和查询物化都可能采用块级 staging，但其选择率、记录大小和后续消费者不同，单看 Global Store Efficiency 无法决定收益。实验应把 global store sectors、共享事务、屏障停顿、occupancy 和总时间随选择率一起绘制，并用带唯一原始下标的记录验证输出严格递增；仅用重复值会掩盖换序。若扇区减少却时间上升，通常说明共享或占用成本更大。这样的因果核对能把“指标看起来更好”与“应用真正更快”区分开。
\variantheading 若一个 8 项块的 flags 为 $[1,0,1,1,0,0,1,0]$，私有数组依次保存原输入
的第 $0,2,3,6$ 项，\code{total=4}，随后四个连续位置从块基址写出。
\practiceheading 用 Nsight Compute 比较原始散射与私有化版本的 Global Store Efficiency、
Shared Memory Usage 和 Occupancy；用 CUB DeviceSelect 校验全保留、全丢弃及交替谓词的稳定次序。
\sourcecheck{https://github.com/NVIDIA/cccl/tree/main/cub}{CUB DeviceSelect 官方稳定选择语义；高置信度}
边界审查：共享缓冲必须至少容纳块可能保留的全部 $B$ 项。若谓词有副作用或依赖处理顺序，
并行求值本身可能改变语义，不能仅凭稳定输出次序弥补。
\end{solutionexercise}

\begin{solutionexercise}{4}
\problemheading 对练习 3 的内核应用线程粗化。

练习 3 的基线先把每块保留元素稳定打包到共享内存，再连续写回全局输出。本题要求
每个线程处理多个输入元素，同时保持全局稳定次序；应明确线程到输入子段的映射、
逐项尾部检查、线程内秩与线程间扫描的组合方式，以及粗化对并行度、寄存器和共享内存
占用的影响。

\approachheading 每线程按输入顺序处理 $C$ 个连续元素，先得到线程保留数，再扫描线程
计数。线程基址加线程内秩给出共享私有数组位置；这同时保持线程内和线程间次序。

\answerheading
\begin{lstlisting}[language=C++]
template<int B,int C>
__global__ void stable_filter_coarse(const int* x,int* y,size_t n,
    unsigned* counter,int* aggregates,int* prefixes,unsigned* states,
    int* outSize) {
  using Scan=cub::BlockScan<int,B>;
  __shared__ typename Scan::TempStorage scan;
  __shared__ int packed[B*C],base_s,total_s;
  __shared__ unsigned bid_s;
  if(threadIdx.x==0) bid_s=atomicAdd(counter,1u);
  __syncthreads();
  unsigned bid=bid_s;
  size_t first=size_t(bid)*B*C+size_t(threadIdx.x)*C;
  int values[C], localCount=0;
#pragma unroll
  for(int c=0;c<C;++c){
    size_t i=first+c; values[c]=(i<n)?x[i]:0;
    localCount += (i<n && values[c]!=0);
  }
  int threadBase,total; Scan(scan).ExclusiveSum(localCount,threadBase,total);
  int q=0;
#pragma unroll
  for(int c=0;c<C;++c)
    if(first+c<n && values[c]!=0) packed[threadBase+q++]=values[c];
  __syncthreads();
  int blockBase=int(publishLookback(total,bid,aggregates,prefixes,states));
  if(threadIdx.x==0){base_s=blockBase;total_s=total;}
  __syncthreads();
  for(int j=threadIdx.x;j<total_s;j+=B) y[base_s+j]=packed[j];
  unsigned G=unsigned((n+B*C-1)/(B*C));
  if(threadIdx.x==0 && bid==G-1) *outSize=base_s+total_s;
}
\end{lstlisting}
每个块覆盖 $BC$ 个不相交输入，尾部逐项检查。线程内循环没有同步或竞态；CUB 只扫描
$B$ 个计数而非 $BC$ 个标志，减少同步和跨块状态数量。

\complexity{$O(N)$ 工作、每线程 $O(C)$ 串行段}{$BC$ 个私有共享元素/块，$O(N/(BC))$ 状态}{$N/C$ 个线程；$C$ 增大时并行度下降}
\conclusionheading 粗化摊薄扫描、同步和回看成本，并保持稳定性；$C$ 应以占用率和输入规模调优。
\discussionheading
\discussionpoint{线程到输入区段的映射决定稳定次序。}blocked 映射让线程 $t$ 按顺序处理 $[tC,(t+1)C)$，线程计数的排他扫描按线程号分配基址，因此线程内次序和线程间区段次序都与原输入一致。若改为 striped 下标 $t+cB$ 却仍按线程号拼接，$B=2,C=2$ 会把顺序 $0,1,2,3$ 变成 $0,2,1,3$，数值相等时很难察觉。要兼得 striped 合并与稳定性，必须经共享交换恢复 blocked 所有权，或直接计算每个原始位置的 rank。这个反例与稳定基数排序中的布局约束完全相同。

\discussionpoint{粗化把块间开销换成线程内串行链。}因子 $C$ 使块数、跨块状态和块扫描参与项约缩小 $C$ 倍，一次同步覆盖更多输入；每线程却要依次求值 $C$ 个谓词并保存局部值，关键路径和寄存器需求随之增长。编译器若不能把数组保留在寄存器，可能把它放入 local memory；共享 \code{packed[BC]} 也随 $C$ 线性增大并压低 occupancy。因而最优点出现在通信摊薄与延迟隐藏的交叉处，而不是最大可编译 $C$。还要测试 $N\bmod (BC)\ne0$ 的尾块，因为高粗化下多数线程只处理虚拟槽时，摊薄收益会迅速消失。这种成本曲线也适用于每线程多像素卷积、向量归约和批量哈希探测。

\discussionpoint{加载合并与谓词负载均衡可能要求不同布局。}blocked 输入让每线程获得连续段，便于 \code{float4} 等向量指令，但一个 warp 的线程起点相隔 $C$，$C$ 较大时事务利用率下降；striped 加载让同一轮 lane 访问连续地址，却要转置才能保持稳定输出。如果谓词成本还随数据变化，一个线程连续遇到多个昂贵记录会拖慢整个 warp，可先用 ballot 分类、使用工作队列或把昂贵项二次调度。每种修正都会增加交换与元数据，因此应把访存布局和计算不均分别建模。它们经常在 ETL 解析、可变长度文本和图邻接过滤中同时出现。

\discussionpoint{真实流水线的最佳粗化因子由消费者共同决定。}过滤输出可能马上进入排序、压缩或图遍历，下游若偏好某种 tile 或对齐，单独让过滤最快的 $C$ 未必令端到端最快；内存池是否复用共享状态也会改变固定成本。测试应覆盖 $N\bmod BC$ 的各种尾部、全保留、全丢弃与交替谓词，还要检查 $BC$ 和全局 rank 的宽度不会溢出。可将多项 BlockScan、单项版本和成熟 DeviceSelect 放在相同输入上比较，再计入消费者时间。这样才能从内核参数调优上升到流水线级设计，而不是把局部峰值当成最终目标。
\variantheading 若 $B=128,C=4,N=513$，共有 2 个逻辑块；第二块只有一个有效输入，必须
逐项掩蔽，其余 511 个槽不参与计数，不能作为 0 值误保留。
\practiceheading 对 $C=1,2,4,8$ 用 Nsight Compute 比较寄存器数、占用率和写事务，并与
CUB BlockScan 的 blocked 多项接口对照；主机需特判 $N=0$，还要验证 $BC$ 乘法、网格数
和 \code{int} 输出秩均不溢出。
\sourcecheck{https://github.com/NVIDIA/cccl/tree/main/cub}{CUB BlockScan 对每线程多项输入与前缀回调的官方说明；高置信度}
对抗审查：若改成每线程处理条带式下标 \code{first=c*B+tid} 却仍按线程顺序打包，会改变
全局相对次序。这里采用连续 blocked 映射；大 $C$ 还会增加寄存器与共享内存压力。
\end{solutionexercise}

\chapter{归并与动态输入数据识别}

\begin{solutionexercise}{1}
\problemheading 假设要归并列表
$A=\{1,7,8,9,10\}$ 与 $B=\{7,10,10,12\}$。$C[8]$ 的两个协同秩是多少？

\approachheading 协同秩 $(i,j)$ 满足 $i+j=k$，且长度为 $k$ 的输出前缀恰由
$A[0..i-1]$ 与 $B[0..j-1]$ 稳定归并而成；相等时令 $A$ 先于 $B$。

\answerheading 完整稳定归并为
\[
[1,7_A,7_B,8,9,10_A,10_{B_1},10_{B_2},12].
\]
前 8 个元素已经耗尽 $A$ 的 5 个元素以及 $B$ 的前 3 个元素，故
\[
 i=5,\qquad j=8-i=3.
\]
边界条件也成立：$A[i-1]=10\le B[j]=12$；$i=m$ 使第二个条件自动成立。

\complexity{二分协同秩为 $O(\log(\min(m,n)+1))$，手算为常数}{$O(1)$}{$i+j=k$ 的候选可二分搜索}
\conclusionheading 两个协同秩为 $\boxed{(5,3)}$。
\discussionheading
\discussionpoint{协同秩把输出位置变成输入切分。}若归并结果的前 $k$ 项消费了 $A$ 的 $i$ 项和 $B$ 的 $j$ 项，就必有 $i+j=k$，所以所有候选都落在归并矩阵的同一条对角线上。合法点还要满足切分左侧元素不大于下一侧候选，保证两段输入前缀合在一起恰是最小的 $k$ 项。找到这个点后，输出区间 $[k_0,k_1)$ 只需读取两端协同秩之间的输入片段，可以与其他区间独立归并。例如本题 $k=8$ 的 $(5,3)$ 表示无论两数组长度如何失衡，前八个输出的消费总数始终精确为八。它把按输入切块造成的数据相关负载不均，转化成按输出数量均匀分工，也是 merge path 几何图像的核心。

\discussionpoint{稳定性隐藏在严格与非严格比较的方向中。}当两个键相等并约定 $A$ 中记录先于 $B$，边界的一侧必须允许等号，另一侧必须严格，否则同一对角线上可能出现两个合法候选或一个也没有。以 $A=[2_A]$、$B=[2_B]$、$k=1$ 为例，稳定切分应消费 $A$ 而不是 $B$；只看输出键值都是 2，错误完全不可见。给记录附上来源与原始下标后，比较规则才成为可观察契约。也可把比较键形式化为 $(key,sourcePriority,index)$，使手算、CPU 参考和 GPU 边界判断共享同一个全序。相同的 tie-break 还决定数据库稳定连接、时间序列同刻事件和外部排序多路归并的顺序，因此不是实现细节。

\discussionpoint{概念哨兵不能变成越界解引用。}数学证明常把数组左外侧写作 $-\infty$、右外侧写作 $+\infty$，代码却应先判断 $i=0,m$ 或 $j=0,n$，再决定是否访问邻居，不能真的形成并解引用越界地址。搜索区间至多包含较短输入长度再加一个边界候选，所以复杂度写为 $O(\!\log(\min(m,n)+1))$ 才在某侧为空时仍有定义；直接写 $O(\log\min(m,n))$ 会出现 $\log0$。当 tile 只有几项，线性推进或复用上一个切分反而比二分的分支更便宜，说明渐近最优不等于每个局部规模都最优。

\discussionpoint{不变量测试能覆盖归并路径的整个应用族。}每个切分都应同时满足 $i+j=k$、下标范围和两条有序边界，所有相邻区间拼接后还必须无重无漏地等于稳定串行归并。测试应包含空数组、长度悬殊、完全分离和大量重复键，因为它们分别触发对角线端点、窄搜索区间和 tie-break。特别是某侧为空时，搜索区间仍有一个合法边界候选，复杂度应写成 $O(\log(\min(m,n)+1))$，不能让公式出现 $\log0$。相同切分方法可用于数据库 sort-merge join、倒排表合并和有序流重分区，但记录负载可能比键大得多。将切分正确性与实际搬运字节分开验证，才能既保证排序语义又评价负载均衡。
\variantheading 若改求 $k=6$ 的边界，前六项为 $[1,7_A,7_B,8,9,10_A]$，故协同秩为
$\boxed{(5,1)}$；这与本题 $k=8$ 的 $(5,3)$ 不同。
\practiceheading 可用 Thrust \code{stable_merge} 或 CPU 稳定归并生成黄金输出，并随机加入
大量重复键检查每个 $k$ 的 $i+j=k$、左右边界条件及相邻分区无重叠。
\sourcecheck{https://doi.org/10.1007/978-3-642-33518-1_25}{Siebert 与 Träff 的稳定并行归并原始论文；高置信度}
假设审查：$C[8]$ 表示“秩为 8 的边界”，协同秩描述它之前的 8 项，并非声称 $C[8]$
本身来自 $A[5]$ 或 $B[3]$。若相等时改让 $B$ 优先，中间三个 10 的来源顺序会变，但该
边界仍是 $(5,3)$。
\end{solutionexercise}

\begin{solutionexercise}{2}
\problemheading 完成原书图 13.6 中线程 2 的协同秩函数计算过程。为使题目可独立理解，
该图归并的输入为 $A=[1,7,8,9,11]$、$B=[7,10,10,12]$，并用 3 个线程、每线程
负责 3 个连续输出。线程 2 负责 $C[6]$ 至 $C[8]$，因此以
$(k,A,m,B,n)=(6,A,5,B,4)$ 调用协同秩函数；请给出二分搜索各轮的候选
$(i,j)$、下界以及最终协同秩。

\approachheading 严格执行书中 \code{co_rank(k,A,m,B,n)}：初始
$i=\min(k,m)$、$j=k-i$，并保持 $i+j=k$。

\answerheading $m=5,n=4,k=6$，初值与下界为
\[
(i,j)=(5,1),\qquad i_{low}=\max(0,6-4)=2,\quad
j_{low}=\max(0,6-5)=1.
\]
第 0 轮：$A[4]=11>B[1]=10$，故
$\delta=\lceil(5-2)/2\rceil=2$，更新为
$(i,j)=(3,3)$、$j_{low}=1$。

第 1 轮：$A[2]=8\le B[3]=12$，但 $B[2]=10\ge A[3]=9$，故
$\delta=\lceil(3-1)/2\rceil=1$，更新为
$(i,j)=(4,2)$、$i_{low}=3$。

第 2 轮：$A[3]=9\le B[2]=10$ 且 $B[1]=10<A[4]=11$，两个停止条件均满足，返回
$i=4$，从而 $j=2$。前 6 项确为 $[1,7_A,7_B,8,9,10_{B_1}]$。

\editorialnote 英文底本和中文正文紧邻该图的一段叙述把 $k=6$ 的结果写成 $(5,1)$，并
声称线程 1 输出 $[8,9,11]$；其后线程才输出 10，会使 $C$ 不是有序数组。按底本给出的
函数逐行执行得到 $(4,2)$，上面的反例也独立证实本题应采用这一结果。

\complexity{$O(\log(\min(m,n)+1))$ 时间}{$O(1)$}{不同线程可独立求各自边界}
\conclusionheading 线程 2 的正确协同秩为 $\boxed{(4,2)}$，共经历三次条件检查轮。
\discussionheading
\discussionpoint{先用不变量审图，而不是让图替算法作证。}一条可信协同秩轨迹中的每个候选都必须在目标对角线上，即 $i+j=k$；终点还要满足稳定归并的左右夹逼条件。若教材箭头指到不满足这些式子的格点，应最小修正图示，而不能为了匹配图而改动算法边界。比如目标 $k=6$ 时，图中任何 $(4,3)$ 都因和为 7 而立即失效，无需继续争论比较结果。只从最终合并序列倒推可能掩盖中间搜索区间曾扩大、比较方向曾写反等问题。逐轮核对不变量能指出究竟是候选坐标、上下界还是严格符号出错，这也是审查论文伪代码和硬件波形的通用思路。

\discussionpoint{重复键是观察隐藏稳定约定的探针。}若两个输入都含键 5，选择 $5_A$ 或 $5_B$ 作为第一个输出在数值数组中完全一样，但会给出不同协同秩，并影响携带载荷的记录次序。测试数据应把元素写成 $(key,source,index)$，期望次序便能明确规定同键时先比较来源或原始下标。严格与非严格不等式必须与这条规则成对设计，不能两侧都严格或都宽松。若载荷分别是“旧版本”和“新版本”，一次隐藏换序甚至会改变下游去重保留哪条记录，而不只是展示顺序。数据库同键事件、日志时间戳和稳定排序都依赖相同约定；通过来源标签验证后，稳定性才从口头承诺变成可执行性质。

\discussionpoint{二分轨迹是压缩而高价值的调试证据。}每轮只需记录搜索下界、上界、候选 $(i,j)$ 以及哪条边界比较失败，就能检查区间是否严格缩小、候选是否始终满足 $i+j=k$。若循环卡住，轨迹会直接显示整数中点取整后上下界未变化；若终点越界，则可追到第一次非法更新。合法实现每轮至少排除一个候选，因此迭代上界可由初始候选数的对数给出，超过上界本身就是错误信号。GPU 上不应让成千线程同时 \code{printf}，可只让失败块领导者原子保留一个紧凑样例，随后在 CPU 重放完整搜索。类似的“设备捕获最小反例、主机确定性复现”也适合哈希探测和并行图算法调试。

\discussionpoint{勘误记录本身应当可复现。}修正底本时应保留原图号、原坐标或比较符号、修正结果和依据，使读者能从同一输入重新走出差异；静默改图会切断证据链。最小复现应同时保存两输入、目标 $k$、稳定规则和逐轮区间，因为只截终点无法区分图错与算法错。生产代码也不能只保存这一个示例，因为固定数据可能恰好没有重复键或端点。应把对角线、范围和夹逼条件写成调试断言，再用属性测试生成空侧、重复键和极端长度比。这样一次教材勘误可以转化为长期回归规则，未来更换索引宽度、tile 大小或比较器时仍能捕获同类错误。
\variantheading 若把 $A[4]$ 从 11 改为 10，$k=6$ 的稳定前缀仍先取 $A$ 的三个 10 之前
可用项，协同秩变为 $(5,1)$；停止条件必须继续让 $A$ 的相等键优先。
\practiceheading 把底本示例和随机重复键数据加入回归测试，逐个输出边界与串行稳定归并
对比；GPU 上再用 Compute Sanitizer memcheck 检查 $i=0,m$ 或 $j=0,n$ 的哨兵分支。
\sourcecheck{https://davidbader.net/publication/2012-gm-ba/2012-gm-ba.pdf}{GPU Merge Path 作者公开版对对角划分/协同秩思想的独立佐证；并以 DOI 10.1145/2304576.2304621 核对正式出版信息；高置信度}
对抗审查：若用不严格条件 $B[j-1]\le A[i]$，相等键可能同时被相邻分区认领，破坏稳定性；
书中第二条件用 $\ge$ 触发调整正是为了让 $A$ 的相等键优先。
\end{solutionexercise}

\begin{solutionexercise}{3}
\problemheading 对原书图 13.12 中载入 $A$、$B$ 输入块的两个 for 循环增加协同秩
函数调用，使其只载入当前 while 循环迭代实际会消费的 $A$、$B$ 元素。原内核每轮最多
生成 \code{tile_size} 个输出，并分别从尚未消费的两个有序输入中各载入至多
\code{tile_size} 项；需要先由本轮输出边界的协同秩确定两个输入的精确消费量，再由块内
线程协作载入，同时正确处理最后一轮和任一输入先耗尽的边界情形。

\approachheading 当前轮最多生成 \code{tile_size} 个输出。先由一个线程对尚未消费的两个
全局子数组求这个输出边界的协同秩，得到精确的 \code{A_take/B_take}，再合作加载。

\answerheading 以下代码替换原书图 13.12 的两段固定上限循环，并相应更新消费计数：
\begin{lstlisting}[language=C++]
__shared__ int A_take_s, B_take_s, C_take_s;

while (C_completed < C_length) {
  if (threadIdx.x == 0) {
    int remainC = C_length - C_completed;
    C_take_s = min(tile_size, remainC);
    int remainA = A_length - A_consumed;
    int remainB = B_length - B_consumed;
    A_take_s = co_rank(C_take_s,
        A + A_curr + A_consumed, remainA,
        B + B_curr + B_consumed, remainB);
    B_take_s = C_take_s - A_take_s;
  }
  __syncthreads();

  for (int p=threadIdx.x; p<A_take_s; p+=blockDim.x)
    A_S[p] = A[A_curr + A_consumed + p];
  for (int p=threadIdx.x; p<B_take_s; p+=blockDim.x)
    B_S[p] = B[B_curr + B_consumed + p];
  __syncthreads();

  int c0 = int((static_cast<unsigned long long>(threadIdx.x) * C_take_s)
               / blockDim.x);
  int c1 = int((static_cast<unsigned long long>(threadIdx.x + 1) * C_take_s)
               / blockDim.x);
  int a0 = co_rank(c0,A_S,A_take_s,B_S,B_take_s), b0=c0-a0;
  int a1 = co_rank(c1,A_S,A_take_s,B_S,B_take_s), b1=c1-a1;
  merge_sequential(A_S+a0,a1-a0,B_S+b0,b1-b0,
                   C+C_curr+C_completed+c0);
  __syncthreads();
  A_consumed += A_take_s; B_consumed += B_take_s;
  C_completed += C_take_s;
  __syncthreads();
}
\end{lstlisting}
协同秩定义保证当前输出前缀恰好消费 \code{A_take+B_take=C_take} 个输入；共享内存
中的线程级归并再把这一前缀无重叠地划分。循环末屏障防止快线程覆盖仍在使用的 tile。

\complexity{$O(m+n)$ 归并工作，每轮另有 $O(\log(m+n))$ 边界搜索}{$O(\text{tile\_size})$ 共享空间}{块间、块内线程区间均可并行}
\conclusionheading 全局加载总量从每轮“至多两个完整 tile”降为恰好当前轮消费的元素数。
\discussionheading
\discussionpoint{输出预算可以精确反推出本轮输入需求。}若共享 tile 本轮只准备产生 $C_{take}$ 个结果，从当前输入游标出发求终止对角线的协同秩，就得到恰好要消费的 $A_{take}$ 与 $B_{take}$，并满足两者之和等于输出数。朴素方案若总从两侧各加载一个完整 tile，在两个数组完全分离或一侧接近结束时，会搬入大量本轮不会比较的元素。按协同秩裁剪后，加载集合与消费集合一致，节省的是全局字节而非归并比较本身。相同“从输出容量反推输入边界”的思想还用于分页 join、流窗口和压缩块解码。

\discussionpoint{一次二分搜索只有在省下的通信更贵时才值得。}领导线程执行 $O(\!\log T)$ 比较，还要把消费数广播给全块并同步；若 tile 很小、数据已命中缓存或两侧消费恰好均衡，控制成本可能超过少量多余加载。tile 较大、全局内存昂贵或输入高度偏斜时，避免半个 tile 的无用传输则容易覆盖搜索成本。可用 $T_{saved}\approx bytes_{avoided}/B_{eff}$ 与二分、广播和屏障时间比较盈亏点，并复用前一轮切分作为下一轮搜索窄区间。这个模型把算法优化从“少加载一定快”变成可在目标设备上验证的条件判断。

\discussionpoint{协作加载必须覆盖任意消费数并维护流水生命周期。}让线程 $t$ 处理 $t,t+B,t+2B,\ldots$ 的局部下标，可覆盖从 0 到任意 \code{count-1} 的全部元素，即使数量不是块大小倍数或最后一轮只有几项，也不会遗漏。例如 $B=128,\code{count}=130$ 时，线程 0 和 1 各再搬一项，其余线程只搬一项，循环无需特殊尾部分支。支持异步全局到共享拷贝时，可用双缓冲让下一轮加载与当前轮归并重叠；但覆盖旧缓冲前必须确认所有消费者完成，提交组和等待点也要与该轮边界元数据配对。若只同步数据不更新消费数，线程可能用新 tile 配旧边界。这个版本化问题与网络环形缓冲和生产者--消费者队列完全相似。

\discussionpoint{输入分布决定裁剪加载的实际收益。}均匀交错时两侧消费大致均衡，完全分离时一轮可能几乎只取一侧，全相等又会把比例交给稳定 tie-break，一侧很短则频繁触发端点；四类数据覆盖了主要路径。若每轮输出 $T$ 项而朴素地两侧各载入 $T$ 项，完全分离时裁剪最多可避免近 $T$ 项无用搬运，均匀交错时收益则小得多。性能报告应同时给出实际加载元素数、全局扇区、二分迭代、屏障等待和结果吞吐，才能确认少搬元素确实转化为少事务。正确性还要累计所有轮的 $A_{take}$、$B_{take}$，最终分别等于输入长度，并与稳定 CPU 归并逐项相同。这样数据分布既是性能变量，也是边界证明的测试生成器。
\variantheading 若当前轮生成 1000 项、块含 128 个线程，则线程 127 的区间为
\[
[\lfloor127\times1000/128\rfloor,1000)=[992,1000),
\]
所以最后 8 项不会因整除截断而丢失。
\practiceheading 用 CPU 的标准归并对照不同长度、全相等键和空侧输入，并重点覆盖 tile
大小不能整除块大小的组合；Nsight Compute 可比较修复前后的实际加载字节与领导
线程二分开销。
\sourcecheck{https://davidbader.net/publication/2012-gm-ba/2012-gm-ba.pdf}{GPU Merge Path 作者公开版的分区与分块归并方法；并以 DOI 10.1145/2304576.2304621 核对正式出版信息；高置信度}
边界审查：末轮 \code{C_take} 可能小于块大小，必须用它限制线程输出；若只缩短加载却仍
把 \code{tile_size} 传给协同秩，会读取未初始化共享内存。
\end{solutionexercise}

\begin{solutionexercise}{4}
\problemheading 考虑并行归并两个长度分别为 1,030,400 和 608,000 的数组。假设每线程
归并 8 个元素，线程块大小为 1024。分别回答：
\begin{enumerate}[label=(\alph*)]
  \item 在原书图 13.9 的基本归并内核中，有多少线程在全局内存数据上执行二分搜索？
  \item 在原书图 13.11 至图 13.13 的分块归并内核中，有多少线程在全局内存数据上执行
  二分搜索？
  \item 在同一分块归并内核中，有多少线程在共享内存数据上执行二分搜索？
\end{enumerate}

\approachheading 先求输出总长、线程数和块数，再区分“执行搜索的不同线程数”与“函数
调用次数”。

\answerheading
\[
N=1{,}030{,}400+608{,}000=1{,}638{,}400,
\]
\[
T=N/8=204{,}800,\qquad G=T/1024=200.
\]

基本内核中每个线程都为输出区间两端调用协同秩，所以在全局内存数据上搜索的是
$\boxed{204{,}800}$ 个不同线程（共 $409{,}600$ 次调用）。

分块内核只让每块一个领导线程求块级起止协同秩，所以在全局数据上搜索的是
$\boxed{200}$ 个不同线程（共 400 次调用）。块内每个线程仍需在共享 tile 上求自己的
起止边界，所以在共享数据上搜索的是 $\boxed{204{,}800}$ 个不同线程（共 409,600 次
调用；本题参数使每线程的 8 项恰在一个块级轮次中完成）。

\complexity{搜索本身每次 $O(\log N)$；随后归并总工作 $O(N)$}{$O(\text{tile size})$ 共享空间/块}{$204800$ 个线程、200 个块}
\conclusionheading (a) 204,800；(b) 200；(c) 204,800。
\discussionheading
\discussionpoint{两级切分把昂贵搜索提升到块级共享。}块领导者先在完整全局数组上求当前输出 tile 的两端协同秩，确定要搬入共享内存的输入片段；块内线程再在这两个小片段上为自己的输出子区间求局部切分。于是高延迟、可能跨许多缓存行的全局二分从每线程一次降为每块一次，增加的是低延迟共享搜索。若 1024 个线程各自做全局二分，搜索起点就有 1024 份；提升后每块只需求少数 tile 边界，再让线程在共享窗口内细分。结构上它像页表和缓存的两级定位：粗层缩小工作集，细层在近处完成个体分工。正确性来自局部坐标加回全局起点后仍满足同一对角线与夹逼不变量。

\discussionpoint{线程数、调用数和比较数必须分开统计。}一个线程可能为左右两个边界调用两次二分，一次二分又可能执行若干比较；“有多少线程搜索”不能直接回答总比较工作。尾块仍会启动完整线程块，但超出实际输出范围的线程可能跳过搜索或只处理空区间，取决于边界代码。若每线程生成多个输出，线程数还会少于输出数。粗略模型可写成 $C\approx Q\bar L+C_{block}$，其中 $Q$ 是有效调用数，$\bar L$ 是每次平均迭代，而不是把启动线程数当作比较数。这个计数纪律同样适用于哈希探测、BVH 遍历和数据库索引查询。

\discussionpoint{tile 尺寸在搜索频率与片上驻留之间取平衡。}较小 tile 需要更多块级全局切分和同步，输入片段复用也有限；较大 tile 减少边界搜索，却消耗更多共享内存，并可能让每 SM 只能驻留很少块。若 tile 从 1024 翻倍到 2048 项，边界搜索数近似减半，但双缓冲共享容量也近似翻倍，驻留块数可能发生阶跃下降。记录若包含大 payload，搬入整个记录会进一步放大容量，可只缓存键和索引，归并确定后再搬 payload。最优点还受比较器成本和输入长度比影响，不能从整数键实验直接外推。通过扫描 tile 大小并记录全局搜索、共享占用和端到端吞吐，才能找到真正的资源拐点。

\discussionpoint{分层归并连接到存储系统中的有序数据维护。}数据库 sort-merge join 用有序键对齐记录，LSM compaction 合并不同层的有序段，倒排索引交集处理长度高度不均的 posting list；它们都可用对角切分分配输出工作。实际记录比较可能访问字符串，payload 搬运也远大于整数，所以键搜索、记录物化与写带宽应分别计量。尤其 LSM 压实常受写放大限制，即使切分让线程工作均衡，额外物化一遍 payload 仍可能吞掉全部收益。基准要改变输入长度比、重复键和 payload 大小，并用稳定 CPU 结果核对来源次序。这样才能判断两级切分究竟缓解了负载不均，还是被比较与搬运成本掩盖。
\variantheading 若每线程改为归并 16 项而块仍为 1024 线程，则共有 102,400 个线程和
100 个块；全局搜索线程数为 100，共享搜索线程数为 102,400。
\practiceheading 在生产基准中分别统计全局与共享内存访问、二分分支效率和块占用率，并用
GPU Merge Path 或 Thrust merge 作基线；长度乘加应使用 \code{size_t}，避免大输入的 32 位溢出。
\sourcecheck{https://davidbader.net/publication/2012-gm-ba/2012-gm-ba.pdf}{GPU Merge Path 作者公开版对层次化分区搜索的分析；并以 DOI 10.1145/2304576.2304621 核对正式出版信息；高置信度}
假设审查：题目问“多少线程”而不是“多少次搜索”。若实现复用相邻线程共享的区间端点，
调用次数可以再减；这里严格按原书两个内核的源代码计数。
\end{solutionexercise}

\chapter{排序}

\begin{solutionexercise}{1}
\problemheading 扩展原书图 14.7 的内核，使用共享内存改进内存合并。该内核的一位稳定
基数迭代为每个输入键分配一个线程，以
\[
  b_i=(\code{key}_i\mathbin{\texttt{>>}}\code{iter})\mathbin{\texttt{\&}}1
\]
提取目标位。对位数组执行网格范围排他扫描后，令 $o_i$ 为键 $i$ 之前的 1 键数、$O$
为 1 键总数，并按
\[
  \code{dst}_i=\begin{cases}i-o_i,&b_i=0,\\N-O+o_i,&b_i=1\end{cases}
\]
求稳定目标下标。
扩展后的实现应先在块内共享内存中形成连续的局部 0 桶和 1 桶，再确定各局部桶的全局
起点并以合并访问写回。

\approachheading 每块先用局部排他扫描把键稳定地排入共享内存的 0、1 两个桶，再合并
写入块级暂存数组。对“桶优先、块次序”的计数表做全局排他扫描，最后把各连续局部桶搬到
全局目的区间。

\answerheading 以下三个阶段可直接与 CUB DeviceScan 组合；\code{counts/offsets} 长度
为 $2G$，\code{packed} 长度为 $N$。
\begin{lstlisting}[language=C++]
#include <cub/block/block_scan.cuh>
#include <cub/device/device_scan.cuh>

template<int B>
__global__ void pack_one_bit(const unsigned* in,unsigned* packed,
    unsigned* counts,size_t n,unsigned shift) {
  using Scan=cub::BlockScan<int,B>;
  __shared__ typename Scan::TempStorage temp;
  __shared__ unsigned local[B];
  size_t base=size_t(blockIdx.x)*B, i=base+threadIdx.x;
  int valid=i<n, bit=valid?((in[i]>>shift)&1u):0;
  int onesBefore,totalOnes;
  Scan(temp).ExclusiveSum(valid?bit:0,onesBefore,totalOnes);
  int validCount=int(min(size_t(B),n>base?n-base:0));
  int zeros=validCount-totalOnes;
  if(valid) {
    int pos=bit ? zeros+onesBefore : int(threadIdx.x)-onesBefore;
    local[pos]=in[i];
  }
  __syncthreads();
  if(threadIdx.x<validCount) packed[base+threadIdx.x]=local[threadIdx.x];
  if(threadIdx.x==0) {
    counts[blockIdx.x]=zeros;
    counts[gridDim.x+blockIdx.x]=totalOnes;
  }
}

__global__ void scatter_one_bit(const unsigned* packed,unsigned* out,
    const unsigned* counts,const unsigned* offsets,size_t n,unsigned G) {
  size_t i=size_t(blockIdx.x)*blockDim.x+threadIdx.x;
  if(i>=n)return;
  unsigned t=threadIdx.x, zeros=counts[blockIdx.x];
  unsigned bucket=t<zeros?0:1, rank=t<zeros?t:t-zeros;
  out[offsets[bucket*G+blockIdx.x]+rank]=packed[i];
}

// G=(n+B-1)/B:
// pack_one_bit<B><<<G,B>>>(in,packed,counts,n,shift);
// cub::DeviceScan::ExclusiveSum(temp,tempBytes,counts,offsets,2*G);
// scatter_one_bit<<<G,B>>>(packed,out,counts,offsets,n,G);
\end{lstlisting}
块扫描保持每个局部桶的原相对顺序；计数表先列全部 0 桶、再列全部 1 桶，且每列按块号
排列，所以全局 offset 同时保证桶顺序、块顺序和稳定性。每个输出区间互不重叠，无竞态。

\complexity{$O(N+G)$ 工作（不含扫描常数）}{$N+4G$ 个显式辅助整数、CUB 扫描查询得到的 \code{tempBytes} 工作区及 $B$ 个共享键/块}{$G$ 块；桶搬运为连续访问}
\conclusionheading 共享内存吸收不规则局部散射，全局读写均按连续桶区间进行。
\discussionheading
\discussionpoint{一位稳定分区是 LSD 基数排序的归纳基石。}当前位为 0 或 1 可视为两个谓词桶，排他秩让同桶元素保持原次序，桶起点再把 0 桶整体放在 1 桶之前。最低位 pass 结束后序列按这一位有序；假设前 $p$ 位已经有序，第 $p+1$ 位稳定分桶会重排不同新位的组，却保留同组内旧的低位顺序，因此归纳得到前 $p+1$ 位有序。任一轮若不稳定，后续高位无法恢复被打乱的低位。例如两个高位相同而低位不同的键一旦在高位 pass 内换序，后续已经没有更高优先级的处理来纠正它们。这个证明也解释了为什么 MSD 排序的递归结构不同，不能把两者的 pass 顺序随意交换。

\discussionpoint{桶地址由局部计数和全局前缀共同组成。}块内先把 0、1 项分别压成连续区段，并记录两个桶计数；设备范围按桶扫描各块计数，得到该块在相应全局桶中的起点；最后局部 rank 加全局起点形成唯一地址。对桶 $b$ 的块 $g$，可写成 $pos=base_b+\sum_{h<g}count_{h,b}+rank_{local}$，三项分别隔离桶、前块和块内记录。这个结构可以推广到 $R$ 个桶，写成“局部直方图、桶主序全局前缀、局部稳定散射”。数据库 hash partition、并行 allocator 和稀疏结构组装也用同一两层地址分配方式。正确性关键不是某段 CUDA 代码，而是每个地址可唯一分解为桶起点、前块数量和块内秩。

\discussionpoint{写合并收益需要用临时空间和额外遍历购买。}块级打包通常要保存 $N$ 项临时键、每块桶计数和扫描工作区，还可能多出计数、扫描、散射多个 launch；对很小输入，这些固定成本会超过合并写节省的事务。工作区大小应先查询并由内存池跨 pass 复用，不能每轮重新分配，否则分配同步会淹没内核差异。一个 32 位键数组每增加一次完整读写就至少搬运约 $8N$ 字节，这给优化回本所需减少的事务提供了量级下界。评估时应累计整次排序的读写字节和时间，而非只看最后 scatter 的 store efficiency。这个原则也适用于任何“先重排以改善后续访问”的优化：布局转换本身必须进入成本账本。

\discussionpoint{按位次序必须先定义键的语义编码。}无符号整数的位字典序与数值序一致，可直接从低位处理；二进制补码有符号数需要翻转符号位，IEEE 浮点还要根据符号对位型作条件变换，并明确 NaN、正负零如何排序。若把有符号的 $-1$ 直接当无符号位型，它会成为接近最大值的全 1 模式，数值顺序立即反转，这就是规范化不可省的反例。复合键可能按字段优先级串接，也可能需要稳定地分多阶段处理。验证时应让重复键携带原始下标，以观察稳定性，并覆盖已排序、逆序、单值、NaN 和高度偏斜数据。与成熟 DeviceRadixSort 在同一位范围和同一语义下比较，才能避免把编码差异误判为算法错误。
\variantheading 若当前位的输入键为 $[3,2,1,4]$，0 桶稳定得到 $[2,4]$、1 桶得到
$[3,1]$，本轮输出为 $[2,4,3,1]$；桶内不能按键值再次排序。
\practiceheading 用 CUB \code{DeviceRadixSort} 作稳定性和吞吐基线，并用重复键携带原始下标
验证桶内次序；Nsight Compute 可检查 Scatter 阶段的 Global Store Efficiency 和共享容量
对占用率的影响。
\sourcecheck{https://github.com/NVIDIA/cccl/tree/main/cub}{NVIDIA CCCL/CUB 官方仓库中的 BlockScan 与 DeviceRadixSort 实现；高置信度}
边界审查：末块的无效 lane 不能计入 0 桶；代码用 \code{validCount} 单独处理。单个合作
内核若没有合法网格级同步，不能在产生 counts 后立即读取 offsets。
\end{solutionexercise}

\begin{solutionexercise}{2}
\problemheading 扩展原书图 14.7 的内核，使其支持多位基数。原内核每轮只提取 1 位并
形成两个稳定桶；扩展后每轮应提取 $r$ 位、形成 $2^r$ 个稳定桶，给出块内局部排序、各块
桶计数的全局定位和写回方法，并保持相等高位键在本轮之前已有的相对次序。

\approachheading 每次取 $r$ 位形成 $R=2^r$ 个桶。块内按桶号依次执行稳定扫描并打包；
全局计数表按“桶主序、块次序”展开后做一次排他扫描。

\answerheading 下列清晰版本以 $R$ 次块扫描换取简单、可验证的实现；实际程序可用 CUB
BlockRadixSort 融合这些扫描。
\begin{lstlisting}[language=C++]
template<int B,int RBITS>
__global__ void pack_multibit(const unsigned* in,unsigned* packed,
    unsigned char* packedDigit,unsigned* counts,size_t n,unsigned shift) {
  static_assert(RBITS>=1 && RBITS<=8,
                "unsigned char digit requires 1 <= RBITS <= 8");
  constexpr int R=1<<RBITS;
  using Scan=cub::BlockScan<int,B>;
  __shared__ typename Scan::TempStorage temp;
  __shared__ unsigned keys[B];
  __shared__ unsigned char digs[B];
  size_t base=size_t(blockIdx.x)*B, i=base+threadIdx.x;
  int valid=i<n; unsigned key=valid?in[i]:0;
  unsigned digit=(key>>shift)&(R-1); int running=0;
  for(int b=0;b<R;++b) {
    int rank,total,flag=valid && digit==unsigned(b);
    Scan(temp).ExclusiveSum(flag,rank,total);
    if(flag){keys[running+rank]=key;digs[running+rank]=digit;}
    if(threadIdx.x==0) counts[b*gridDim.x+blockIdx.x]=total;
    running+=total;
    __syncthreads(); // required before reusing TempStorage
  }
  int validCount=int(min(size_t(B),n>base?n-base:0));
  if(threadIdx.x<validCount){
    packed[base+threadIdx.x]=keys[threadIdx.x];
    packedDigit[base+threadIdx.x]=digs[threadIdx.x];
  }
}

template<int RBITS>
__global__ void scatter_multibit(const unsigned* packed,
    const unsigned char* digit,unsigned* out,const unsigned* counts,
    const unsigned* offsets,size_t n,unsigned G) {
  size_t i=size_t(blockIdx.x)*blockDim.x+threadIdx.x;
  if(i>=n)return;
  unsigned d=digit[i], start=0;
  for(unsigned b=0;b<d;++b) start+=counts[b*G+blockIdx.x];
  unsigned rank=threadIdx.x-start;
  out[offsets[d*G+blockIdx.x]+rank]=packed[i];
}
// Scan R*G bucket-major counts, scatter, then ping-pong input/output.
// Repeat ceil(key_bits/RBITS) passes; optionally shrink the final bucket set.
\end{lstlisting}
每个桶内按输入下标稳定；桶按数值升序串接。LSD 基数排序归纳地保持先前低位的次序，故
完成全部 digit 后按完整键排序。

\complexity{本教学实现 $O(RN)$，优化实现可达每遍 $O(N+RG)$}{$O(N+RG)$ 辅助空间，$B$ 个共享键/块}{$G$ 块与 $R$ 个桶并行层次}
\conclusionheading $r$ 位把遍数缩为原来的约 $1/r$，但桶表、共享资源和局部扫描成本增大。
\discussionheading
\discussionpoint{多位 digit 没有改变稳定归纳，只是扩大了桶数。}每轮提取 $r$ 位得到 $R=2^r$ 个桶，并按 digit 数值升序稳定连接；同一 digit 内仍保持此前所有低位建立的顺序。假设进入本轮时已按低 $p$ 位有序，稳定分桶后，不同 digit 由新位决定先后，同 digit 则沿用旧序，所以得到低 $p+r$ 位有序。以 $r=2$ 为例，一轮有 0--3 四桶，但桶内仍必须按进入本轮的次序排列，不能用四个无序原子追加队列替代。局部打包和跨块桶前缀任一处换序都会破坏这条链，错误可能只在重复高 digit 时出现。逐 pass 保存键--原始下标结果，比只检查最终随机键更容易定位是哪一层失去稳定性。

\discussionpoint{radix 宽度以指数资源换取线性减少 pass。}32 位键取 $r=4$ 需要 8 轮和 16 桶，取 $r=8$ 只需 4 轮却有 256 桶；桶计数、清零、全局扫描和共享布局都按 $2^r$ 增长。宽 digit 还会提高寄存器需求并让局部直方图出现 bank 冲突，而窄 digit 则重复读写完整键数组。粗略成本可写成 $\lceil w/r\rceil(C_{stream}+2^rC_{bucket})$，它说明减少 pass 与增加桶工作并非同一比例。最佳点通常位于中等位宽，但会随架构、每线程键数和键范围改变。这个指数--线性交换与查表算法、分支因子和缓存分块相似，必须用总字节和资源占用共同评估。

\discussionpoint{最后一轮的完整掩码在规范化无符号键上仍然正确。}对 32 位 \code{unsigned} 且 \code{shift<32}，右移会在高位补 0；若最后只剩 2 个有效位而仍用 8 位掩码，digit 仍落在 0--3，只是其余 252 个桶永远为空。例如在 \code{shift=30} 时，\code{(key >> 30) \& 0xff} 与掩码 3 得到同一 0--3 值，不会凭空读取第 32 位之外的信息。因此缩小掩码和桶数是减少清零、扫描工作的优化，不是该前提下的正确性补丁。只有逻辑键宽小于存储宽且未清掉高位，或键是有符号、浮点、复合字段时，才要先做顺序保持规范化并明确有效位。把这两个场景分开，可避免把无害空桶误写成越界位错误。

\discussionpoint{桶分布改变冲突位置而不改变总键数。}均匀 digit 会使用全部桶，局部计数较分散，但全局前缀要处理大量非零项；单一 digit 让同一局部计数器高度竞争，散射结果却是一个连续区段。无论分布如何，所有桶计数之和都应严格等于 $N$，这条守恒式可把性能差异与丢键、重键的正确性错误分开。重复模式、已排序键和随机键因此可能有完全不同瓶颈。比较 $r$ 时应报告 pass 数、临时字节、共享占用、直方图冲突与端到端键/秒，并保持键范围和稳定语义一致。将这些结果与 BlockRadixSort、DeviceRadixSort 对照后，才能判断教学实现受限于算法选择、局部实现还是工作区管理。
\variantheading 若排序 32 位键且 $r=8$，共有 256 个桶、需要 4 个 pass，\code{digit} 范围
为 0--255，恰好可由 \code{unsigned char} 表示；$r=9$ 则必须换更宽类型且不被本实现允许。
\practiceheading 比较 CUB \code{BlockRadixSort} 的 4 位与 8 位配置，记录寄存器、共享内存和
吞吐；用全 0、全 1、重复高位及最高位设置的键验证每个 pass，运行前校验
\code{shift < 32}。
\sourcecheck{https://github.com/NVIDIA/cccl/tree/main/cub}{NVIDIA CUB BlockRadixSort/DeviceRadixSort 的官方源代码；高置信度}
反例审查：若局部桶不是稳定的，LSD 多遍排序会破坏已经排好的低位。$r$ 过大还会使
$R=2^r$ 的计数空间和扫描成本失控；有符号数与浮点数需采用 CUB 所定义的位变换语义。
\end{solutionexercise}

\begin{solutionexercise}{3}
\problemheading 扩展原书图 14.7 的内核，应用线程粗化改进内存合并。原内核每线程处理
一个键并直接写目标位置；粗化版本应让每线程处理多个连续键，说明如何合并读取、在块内
保持稳定顺序、计算局部桶与全局桶偏移，并以连续区段写回。

\approachheading 每线程处理四个连续键并用 \code{uint4} 向量加载。先扫描“每线程 0
键数”，再按线程次序与线程内次序写入共享桶；全局桶定位沿用第 1 题的计数扫描。

\answerheading
\begin{lstlisting}[language=C++]
template<int B>
__global__ void pack_one_bit_x4(const unsigned* in,unsigned* packed,
    unsigned* counts,size_t n,unsigned shift) {
  constexpr int C=4;
  using Scan=cub::BlockScan<int,B>;
  __shared__ typename Scan::TempStorage temp;
  __shared__ unsigned local[B*C];
  size_t blockBase=size_t(blockIdx.x)*B*C;
  size_t first=blockBase+size_t(threadIdx.x)*C;
  unsigned v[C]={0,0,0,0}; int valid=0;
  if(first+3<n && ((reinterpret_cast<size_t>(in+first)&15)==0)) {
    uint4 z=*reinterpret_cast<const uint4*>(in+first);
    v[0]=z.x;v[1]=z.y;v[2]=z.z;v[3]=z.w;valid=4;
  } else for(int c=0;c<C;++c) if(first+c<n){v[c]=in[first+c];++valid;}
  int zerosLocal=0;
  for(int c=0;c<valid;++c) zerosLocal+=(((v[c]>>shift)&1u)==0);
  int zeroBase,totalZeros;
  Scan(temp).ExclusiveSum(zerosLocal,zeroBase,totalZeros);
  int validTotal=int(min(size_t(B*C),n>blockBase?n-blockBase:0));
  int validBefore=min(int(threadIdx.x)*C,validTotal);
  int oneBase=validBefore-zeroBase, z=0,o=0;
  for(int c=0;c<valid;++c) {
    if(((v[c]>>shift)&1u)==0) local[zeroBase+z++]=v[c];
    else local[totalZeros+oneBase+o++]=v[c];
  }
  __syncthreads();
  for(int p=threadIdx.x;p<validTotal;p+=B) packed[blockBase+p]=local[p];
  if(threadIdx.x==0){
    counts[blockIdx.x]=totalZeros;
    counts[gridDim.x+blockIdx.x]=validTotal-totalZeros;
  }
}

template<int B>
__global__ void scatter_one_bit_x4(const unsigned* packed,unsigned* out,
    const unsigned* counts,const unsigned* offsets,size_t n,unsigned G) {
  size_t base=size_t(blockIdx.x)*B*4;
  unsigned valid=unsigned(min(size_t(B*4),n>base?n-base:0));
  unsigned zeros=counts[blockIdx.x];
  for(unsigned p=threadIdx.x;p<valid;p+=B) {
    unsigned bucket=p<zeros?0:1, rank=p<zeros?p:p-zeros;
    out[offsets[bucket*G+blockIdx.x]+rank]=packed[base+p];
  }
}
// Scan 2*G counts, then launch scatter_one_bit_x4<B><<<G,B>>>.
\end{lstlisting}
线程 $t$ 之前的有效项数为 \code{validBefore}；减去此前 0 数就是此前 1 数，故两个桶的
目的位置都唯一且稳定。向量路径只在完整且 16 字节对齐时使用。

\complexity{$O(N)$ 工作，每线程 4 项串行段}{$4B$ 个共享键/块和 $O(G)$ 计数}{$N/4$ 个线程，块数约降为四分之一}
\conclusionheading 粗化减少扫描参与线程和块计数，并用 16 字节访问保持连续输入传输。
\discussionheading
\discussionpoint{粗化后的线程所有权必须与稳定次序一致。}若线程 $t$ 按输入顺序处理连续四键，线程内局部 rank 保持这四项次序，线程保留数的排他扫描又按 $t$ 递增分配区段，因此块内整体稳定。若为了全局合并改成 striped 读取 $t,t+B,t+2B,t+3B$，却仍按线程号串接，原序会被改成条带分组，相等键的载荷顺序随之改变。以 $B=2,C=2$ 为例，所有权次序会从 $0,1,2,3$ 变成 $0,2,1,3$，即使四个键相等，载荷标签也能揭示换序。解决方法是在共享内存中真正做线程间转置以恢复 blocked 所有权，或按原始下标计算 rank。这个所有权证明与粗化过滤完全相同，展示了稳定原语之间可复用的核心不变量。

\discussionpoint{十六字节指令只解决单线程指令数。}\code{uint4} 加载要求实际地址满足 16 字节对齐，四个键都属于合法对象；子数组从奇数键偏移开始或尾部不足四项时，应回退为标量受保护访问。即使每个线程都成功执行一条向量加载，warp 中线程起点若相隔十六字节以上，覆盖的 DRAM 扇区仍可能多于 striped 访问，因此指令减少不等于事务减少。可将 load instruction 与 global sectors 同时比较，并用对齐、错位和各尾宽输入分别验证。类似区分也适用于 SIMD 向量化、结构体读取和二维 pitch 子视图。

\discussionpoint{四键粗化减少元数据却延长每线程生命期。}每线程处理四项后，参与扫描的线程数和块计数大致降为四分之一，固定同步与全局桶状态被更多键摊薄；线程却要保存四个键、digit、标志和局部 rank，串行循环也更长。寄存器不足会触发 local memory，输入较小又可能只剩少量块，二者都会削弱延迟隐藏。因而粗化不是无条件四倍优化，应把 $C=1,2,4,8$ 当作策略空间，并在 ptxas 资源报告和实际 occupancy 下找拐点。尾部还需单独测试 $N\bmod C=1,2,3$，否则向量加载掩盖的越界只会在真实长度出现。这个规律也适用于每线程多输出的 GEMM 与卷积。

\discussionpoint{生产排序把多个局部优化放进一条多 pass 流水。}真实实现常融合加载、局部重排、直方图和散射，并让输入输出缓冲 ping-pong 交替，避免每轮拷回固定位置；工作区也跨轮复用。评估 x4 版本不能只计一个 pack kernel，应覆盖对齐与未对齐首址、$N\bmod4$ 的所有尾宽、重复键和不同 digit 偏斜，再观察完整排序的字节与时间。若 pass 数为奇数，最终结果所在缓冲与偶数时相反，API 必须返回实际所有权而非悄悄多拷一遍。若加载指令下降而总吞吐不变，瓶颈可能在桶扫描或写回；若 occupancy 大降，则寄存器是主因。这种逐阶段证据链能避免把局部汇编改善误当成端到端收益。
\variantheading 若 $N=1025,B=256$，首块处理 1024 项，第二块只处理 1 项；尾块不能执行
越界 \code{uint4} 加载，其标量回退只令 \code{valid=1}。
\practiceheading 用 Nsight Compute 比较标量版与 x4 版的加载指令、DRAM sectors 和占用率，
并用 Compute Sanitizer memcheck 覆盖未对齐子数组指针与 $N\bmod4=1,2,3$。
\sourcecheck{https://docs.nvidia.com/cuda/cuda-c-best-practices-guide/}{NVIDIA 对合并全局内存访问的官方建议；高置信度}
边界审查：若用条带式下标却按线程次序打包，会破坏稳定性；若无对齐检查强制解引用
\code{uint4}，行为不受保证。粗化也降低并行度并提高寄存器用量。
\end{solutionexercise}

\begin{solutionexercise}{4}
\problemheading 使用第 13 章的并行归并实现并行归并排序。实现应从长度为 1 的有序
run 开始逐轮两两归并，直到只剩一个有序 run；同时处理数组长度不是 2 的幂、末尾 run
不完整以及键值相等时的稳定次序。

\approachheading 自底向上把有序 run 宽度从 1 倍增到 2、4、8。每一遍的每对 run 再
切成固定输出片；线程用两个输出边界的协同秩确定输入子区间，并顺序归并少量元素。

\answerheading 下面的核心代码支持任意非 2 的幂 $N$、末尾不完整 run 和稳定的相等键处理；
长度与秩统一使用 \code{size_t}，适用范围还受设备内存和单次启动网格上限约束：
\begin{lstlisting}[language=C++]
__device__ size_t co_rank(size_t k,const int* A,size_t m,
                          const int* B,size_t n){
  size_t lo=k>n?k-n:0,hi=k<m?k:m;
  while(lo<hi){
    size_t i=lo+(hi-lo)/2,j=k-i;
    if(j>0 && i<m && B[j-1]>=A[i]) lo=i+1; else hi=i;
  }
  return lo; // ties from A precede ties from B
}
__device__ void serial_merge(const int* A,size_t m,const int* B,
                             size_t n,int* C){
  size_t i=0,j=0,k=0;
  while(i<m&&j<n) C[k++]=(A[i]<=B[j])?A[i++]:B[j++];
  while(i<m)C[k++]=A[i++]; while(j<n)C[k++]=B[j++];
}

template<int ITEMS>
__global__ void merge_pass(const int* src,int* dst,size_t N,size_t width,
                           size_t pairSpan,size_t tilesPerPair){
  size_t task=blockIdx.x,pair=task/tilesPerPair,tile=task%tilesPerPair;
  size_t base=pair*pairSpan; if(base>=N)return;
  size_t m=width<N-base?width:N-base;
  size_t remain=N-base-m, n=width<remain?width:remain;
  size_t total=m+n;
  size_t k0=tile*blockDim.x*ITEMS+threadIdx.x*ITEMS;
  size_t k1=min(k0+ITEMS,total); if(k0>=total)return;
  const int* A=src+base; const int* B=A+m;
  size_t i0=co_rank(k0,A,m,B,n),j0=k0-i0;
  size_t i1=co_rank(k1,A,m,B,n),j1=k1-i1;
  serial_merge(A+i0,i1-i0,B+j0,j1-j0,dst+base+k0);
}

// Host, with two N-element buffers src/dst and B=256:
// for (size_t w=1; w<N; w*=2) {
//   size_t span=(w>N-w)?N:2*w;
//   size_t pairs=(N-1)/span+1;
//   size_t tileItems=size_t(B)*ITEMS;
//   size_t tpp=(span-1)/tileItems+1;
//   // Split launches if pairs*tpp exceeds the device grid.x limit.
//   merge_pass<8><<<pairs*tpp,B>>>(src,dst,N,w,span,tpp);
//   std::swap(src,dst);
//   if(w>N/2) break; // result is now in src; avoids size_t overflow
// }
// After the loop, src points to the sorted buffer.  If the API requires the
// original allocation to own the result and src differs, copy N items back.
\end{lstlisting}
每遍开始时各 run 已有序；协同秩把一对 run 的输出划成不重叠且完整的片，每片稳定归并，
因此整对 run 有序。归纳到宽度不小于 $N$ 时，全数组有序。不同线程只写各自输出片。

\complexity{$O(N\log N)$ 工作，约 $\log_2N$ 次 launch 层}{$O(N)$ 辅助数组}{$O(N/\text{ITEMS})$ 个线程/遍}
\conclusionheading 该实现是稳定、支持任意非 2 的幂长度的并行自底向上归并排序；
\code{size_t} 消除了原代码的 32 位秩窄化，超大任务还须由主机拆分启动。
\discussionheading
\discussionpoint{自底向上归并把有序 run 逐轮翻倍。}初始每个元素是长度 1 的有序段，第 $p$ 轮把相邻长度 $2^p$ 的 run 归并为长度至多 $2^{p+1}$ 的 run，经过 $\lceil\log_2N\rceil$ 轮便覆盖全数组。每轮读取和写入 $O(N)$ 项，总工作为 $O(N\log N)$。例如 $N=10$ 时，宽度依次为 1、2、4、8，最后一轮的右 run 只有两项，截断端点仍保持证明成立。相等键时始终选择左 run 的记录，可由归纳证明每轮都保持原相对次序；若某轮改成右侧优先，后续轮无法恢复。这个层次同时对应外部排序把小文件段合成大段的过程，只是存储介质不同。

\discussionpoint{ping-pong 缓冲用空间换取简单的数据所有权。}每轮只从缓冲 A 读、向缓冲 B 写，下一轮交换角色，读写集合完全分离，因此不需要证明原地覆盖是否会毁掉尚未消费的键；代价是额外 $O(N)$ 存储。最终结果在哪个缓冲由 pass 奇偶决定，API 可以返回实际指针、只在必要时拷回，或安排初始角色使常见轮数落到目标位置。若无条件拷回，一个 1 GiB 数组就会额外产生约 2 GiB 的读写流量，而交换句柄几乎没有数据成本；容量规划还必须同时容纳两份峰值数据。双缓冲思想也用于迭代 stencil、图前沿和神经网络激活，所有权契约应与数据本身一起返回。

\discussionpoint{不同 run 尺度需要不同并行粒度。}早期 run 很短且数量多，为每个小归并执行二分切分会让控制开销占主导，可先在单块共享内存中排序出较大的初始 run；后期 run 很长，单线程顺序归并又会负载不足，应使用 merge path 或 co-rank 按输出对角线均匀分给线程块。最后一个右 run 可能为空或不完整，只需把端点截到 $N$，算法不应假定长度为 2 的幂。切换阈值应由共享容量、每次搜索成本和可用块数共同决定，并在阈值两侧测量，避免分派抖动。按 run 尺寸切换 kernel family，与 BLAS 按矩阵形状分派实现具有同样的分段优化思想。

\discussionpoint{归并排序的优势常来自比较器与外部有序段。}固定宽整数键通常适合吞吐更高的基数排序，但字符串、复合比较器或必须稳定搬运记录时，比较归并更自然；已有有序 run 时更无需从头基数分桶。若比较器只查看变长字符串到首个差异字节，其成本随数据共同前缀变化，固定“每次比较一次操作”的模型就会失真。数据库 merge join、LSM compaction 和外部排序的瓶颈还可能是 payload 搬运、压缩或存储带宽，而不是键比较。基准应改变重复键、非幂长度、记录宽度和超大索引，并与合适语义的 CUB/Thrust 基线比较端到端时间。这样才能从算法复杂度走向真实存储系统的选择依据。
\variantheading 若 $N=8$，宽度依次为 1、2、4，共 3 个 pass；交换三次后结果位于初始
辅助缓冲。若 $N=4$，只有 2 个 pass，结果回到初始输入缓冲。
\practiceheading 用 CUB \code{DeviceMergeSort} 或 Thrust \code{stable_sort} 作为基线，测试
$N=0,1$、非 2 的幂和重复键；主机封装应返回最终指针或按 pass 奇偶明确拷回，并检查
\code{width*=2} 的 \code{size_t} 溢出。
\sourcecheck{https://davidbader.net/publication/2012-gm-ba/2012-gm-ba.pdf}{GPU Merge Path 作者公开版的并行归并分区；并以 DOI 10.1145/2304576.2304621 核对正式出版信息；高置信度}
反例审查：每遍不能原地覆盖尚未读取的 run，必须 ping-pong。\code{width*=2} 在极大
\code{size_t} 上还应加入溢出保护；相等比较若改成 \code{<} 会改变跨 run 稳定次序。
\end{solutionexercise}

\chapter{矩阵乘法的高级优化}

\begin{solutionexercise}{1}
\problemheading 修改原书图 15.5 的设备函数，用向量加载替代标量加载。原
\code{loadTile} 依次接收全局 tile 指针 \code{T}、主维度 \code{lda}、有效边界
\code{maxRow/maxCol}、共享 tile 指针 \code{T_s}、共享主维度 \code{ldas} 以及
\code{height/width}。块内线程遍历 \code{height}\,$\times$\,\code{width} 输入 tile，
把合法的
\code{T[row*lda+col]} 写入共享内存
\code{T_s[row*ldas+col]}，越界项填 0。向量版本须保持这些边界语义，并处理向量宽度、
源地址对齐和不足一个向量的行尾。

\approachheading 把 tile 每行按四个 \code{float} 分组；线程遍历向量任务而非标量任务。
只有源地址 16 字节对齐且四项均在矩阵边界内时才解引用 \code{float4}，否则逐项加载并
为边界外位置填 0。

\answerheading
\begin{lstlisting}[language=C++]
__device__ __forceinline__ void loadTileVec4(
    const float* T,unsigned lda,unsigned maxRow,unsigned maxCol,
    float* S,unsigned ldas,unsigned height,unsigned width) {
  unsigned vecCols=(width+3)/4;
  unsigned tasks=height*vecCols;
  for(unsigned v=threadIdx.x;v<tasks;v+=blockDim.x) {
    unsigned row=v/vecCols, col=4*(v%vecCols);
    float a[4]={0,0,0,0};
    bool full=row<maxRow && col+3<maxCol && col+3<width;
    const float* src=full ? T+size_t(row)*lda+col : nullptr;
    if(full && ((reinterpret_cast<size_t>(src)&15)==0)) {
      float4 x=*reinterpret_cast<const float4*>(src);
      a[0]=x.x;a[1]=x.y;a[2]=x.z;a[3]=x.w;
    } else {
#pragma unroll
      for(unsigned j=0;j<4;++j)
        if(row<maxRow && col+j<maxCol && col+j<width)
          a[j]=T[size_t(row)*lda+col+j];
    }
#pragma unroll
    for(unsigned j=0;j<4;++j)
      if(col+j<width) S[row*ldas+col+j]=a[j];
  }
}
\end{lstlisting}
线性向量任务覆盖 tile 中每个逻辑元素恰好一次；完整向量与标量回退写入相同共享下标，
所以转换不改变数值。函数不含屏障，调用者必须在所有线程调用后执行一次
\code{__syncthreads()} 才能使用 tile。

\complexity{$O(hw)$ 总搬运工作，每线程约 $hw/B$ 项}{$O(1)$ 每线程寄存器；共享 tile 由调用方提供}{$B$ 线程合作加载}
\conclusionheading 内部完整区域使用一条 128 位加载搬运四个标量，边缘和未对齐行保持正确。
\discussionheading
\discussionpoint{向量指令减少发射，不保证减少内存事务。}一次 \code{float4} 可以替代四条标量加载，少做地址计算并降低指令发射压力；但 DRAM 看到的是整个 warp 覆盖的地址扇区，而不是 C++ 类型名称。若相邻线程的标量访问本来连续且对齐，四轮标量加载也可能完全合并；反之，各线程向量起点分散时，一条 \code{float4} 仍会触及许多扇区。性能验证必须同时看 load instructions 和 global sectors，再结合时间判断瓶颈在指令前端还是内存系统。这个区分也适用于 CPU SIMD、网络 scatter-gather 与存储批量 I/O。

\discussionpoint{二维 pitch 不能替子视图证明向量对齐。}行首 pitch 可能按 128 字节对齐，但子矩阵从列 1 开始时，\code{float} 地址只偏移 4 字节，已经不满足 \code{float4} 的 16 字节要求。实现应先用基矩阵尺寸验证绝对行、列以及后续四项都在对象内，再计算地址并选择向量或标量路径；不能先形成越界派生指针，期待条件为假时不解引用便自动合法。还要遵守 C++ 对齐与别名规则，可使用正确向量类型或受支持的拷贝原语。类似风险广泛存在于张量 view、图像 ROI 和带 leading dimension 的 BLAS 接口。

\discussionpoint{异步 tile 搬运把边界正确性扩展为流水协议。}协作线程可在计算当前共享 tile 时预取下一 tile，用双缓冲隐藏全局延迟；异步拷贝需要提交组、等待点和复用屏障，确保消费者只读取已经完成的数据。边缘 tile 的越界位置仍要显式填零或掩蔽，不能因为拷贝异步便省略容量检查。若第 $k+1$ 轮过早复用仍被第 $k-1$ 轮消费者读取的缓冲，即使拷贝本身不越界，也会出现跨迭代数据竞态。较新的 bulk 或 tensor memory copy 还依赖布局描述和对齐契约，描述符错误可能让整个 tile 错位。这个机制与软件流水、DMA 双缓冲和网络收发环相同：性能来自重叠，正确性来自每一级缓冲的所有权与完成事件。

\discussionpoint{矩阵内核应把主区域与边界路径分别验收。}GEMM、卷积和注意力通常让完整对齐 tile 走快速向量路径，不完整边缘用谓词、标量回退或专门 kernel；若只测尺寸恰为 tile 倍数，最危险路径从未执行。以 $128\times128$ tile 为例，$256\times256$ 只覆盖主路径，而 $257\times129$ 会同时触发右边缘、下边缘和角落。基准应覆盖对齐主矩阵、错位子视图、窄列和各种尾宽，记录指令、扇区、有效带宽和总时间，并与标量 CPU 参考逐元素比较。还要用非单位 pitch 检查行寻址，而不仅是紧密数组。这样才能区分“快路径确实更快”和“为了快路径牺牲通用接口正确性”两种完全不同结果。
\variantheading 若 \code{width=10}，每行前两个四元组可在对齐时向量加载，最后一组只含列
8、9，必须标量加载两项并把另外两项填 0。
\practiceheading 用 CUTLASS 的协作 tile copy 或 \code{cuda::memcpy_async} 实现作性能基线，
并以 Compute Sanitizer memcheck 覆盖偏移子矩阵、\code{lda\%4!=0} 和尾行；Nsight Compute
可统计向量与标量加载指令比例。
\sourcecheck{https://docs.nvidia.com/cuda/cuda-c-programming-guide/}{NVIDIA CUDA C++ Programming Guide 的内建向量类型与对齐要求；高置信度}
对抗审查：\code{cudaMalloc} 的基址对齐并不保证任意行地址仍对齐；若 \code{lda} 不是
4 的倍数，部分行会走标量回退。不能把未对齐 \code{float*} 无条件强转成 \code{float4*}。
\end{solutionexercise}

\begin{solutionexercise}{2}
\problemheading 修改原书图 15.7 的设备函数，用向量存储替代标量存储。原
\code{writeTile} 依次接收全局输出指针 \code{c}、主维度 \code{ldc}、有效边界
\code{maxRow/maxCol}、寄存器 tile \code{C_r} 以及尺寸 \code{m/n}。它遍历寄存器中的
$m\times n$ 输出 tile，并仅在行、列分别小于 \code{maxRow}、\code{maxCol} 时把
\code{C_r[row][col]} 写到
\code{c[row*ldc+col]}。向量版本须保留逐元素边界语义，并正确处理目的地址对齐和不完整
行尾。

\approachheading 沿列每四项打包成 \code{float4}；整组在输出边界内且目的地址对齐时
执行向量存储，否则标量写合法尾部。

\answerheading
\begin{lstlisting}[language=C++]
template<int TN>
__device__ __forceinline__ void writeTileVec4(float* C,size_t ldc,
    size_t M,size_t N,size_t tileRow,size_t tileCol,
    const float Creg[][TN],unsigned m,unsigned n) {
  if(tileRow>=M || tileCol>=N)return;
#pragma unroll
  for(unsigned row=0;row<m;++row) {
    if(size_t(row)>=M-tileRow)break;
    size_t globalRow=tileRow+row;
#pragma unroll
    for(unsigned col=0;col<n;col+=4) {
      if(size_t(col)>=N-tileCol)break;
      size_t globalCol=tileCol+col;
      size_t valid=N-globalCol;
      bool full=(n-col>=4 && valid>=4);
      float* dst=C+globalRow*ldc+globalCol;
      if(full) {
        if((reinterpret_cast<size_t>(dst)&15)==0) {
          float4 v=make_float4(Creg[row][col],Creg[row][col+1],
                               Creg[row][col+2],Creg[row][col+3]);
          *reinterpret_cast<float4*>(dst)=v;
        } else {
#pragma unroll
          for(unsigned j=0;j<4;++j) dst[j]=Creg[row][col+j];
        }
      } else {
#pragma unroll
        for(unsigned j=0;j<4;++j)
          if(col+j<n && size_t(j)<valid) dst[j]=Creg[row][col+j];
      }
    }
  }
}
\end{lstlisting}
每个输出元素仍由原拥有线程写一次，故没有新增竞态。向量化只合并单线程相邻四项；跨线程
是否形成完整合并事务还取决于线程级 tile 的排列，下一题会解决这一点。

\complexity{$O(mn)$ 存储工作，完整区约 $mn/4$ 条存储指令}{$O(1)$ 额外空间}{所有输出 tile 所有者并行写回}
调用方必须传矩阵基址与线程 tile 的绝对起始坐标，不能预先计算可能越界的
\code{C+tileRow*ldc+tileCol} 再传入。函数先证明起始坐标和当前行列合法，之后才形成
目的指针。

\conclusionheading 该函数在满足边界与对齐时使用 \code{float4}，其余路径逐项安全写回；
边界检查发生在任何派生目的指针形成之前。
\discussionheading
\discussionpoint{连续的线程私有片段还不等于合并的 warp 写回。}线程拥有四个相邻结果，只说明它可发出一条 \code{float4}；要让 warp 覆盖少量连续扇区，线程 $t+1$ 的四元组还应紧接线程 $t$。若逻辑 tile 把相邻线程映射到不同行、隔列或不同批次，每条向量都合法，整体地址却仍是分散的。例如 lane $t$ 写列 $4t$ 时相邻，而写列 $4t$、行 $t$ 时地址间隔约为 leading dimension，事务无法合并。可以先画出 lane 到绝对行列的表，再计算一个 warp 的地址跨度和扇区数量，这比只检查每线程指针更可靠。相同所有权分析决定 GEMM epilogue、图像通道写回和结构数组存储能否真正合并。

\discussionpoint{尾部写回必须尊重对象边界而非事后修补。}只有目的绝对行列合法、连续四项均存在且地址对齐时，才能形成并解引用 \code{float4*}；未对齐完整组和不足四项尾部应逐项谓词写。先覆盖邻接对象再把它恢复，在并发线程存在时会产生竞态，也可能破坏另一个逻辑结果；先计算对象外指针再避免 store，同样不是稳健的 C++ 接口。比如一行只剩 6 项时，可安全写一个四元组再标量写 2 项，绝不能让第二个四元组跨到下一行。用宽度 0--7、错位列和最后一行做测试，可以覆盖每种标量回退。这个原则也适用于 SIMD 掩码存储和文件尾部块，边界必须在形成访问之前成立。

\discussionpoint{缓存策略取决于写后是否很快复用。}向量 store 减少指令，但总吞吐还受事务合并、L2 写分配、ECC 和内存控制器影响；纯流式输出若之后很久不用，缓存驻留可能挤掉更有价值数据。若下一 kernel 立即读取同一矩阵，让结果保留在 L2 又可能明显缩短端到端时间。可用“写后立刻读”和“写后先扫过大于 L2 的干扰数组”两组实验隔离复用收益，而不能把冷热场景混成一个平均值。不同架构提供的缓存或流式提示语义并不完全相同，应把生产者和消费者一起测量，而不是只计 store kernel。这个生产--消费距离的视角也适用于数据库物化、图中间 frontier 和深度学习算子融合。

\discussionpoint{epilogue 融合把向量写回变成数值转换边界。}高性能 GEMM 常在累加器仍位于寄存器时加入偏置、激活、缩放或量化，再一次性写出，从而避免中间张量往返全局内存。融合会引入新的语义：低精度转换需要舍入与饱和规则，ReLU 要明确 NaN 行为，按通道缩放又改变广播索引。测试应覆盖完整向量和每种尾宽、任意 leading dimension、NaN、极值与非单位缩放，并比较输出误差。只有在数值契约一致时，store instructions 与 global sectors 的下降才代表真正可替代的行业实现。
\variantheading 若 \code{n=6,maxCol=5}，列 0--3 在对齐时用一次 \code{float4} 写，列 4
由尾部标量路径写出，列 5 因超出 \code{maxCol} 不写。
\practiceheading 用 Nsight Compute 检查 Global Store Efficiency 和 store 指令数，并以
cuBLAS 输出作数值基线；Compute Sanitizer memcheck 应覆盖未对齐列偏移、部分行以及
\code{maxCol<n} 的裁剪视图。
\sourcecheck{https://docs.nvidia.com/cuda/cuda-c-best-practices-guide/}{NVIDIA Best Practices Guide 的向量加载/存储建议；高置信度}
边界审查：线程自己的四项连续不代表 warp 的向量起点连续；旧的 $8\times8$ 映射仍可能
跨步。向量写也不能越过 \code{maxCol} 后再依赖掩码补救，因为写指令本身没有逐 lane 掩码。
\end{solutionexercise}

\begin{solutionexercise}{3}
\problemheading 按原书第 15.5 节所述的线程级输出 tile 重排方式，重新实现本章的分块
矩阵乘法内核。具体映射是：一个 256 线程块计算 $128\times128$ 输出 tile，8 个 warp
按 $2\times4$ 排列、每个负责 $64\times32$；每个 warp 的 32 个线程按 $8\times4$
排列，每线程从四个 $32\times16$ 象限各取得一个 $4\times4$ 物理 tile，合成逻辑
$8\times8$ 输出 tile。实现应使相邻线程逐行写出相邻的四元素向量，并处理矩阵边界。

\approachheading 一个 256 线程块由 8 个 warp 构成，按 $2\times4$ 覆盖
$128\times128$ 输出 tile；每 warp 覆盖 $64\times32$。lane 按 $8\times4$ 排列，
从四个 $32\times16$ 象限各取一个 $4\times4$ tile。每行的四个 lane 因而写相邻
\code{float4}。

\answerheading 以下内核计算行主序 $C_{M\times N}=A_{M\times K}B_{K\times N}$：
\begin{lstlisting}[language=C++]
template<int BM=128,int BN=128,int BK=8>
__global__ void gemm_rearranged(const float* A,const float* B,float* C,
                                int M,int N,int K) {
  static_assert(BM==128 && BN==128 && BK==8,"mapping constants");
  __shared__ float As[BM][BK+1]; // padding removes A bank conflicts
  __shared__ float Bs[BK][BN];
  float acc[4][4][4]={0};
  int tid=threadIdx.x, warp=tid>>5, lane=tid&31;
  int warpRow=warp/4, warpCol=warp%4;
  int laneRow=lane/4, laneCol=lane%4;
  int localRow=warpRow*64+laneRow*4;
  int localCol=warpCol*32+laneCol*4;
  int blockRow=blockIdx.y*BM, blockCol=blockIdx.x*BN;

  for(int kb=0;kb<K;kb+=BK) {
    for(int p=tid;p<BM*BK;p+=blockDim.x) {
      int r=p/BK,k=p%BK,gr=blockRow+r,gk=kb+k;
      As[r][k]=(gr<M && gk<K)?A[size_t(gr)*K+gk]:0.0f;
    }
    for(int p=tid;p<BK*BN;p+=blockDim.x) {
      int k=p/BN,c=p%BN,gk=kb+k,gc=blockCol+c;
      Bs[k][c]=(gk<K && gc<N)?B[size_t(gk)*N+gc]:0.0f;
    }
    __syncthreads();
#pragma unroll
    for(int k=0;k<BK;++k) {
#pragma unroll
      for(int q=0;q<4;++q) {
        int ro=(q/2)*32,co=(q%2)*16;
#pragma unroll
        for(int r=0;r<4;++r) {
          float av=As[localRow+ro+r][k];
#pragma unroll
          for(int c=0;c<4;++c)
            acc[q][r][c]+=av*Bs[k][localCol+co+c];
        }
      }
    }
    __syncthreads();
  }

#pragma unroll
  for(int q=0;q<4;++q) {
    int ro=(q/2)*32,co=(q%2)*16;
#pragma unroll
    for(int r=0;r<4;++r) {
      int gr=blockRow+localRow+ro+r, gc=blockCol+localCol+co;
      if(gr<M && gc<N) {
        float* dst=C+size_t(gr)*N+gc;
        if(gc+3<N && ((reinterpret_cast<size_t>(dst)&15)==0)) {
          *reinterpret_cast<float4*>(dst)=make_float4(
              acc[q][r][0],acc[q][r][1],acc[q][r][2],acc[q][r][3]);
        } else for(int c=0;c<4 && gc+c<N;++c) dst[c]=acc[q][r][c];
      }
    }
  }
}
// dim3 block(256), grid((N+127)/128,(M+127)/128);
\end{lstlisting}
每个 $(i,j)$ 输出由唯一的块、warp、lane、象限和片内坐标确定；内层遍历全部 $k$ tile，
边缘用 0 填充，故累加值就是点积。两次块屏障分别保护加载完成与下一轮覆盖。四个同
\code{laneRow} 的相邻 lane 写连续 16 字节段，形成连续 64 字节区域；四个象限完整覆盖
warp tile 且互不重叠。

\complexity{$O(MNK)$ 工作；每线程每 $K$ tile 做 512 次 FMA}{$64$ 个累加寄存器/线程，约 8.5 KiB 共享空间/块}{$\lceil M/128\rceil\lceil N/128\rceil$ 块，256 线程/块}
\conclusionheading 重排保持每线程 $8\times8$ 的总输出量，同时实现 warp 级合并向量写。
\discussionheading
\discussionpoint{三级 tile 是描述 GEMM 数据所有权的共同语言。}线程块 tile 决定从全局内存搬入的 $A/B$ 区域以及多少输出复用它们，warp tile 决定一组 32 线程怎样协作读取共享数据，线程 tile 则规定每个线程在寄存器中累加哪些 $C_{ij}$。本题的 $128\times128$、$64\times32$、$8\times8$ 并非三个独立数字：块 tile 必须能被 warp tile 覆盖，warp tile 又要由线程输出无重无漏地铺满。改变任一层都会连带改变共享布局、加载分工和 epilogue 写回。建立这种层次地图后，才能系统推导索引，而不是从大量魔数中猜测所有权。

\discussionpoint{寄存器复用提高算术强度，也可能减少可用并发。}每线程保存 $8\times8=64$ 个累加器，使一小片输入被重复用于许多 FMA，减少共享读取；加上地址、$A/B$ 片段和双缓冲状态后，寄存器总数可能很高。若每 SM 因此只能驻留一个块，occupancy 下降；若编译器溢出到 local memory，原本片上的累加器甚至会产生全局流量。不过低 occupancy 不必然慢，足够的指令级并行和异步预取仍可隐藏延迟。应结合 ptxas 寄存器/local 报告、驻留块数和时间判断，而不是追求某个固定 occupancy 百分比。

\discussionpoint{Tensor Core 路径重新定义了 warp 内数据布局。}标量 FMA 版本把共享元素加载到普通寄存器后逐项乘加，是理解 tile 边界和正确性最透明的基线；现代实现会用多级异步预取，并让 warp 发出 MMA 指令处理专门的矩阵 fragment。fragment 的 lane 映射、允许类型、累加精度和对齐要求都由指令族规定，不能从标量 $8\times8$ 所有权直接推断。CUTLASS 将块、warp、instruction tile 和 epilogue 参数化，正是为了管理这些耦合约束。迁移时应先保持数学与边界参考，再逐层替换搬运和计算机制。

\discussionpoint{矩阵形状和数学模式决定应该使用哪个 kernel family。}大方阵能充分填满固定 $128\times128$ tile，瘦高、瘦宽、小批量或任一维小于 tile 时会产生大量掩码工作，甚至只有少数 warp 有效；生产库通常按形状分派不同 tile。比较 cuBLAS 或 CUTLASS 时还要固定数据类型、Tensor Core/TF32 模式、转置和累加精度，否则 TFLOP/s 与误差都不可比。测试需覆盖 $M,N,K$ 各自的尾部、非单位 leading dimension 和量级悬殊输入，并报告有效而非填充 FLOP。这样性能数字才同时反映形状适配和数值契约。
\variantheading 若 $N=130$，最右输出块只有两列有效，多数 \code{float4} 写回退化为标量；
数值仍正确，但合并写优势显著下降，适合为窄尾块另设内核。
\practiceheading 以 cuBLAS SGEMM 和 CUTLASS 相同数据类型配置作基线，用 ptxas 报告核对
每线程寄存器与 spill，再用 Nsight Compute 的 Roofline、Occupancy 和内存事务选择
\code{BK} 与输出 tile。
\sourcecheck{https://github.com/NVIDIA/cutlass}{NVIDIA CUTLASS 官方仓库的层次化 GEMM tile 与合并复制实现；高置信度}
对抗审查：大量累加器可能造成寄存器溢出或低占用率，性能必须实测。边缘矩阵或 $N$ 非
4 倍数会退回标量写；缺少任一轮末屏障都会让部分线程在其他线程仍计算时覆盖共享 tile。
\end{solutionexercise}


\part{高级模式与应用}
\chapter{动态规划与波前并行}

\begin{solutionexercise}{1}
\problemheading 实现使用正方形线程块的 Floyd--Warshall 内核，而不是原书第 16.4 节
的一维线程块实现。把距离表划分成正方形数据片并分配给各线程块；块内线程把第 $k$ 列和
第 $k$ 行的相关部分存入共享内存以供复用。比较两种实现的内存访问、数据复用和性能。
两种实现都应按
$d(i,j,k)=\min\{d(i,j,k-1),d(i,k,k-1)+d(k,j,k-1)\}$ 更新，并正确处理不可达距离
的哨兵值和距离相加的溢出边界。

\approachheading 一个 $B\times B$ 块覆盖距离表的方形 tile。tile 只需从第 $k$ 列
加载 $B$ 项、从第 $k$ 行加载 $B$ 项，供 $B^2$ 个输出复用。每个 $k$ 仍是一次独立内核
启动，形成全局波前屏障。

\answerheading 使用 64 位距离表。取 \code{INF=LLONG_MAX/4}，假定输入无负环、对角线为 0，
且算法过程中每个有限距离的绝对值始终小于 \code{INF/2}；于是两个有限距离相加既不会碰到
哨兵，也不会发生 64 位有符号溢出：
\begin{lstlisting}[language=C++]
template<int B>
__global__ void fw_square(long long* d,int n,int k,long long INF) {
  __shared__ long long dik[B], dkj[B];
  int i=blockIdx.y*B+threadIdx.y;
  int j=blockIdx.x*B+threadIdx.x;
  if(threadIdx.x==0)
    dik[threadIdx.y]=(i<n)?d[size_t(i)*n+k]:INF;
  if(threadIdx.y==0)
    dkj[threadIdx.x]=(j<n)?d[size_t(k)*n+j]:INF;
  __syncthreads();
  if(i<n && j<n && i!=k && j!=k && dik[threadIdx.y]!=INF &&
     dkj[threadIdx.x]!=INF) {
    long long via=dik[threadIdx.y]+dkj[threadIdx.x];
    long long old=d[size_t(i)*n+j];
    if(via<old) d[size_t(i)*n+j]=via;
  }
}
// dim3 block(B,B), grid((n+B-1)/B,(n+B-1)/B);
// for(int k=0;k<n;++k) fw_square<16><<<grid,block>>>(d,n,k,INF);
\end{lstlisting}
第 $k$ 轮开始时 \code{d[i,k]} 与 \code{d[k,j]} 是递推式的两个输入。无负环且
\code{d[k,k]=0} 时第 $k$ 行、列本轮不变；代码显式跳过对它们的写入，从而让共享加载
与全局访问也不存在读写竞态。不同 $(i,j)$ 各有唯一写者。距离表与中间和都采用
\code{long long}，因此不存在把 64 位结果无范围检查地窄化回 \code{int}；上述输入契约
同时排除了 64 位加法溢出。

一维版本每块复用一个 \code{d[i,k]}，但每个线程仍读自己的 \code{d[k,j]}；方块版本
用约 $2B$ 次全局加载服务 $B^2$ 个输出，把两条输入边都复用。输出读写量不变，且二维块
多一次共享加载/屏障；实际收益取决于缓存、$B$、占用率和图规模，应计时比较。

\complexity{$O(n^3)$ 工作、$n$ 个全局波前}{$O(n^2)$ 距离表，$2B$ 个共享整数/块}{$O(n^2)$ 线程任务/轮}
\conclusionheading 正方形 tile 同时复用第 $k$ 行和列，理论全局输入流量较一维版本更低。
\discussionheading
\discussionpoint{Floyd--Warshall 是 min-plus 代数上的闭包。}普通矩阵乘法用乘法连接路径、用加法汇总候选；最短路把“连接”换成距离相加，把“汇总”换成取最小，形成 min-plus 运算。第 $k$ 阶段只允许顶点 $0\ldots k$ 作为中间点，递推 $D^{(k)}_{ij}=\min(D^{(k-1)}_{ij},D^{(k-1)}_{ik}+D^{(k-1)}_{kj})$ 在旧距离与经 $k$ 的两段距离之间取小者，因此 $k$ 必须按顺序推进。固定 $k$ 时所有 $(i,j)$ 更新相互独立，才是 GPU 可以展开的二维并行面。这个代数视角还连接到传递闭包、Viterbi 和热带代数，说明替换底层运算就能复用矩阵式依赖结构。

\discussionpoint{二维 tile 只是复用起点，完整分块还要处理阶段依赖。}固定 $k$ 时，一个块内多项 $(i,j)$ 会反复读取第 $k$ 行和第 $k$ 列，把它们放入共享内存可减少全局访问。更深的 blocked Floyd--Warshall 一次处理一组 $k$：先更新包含该组的对角 tile，再更新同排和同列 tile，最后更新其余 tile；后三类所需数据由前一阶段产生，不能任意并行。这样能跨多个 $k$ 复用片上数据，却增加全局阶段同步和共享容量。它与块 LU 分解的面板、行列、尾部更新结构相似，分块优化必须尊重数据依赖而非只追求大 tile。

\discussionpoint{不可达、溢出和负环是三种不同语义边界。}若 $d_{ik}$ 或 $d_{kj}$ 为不可达哨兵，就不能直接相加，否则有限整数哨兵可能溢出成负数；两个有限距离相加也应在足够宽类型中检查范围。例如以 \code{INT_MAX} 表示不可达时，再加正权会发生有符号溢出，错误的负数反而会被 \code{min} 选中。负边本身允许存在，负环则只影响能够先到达该环、又能从环到达目标的点对，这些点对可无限绕环降低代价；图中其他点对仍可能有有限最短路或不可达。对角线负值可发现负环，但还需结合前后可达性传播受影响范围。这个精细分类比笼统说“有负环就全图无解”更适合实际路由与约束系统。

\discussionpoint{图密度决定全源最短路的算法家族。}Floyd--Warshall 固定做 $O(n^3)$ 更新，控制规则且适合稠密邻接矩阵；稀疏大图若只有 $m\ll n^2$ 条边，重复 Dijkstra、Johnson 或专用稀疏 GPU 算法通常少做许多工作。比较时应报告真正执行的 min-plus 更新率、全局字节、tile 复用与时间，而不仅是理论 FLOP；负边又会限制 Dijkstra 的适用性。测试需包含不可达点、负边、局部负环和无负环图，并与宽整数 CPU 参考比对。这样算法选择由密度、权值语义和硬件访问共同决定。
\variantheading 若 $B=16,n=1000$，网格为 $63\times63$，每轮启动 3969 个块；最右和最下
tile 只有 8 个有效行或列，越界线程仍需参加共享加载后的屏障。
\practiceheading 与 cuGraph 或 CPU Floyd--Warshall 结果对照含不可达点和负边但无负环的图，
用 Compute Sanitizer 检查边缘 tile；启动前验证 $B^2$ 不超过每块线程上限并以 Nsight
Compute 比较 $B=8,16,32$ 的占用率与全局流量。
\sourcecheck{https://docs.nvidia.com/cuda/cuda-c-programming-guide/}{NVIDIA 对共享内存协作加载与块级同步的官方说明；高置信度}
反例审查：存在负环或 \code{d[k,k]<0} 时第 $k$ 行列可能改变，原地并行更新不再符合该
证明，应双缓冲。\code{B*B} 不得超过设备每块线程上限，通常取 16 或 32。
\end{solutionexercise}

\begin{solutionexercise}{2}
\problemheading 修改采用线程块级分块的 Smith--Waterman 内核，使其使用矩形数据片。
设数据片高、宽分别为 $R,C$，片内仍须沿 $R+C-1$ 条反对角线推进；每个单元依赖北、
西和西北单元，并按
$H_{r,q}=\max(0,H_{r-1,q-1}+s,H_{r,q-1}+g,H_{r-1,q}+g)$ 更新。实现须处理矩阵边界、
片内同步和相邻数据片的依赖。

\approachheading 对 $R\times C$ tile 使用带一行、一列 halo 的共享数组。块内按
$R+C-1$ 条反对角线推进；令一个线程对应局部列，局部行为 $d-q$。主机仍按 tile 反对角线
依次启动内核。

\answerheading 以下设备例程还会被后续同步版本复用；\code{rows/cols} 包含得分矩阵的
全 0 第 0 行/列。
\begin{lstlisting}[language=C++]
template<int R,int C>
__device__ void swTile(int* H,const char* rea,const char* ref,
    int rows,int cols,int tr,int tc,int MATCH,int MISMATCH,int GAP,int* s) {
  constexpr int LD=C+1;
  int row0=1+tr*R,col0=1+tc*C;
  for(int q=threadIdx.x;q<C;q+=blockDim.x)
    s[q+1]=(col0+q<cols)?H[size_t(row0-1)*cols+col0+q]:0;
  for(int r=threadIdx.x;r<R;r+=blockDim.x)
    s[(r+1)*LD]=(row0+r<rows)?H[size_t(row0+r)*cols+col0-1]:0;
  if(threadIdx.x==0)s[0]=H[size_t(row0-1)*cols+col0-1];
  __syncthreads();
  for(int wave=0;wave<R+C-1;++wave) {
    int q=threadIdx.x,r=wave-q,gr=row0+r,gc=col0+q;
    if(q<C && r>=0 && r<R && gr<rows && gc<cols) {
      int nw=s[r*LD+q],n=s[r*LD+q+1],w=s[(r+1)*LD+q];
      int sub=(rea[gr-1]==ref[gc-1])?MATCH:MISMATCH;
      int v=max(0,max(nw+sub,max(w+GAP,n+GAP)));
      s[(r+1)*LD+q+1]=v;
    }
    __syncthreads();
  }
  for(int p=threadIdx.x;p<R*C;p+=blockDim.x) {
    int r=p/C,q=p%C,gr=row0+r,gc=col0+q;
    if(gr<rows&&gc<cols)H[size_t(gr)*cols+gc]=s[(r+1)*LD+q+1];
  }
  __syncthreads();
}

template<int R,int C>
__global__ void sw_rect_wave(int* H,const char* rea,const char* ref,
    int rows,int cols,int tileWave,int MATCH,int MISMATCH,int GAP) {
  static_assert(C>0 && C<=1024,"one thread is required per tile column");
  __shared__ int s[(R+1)*(C+1)];
  int tr=blockIdx.x,tc=tileWave-tr;
  int ntr=(rows-1+R-1)/R,ntc=(cols-1+C-1)/C;
  if(tr<ntr && tc>=0 && tc<ntc)
    swTile<R,C>(H,rea,ref,rows,cols,tr,tc,MATCH,MISMATCH,GAP,s);
}
// If ntr==0 or ntc==0, return on the host without launching.
// Otherwise, after checking C<=maxThreadsPerBlock and shared-memory use:
// for wave=0..ntr+ntc-2: sw_rect_wave<R,C><<<ntr,C>>>(...) in one stream.
\end{lstlisting}
halo 提供 tile 外依赖，内部每个单元只读取前两条已经完成的反对角线。主机的相邻 launch
在同一 stream 中有序，故上方、左方 tile 已写回。每单元唯一写者，无原子操作。

\complexity{$O((rows-1)(cols-1))$ 工作；每 tile $R+C-1$ 轮}{$(R+1)(C+1)$ 个共享整数/块}{每条 tile 波前多个块，每块 $C$ 线程}
\conclusionheading 矩形 $R,C$ 可分别适配序列形状与硬件资源，正确性不依赖 $R=C$。
\discussionheading
\discussionpoint{二维动态规划暴露出反对角线波前。}Smith--Waterman 单元 $(r,q)$ 只依赖北 $(r-1,q)$、西 $(r,q-1)$ 和西北 $(r-1,q-1)$，这些前驱的坐标和都小于当前点；具有相同 $r+q$ 的单元彼此没有依赖，可以同时计算。矩形 $R\times C$ tile 不改变这个偏序，只让片内波前先增长到 $\min(R,C)$，再保持或缩短，并改变相邻 tile 的边界形状。正确实现应先画依赖 DAG，再安排线程与屏障，而不是从二维线程坐标直接猜并行性。编辑距离、动态时间规整和许多 stencil 都共享这种波前结构。

\discussionpoint{矩形 tile 用形状适配序列与硬件资源。}片内共有 $R+C-1$ 条反对角线，共享表若保存边界通常按 $(R+1)(C+1)$ 增长；增大 $C$ 能提高中段同波前宽度，却可能需要更多线程或让每线程处理多列。增大 $R$ 会让一个块经历更多串行波前，并减少可并行 tile 行数。例如从 $16\times16$ 改为 $8\times32$ 保持 256 个单元，却把波前数从 31 增至 39，并改变峰值活动宽度与边界形状。长短序列高度不对称时，方形 tile 未必有效，矩形可以沿长维提供吞吐、沿短维节约无效槽。最优 $R,C$ 因线程上限、共享容量和序列长宽比分段变化，与矩阵乘法按形状选择 tile 的原因相同。

\discussionpoint{最高分与完整回溯对应不同的存储问题。}局部比对用零截断使表值非负，却不防止长序列在持续匹配奖励下超过 16 位甚至 32 位范围；可用最大可能匹配数乘奖励上界选择类型。若只求最高分，每轮只需保存有限边界或滚动行；要恢复对齐路径，则需记录每格来源方向、周期检查点，或在选定区域重新计算。即使每格方向只占 2 位，$m\times n$ 表的存储仍随面积增长，而滚动分数仅随较短边增长。方向可压缩成少量位，但仍会改变内存带宽和 tile 设计。因而一个“更快的分数 kernel”不一定是完整比对器的最佳核心，API 应先明确是否需要位置和 traceback。

\discussionpoint{生产比对通过约束搜索区域而非总填满矩阵。}已知序列大致对齐时可用带状限制只计算主对角附近，X-drop 在分数落后当前最佳过多时提前剪枝，批处理又把相近长度序列放在一起减少尾部浪费。若长度约为 $n$、半带宽为 $w\ll n$，候选单元可从 $O(n^2)$ 降到约 $O(nw)$，但前提是最优路径确实没有离开该带。多 GPU 系统还需切分批次或波前并交换边界，支持 DPX 的设备则可压缩单元整数指令。验证应覆盖空序列、全失配、多个同分最优位置和不完整 tile，并区分只核对最高分还是连路径也一致。性能报告用 GCUPS 时还应说明实际计算单元数，不能把被剪枝单元也计入有效工作。
\variantheading 若 $R=32,C=16$，每 tile 有 $47$ 条反对角线、需 561 个共享整数并以
16 个线程启动；相比 $16\times16$，行方向复用增加而单线程承担更多行波次。
\practiceheading 用 NCBI 风格已知比对样例和 CPU Smith--Waterman 作黄金结果，覆盖空序列、
末尾部分 tile 和不同 gap 分数；启动前查询每块线程与共享内存上限，再用 Nsight Compute
调优 $R,C$。
\sourcecheck{https://doi.org/10.1016/0022-2836(81)90087-5}{Smith 与 Waterman 的局部序列比对原始论文；高置信度}
边界审查：若 $C>1024$ 就不能一列一线程；应让线程跨步处理列。所有线程必须执行每轮
屏障，故不能在 tile 内因全局越界而提前返回。
\end{solutionexercise}

\begin{solutionexercise}{3}
\problemheading 实现采用线程块级分块、并用协作组进行跨块同步的 Smith--Waterman
内核。每个数据片内部沿反对角线计算，数据片之间也按反对角波前推进；合作启动的网格必须
保证所有驻留块都能到达每次网格范围同步，并正确处理不完整的边缘数据片。

\approachheading 把主机的 tile 波前循环移入合作内核。每个物理块固定对应一个 tile 行；
每轮选择相应 tile 列，活动块调用完整 tile 例程，随后全体块无条件执行 \code{grid.sync()}。

\answerheading
\begin{lstlisting}[language=C++]
#include <cooperative_groups.h>
namespace cg=cooperative_groups;

template<int T>
__global__ void sw_cooperative(int* H,const char* rea,const char* ref,
    int rows,int cols,int MATCH,int MISMATCH,int GAP) {
  static_assert(T>0 && T<=1024,"one thread is required per tile column");
  cg::grid_group grid=cg::this_grid();
  __shared__ int s[(T+1)*(T+1)];
  int ntr=(rows-1+T-1)/T,ntc=(cols-1+T-1)/T;
  int tr=blockIdx.x;
  for(int wave=0;wave<ntr+ntc-1;++wave) {
    int tc=wave-tr;
    if(tr<ntr && tc>=0 && tc<ntc)
      swTile<T,T>(H,rea,ref,rows,cols,tr,tc,MATCH,MISMATCH,GAP,s);
    grid.sync(); // every block, including idle blocks, must participate
  }
}

// Host must reject ntr==0 or ntc==0, then query cooperative support,
// maxThreadsPerBlock, shared-memory capacity, and occupancy first:
// int perSM; cudaOccupancyMaxActiveBlocksPerMultiprocessor(
//   &perSM,sw_cooperative<32>,32,0);
// require(ntr <= perSM * deviceProp.multiProcessorCount);
// void* args[]={&H,&rea,&ref,&rows,&cols,&match,&mismatch,&gap};
// cudaLaunchCooperativeKernel((void*)sw_cooperative<32>,ntr,32,args,0,stream);
\end{lstlisting}
第 $w$ 次网格同步之前只计算满足 $tr+tc=w$ 的 tile；其依赖 tile 位于更小波前，已在
上一次或更早同步前写完。\code{grid.sync()} 同时提供执行与内存可见性屏障。

\complexity{$O(RC)$ 得分工作，加 $ntr+ntc-1$ 次网格屏障}{$O(RC)$ 得分表，$(T+1)^2$ 共享整数/块}{$ntr$ 个常驻合作块}
\conclusionheading 单个合作内核消除逐波前 launch，但网格必须整体可同时驻留。
\discussionheading
\discussionpoint{网格屏障是一项带启动前提的全设备契约。}\code{grid.sync()} 让合作网格中所有线程在同一点会合，并为屏障前后的内存通信提供所需可见性，因此一个 kernel 内相邻 DP 波前可以安全衔接。普通启动没有全网格屏障，块完成顺序也未定义，不能因为一次实验恰按编号运行就依赖该时序。使用合作组还要查询设备支持并通过 cooperative launch 调用，错误启动应在主机端立即检查。这个显式契约类似 MPI collective：所有参与者、调用顺序和上下文必须一致，少一个参与者都不能靠普通同步替代。

\discussionpoint{合作网格必须整体可驻留才能保证所有参与者到达。}允许的块数由每 SM 活跃块数乘 SM 数决定，而活跃块数又受线程、寄存器、静态与动态共享内存共同限制；只按最大线程数估算会过量。若未驻留块等已驻留块越过 \code{grid.sync()} 才能获得资源，系统会形成死锁，所以 cooperative launch 会限制网格容量。问题逻辑 tile 更多时，可让有限持久块循环领取任务，或退回每个波前一个 kernel/图节点。这个约束说明“少 launch”并非免费，必须用常驻资源换取 kernel 内全局同步。

\discussionpoint{全网格同步会让无工作块也支付等待成本。}在 tile 波前的开头和结尾，只有一两个对角线块有有效单元，其余合作块仍必须执行控制骨架并到达每个 \code{grid.sync()}；狭长矩阵的空闲比例尤其高。若合作网格有 64 块而首轮只有 1 块有效，其余 63 块仍需驻留并参与屏障，活动比例仅 $1/64$。持久任务队列可以让块领取当前可执行 tile，单向依赖则只唤醒真实后继，多 kernel 方法又能只启动当轮块数，但三者分别增加调度状态、内存序或 launch。应比较同步粒度带来的闲置与替代协议的管理成本。相同权衡出现在图层次遍历、全局迭代 stencil 和多阶段求解器中。

\discussionpoint{部署环境会改变昨天可行的合作容量。}启动前应调用占用率 API 以实际 kernel 的寄存器、共享量和块形状计算容量，再检查合作启动返回值；MIG 分区、多进程共享或编译器资源变化都可能减少可用 SM 与驻留块。同一源码若编译器多分配几个寄存器，单 SM 驻留块数可能从 2 阶跃降到 1，使原本合法的合作网格超出容量。基准需要把省下的 launch 时间、网格屏障停顿和有效 DP 单元比例分开，否则总加速无法归因。正确性还应在矩阵不足一个 tile、最后行列不完整和极端长宽比下，与多 kernel 版本及 CPU 表逐项比较。这样才能把硬件配置假设纳入持续验证，而不是写死一次设备查询结果。
\variantheading 若占用率查询给出每 SM 2 个常驻块、设备有 80 个 SM，则 $ntr\le160$ 才满足
该启动的驻留必要条件；$ntr=161$ 必须改回多内核波前或分批设计。
\practiceheading 在启动前查询设备属性和占用率，并检查合作启动 API 的返回值；用
Nsight Systems 比较多次小内核 launch 与合作
内核网格屏障的时间分布。
\sourcecheck{https://docs.nvidia.com/cuda/cuda-c-programming-guide/}{NVIDIA Cooperative Groups 网格同步与合作启动要求；高置信度}
对抗审查：普通 \code{<<<>>>} 启动不保证 \code{this_grid().sync()} 合法；网格超过并发
驻留上限会启动失败。把无任务块提前 return 会让其余块永久停在网格屏障。
\end{solutionexercise}

\begin{solutionexercise}{4}
\problemheading 实现采用线程块级分块、并使用单向同步的 Smith--Waterman 内核。可参考
第 11 章的单次回看同步。每个块开始一个数据片前只等待其真实依赖的数据片完成标志，而不是
执行全网格屏障；实现须说明完成标志的发布/获取语义，并避免持久块调度导致的死锁。

\approachheading 为每个 tile 行使用一个持久块，块内从左到右处理。第 $r$ 行块计算列
$c$ 前只等待上方 $(r-1,c)$ 的完成标志；自己的左 tile 已顺序完成，上左依赖也由上方块
的顺序隐含保证。动态分配逻辑行号避免后继行块先占满设备造成死锁。

\answerheading
\begin{lstlisting}[language=C++]
#include <cuda/atomic>
template<int T>
__global__ void sw_unidirectional(int* H,const char* rea,const char* ref,
    int rows,int cols,int MATCH,int MISMATCH,int GAP,
    unsigned* rowCounter,unsigned* done) {
  static_assert(T>0 && T<=1024,"launch exactly T threads per block");
  __shared__ int s[(T+1)*(T+1)]; __shared__ unsigned tr_s;
  if(threadIdx.x==0) tr_s=atomicAdd(rowCounter,1u);
  __syncthreads();
  unsigned tr=tr_s,ntr=(rows-1+T-1)/T,ntc=(cols-1+T-1)/T;
  if(tr>=ntr)return;
  for(unsigned tc=0;tc<ntc;++tc) {
    if(tr>0 && threadIdx.x==0) {
      cuda::atomic_ref<unsigned,cuda::thread_scope_device> above(
          done[(tr-1)*ntc+tc]);
      while(above.load(cuda::memory_order_acquire)==0){}
    }
    __syncthreads();
    swTile<T,T>(H,rea,ref,rows,cols,tr,tc,MATCH,MISMATCH,GAP,s);
    if(threadIdx.x==0) {
      cuda::atomic_ref<unsigned,cuda::thread_scope_device> mine(
          done[tr*ntc+tc]);
      mine.store(1,cuda::memory_order_release);
    }
    __syncthreads();
  }
}
// If ntr==0 or ntc==0, return on the host.  Otherwise check T threads and
// (T+1)^2*sizeof(int) shared bytes, zero state, then launch:
// sw_unidirectional<T><<<ntr,T>>>(...);
\end{lstlisting}
release 发布发生在全块写回后；上方 acquire 成功后，其得分对当前块可见。逻辑行 0 最先
被分配且不等待，必能推进并逐列唤醒后继行，因而不会形成“等待尚未调度前驱”的环。

\complexity{$O(RC)$ 有效工作与 $O(ntr\,ntc)$ 发布；轮询加载次数取决于调度且无固定渐近上界}{$O(ntr\,ntc)$ 标志，$(T+1)^2$ 共享整数/块}{$ntr$ 个流水块，可沿反对角线重叠}
\conclusionheading 单向标志只等待真正依赖的上方 tile，避免全网格屏障的无关等待。
\discussionheading
\discussionpoint{局部依赖事件可以替代过强的全局阶段。}Smith--Waterman tile 形成 DAG，每个 tile 只需等待提供北、西和西北边界的真实前驱；全网格 barrier 却强迫同一阶段所有无关块也一起等待。把完成标志放在 DAG 边上，前驱结束后只唤醒后继，多个相邻波前就能流水重叠。例如 tile $(r,c)$ 的就绪条件可明确写成 $(r-1,c)$、$(r,c-1)$ 与 $(r-1,c-1)$ 已发布，而不必等待同对角线远处的无关 tile。正确性证明应列出每个片外输入由哪个事件保护，并确认片内映射已经隐含满足的依赖无需重复等待。这个从 bulk-synchronous 到数据流调度的转变也适用于任务图、流水 stencil 和流式查询执行。

\discussionpoint{完成标志必须发布它所保护的边界载荷。}前驱先写完矩阵边界或共享到全局的摘要，再用 device-scope release 原子存储完成状态；后继以 acquire 观察到该状态后，才可读取普通载荷。若先发布标志再写数据，另一个 SM 可能读到就绪和旧边界；\code{volatile} 不建立 happens-before，只对标志使用原子但采用 relaxed 顺序也未必满足协议。可把边界先填唯一哨兵并随机延迟载荷写入，主动放大“标志新、数据旧”的非法观察，但压力测试仍不能替代证明。竞争检查可发现部分冲突，内存序审查则解释为何所有合法执行都正确。该模式与无锁队列生产者写槽后发布尾标完全相同。

\discussionpoint{等待协议还必须证明调度能够前进。}若静态后继块先占满驻留槽并自旋等待尚未调度的前驱，前驱拿不到资源便会死锁；即使前驱已经运行，长时间自旋也可能占据执行与内存带宽，造成近似活锁。设设备总共只能驻留 8 块，若调度器先放入 8 个都依赖未启动前驱的后继，它们共同占满资源，等待多久都不会让状态改变。可让有限持久块按依赖方向领取可执行任务、动态分配单调逻辑编号，或把网格限制在可驻留容量内；轮询还可加入退避以减少流量，但退避本身不能修复资源环。自旋读取次数由调度和前驱耗时决定，不能写成一次或固定渐近常数。并发算法的正确性因此包含安全性与前向进度两部分，缺一不可。

\discussionpoint{更易证明的多 kernel 路线仍是重要基线。}主机逐波前 launch、CUDA Graphs 或库级任务图用 kernel 边界提供清晰全局同步，launch 更多但不需要自定义自旋和原子内存序；图执行还能摊薄重复启动开销。若同一矩阵形状反复执行，预先捕获上百个波前节点可把图构建成本摊到多次运行，而单次小问题未必回本。比较时应记录标志流量、自旋指令、前驱等待、有效单元率和总时间，而不能只数 kernel 个数。随机延迟注入、低 occupancy 和极端长宽比会主动放大隐藏死锁与等待不均。若局部协议只在理想调度下更快，它尚不足以替代可预测的阶段实现。
\variantheading 若 $ntr=3,ntc=4$，逻辑行 0 无需等待并依次发布 4 个标志；行 1 在每列等
行 0，同理唤醒行 2，共有 12 次发布，但轮询 load 次数不是 12 的常数倍保证。
\practiceheading 用 Compute Sanitizer racecheck 验证 release/acquire 载荷，并在人为降低每
SM 驻留块数时做前向进度压力测试；启动必须为 \code{<<<ntr,T>>>}，同时查询 $T$ 线程和
$(T+1)^2$ 共享整数是否可用。
\sourcecheck{https://docs.nvidia.com/cuda/cuda-c-programming-guide/}{NVIDIA 对设备作用域内存排序与跨块可见性的官方说明；高置信度}
对抗审查：普通 \code{volatile} 标志不能替代 release/acquire。若用静态行号且调度器先让
后继行常驻自旋，可能死锁；动态单调行号是前向进度证明的一部分。
\end{solutionexercise}

\begin{solutionexercise}{5}
\problemheading 实现采用超数据片和单向同步的 Smith--Waterman 内核。可参考第 11 章
的单次回看同步。超数据片采用原书的剪切变换
$q=q_{\mathrm{tile}}-r_{\mathrm{tile}}$（剪切因子 $-1$），使片内每轮都有固定数量的
活动线程；实现须据此推导相邻超数据片依赖、等待/发布标志和矩阵边界。

\approachheading 超数据片 $(r,c)$ 的全局列为 $cT+q_{tile}-r_{tile}+1$。它依赖本行
前一片以及上一行的 $c$、$c+1$ 两片，所以持久行块在处理 $c$ 前等待上一行的 $c+1$
标志（最右边界退化为 $c$）；这也隐含上一行的 $c$ 已完成。

\answerheading
\begin{lstlisting}[language=C++]
template<int T>
__global__ void sw_hyper_oneway(int* H,const char* rea,const char* ref,
    int rows,int cols,int MATCH,int MISMATCH,int GAP,
    unsigned* rowCounter,unsigned* done) {
  static_assert(T>0 && T<=1024,"launch exactly T threads per block");
  __shared__ int tile[T*T]; __shared__ unsigned tr_s;
  if(threadIdx.x==0)tr_s=atomicAdd(rowCounter,1u);
  __syncthreads();
  unsigned tr=tr_s,ntr=(rows-1+T-1)/T,ntc=(cols-1+T-1)/T;
  if(tr>=ntr)return;
  unsigned pitch=ntc+1; // includes the right-edge partial hypertile
  for(unsigned tc=0;tc<=ntc;++tc) {
    if(tr>0 && threadIdx.x==0) {
      unsigned need=(tc<ntc)?tc+1:tc;
      cuda::atomic_ref<unsigned,cuda::thread_scope_device> above(
          done[(tr-1)*pitch+need]);
      while(above.load(cuda::memory_order_acquire)==0){}
    }
    __syncthreads();
    for(int p=threadIdx.x;p<T*T;p+=T)tile[p]=0;
    __syncthreads();
    int rt=threadIdx.x,gr=1+int(tr)*T+rt;
    for(int qt=0;qt<T;++qt) {
      int gc=1+int(tc)*T+qt-rt;
      if(gr<rows && gc>=1 && gc<cols) {
        int n=(rt==0||qt==0)?H[size_t(gr-1)*cols+gc]
                             :tile[(rt-1)*T+qt-1];
        int w=(qt==0)?H[size_t(gr)*cols+gc-1]:tile[rt*T+qt-1];
        int nw=(rt==0||qt<=1)?H[size_t(gr-1)*cols+gc-1]
                              :tile[(rt-1)*T+qt-2];
        int sub=(rea[gr-1]==ref[gc-1])?MATCH:MISMATCH;
        tile[rt*T+qt]=max(0,max(nw+sub,max(w+GAP,n+GAP)));
      }
      __syncthreads();
    }
    for(int rr=0;rr<T;++rr) {
      int r=1+int(tr)*T+rr,gc=1+int(tc)*T+threadIdx.x-rr;
      if(r<rows&&gc>=1&&gc<cols)H[size_t(r)*cols+gc]=tile[rr*T+threadIdx.x];
    }
    __syncthreads();
    if(threadIdx.x==0) {
      cuda::atomic_ref<unsigned,cuda::thread_scope_device> mine(
          done[tr*pitch+tc]);
      mine.store(1,cuda::memory_order_release);
    }
    __syncthreads();
  }
}
// If ntr==0 or ntc==0, return on the host.  Otherwise validate T threads and
// T*T*sizeof(int) shared bytes, zero ntr*(ntc+1) flags, then launch:
// sw_hyper_oneway<T><<<ntr,T>>>(...);
\end{lstlisting}
片内每轮所有 $T$ 个线程对应同一条剪切后的反对角线，故只需 $T$ 轮。上述北、西、西北
下标分别按仿射映射逆推；标志依赖恰覆盖所有片外来源。每个合法得分单元属于唯一
$(tr,tc,rt,qt)$，不会重复写。

\complexity{$O(RC)$ 有效工作及边界掩蔽；发布为 $O(ntr\,ntc)$，轮询加载调度相关且无固定渐近上界}{$T^2$ 共享整数/块及 $ntr(ntc+1)$ 标志}{$ntr$ 个持久行块，片内恒定 $T$ 路并行}
\conclusionheading 超数据片把片内轮数从 $2T-1$ 降为 $T$，单向两列领先协议形成跨行流水。
\discussionheading
\discussionpoint{剪切通过改变坐标轴而不是改变依赖来整形波前。}令新列坐标 $q_{tile}=q+r_{tile}$，原来在二维表上倾斜的反对角线会在新坐标中接近固定宽竖带，使每轮可安排 $T$ 个线程处理不同逻辑行。正确性不能靠“图看起来规整”证明，必须把北、西、西北三个原坐标前驱逐一代入逆变换，确认它们在新调度中都早于当前点，并推导越界条件。坐标变换只重排执行顺序，不改变每个 DP 单元公式。编译器 polyhedral 优化中的 loop skewing 正是同一数学动作，用于把跨迭代依赖变成可 tile 的方向。

\discussionpoint{两列领先源于剪切后相邻行的横向错位。}后继逻辑行处理超片 $c$ 时，其顶部边界除首项外可能来自上一行的 $c+1$，所以只等待“上一行同列完成”会读取陈旧数据；协议必须等待足够领先的前驱标志。把 $T=2$ 的四个单元逐个代入剪切坐标，即可看见同列标志尚不足以覆盖右上来源，这比凭图估计等待距离可靠。最右侧额外超片承接剪切移出矩形边界的有效单元，大部分 lane 虽被掩蔽，仍需参加块内屏障。换来的收益是片内轮次从普通方形波前的 $2T-1$ 缩到 $T$，并让各行形成稳定流水。这个例子说明同步距离应由变换后的依赖推导，不能凭 tile 名称猜测。

\discussionpoint{超片减少同步深度但不改变有效工作阶数。}矩阵中每个合法 DP 单元仍需计算一次，所以有效工作保持 $O(mn)$；剪切改善的是每轮活动线程比例和同步轮数，而不是把表格算法变成次二次复杂度。额外超片、边界谓词、完成标志和自旋都增加常数开销，短序列可能无法摊薄。若有效单元为 $mn$ 而变换额外执行 $E$ 个掩蔽槽，利用率应报告为 $mn/(mn+E)$，而不能把虚拟槽算入 GCUPS 分子。增大 $T$ 可扩大流水宽度，却让 $T^2$ 共享 tile 和寄存器状态迅速增长，并减少同时推进的行块。应把有效单元率、屏障数、标志流量与驻留量一起建模，避免只用理论轮数宣称加速。

\discussionpoint{小尺寸依赖图是坐标优化最有力的回归工具。}可先用 $T=2$ 或 3 画出每个 $(r,q)$ 在超片中的位置、三个前驱和发布标志，再与 CPU 完整表逐点比较；这会迅速暴露一列偏移、右边界遗漏和等待距离错误。测试还应验证坐标变换在合法域内一一对应：每个原单元恰好映射一次，逆变换回去仍是原坐标，不能只看最终最高分碰巧相同。相同仿射剪切可用于 stencil、动态时间规整、编辑距离和其他波前 DP，但每种依赖向量不同，变换必须重新验证。性能分析还要观察活动 lane、同步等待和附加超片无效工作。只有依赖证明与实测利用率同时成立，坐标整形才真正开阔到可复用的优化方法。
\variantheading 若 $T=32,ntc=10$，每行分配 pitch 11 并处理 $tc=0\ldots10$；处理
$tc=9$ 前等待上一行标志 10，最右附加片 $tc=10$ 则等待上一行同号 10。
\practiceheading 用 CPU 串行参考对照短序列穷举，并以竞争检查工具检查标志发布；启动固定为
\code{<<<ntr,T>>>}，用 Nsight Systems 观察自旋热点，
同时验证 $T^2$ 共享空间和标志乘法不溢出。
\sourcecheck{https://www.cri.minesparis.psl.eu/classement/doc/A-179.pdf}{Irigoin 与 Triolet 的 Supernode Partitioning 机构公开版；并以 DOI 10.1145/73560.73588 核对 POPL'88 正式出版信息；高置信度}
\sourcecheck{https://docs.nvidia.com/cuda/cuda-c-programming-guide/}{NVIDIA CUDA C++ Programming Guide 对设备作用域原子引用及 release/acquire 内存序的官方说明；高置信度}
对抗审查：只等待上一行同列 $c$ 不够，因为超片顶部除首项外来自上一行 $c+1$；这是一个
可构造的陈旧读取反例。最右附加片多数 lane 越界，但仍必须参与全部屏障。
\end{solutionexercise}

\begin{solutionexercise}{6}
\problemheading 用功能相同的 DPX 指令替换本章 Smith--Waterman 内核所用的
\code{max4()} 设备函数。原函数计算
$\max(0,H_{r-1,q-1}+s,H_{r,q-1}+g,H_{r-1,q}+g)$；实现须在支持 DPX 的 NVIDIA
Hopper 或更新架构上使用等价指令，并为较早架构保留语义相同的回退路径。

\approachheading 原表达式是 0 与三个 32 位有符号候选的最大值，恰对应 DPX intrinsic：
\[
\code{__vimax3_s32_relu(a,b,c)}.
\]
为非 Hopper 目标保留普通指令回退。

\answerheading
\begin{lstlisting}[language=C++]
__device__ __forceinline__ int sw_max(int diagonal,int west,int north) {
#if defined(__CUDA_ARCH__) && __CUDA_ARCH__ >= 900
  return __vimax3_s32_relu(diagonal,west,north);
#else
  return max(0,max(diagonal,max(west,north)));
#endif
}
// Replacement at each cell:
int value=sw_max(nw+substitution,w+DELETION,n+INSERTION);
\end{lstlisting}
DPX 的 ReLU 后缀表示再与 0 取最大值，故与
\code{max4(0, diagonal, west, north)} 逐个 32 位整数等价。它只替换单元的最大值计算，
不改变波前依赖、边界或同步。

\complexity{每单元 $O(1)$；硬件 DPX 可用更少指令}{$O(1)$}{各波前单元照常并行}
\conclusionheading Hopper（计算能力 9.0+）走 DPX，旧架构走语义相同的可移植回退。
\discussionheading
\discussionpoint{不同 DPX intrinsic 的融合边界不同。}本题调用 \code{__vimax3_s32_relu(a,b,c)}，其语义只是 $\max(a,b,c,0)$；对角、北、西三个候选中的 substitution 或 gap 加法已经在调用前分别完成。\code{__viaddmax_s32_relu(x,y,z)} 才计算 $\max(x+y,z,0)$，它融合一次加法、一次二选最大和零截断，也不能单条覆盖三个各带加法的候选。名称都含 max 与 ReLU 不代表可互换，替换时要逐个写出数学式、位宽和溢出行为。这个纪律同样适用于 Tensor Core 与 SIMD 复合指令。

\discussionpoint{源码分支存在不等于目标二进制含有硬件路径。}设备代码要在 \code{__CUDA_ARCH__>=900} 时选择 Hopper DPX，并为旧目标保留普通 max 回退；构建还必须真正生成相应 SM 90 映像，否则运行时只能执行兼容代码。fat binary 可同时包含多个目标，驱动按设备选择，但编译宏是在各目标编译阶段决定而非主机运行时判断。应使用 \code{cuobjdump} 或 SASS 反汇编确认目标指令，并在旧设备验证回退。这个从源、编译目标到运行选择的三层审查也适用于架构专用异步拷贝和矩阵指令。

\discussionpoint{Amdahl 定律限制单元算术优化的端到端收益。}DPX 能减少每个 DP 单元中的整数指令，但波前屏障、边界加载、完成标志、自旋和 traceback 流量不会按同一比例下降。若算术只占总时间的 30\%，即使理想地无限加速它，整体也最多提升到约 $1/(1-0.3)=1.43$ 倍。具体 intrinsic 的融合边界也必须准确：\code{__vimax3_s32_relu} 融合三路最大值与 ReLU，候选值的加法仍要先算；\code{__viaddmax_s32_relu} 只融合一次加法、两路最大值与 ReLU，不能声称整条多候选递推成为一条指令。使用 16 位 packed 指令可能增加吞吐，却必须证明最大可能得分不溢出，并确认指令采用普通、饱和还是其他算术语义。因而优化报告应先分解时间和数值范围，再解释 DPX 加速，而不是只展示内层指令数。

\discussionpoint{行业验证必须覆盖性能分布、数值路径和构建矩阵。}基因组序列的长度、相似度、带宽限制与批大小都会改变有效 DP 单元和同步占比，所以应分组报告 GCUPS、能耗以及实际单元数。GPU 结果至少要与 CPU 参考核对最高分，支持 traceback 时还要处理多个同分最优路径的约定；DPX 与回退路径应使用同一黄金样例。还应加入恰好接近 16 位上限的合成长序列，验证 packed 路径的溢出策略与 32 位路径一致，而不只测短读段。持续集成可在无 Hopper 主机上仍编译 SM 90 对象、反汇编检查指令并运行 CPU/回退测试，在 Hopper 环境再执行硬件路径。这样专用指令不会成为无法验证的条件分支。
\variantheading 若三个候选为 $-7,-2,-9$，\code{__vimax3_s32_relu} 返回 0；若为
$4,-2,3$ 则返回 4，分别验证 ReLU 与三路最大值语义。
\practiceheading 构建时同时生成 \code{sm_90} 与受支持的回退目标，用 \code{cuobjdump} 或
Nsight Compute 核对 DPX 指令并与 CPU 参考逐单元比较；另测接近 \code{INT_MAX} 的分数，
因为 intrinsic 不会修复候选加法溢出。
\sourcecheck{https://docs.nvidia.com/cuda/cuda-c-programming-guide/}{NVIDIA CUDA C++ Programming Guide 的 DPX intrinsic 官方说明；高置信度}
边界审查：DPX 不是新的同步机制，也不会解决加法溢出；候选分数若可能越过 32 位范围，
应先采用饱和策略或更宽类型。编译还要求提供该 intrinsic 的 CUDA 12 代工具链。
\end{solutionexercise}

\chapter{稀疏矩阵计算}

\begin{solutionexercise}{1}
\problemheading 考虑下面这个稀疏矩阵：
\[
\begin{bmatrix}
1&0&7&0\\0&0&8&0\\0&4&3&0\\2&0&0&1
\end{bmatrix}
\]
分别写出它在以下格式下的表示：\textup{(a)} COO；\textup{(b)} CSR；
\textup{(c)} ELL；\textup{(d)} JDS；\textup{(e)} CSC。表示中须明确值数组、
索引/指针数组、填充规则和 JDS 的行置换；统一采用零基下标。第 \textup{(e)} 项来自
中文译稿的扩展，底本范围差异见本题后的题解编者注。

\approachheading 统一采用零基下标；行内元素按列号升序排列。ELL 按“同一迭代的各行
元素连续”方式存放，JDS 则先按行长降序稳定排序，再按锯齿对角线存放。

\answerheading 矩阵共有 7 个非零元素。各表示如下；星号是不会参与计算的填充槽，
其列下标可任取合法值。
\begin{align*}
\text{COO:}\quad
 \code{rowIdx}&=[0,0,1,2,2,3,3],\\
 \code{colIdx}&=[0,2,2,1,2,0,3],\\
 \code{value}&=[1,7,8,4,3,2,1].
\end{align*}
\begin{align*}
\text{CSR:}\quad
 \code{rowPtrs}&=[0,2,3,5,7],\\
 \code{colIdx}&=[0,2,2,1,2,0,3],\\
 \code{value}&=[1,7,8,4,3,2,1].
\end{align*}
最长行含 2 个非零元素，所以 ELL 的两个“列”按列主序铺开为
\begin{align*}
 \code{nnzPerRow}&=[2,1,2,2],\\
 \code{colIdx}&=[0,2,1,0,\;2,*,2,3],\\
 \code{value}&=[1,8,4,2,\;7,*,3,1].
\end{align*}
JDS 的一种稳定排序为原行 $[0,2,3,1]$，相应行长为 $[2,2,2,1]$：
\begin{align*}
 \code{row}&=[0,2,3,1],&
 \code{iterPtr}&=[0,4,7],\\
 \code{colIdx}&=[0,1,0,2,\;2,2,3],&
 \code{value}&=[1,4,2,8,\;7,3,1].
\end{align*}
第一段是各排序行的第 0 个元素，第二段是前三行的第 1 个元素。长度相同的三行允许
任意置换，但必须同步置换 \code{row} 及各段中的元素。

CSC 按列分组：
\begin{align*}
 \code{colPtrs}&=[0,2,3,6,7],\\
 \code{rowIdx}&=[0,3,2,0,1,2,3],\\
 \code{value}&=[1,2,4,7,8,3,1].
\end{align*}

\editorialnote 英文底本第 17 章练习 1 只列 COO、CSR、ELL、JDS 四项；中文学习译稿
另加入了 CSC。本题解保留并回答中文译稿的扩展子问，但不把它误称为英文底本原题。

\conclusionheading 逐一由这些数组还原坐标和值，均得到题给矩阵；JDS 的 \code{row}
数组负责把排序后的输出恢复到原行号。
\discussionheading
\discussionpoint{格式名称背后是遍历方向。}
把同一个矩阵换成 COO、CSR 或 CSC，看起来只是重新排列三组数组，实际却是在提前回答“计算将沿哪个方向推进”。例如 SpMV 逐行形成 $y_i$ 时，CSR 的两个行指针直接给出一段连续区间；若要反复访问某列，CSC 才能避免扫描全部非零元。COO 没有压缩某一维，因而适合流式生成、拼接和按键排序，却要为每项同时携带行列坐标。ELL 与 JDS 又把规则性编码进布局：前者接受填充，后者接受行置换，以便硬件按统一节拍取数。因此选择格式不是无代价的“换皮”，而是在存储空间、构造成本和预期算子之间签订一份数据访问契约。

\discussionpoint{规则布局怎样转化为并行效率。}
设一个 warp 的 32 个线程各负责一行，当它们同时读取各自行中的第 $t$ 个非零元时，按迭代列主序存放的 ELL 可以给出连续地址，这正是合并访问希望看到的形状。然而只要一条行很长而其余行很短，统一宽度就会制造大量填充，短行线程也会提前失活。JDS 先按行长排序，使同一锯齿段中的活动线程数逐段递减，通常比原始次序更整齐，但输出必须经置换恢复原行。CPU 的缓存线、GPU 的 warp 和传统向量机的向量长度并不相同，所以“规则”必须相对于执行平台定义；仅凭格式名称，无法推出某个矩阵一定更快。

\discussionpoint{转换成本要放进矩阵的整个寿命。}
如果矩阵只参与一次乘法，先统计行长、排序行、分配填充并搬运数据，很可能比直接使用 COO 或 CSR 更贵。相反，Krylov 迭代可能对同一矩阵执行数百次 SpMV，这时一次性的 $T_{\mathrm{convert}}$ 可以由每次节省的时间摊销；简单的盈亏点是 $k>T_{\mathrm{convert}}/(T_{\mathrm{old}}-T_{\mathrm{new}})$。这个模型还应加入输入向量和输出置换、临时缓冲以及多 GPU 分发，因为它们同样属于格式生命周期。工程报告若只展示稳定阶段的内核带宽，就会掩盖一次性任务、小矩阵或结构频繁变化时真正占主导的成本。

\discussionpoint{领域结构促成混合与分桶格式。}
规则有限差分网格的行长往往集中在一个狭窄区间，统一 ELL 宽度通常不会浪费太多；推荐系统和幂律图则可能同时出现大量短行与少数超长行，用全局最大行长填充会迅速失控。实际库因而常按行长分桶：短行交给 ELL 或 SELL，中等行使用 warp 协作 CSR，极长行再由单独的块级归约处理。这样做增加了描述符、内核启动和构造复杂度，却能让每一类行匹配适当的并行粒度。评价方案时应同时给出行长分位数、非零元字节数、填充率、转换时间与多次迭代后的总时间，才能把领域结构与性能结论真正联系起来。
\variantheading 若在零基坐标 $(1,1)$ 再加入值 5，则 $z=8$，CSR 的 \code{rowPtrs} 变为 $[0,2,4,6,8]$；四行均长 2，ELL 不再有填充，JDS 的两个锯齿段也都含 4 项。
\practiceheading 在 32-lane warp、FP32 值和 32 位索引的 NVIDIA GPU 上，用 cuSPARSE \code{cusparseSpMV} 的 CSR 结果校验自写 ELL/JDS；取 $x=[1,1,1,1]^T$ 时本题原矩阵应得 $y=[8,8,7,3]^T$，再用 Nsight Compute 比较 DRAM 字节数。
\sourcecheck{https://ntrs.nasa.gov/api/citations/19890017045/downloads/19890017045.pdf}{Saad 1989 年原始论文直接定义了按行长降序重排、依次存放锯齿对角线以及 \code{IDIAG}/\code{JDIAG} 指针与列号；对抗检查：零基/一基下标和同长行的排序都不是格式本身的唯一选择，本解已把约定写明}
\end{solutionexercise}

\begin{solutionexercise}{2}
\problemheading 设有一个整数稀疏矩阵，含 $m$ 行、$n$ 列、$z$ 个非零元素。分别写出用
\textup{(a)} COO、\textup{(b)} CSR、\textup{(c)} ELL、\textup{(d)} JDS 表示该矩阵
各需要多少个整数。计数应同时包括非零值和格式所需的整数索引/指针；若已知信息不足以
作答，请指出还缺少哪些信息。

\approachheading 同时计入非零值和整数索引。令
$r=\max_i\ell_i$ 为最大行长，其中 $\ell_i$ 是第 $i$ 行非零元素数；采用本章 ELL 的
\code{nnzPerRow} 数组以及 JDS 的行置换和带末端哨兵的 \code{iterPtr}。

\answerheading
\begin{center}
\begin{tabular}{lll}
\toprule
格式 & 整数数目 & 组成\\
\midrule
COO & $3z$ & 值、行号、列号各 $z$\\
CSR & $2z+m+1$ & 值、列号各 $z$，行指针 $m+1$\\
ELL & $2mr+m$ & 值、列号各 $mr$，每行长度 $m$\\
JDS & $2z+m+r+1$ & 值、列号各 $z$，行置换 $m$，段指针 $r+1$\\
\bottomrule
\end{tabular}
\end{center}
所以 COO、CSR 只需题给 $m,n,z$。ELL 和 JDS 还需要最大行长 $r$；若要真正
构造数组而不只是求总长度，还需要全部行长分布和各非零元素坐标。某些 ELL 实现把填充值
固定为 0，省略 \code{nnzPerRow}，此时为 $2mr$；某些 JDS 实现显式保存每个排序行的
长度，则还要增加 $m$。这些是接口约定差异，不能静默混算。

\editorialnote 英文底本和中文译稿的第 2 题都只问 COO、CSR、ELL、JDS；CSC 只出现在
第 1 题中。旧稿误把 CSC 加进本题题面，现已按练习原文删去。

\conclusionheading 信息不足只影响依赖最大行长或完整行分布的格式；不能仅由平均行长
$z/m$ 反推出 ELL/JDS 的大小。
\discussionheading
\discussionpoint{整数个数还不是内存流量。}
题目要求数“多少个整数”，这能比较数组长度，却不能直接回答矩阵能否装入显存或一次 SpMV 要搬多少字节。假设值为 FP64、索引为 32 位，CSR 每个非零元至少携带 8 字节值和 4 字节列号；若索引改为 64 位，元数据流量便从每项 4 字节增到 8 字节，而乘加次数没有变化。SpMV 通常受内存带宽约束，因而少读一个索引可能比省下一条整数指令更重要。进一步的字节模型还要估计输入向量 $x$ 的缓存复用和输出写回，区分“数组静态大小”与“运行时实际传输”，否则不同格式之间的容量比较会被误当成性能预测。

\discussionpoint{元数据约定必须先于公式统一。}
两个实现都自称 ELL，并不保证它们的数组长度相同：一个接口可能保存每行有效长度，另一个把填充列号设为哨兵而省略长度。JDS 的 \code{iterPtr} 是否含末端哨兵、空行是否仍进入置换、同长行是否稳定排序，也都会改变 $m$ 或 $r$ 级别的常数项。若论文 A 把这些辅助数组计入、论文 B 却只统计值和列号，“B 更省内存”的结论只是口径差异。格式描述因此应像 API 契约一样写清零基下标、指针长度、填充值和空输入语义；这些细节不仅影响计数，也决定内核能否在边界处无歧义地停止。

\discussionpoint{索引宽度连接规模上限与带宽。}
32 位无符号位置最多表达 $2^{32}-1$，但许多库接口使用有符号整数，实际安全上限常是 $2^{31}-1$；一旦 $z$ 或 \code{rowPtrs[m]} 超界，截断可能把合法区间变成负数或较小地址。全部改成 64 位可以容纳更大的矩阵，却让每个列号和指针多占 4 字节，在内存受限算子中并非免费。一个折中是分块保存：块间基址使用 64 位，块内列号或偏移使用 16/32 位，只要显式证明局部范围不会溢出。这种分层索引也见于图分区、文件格式和数据库页，核心思想都是用较宽地址定位区域，再用较窄偏移描述区域内部。

\discussionpoint{峰值驻留量才决定能否运行。}
矩阵本体占 12~GiB 并不意味着它能放进 16~GiB GPU，因为格式转换可能同时保留源 COO、目标 CSR、直方图、扫描临时区和排序工作区。迭代求解器还要保存多个向量，检查点或多 GPU 通信又会增加暂存缓冲；分配器对齐与碎片使可用连续空间进一步小于标称容量。可靠预算应画出各阶段对象的生命周期，取同一时刻仍存活对象之和的峰值，而不是把所有数组总和或单个稳态快照当答案。生产系统还会预留错误恢复和库 workspace 的余量，这把一道格式计数题延伸成了资源规划与内存生命周期分析问题。
\variantheading 若 $m=4,z=7,r=2$ 且沿用本题元数据约定，则 COO、CSR、ELL、JDS 分别需要 $21,19,20,21$ 个整数；这组数可直接由四个公式交叉复算。
\practiceheading 假定值和索引均为 4 字节、GPU 显存为 16 GiB，用 cuSPARSE 的 32 位 CSR/COO 描述符建立同一矩阵，并以 Nsight Compute 的 \code{dram__bytes} 检查元数据流量；若 $z\ge2^{31}$，改用 64 位索引后重新核算容量。
\sourcecheck{https://ntrs.nasa.gov/api/citations/19890017045/downloads/19890017045.pdf}{Saad 1989 年原始论文直接并列讨论 ITPACK/ELL 与 JDS 的值、列号、行重排和段指针结构；对抗检查：不同接口是否另存 ELL/JDS 行长会改变常数项，本解因此显式声明计数模型，且没有用不含 JDS 的库文档代替证据}
\end{solutionexercise}

\begin{solutionexercise}{3}
\problemheading 用直方图和前缀和这两个基础并行原语，实现从 COO 到 CSR 的格式转换。
输入 COO 三元组可以是任意顺序；输出须给出长度 $m+1$ 的合法 CSR 行指针，并保证同一行
中的每个非零元素被写入唯一槽位。还须说明是否保持 COO 原有的行内次序。

\approachheading 先按行统计元素数，再对计数作排他前缀和得到每行区间；散布阶段必须
使用独立游标，不能原地破坏 \code{rowPtrs}。同一行的多个线程以原子加领取唯一槽位。

\answerheading 下面代码省略常规 CUDA 错误检查，但给出了转换的全部设备步骤：
\begin{lstlisting}[language=C++]
__global__ void count_rows(const int* cooRow, const int* cooCol,
                           int z, int m, int n,
                           unsigned* counts, unsigned* bad) {
    int e = blockIdx.x * blockDim.x + threadIdx.x;
    if (e >= z) return;
    int r = cooRow[e], c = cooCol[e];
    if (r < 0 || r >= m || c < 0 || c >= n) {
        atomicExch(bad, 1u); return;
    }
    atomicAdd(&counts[r], 1u);
}

__global__ void scatter_csr(const int* cooRow, const int* cooCol,
                            const int* cooVal, int z, int m, int n,
                            unsigned* next, int* csrCol, int* csrVal,
                            unsigned* bad) {
    int e = blockIdx.x * blockDim.x + threadIdx.x;
    if (e >= z) return;
    int r = cooRow[e], c = cooCol[e];
    if (r < 0 || r >= m || c < 0 || c >= n) {
        atomicExch(bad, 1u); return;
    }
    unsigned p = atomicAdd(&next[r], 1u);
    csrCol[p] = c;
    csrVal[p] = cooVal[e];
}

// Require m,n,z >= 0; counts[0..m-1] = 0 and bad = 0.
if (z > 0) {                         // CUDA cannot launch a zero-block grid
    count_rows<<<(z + 255) / 256, 256>>>(row, col, z, m, n,
                                         counts, bad);
    CHECK(cudaGetLastError());
}
unsigned hBad = 0;
CHECK(cudaMemcpy(&hBad, bad, sizeof(hBad), cudaMemcpyDeviceToHost));
if (hBad) throw std::out_of_range("COO coordinate");

// For m>0, run ExclusiveScan(counts,rowPtrs,m), for example with CUB.
// In every case write the sentinel rowPtrs[m] = z; thus m=z=0 gives [0].
unsigned uz = unsigned(z);
CHECK(cudaMemcpy(rowPtrs + m, &uz, sizeof(uz), cudaMemcpyHostToDevice));
if (m > 0)
    CHECK(cudaMemcpy(next, rowPtrs, size_t(m) * sizeof(unsigned),
                     cudaMemcpyDeviceToDevice));
if (z > 0) {
    scatter_csr<<<(z + 255) / 256, 256>>>(row, col, val, z, m, n,
                                         next, csrCol, csrVal, bad);
    CHECK(cudaGetLastError());
}
\end{lstlisting}
散布内核显式接收 \code{m,n}；各内核在同一 stream 中依次启动，内核边界就是全网格
依赖点。第二次行列检查使散布本身也不会因调用者错误地跳过主机检查而越界。

排他扫描给出 \code{rowPtrs[r]}，而最后的哨兵必须显式写成 $z$。代码先验证全部行、列
坐标并在发现 \code{bad} 后停止，散布时又作防御性复查；所以非法输入只能报错，不会用
非法行号访问 \code{next}。$z=0$ 时跳过两个元素内核，仍产生全零行指针（特别地，
$m=0$ 时 CSR 行指针就是单元素数组 $[0]$）。

散布后同一行内部的次序不确定。CSR 本身不要求列有序；若下游接口要求列升序，还需对每个
\code{rowPtrs[r]:rowPtrs[r+1]} 段排序。重复坐标会作为多个条目保留；若要合并重复项，
再做分段归约。

\complexity{$O(m+z)$ 工作；并行扫描深度 $O(\log m)$，热点行的原子散布可能串行化}{$O(m+z)$ 输出及 $O(m)$ 计数/游标临时空间}{$z$ 个统计与散布任务，扫描由库原语并行执行}
\conclusionheading 直方图决定每行长度，排他扫描决定区间，原子游标保证每个 COO 条目
恰好占据该行的一个不同 CSR 槽位。
\discussionheading
\discussionpoint{计数、扫描、散布是一条可迁移的管线。}
怎样让并行线程在不知道前面有多少输出时仍得到互不重叠的位置？先统计每个桶的大小，再对桶大小做排他扫描，便能把第 $r$ 个桶映射到区间 $[p_r,p_{r+1})$；最后的散布只需在这个区间内领取槽位。COO 转 CSR 中桶是矩阵行，而图邻接表构造、粒子空间分箱和基数排序只是换了“桶键”的含义。三个阶段各有清楚不变量：计数总和为 $z$，扫描区间首尾相接，散布对每个输入恰写一次。掌握这些不变量，比记住某个库调用更有价值，因为它们还能指导错误检测和并行算法的组合证明。

\discussionpoint{唯一槽位不等于稳定顺序。}
原子加游标能保证同一行的两个线程不会写到同一位置，但哪个线程先获得较小下标取决于调度，所以行内次序并不稳定。CSR 的数学语义通常不要求列有序，然而稀疏三角求解、重复项合并或逐位可复现输出可能依赖确定顺序。若必须保持 COO 输入顺序，可以先为每个元素计算稳定的行内秩；若必须按列排序，则可按 $(row,col)$ 做稳定基数排序，再由行边界生成指针。稳定性要付出额外搬运和临时空间，却也让重复坐标相邻、搜索更容易、测试结果更可预测，因此它应作为明确契约选择，而不是原子散布的默认副产品。

\discussionpoint{平均行长会掩盖真正的原子热点。}
设 $z/m=8$，这个平均值既可能来自每行恰有 8 项，也可能来自一条百万项长行和大量空行；两种输入对同一原子游标的竞争完全不同。幂律图中的超级节点会让许多线程争用一个计数器，使散布阶段局部串行化，即使全局线程数很大也无法消除这个瓶颈。可按线程块建立局部直方图后再合并，对超长行分段并分配独立区间，或直接排序所有键来换掉热点原子。选择策略时应查看最大值、分位数和实际冲突率，而不能只报告 $z/m$；这也是并行负载均衡中“均值不足以描述偏斜”的典型例子。

\discussionpoint{转换接口还要定义坏数据的命运。}
真实 COO 数据可能含越界坐标、重复坐标、空矩阵，甚至在累计计数时发生整数溢出；若这些情况没有契约，转换程序即使在正常样例上正确，也可能静默写坏内存。重复坐标可以原样保留，使 SpMV 分别累加，也可以排序后做分段归约，把它们合成一个值；两种做法在浮点舍入顺序和后续非零元计数上都不同。稳健管线应先验证坐标范围与计数容量，并记录拒绝、丢弃或合并的数量，使结果能够追溯。这个原则同样适用于数据清洗和图导入：格式转换不仅是内存重排，还是从外部数据到内部不变量的一道可信边界。
\variantheading 若 $m=3$ 且 COO 行号依次为 $[2,0,2,1,2]$，直方图为 $[1,1,3]$、排他结果含哨兵为 \code{rowPtrs} $=[0,1,2,5]$；原子只会改变最后一段三个条目的内部次序。
\practiceheading 在支持 32 位全局原子的 NVIDIA GPU（计算能力 7.0 或更新、32-lane warp）上，用 CUB \code{DeviceScan::ExclusiveSum} 和 \code{DeviceRadixSort} 实现两条路径，并以 Nsight Compute 的原子吞吐和 L2 sector 指标定位热点行。
\sourcecheck{https://github.com/NVIDIA/cccl/tree/main/cub}{NVIDIA CCCL/CUB 官方源代码与文档；对抗检查：扫描只生成区间，不能单独解决同一行散布竞态，故本解另设原子游标并保留行指针}
\end{solutionexercise}

\begin{solutionexercise}{4}
\problemheading 实现构造混合 ELL--COO 表示的主机端代码，并用它完成一次 SpMV：
每行前至多 $K$ 个非零元素放入 ELL，超出 $K$ 的余项放入 COO；ELL 部分在设备上通过
内核计算，COO 部分的贡献在主机上计算。须说明两部分的输出如何合并，并处理空行、短行和
超过 $K$ 的长行。

\approachheading 输入采用按行分组的三元组。每行前 $K$ 项进入列主序 ELL，其余项进入
COO。设备输出先清零并计算 ELL，回传后主机再累加 COO；这样不存在 CPU/GPU 同时写同一
输出的竞态。

\answerheading 核心实现如下（\code{Triplet} 含整数 \code{row,col} 和浮点
\code{value}）：
\begin{lstlisting}[language=C++]
struct Triplet { int row, col; float value; };

__global__ void ell_spmv(const float* value, const int* col,
                         const int* len, int m, int K,
                         const float* x, float* y) {
    int r = blockIdx.x * blockDim.x + threadIdx.x;
    if (r >= m) return;
    float s = 0.0f;
    for (int t = 0; t < len[r]; ++t) {
        int p = t * m + r;              // column-major ELL
        s += value[p] * x[col[p]];
    }
    y[r] = s;
}

void hybrid_spmv(const std::vector<Triplet>& a, int m, int n, int K,
                 const std::vector<float>& x, std::vector<float>& y) {
    if (m < 0 || n < 0 || K < 0 || x.size() != size_t(n))
        throw std::invalid_argument("bad shape");
    std::vector<std::vector<Triplet>> rows(m);
    for (const auto& e : a) {
        if (e.row < 0 || e.row >= m || e.col < 0 || e.col >= n)
            throw std::out_of_range("COO coordinate");
        rows[e.row].push_back(e);
    }

    y.assign(size_t(m), 0.0f);
    if (m == 0 || n == 0 || K == 0) {
        // K==0 is the pure-COO endpoint; the other two cases have no entries.
        for (const auto& e : a) y[e.row] += e.value * x[e.col];
        return;
    }

    std::vector<float> ev(size_t(m) * K, 0.0f);
    std::vector<int> ec(size_t(m) * K, 0), len(m, 0);
    std::vector<Triplet> overflow;
    for (int r = 0; r < m; ++r) {
        len[r] = std::min<int>(K, rows[r].size());
        for (int t = 0; t < len[r]; ++t) {
            ev[size_t(t) * m + r] = rows[r][t].value;
            ec[size_t(t) * m + r] = rows[r][t].col;
        }
        overflow.insert(overflow.end(), rows[r].begin() + len[r],
                         rows[r].end());
    }

    float *dv = nullptr, *dx = nullptr, *dy = nullptr;
    int *dc = nullptr, *dl = nullptr;
    // CHECK denotes the usual cudaError_t checker.
    CHECK(cudaMalloc(&dv, ev.size() * sizeof(float)));
    CHECK(cudaMalloc(&dc, ec.size() * sizeof(int)));
    CHECK(cudaMalloc(&dl, len.size() * sizeof(int)));
    CHECK(cudaMalloc(&dx, x.size() * sizeof(float)));
    CHECK(cudaMalloc(&dy, size_t(m) * sizeof(float)));
    CHECK(cudaMemcpy(dv, ev.data(), ev.size()*sizeof(float), cudaMemcpyHostToDevice));
    CHECK(cudaMemcpy(dc, ec.data(), ec.size()*sizeof(int), cudaMemcpyHostToDevice));
    CHECK(cudaMemcpy(dl, len.data(), len.size()*sizeof(int), cudaMemcpyHostToDevice));
    CHECK(cudaMemcpy(dx, x.data(), x.size()*sizeof(float), cudaMemcpyHostToDevice));
    ell_spmv<<<(m + 255) / 256, 256>>>(dv, dc, dl, m, K, dx, dy);
    CHECK(cudaGetLastError());
    CHECK(cudaMemcpy(y.data(), dy, size_t(m)*sizeof(float), cudaMemcpyDeviceToHost));

    for (const auto& e : overflow) y[e.row] += e.value * x[e.col];
    CHECK(cudaFree(dy)); CHECK(cudaFree(dx)); CHECK(cudaFree(dl));
    CHECK(cudaFree(dc)); CHECK(cudaFree(dv));
}
\end{lstlisting}
生产代码应使用 RAII 包装设备分配，使中途异常也能释放资源。这里在任何设备分配和启动
之前处理 $m=0$、$n=0$ 和 $K=0$：既不会把零字节分配当成正常设备缓冲区，也不会启动
零块网格；$K=0$ 时全部合法条目恰由主机 COO 路径处理。其余情况下 ELL 与 COO 恰好划分
全部条目，因此每个乘积贡献一次且仅一次。

\complexity{构造和 SpMV 均为 $O(z+mK)$ 上界，实际乘加总数为 $z$}{ELL 为 $O(mK)$，COO 余项为 $O(\sum_i\max(\ell_i-K,0))$，另有向量空间}{设备并行度为 $m$ 行；主机余项顺序执行，是题目指定的可扩展性瓶颈}
\conclusionheading 必须先完成设备计算和回传，再由主机加入溢出贡献；若改成设备 COO
内核，则需原子加或第二输出向量后再归并。
\discussionheading
\discussionpoint{混合格式先要证明划分完整。}
把每行前 $K$ 项放入 ELL、其余项放入 COO，直觉上是在让常见的短前缀走规则快路径，让不规则长尾走旁路。但正确性的核心不是“两个内核都算了一些数”，而是每个原始非零元恰好属于一部分：集合必须互斥，二者并集又必须等于原矩阵。空行、长度恰为 $K$ 的行和含重复坐标的行都是检验这一不变量的边界。只要划分成立，最终输出就是 $y=y_{\mathrm{ELL}}+y_{\mathrm{COO}}$；若某项被遗漏或重复，线性性也无法挽救。类似的常见路径与异常路径分解广泛用于数据库执行和网络处理，但都必须先建立完备、无重叠的分类规则。

\discussionpoint{不规则尾部可以有多种执行位置。}
题目让 CPU 计算 COO 溢出区，是为了把所有权和合并过程讲清楚，并不表示主机一定是最佳位置。GPU 可以让每个 COO 项原子累加到 $y$，也可以先按行排序后做分段归约，或者生成独立的 $y_{\mathrm{COO}}$ 再执行一次向量相加。原子方案省临时空间但会在长行上争用，分段方案更规则却增加排序和 workspace，双输出方案避免写冲突但多一次全向量流量。CPU 路径实现简单，却需要设备完成点、输出回传和串行尾部。选项之间没有脱离数据分布的优胜者，真正的比较对象应是包含同步与合并在内的端到端路径。

\discussionpoint{互连拓扑会改变 CPU 与 GPU 的分工。}
在离散 GPU 上，把长度为 $m$ 的输出向量经 PCIe 回主机，往往比处理少量 COO 项更昂贵；而 CPU 与 GPU 共享物理内存的系统可能省掉显式复制，却仍有缓存一致性、页迁移和同步代价。若溢出项来自另一张 GPU，GPU-aware MPI、NVLink 或主机中转又会形成不同数据路径。因此阈值 $K$ 不能只由“还有多少溢出项”决定，还要估计链路固定延迟、有效带宽、输出大小以及是否能与其他工作重叠。这个例子说明异构算法的切分点不仅由 FLOP 决定，数据放在哪里、何时必须可见，往往才是系统级性能的主导因素。

\discussionpoint{结构变化要求把重建也纳入控制回路。}
静态有限元矩阵可以一次选择 $K$ 后反复使用，但自适应网格或动态图的行长分布会随时间变化，旧阈值可能逐渐制造过多填充或溢出。每轮都重建 ELL/COO 又可能让转换成本吞掉计算收益，因此更实际的做法是周期采样行长分位数和溢出比例，只在预计节省超过重建代价时切换。判断条件可以比较未来预计迭代数乘以单次收益与重建、传输和缓存失效成本，并为观测噪声预留余量。这样的反馈式格式选择接近在线优化问题：既要适应新数据，也要设置回滞，避免在两个接近的阈值之间反复震荡。
\variantheading 若三行长度为 $[1,3,5]$ 且 $K=2$，ELL 有 6 个槽位、其中 5 个是真实项，COO 溢出恰为 $0+1+3=4$ 项；两部分合计仍覆盖全部 9 个非零元一次。
\practiceheading 假定离散 PCIe GPU、FP32 值和 32 位索引，用 cuSPARSE CSR SpMV 作数值与吞吐基线，再由 Nsight Systems 测量本实现的 D2H 复制和 CPU 溢出区间；测试时矩阵与 $x$ 常驻设备以免混入重复上传成本。
\sourcecheck{https://docs.nvidia.com/cuda/cuda-runtime-api/group__CUDART__MEMORY.html}{CUDA Runtime API 官方内存复制与分配语义；对抗检查：主机在异步设备写尚未完成时读取输出会产生生命周期错误，本解使用同步的设备到主机复制作为完成点}
\end{solutionexercise}

\begin{solutionexercise}{5}
\problemheading 实现一个使用 JDS 格式的并行 SpMV 内核。输入包括按行长降序排列后的
行置换、按锯齿对角线存放的非零值和列号，以及每个锯齿段的起始指针；输出须按原矩阵行号
写回，并正确处理各段活动行数递减的边界。

\approachheading 排序后每个线程负责一个行位置。记 $p_t=\code{iterPtr[t]}$，则第 $t$ 个
锯齿段的活动行数为 $p_{t+1}-p_t$。行位置超过该值时，由于活动数单调不增，可以立即结束。

\answerheading
\begin{lstlisting}[language=C++]
__global__ void spmv_jds(const float* value, const int* colIdx,
                         const int* row, const int* iterPtr,
                         int numRows, int numIters,
                         const float* x, float* y) {
    int r = blockIdx.x * blockDim.x + threadIdx.x; // sorted row position
    if (r >= numRows) return;
    float sum = 0.0f;
    for (int t = 0; t < numIters; ++t) {
        int active = iterPtr[t + 1] - iterPtr[t];
        if (r >= active) break;
        int p = iterPtr[t] + r;
        int c = colIdx[p];
        sum += value[p] * x[c];
    }
    y[row[r]] = sum;                    // restore original row order
}
// iterPtr has numIters + 1 entries; row is a permutation of 0..numRows-1.
\end{lstlisting}
构造时应验证 \code{iterPtr[0]==0}、末项等于 $z$、各段长度单调不增、列下标合法且
\code{row} 是排列。每个输出行只有一个写者，所以无需原子操作或块间同步。

\complexity{$O(z)$ 总乘加，关键路径为最大行长 $r$}{$O(z+m+r)$ JDS 元数据与数值空间，每线程 $O(1)$ 寄存器}{最多 $m$ 个线程；排序使同一 warp 的退出位置接近，降低分歧}
\conclusionheading \code{iterPtr} 同时给出段起点和活动行数，\code{row} 保证重排不会
改变 SpMV 输出的逻辑行号。
\discussionheading
\discussionpoint{JDS 把短行退出变成长向量收缩。}
设排序后的行长是 $[5,5,3,1]$，第一条锯齿对角线含 4 项，第二、三条仍各含 3 项，最后两条只剩 2 项；活动宽度因此单调缩短。传统向量机可以把每一段当作一条连续向量，GPU 则可让相邻线程读取相邻项，避免原始行次序下某些 lane 很早退出。这个规则性来自预先按行长排序，而不是矩阵天然规则，所以还要支付置换和段指针的成本。现代 CSR 自适应内核与 SELL 也能通过分桶或局部排序处理不均匀行，JDS 不是唯一答案；它更重要的启示是，改变数据次序可以把控制流不规则性转化为可管理的段长度。

\discussionpoint{行置换必须贯穿算子的语义。}
JDS 只重排矩阵行，不改变每个非零元的列号，因此输入仍按原始 $x_j$ 读取；线程位置 $r$ 算出的和却属于原行 \code{row[r]}。如果把结果直接写到 $y[r]$，程序不会越界，数值也可能看似平滑，但它实际计算了 $PAx$ 而不是 $Ax$，其中 $P$ 是行置换矩阵。连续执行多个算子时，可以让中间向量暂时保持置换顺序，以省掉来回搬运，但下一个算子的行列布局必须同步变换。这个问题与数据库列重排、图顶点重编号相同：置换能改善局部性，却必须作为显式元数据传播，不能在单个内核结束时凭记忆恢复。

\discussionpoint{规则访存并不保证逐位可重复。}
一个线程顺序累加一行时，加法顺序通常由存储顺序固定；若为了超长行并行度而让多个 lane 先算部分和，再用树形归约合并，舍入路径就发生了变化。对于 $[10^{20},-10^{20},3]$ 这样的 FP32 数据，不同配对可以得到 0 或 3，尽管两种执行都没有数据竞态。固定归约树能让结果跨运行重复，却不自动让它更接近实数精确和；FP64 累加、pairwise 或补偿求和解决的是另一维度的准确性。稀疏库因而应分别声明数据竞态、确定性和误差容限，而不能用“访问规则”笼统替代数值契约。

\discussionpoint{预处理只有被复用时才有经济意义。}
把行排序、重排全部非零元并建立 \code{iterPtr} 都消耗带宽，最终写回还要经过行置换；一次性 SpMV 很可能在内核加速之前就已亏损。若构造耗时 20~ms，而 JDS 相比 CSR 每次节省 0.2~ms，那么至少约 100 次复用才刚达到盈亏点，尚未计入额外内存。迭代求解器、固定拓扑图传播和多右端项计算较容易满足这一条件，动态图或临时查询则未必。真实决策应在代表性行长分布上同时测构造与稳态，并把预热、缓存和转换生命周期写进报告，这也是所有“先编译后执行”型优化的共同评估方式。
\variantheading 若 \code{iterPtr} $=[0,5,8,9]$，三个段的活动行数为 $5,3,1$；排序位置 $r=3$ 的线程只执行第 0 段一次乘加，并在第 1 段因 $3\not<3$ 退出。
\practiceheading 在 32-lane warp、FP32/32 位索引的 NVIDIA GPU 上，把自写 JDS 与 cuSPARSE CSR 对同一矩阵重复执行 100 次；用 Nsight Compute 的 Branch Efficiency、Global Load Efficiency 和 DRAM Throughput 判断排序成本能否被稳态 SpMV 摊薄。
\sourcecheck{https://ntrs.nasa.gov/api/citations/19890017045/downloads/19890017045.pdf}{Saad 1989 年原始论文直接给出 JDS 的行置换、逐渐缩短的连续锯齿段、列下标和段起点指针；对抗检查：若活动行数不单调或 row 含重复项，提前退出和无原子写均不再成立，故本解把它们列为构造不变量}
\end{solutionexercise}

\begin{solutionexercise}{6}
\problemheading 原书第 17.6 节指出，SpMV/JDS 各次迭代的起始地址难以对齐到架构规定
的边界。请说明：如果把 JDS 改造成该节提到的“分段 ELL”变体，也就是把排序后的行分组、
每组分别按 ELL 方式填充并对齐各段，这个对齐问题能否得到缓解？代价是什么？

\editorialnote 所核英文第 5 版底本的第 17 章练习止于第 5 题；中文译稿增加了第 6--8
题。本题解依“中文译稿为题目主要来源”的口径保留这三题，并在此统一标明版本差异。

\approachheading JDS 每段长度随活动行数变化，所以下一段自然紧跟上一段，起点可能落在
任意元素位置。分段 ELL 可以人为选择每段的基址和行距。

\answerheading 可以缓解。把行长相近的排序行切成段，对每段独立填充到该段最大行长；
再令段基址按交易边界对齐，并把段的行距（以字节计）补到该边界的整数倍。这样该段内
第 $t$ 次迭代的起点
\[
  \text{base}_{s}+t\,\text{pitch}_{s}
\]
都保持对齐，而纯 JDS 的 \code{iterPtr[t]} 没有这一保证。

代价包括：段内填充元素；为了对齐基址和行距而增加的额外空洞；段偏移、行距、段长等
元数据；更多预处理和行置换；以及按段启动内核或在内核中查段的控制开销。段太大趋近纯
ELL，填充严重；段太小则元数据、启动和尾部 warp 浪费增加。因此答案不是“完全消除且
没有代价”，而是在交易效率和空间/调度开销之间折中。

\complexity{SpMV 的有效工作仍为 $O(z)$，但会增加 $O(p)$ 填充工作，其中 $p$ 由分段选择决定}{比 JDS 多 $O(p+s)$，$s$ 为段数}{段内访问规则且合并；段间需要独立调度，短段可能降低占用}
\conclusionheading 分段 ELL 能以可控填充换取可控对齐，最优段宽需结合行长分布和目标
GPU 的内存交易粒度实测。
\discussionheading
\discussionpoint{pitch 防止每轮访问逐渐失去对齐。}
某段含 30 个 FP32 元素时，第一轮即使从 128 字节边界开始，下一轮起点也会前进 120 字节，于是地址相位相对交易边界移动了 120、112、104 字节。把段宽补到 32 项后，pitch 变为 128 字节，每一轮都重新落在同样的对齐位置；代价是每轮多读或跳过两个填充槽。这里要区分“首地址对齐”和“步长对齐”，前者只保证第一轮，后者才保证整个迭代序列。二维图像、张量行距和网络环形缓冲也有同样问题，因此 pitch 是一种跨行维持访问相位的布局参数，而不只是 CUDA 分配接口返回的数字。

\discussionpoint{局部排序把 JDS 延伸为 SELL。}
全局 JDS 可以获得最长的规则段，却可能把原本相邻的行搬到很远位置，损害输出局部性并增加大规模置换。SELL-C-$\sigma$ 把矩阵切成高度为 $C$ 的 slice，只在大小为 $\sigma$ 的窗口内按行长重排，再把每个 slice 填到局部最大行长。$C$ 接近 warp 宽度时便于线程协作，较大的 $\sigma$ 通常减少填充，却让置换范围扩大。于是窗口、slice 高度和 pitch 共同构成三方取舍：规则向量长度、填充空间与原始行局部性不能同时无限改善。理解这一关系有助于读懂现代稀疏格式为何常带多个调参，而不是只有一个缩写名称。

\discussionpoint{段宽增大并不会单调提速。}
更宽的段能减少段描述符和短尾 warp，似乎应该越大越好；反例是一个 slice 中只有一条超长行，其长度会成为其余短行的填充上限。若行长为 $[2,2,2,100]$，高度 4 的统一 ELL 需要 400 个槽位却只有 106 项有效，而把长行单独分桶可大幅降低浪费。反过来，段切得过细又会产生更多边界、较短连续请求和额外调度。最优段宽因此取决于行长方差、元素与索引字节数、缓存容量及硬件并行宽度，通常需要按矩阵族调参；这类非单调关系也是自动调优比固定经验常数更可靠的原因。

\discussionpoint{有用负载与硬件交易必须分开核算。}
用 $z$ 乘以值和索引宽度只能得到算法真正需要的数据量，却看不到填充槽、未对齐请求和缓存重复取数。一个内核可能报告很高的 DRAM 吞吐量，只因为搬运了大量无用填充；另一个内核吞吐数字较低，却以更少实际字节完成同样的 $z$ 次乘加。应先计算逻辑负载，再通过 sector/request、L2 命中和实际 DRAM 字节解释放大系数，并把格式转换时间按复用次数摊销。这样才能区分“总线很忙”和“算法有效”，也能把稀疏矩阵分析连接到压缩、数据库扫描等所有需要区分有效字节与物理流量的领域。
\variantheading 假定 FP32 数值数组按 128 字节交易边界对齐，某段有 30 行，则自然 pitch 为 120 字节；把段宽填充到 32 行后 pitch 为 128 字节，每次迭代多 2 个槽，却使所有迭代起点继续对齐。
\practiceheading 在 32-lane warp、4 字节元素且全局交易按 128 字节对齐分析的 NVIDIA GPU 上，实现 SELL-C-$\sigma$ 风格分段并与 cuSPARSE CSR 比较；用 Nsight Compute 的 sectors/request、L2 命中率和有效 DRAM 字节扫描段宽 16、32、64。
\sourcecheck{https://docs.nvidia.com/cuda/cuda-c-best-practices-guide/}{CUDA C++ Best Practices Guide 的合并访问与对齐规则；对抗检查：仅对齐段首而不让 pitch 同样对齐，后续迭代仍会失去对齐}
\end{solutionexercise}

\begin{solutionexercise}{7}
\problemheading 原书第 17.2 节提到，SpMV/COO 中原子操作的执行顺序不确定，可能带来
浮点数值稳定性问题。请说明：本章所示每行由一个线程顺序累加的 SpMV/CSR 是否也存在同样
的问题？为什么？回答须区分一次固定内核执行中的行内加法次序和不同实现/优化可能改变的
归约次序。

\approachheading 区分“浮点加法本来不满足结合律”和“多个线程对同一输出进行无序原子
更新”这两件事。

\answerheading 本章的基础 CSR 内核没有与 COO \emph{同样的}非确定性：第 $r$ 个线程
独占 $y_r$，并按 \code{rowPtrs[r]} 到 \code{rowPtrs[r+1]-1} 的固定顺序串行累加；不存在
同一位置的跨线程竞态，固定二进制、编译选项和输入下通常可重复。

但这不表示 CSR 的数值结果精确或与其他实现逐位相同。若一行有 $k$ 项，采用逐次的
round-to-nearest 加法且无上溢/下溢，常用前向误差界为
\[
 |\widehat{s}-s|\le \gamma_{k-1}\sum_{j=1}^{k}|a_jx_j|,
 \qquad
 \gamma_{k-1}=\frac{(k-1)u}{1-(k-1)u},
\]
其中 $u$ 是 unit roundoff。改变列内顺序、采用 FMA、用 warp 对一行做树形归约，都会改变
舍入路径和 ULP 差异。若所谓 CSR 内核让多个线程协作同一行并以原子方式汇总，它又会重新
具有 COO 类似的顺序不确定性。

\conclusionheading 基础 CSR 消除了原子更新顺序的不确定性，却没有消除浮点舍入误差；
“可重复”也不等于“数学精确”。
\discussionheading
\discussionpoint{没有竞态不代表结果唯一或准确。}
若每一行只由一个线程写，确实消除了两个线程同时更新 $y_i$ 的数据竞态，但该线程仍按某个有限精度顺序累加。把输入重新排序后，执行依然合法，却可能因浮点加法不结合而得到不同末位；固定顺序只能保证这次实现较可重复，不能保证它接近实数精确和。反过来，一个确定的树形归约也可能稳定地产生较大误差。因此测试必须先说明契约：要求逐位相同、允许若干 ULP、比较范数误差，还是只要求迭代求解最终达到目标残差。把这三种性质拆开，是理解所有并行数值程序而不仅是 SpMV 的起点。

\discussionpoint{长行并行化实质上是在选择加法树。}
一条含百万项的行若由单线程处理，延迟无法接受；让 warp 计算部分和再归约，或让多个块原子累加，可以释放并行度，却改变项的配对顺序。固定二叉树能让相同分区重复，动态原子顺序则可能随调度变化；二者都不保证最小误差。提高累加到 FP64、按量级做 pairwise summation，或引入 Kahan/Neumaier 补偿，可以减少舍入损失，但会增加指令、寄存器或合并状态。工程选择应把行长分布和误差预算一起考虑，因为“最快的归约树”与“最稳定的归约树”通常不是同一个目标。

\discussionpoint{工具链也是数值实验的一部分。}
同一索引顺序在不同编译选项下仍可能变化：编译器把乘加收缩成 FMA 时只舍入一次，禁用收缩则在乘法和加法后各舍入一次；快速数学可能使用近似倒数，FTZ 又会把反规格化结果冲成零。库还可能根据 GPU 架构和 workspace 自动选择不同分块，从而改变归约树。需要审计或跨机器复现的计算应记录编译器版本、架构目标、数学选项、库算法标识和输入哈希，并用误差统计而非单个样例作比较。这个要求类似实验科学记录仪器配置：硬件与软件环境不是背景信息，而是结果生成过程的一部分。

\discussionpoint{最终应用决定什么误差才重要。}
Krylov 求解器通常更关心残差曲线和迭代次数，即使单次 SpMV 相差几个 ULP，只要没有改变收敛结论，逐位差异未必有应用意义。整数 BFS 的访问标记不存在同类舍入问题，机器学习训练则可能接受统计波动，却要求损失曲线和最终指标处于可接受分布。因而完整评测可以同时报告逐位一致率、最大与分位 ULP、向量范数误差以及端到端收敛或任务质量。这样做不是放松正确性，而是把局部算术误差连接到问题条件数和业务容限，避免用一个脱离领域的“相等/不等”标签概括所有风险。
\variantheading 对 IEEE FP32 序列 $[10^{20},-10^{20},3]$，顺序 $((10^{20}-10^{20})+3)$ 得 3，而 $10^{20}+(-10^{20}+3)$ 得 0；二者都无数据竞态，却展示了重排造成的可检查差异。
\practiceheading 固定同一 NVIDIA GPU、CUDA 工具链、FP32 FMA 设置和输入顺序，以 cuSPARSE CSR 为性能基线，并用 CUB 固定树归约实现协作行版本；重复运行后同时报告 bitwise 一致率、最大 ULP 和 Nsight Compute 吞吐。
\sourcecheck{https://docs.nvidia.com/cuda/floating-point/index.html}{NVIDIA CUDA Floating Point 官方说明；对抗检查：把无竞态误解为零 ULP 误差、或把任何 CSR 变体都视为确定，都会得到过强结论}
\end{solutionexercise}

\begin{solutionexercise}{8}
\problemheading 混合 ELL--COO 格式中，每行保留的元素上限 $K$ 是一个可调参数。请分析：
$K$ 取得过大或过小分别会带来什么后果？如果已知矩阵各行长度的分布，你会如何选择 $K$？
分析须同时考虑 ELL 填充浪费、规则合并访问、COO 余项数量及原子累加代价。

\approachheading 对行长 $\ell_i$，ELL 槽位数是 $mK$，真实 ELL 项数是
$\sum_i\min(\ell_i,K)$，填充量与 COO 溢出量分别为
\[
P(K)=\sum_i\max(K-\ell_i,0),\qquad
C(K)=\sum_i\max(\ell_i-K,0).
\]

\answerheading $K$ 过大时，$P(K)$ 和设备存储增大，长循环的尾部只有少数线程活跃，
控制分歧与无效内存跨度上升；极端情况下退化为纯 ELL。$K$ 过小时，$C(K)$ 增大，COO
需要更多行/列元数据和原子累加；热点长行会产生原子争用。若按第 4 题在主机处理 COO，
还会增加回传后的串行工作，甚至让 CPU/PCIe 路径主导总时间。$K=0$ 和
$K\ge\max_i\ell_i$ 分别是纯 COO 与纯 ELL 两个边界情形。

选择时先由行长直方图计算 $P(K),C(K)$，再在内存预算内最小化一个面向目标设备的代价模型，
例如
\[
 T(K)\approx \alpha\sum_i\min(\ell_i,K)
       +\beta P(K)+\chi C(K)+\rho\,R(K),
\]
其中 $R(K)$ 表示 COO 中按行集中的原子争用，系数由微基准测得。实用初值可取行长的高分位数
（如 90--99 分位），再在附近枚举；还应把 $mK$ 的字节数、对齐、warp 分组后的行长分布和
矩阵将被复用的迭代次数纳入测量。仅取平均行长会被少量超长行严重误导。

\complexity{评估所有候选 $K$ 可由行长直方图的前后缀和在线性时间完成；每个 SpMV 有 $z$ 个有效乘加}{$2mK$ 个 ELL 值/列槽加 $3C(K)$ 个 COO 值/坐标槽（不计行长元数据）}{ELL 按行并行且访问合并；COO 按溢出元素并行但需要原子汇总}
\conclusionheading 行长分布给出候选范围，硬件代价和真实矩阵复用场景决定最终 $K$；没有
对所有矩阵都最优的固定百分位。
\discussionheading
\discussionpoint{$K$ 可以理解为规则快路径的预算。}
增大 $K$ 会把更多非零元纳入 ELL，使 COO 尾部更短、原子争用更少；与此同时，每一行都必须为更宽的 ELL 预留槽位，短行产生填充，缓存工作集也随之扩大。题中 $P(K)$ 与 $C(K)$ 分别刻画填充和溢出项数，却不能直接相加后宣称时间最优，因为一个填充槽主要消耗带宽，一个 COO 项还可能引入随机访问和原子串行化。更现实的模型是 $T(K)=aP(K)+bC(K)+T_{\mathrm{launch}}$，系数要由设备和数据测得。于是 $K$ 不是纯数学阈值，而是把可预测工作与不规则工作分配给两条执行路径的资源预算。

\discussionpoint{单阈值可以继续发展成多桶调度。}
全矩阵共用一个 $K$，隐含假设是所有行都适合相同并行策略；若行长跨越几个数量级，这个假设很快失效。可以把 0--4 项的行放进窄 ELL，5--32 项的行放进 SELL，极长行交给一个 warp 或线程块协作归约，从而分别减少短行填充和长行串行。多桶方案需要额外的行置换、描述符和多个内核启动，还可能让小桶不足以占满 GPU，所以桶越多也不一定越好。这个演化与任务调度中的 size class 相似：按规模选择专门执行器能降低内部浪费，但分类和调度本身也必须计入总成本。

\discussionpoint{硬件资源会移动最佳阈值。}
若目标 GPU 的 FP32 原子吞吐较高且 L2 能缓存热点输出，较小 $K$ 留下的 COO 尾部可能并不昂贵；在原子争用严重或显存带宽紧张的设备上，更大的规则区反而可能有利。索引从 32 位变为 64 位会放大每个填充槽的字节代价，输入向量复用提高又可能让原子而非读取成为瓶颈。具体地，$K$ 每增加 1 都为全部 $m$ 行新增一个值槽和一个列索引槽，却只把行长超过旧 $K$ 的那些 COO 项移入规则区；短行占多数时，这笔固定带宽支出会迅速压过尾部收益。甚至相同 GPU 上，矩阵常驻与每轮重传也会给出不同最佳点，批量大小变化也会移动平衡位置。因此阈值应附带设备、数据类型和运行方式的假设，不能从一个矩阵或一个架构的曲线推广成普遍常数。

\discussionpoint{动态负载可以用反馈而非固定经验值。}
对结构稳定的矩阵，可以离线扫描候选 $K$，把格式构造成本按预计迭代次数摊销后选择端到端最小值。动态图或自适应网格中，行长分布持续变化，更合适的做法是周期采样填充率、溢出比例和实际计时，仅当未来收益预计覆盖重建成本时才切换，并设置回滞避免频繁往返。例如重建格式耗时 20~ms、预计每轮只节省 0.2~ms 时，至少还要运行约 100 轮才刚能回本，剩余迭代不足便应维持原布局。控制器还应保留一个稳定 CSR 基线，以便模型预测失准时回退。这样的反馈式布局选择把稀疏格式问题连接到在线 autotuning：性能参数不是一次决定永久有效，而是受工作负载漂移和观测成本共同约束。
\variantheading 对行长 $[0,1,2,9]$，$K=2$ 时 $P=3,C=7$，而 $K=3$ 时 $P=6,C=6$；后者只少 1 个溢出项却多 3 个填充槽，若两类槽成本相同则 $K=2$ 的总额外量更小。
\practiceheading 假定 32-lane NVIDIA GPU、FP32 值、32 位索引且矩阵常驻显存，扫描 $K=0\ldots\max\ell_i$，用 cuSPARSE CSR 作基线，并以 Nsight Compute 的 DRAM 字节、原子吞吐和 Warp Execution Efficiency 拟合 $T(K)$。
\sourcecheck{https://mgarland.org/files/papers/nvr-2008-004.pdf}{Bell 与 Garland 对混合 ELL/COO（HYB）格式的原始性能研究；对抗检查：只最小化填充量会忽略 COO 原子争用，只最小化溢出量又会退化成过大的 ELL}
\end{solutionexercise}

\chapter{图遍历}

\begin{solutionexercise}{1}
\problemheading 考虑原书练习 1 所示的有向无权图。为避免复制原图，按图重述其顶点集
$V=\{0,1,2,3,4,5,6,7\}$ 和有向边集
\[
\begin{split}
E=\{&(0,5),(0,2),(5,1),(5,7),(1,7),(1,0),(1,4),\\
    &(7,2),(7,6),(7,4),(2,3),(3,6),(3,0),(6,4),(4,3)\}.
\end{split}
\]
分别完成以下各项：
\begin{enumerate}[label=(\alph*)]
  \item 用邻接矩阵表示该图。
  \item 用 CSR 格式表示该图。每个顶点的邻居列表必须有序。
  \item 从顶点 0 开始在该图上执行并行 BFS，也就是说顶点 0 位于第 0 层。对于 BFS
  遍历的每一次迭代，分别回答：
  \begin{itemize}
    \item 若采用顶点中心推送式实现：\textup{A.} 启动多少个线程？
    \textup{B.} 有多少个线程会遍历自己顶点的邻居？
    \item 若采用顶点中心拉取式实现：\textup{A.} 启动多少个线程？
    \textup{B.} 有多少个线程会遍历自己顶点的邻居？
    \textup{C.} 有多少个线程会标记自己的顶点？
    \item 若采用边中心实现：\textup{A.} 启动多少个线程？
    \textup{B.} 有多少个线程可能标记一个顶点？
    \item 若采用基于前沿的顶点中心推送式实现：\textup{A.} 启动多少个线程？
    \textup{B.} 有多少个线程会遍历自己顶点的邻居？
  \end{itemize}
\end{enumerate}

\approachheading 从图中逐边抄录并按源顶点排序。BFS 的前沿依次决定推送、边中心和
前沿实现的活动线程；拉取实现则检查所有尚未访问顶点的入邻居。

\answerheading 边集合为
\[
\begin{split}
&0\to\{2,5\},\quad 1\to\{0,4,7\},\quad 2\to\{3\},
\quad 3\to\{0,6\},\\
&4\to\{3\},\quad 5\to\{1,7\},\quad 6\to\{4\},
\quad 7\to\{2,4,6\}.
\end{split}
\]
所以邻接矩阵（行是源、列是目的）为
\[
\begin{bmatrix}
0&0&1&0&0&1&0&0\\
1&0&0&0&1&0&0&1\\
0&0&0&1&0&0&0&0\\
1&0&0&0&0&0&1&0\\
0&0&0&1&0&0&0&0\\
0&1&0&0&0&0&0&1\\
0&0&0&0&1&0&0&0\\
0&0&1&0&1&0&1&0
\end{bmatrix}.
\]
排序邻接表的 CSR 是
\[
 \code{srcPtrs}=[0,2,5,6,8,9,11,12,15],
\]
\[
 \code{dst}=[2,5,\;0,4,7,\;3,\;0,6,\;3,\;1,7,\;4,\;2,4,6].
\]

从根 0 得到
\[
L_0=\{0\},\quad L_1=\{2,5\},\quad
L_2=\{1,3,7\},\quad L_3=\{4,6\}.
\]
下表把产生 $L_1,L_2,L_3$ 的三轮以及用于确认没有新顶点的最后一轮都列出。每格按题目
的子问顺序排列；“边可标记”统计同时满足“源在上一层、目的在本轮开始时尚未访问”的边，
所以同一目的可能被计多次。
\begin{center}
\small
\setlength{\tabcolsep}{3.5pt}
\begin{tabular}{c|c|c|c|c|c}
\toprule
轮次 & 新前沿 & \shortstack{推送\\启动/遍历} &
\shortstack{拉取\\启动/遍历/标记} & \shortstack{边中心\\启动/可标记} &
\shortstack{前沿推送\\启动/遍历}\\
\midrule
1 & $\{2,5\}$ & $8/1$ & $8/7/2$ & $15/2$ & $1/1$\\
2 & $\{1,3,7\}$ & $8/2$ & $8/5/3$ & $15/3$ & $2/2$\\
3 & $\{4,6\}$ & $8/3$ & $8/2/2$ & $15/4$ & $3/3$\\
4 & $\varnothing$ & $8/2$ & $8/0/0$ & $15/0$ & $2/2$\\
\bottomrule
\end{tabular}
\end{center}
第 3 轮的 4 条候选边是 $1\to4$、$3\to6$、$7\to4$、$7\to6$；它们只产生两个
不同顶点，体现了推送/边中心实现需要原子“首次访问”或幂等写。若实现通过已知全图访问数
提前判断结束而不启动空轮，可删去表中第 4 轮；本章依赖“本轮是否发现新顶点”的实现必须
执行它才能终止。

\complexity{邻接矩阵构造 $O(V^2+E)$，CSR 构造 $O(V+E)$；BFS 有效工作 $O(V+E)$，全扫描实现可达 $O((V+E)D)$}{$O(V^2)$ 邻接矩阵或 $O(V+E)$ CSR，另有 $O(V)$ 层号/前沿}{$V$ 个顶点线程、$E$ 个边线程或仅当前前沿线程，取决于实现}
\conclusionheading 矩阵含 15 个 1，CSR 末项也是 15，BFS 层集合并集恰为全部 8 个顶点，
三项交叉检查一致。
\discussionheading
\discussionpoint{BFS 也可以看成布尔稀疏线性代数。}
从根开始的第 $k$ 层前沿可写成布尔向量 $f_k$。按照本题“行是源、列是目的”的邻接矩阵约定，列向量形式的推送应在布尔半环上计算 $A^{\mathsf T}f_k$；等价地，若把前沿写成行向量，则计算 $f_k^{\mathsf T}A$，之后再用未访问集合掩码。这里的“加法”是逻辑或，“乘法”是逻辑与，重复父节点不会改变集合语义，但实现仍要避免重复入队。只有 frontier-sparse 的 CSR 推送仅枚举当前前沿各源顶点的出边时，才可因每个顶点至多扩展一次而把全程累计为 $O(V+E)$；若每层都对完整稀疏矩阵做 SpMV，则可能扫描全部 $E$ 条边 $D$ 次，并不自动具有该界。这个视角既连接了 GraphBLAS 原语，也说明存储格式与执行方式共同决定复杂度。

\discussionpoint{前沿的形状会改变遍历方式。}
当一层只有几个顶点时，队列只枚举这些顶点的出边，几乎不碰无关数据；当前沿覆盖大半个图时，维护和争抢队列可能不如用位图扫描未访问顶点。位图让集合交、并和成员测试变得规则，也便于 pull 阶段查询“某个邻居是否在前沿”，但每层要扫描更多位置并承担清零或代际标记成本。生产图框架常在队列、位图和稠密布尔数组之间切换，阈值应把表示转换的时间也完整算进去。这个例子说明数据结构不只是保存同一集合，它还能改变算法可执行的操作以及每轮应启动多少工作。在方向优化 BFS 中，切换本身因此也是被计量的算法步骤。

\discussionpoint{并发发现需要把集合语义落实成唯一写者。}
图中两个父顶点可能在同一层同时到达顶点 $v$，若线程都先普通读取“未访问”再写入，它们可能重复增加下一前沿计数。用 CAS 把距离从未访问哨兵原子地改成 $k+1$，只有成功者获得入队权，便把“集合中只有一个 $v$”落实到了并发执行。某些算法允许重复入队，最终最短层号仍可能正确，但前沿容量、后续边扫描和执行不确定性都会放大，甚至触发容量溢出。原子唯一化因此不是装饰性的同步细节，而是把数学集合、内存容量和终止条件连接起来的实现不变量。若距离数组和队列计数使用不同原子协议，还需检查发布次序。

\discussionpoint{线程计数必须先说明统计边界。}
题目中的小图可以逐层数线程，但“执行了多少”至少有三种口径：只计非空前沿的有效顶点，计入检测终止的空轮，或统计内核实际启动的全部 lane，包括边界检查后退出者。边级映射还会让高出度顶点产生更多任务，所以顶点数与边访问数也不能混为一个指标。两张图即使拥有相同的 $V$ 和 $E$，道路网可能有深而窄的前沿，社交图却会在少数层突然膨胀，运行形状完全不同。严谨实验应给出每层前沿、访问边数、启动线程和耗时，先固定口径再比较算法，避免用单一总数掩盖遍历结构。
\variantheading 若根改为 4，有向 BFS 层依次为 $\{4\},\{3\},\{0,6\},\{2,5\},\{1,7\}$，下一轮为空；五个非空层仍覆盖全部 8 个顶点，但深度由 3 增为 4。
\practiceheading 在顶点和边均用 32 位编号、CSR 可完整放入 NVIDIA GPU HBM 的前提下，用 RAPIDS cuGraph \code{bfs} 校验两组根的 distance；再以 Nsight Systems 核对是否执行终止空轮及每轮内核启动数。
\sourcecheck{https://github.com/rapidsai/cugraph}{RAPIDS cuGraph 官方实现仓库；对抗检查：把有向边当成无向边会改变矩阵、CSR 和所有 BFS 层，漏掉终止空轮则会少报一次全扫描内核}
\end{solutionexercise}

\begin{solutionexercise}{2}
\problemheading 实现原书第 18.3 节所述方向优化 BFS 的主机端代码。主机端须维护层次、
当前/下一前沿及模式状态，根据前沿边数或未访问工作量在顶点中心推送和拉取内核之间切换，
逐层启动正确内核，直到前沿为空；还须说明设备缓冲区清零、交换、同步和终止判定。

\approachheading 同时保留出边 CSR 和入边 CSC。用“本轮前沿规模乘平均度”与剩余未访问
顶点数作启发式比较，并使用滞回阈值避免频繁切换。发现计数必须只统计第一次标记成功的顶点。

\answerheading 下面的两个内核展示关键的竞态处理，主机函数接收已经上传到设备且生命周期
覆盖整个调用的 \code{CSRGraph}/\code{CSCGraph}：
\begin{lstlisting}[language=C++]
constexpr unsigned UNVISITED = 0xffffffffu;

__global__ void bfs_push_cas(CSRGraph g, unsigned* level,
                             unsigned* newCount, unsigned curr) {
    unsigned v = blockIdx.x * blockDim.x + threadIdx.x;
    if (v >= g.numVertices || level[v] != curr - 1) return;
    for (unsigned e = g.srcPtrs[v]; e < g.srcPtrs[v + 1]; ++e) {
        unsigned w = g.dst[e];
        if (atomicCAS(&level[w], UNVISITED, curr) == UNVISITED)
            atomicAdd(newCount, 1u);
    }
}

__global__ void bfs_pull(CSCGraph g, unsigned* level,
                         unsigned* newCount, unsigned curr) {
    unsigned v = blockIdx.x * blockDim.x + threadIdx.x;
    if (v >= g.numVertices || level[v] != UNVISITED) return;
    for (unsigned e = g.dstPtrs[v]; e < g.dstPtrs[v + 1]; ++e) {
        if (level[g.src[e]] == curr - 1) {
            // Exactly one thread owns v in this kernel.
            level[v] = curr;
            atomicAdd(newCount, 1u);
            break;
        }
    }
}

std::vector<unsigned> direction_bfs(CSRGraph csr_d, CSCGraph csc_d,
                                    unsigned V, unsigned E,
                                    unsigned root) {
    if (V == 0 || root >= V) throw std::invalid_argument("BFS root");
    unsigned *dLevel = nullptr, *dNew = nullptr;
    CHECK(cudaMalloc(&dLevel, size_t(V) * sizeof(unsigned)));
    CHECK(cudaMalloc(&dNew, sizeof(unsigned)));
    CHECK(cudaMemset(dLevel, 0xff, size_t(V) * sizeof(unsigned)));
    unsigned zero = 0;
    CHECK(cudaMemcpy(dLevel + root, &zero, sizeof(unsigned),
                     cudaMemcpyHostToDevice));

    constexpr unsigned BLOCK = 256;
    unsigned blocks = (V + BLOCK - 1) / BLOCK;
    unsigned frontier = 1, visited = 1, curr = 1;
    bool pull = false;
    const double avgDegree = double(E) / V;
    const double alpha = 14.0; // Tune on the target graph/device.
    const double beta  = 24.0;

    while (frontier != 0) {
        unsigned remaining = V - visited;
        if (!pull && double(frontier) * avgDegree > remaining / alpha)
            pull = true;
        else if (pull && frontier < double(V) / beta)
            pull = false;

        CHECK(cudaMemset(dNew, 0, sizeof(unsigned)));
        if (pull)
            bfs_pull<<<blocks, BLOCK>>>(csc_d, dLevel, dNew, curr);
        else
            bfs_push_cas<<<blocks, BLOCK>>>(csr_d, dLevel, dNew, curr);
        CHECK(cudaGetLastError());
        CHECK(cudaMemcpy(&frontier, dNew, sizeof(unsigned),
                         cudaMemcpyDeviceToHost)); // completion point
        visited += frontier;
        ++curr;
    }

    std::vector<unsigned> level(V);
    CHECK(cudaMemcpy(level.data(), dLevel, size_t(V)*sizeof(unsigned),
                     cudaMemcpyDeviceToHost));
    CHECK(cudaFree(dNew)); CHECK(cudaFree(dLevel));
    return level;
}
\end{lstlisting}
启发式常数不是 BFS 正确性的一部分，合理答案可以改用前沿边数与未访问边数、或离线调优的
切换规则。正确性只要求一轮完整处理第 $d-1$ 层后才开始标记第 $d+1$ 层；默认 stream 中
逐轮内核和同步回传满足该依赖。推送用 CAS 保证多父顶点争抢同一邻居时只计数一次；拉取的
每个目的顶点只有一个线程，并在首个匹配入邻居处退出。

\complexity{取决于切换点；推送检查前沿出边，拉取检查未访问顶点直到首个匹配入边，最坏仍为 $O((V+E)D)$}{$O(V+E)$ 的 CSR 与 CSC 各一份，另有 $O(V)$ 层号}{每轮启动 $V$ 个线程；可进一步用显式前沿减少推送阶段的无效线程}
\conclusionheading 方向优化改变每轮检查的边集合，不改变层次语义。
CAS、轮间完成点以及 CSR 与 CSC 的生命周期，都是主机代码不可省略的正确性条件。
\discussionheading
\discussionpoint{push 与 pull 利用了前沿密度的变化。}
遍历早期只有少数活动顶点，push 沿它们的出边寻找新节点，工作量接近当前前沿出边数；当前沿膨胀后，许多边会指向已经访问的顶点，这些尝试大多浪费。pull 改让每个未访问顶点检查入边，只要找到一个前沿父节点就能早停，因此在稠密阶段可能少看很多无效边。二者的交叉点取决于度分布、剩余未访问集合和命中位置，不是固定的第几层。方向优化的广泛意义在于：同一个可达性问题可以从生产者或消费者一侧求解，系统应根据当前活动集合选择较少的候选工作。高偏斜图中，切换判据应估计边数而不能只数前沿顶点。

\discussionpoint{双向邻接索引用空间换取遍历方向。}
push 需要快速枚举出边，CSR 按源顶点分段最自然；pull 需要枚举入边，CSC 或转置图才能给出同样的顺序访问。两种格式同时驻留几乎复制了全部边索引，并各自带一组指针，但切换方向时不必临时扫描或转置。以 32 位顶点编号为例，额外方向至少再保存约 $4E$ 字节的邻接索引和 $4(V+1)$ 字节的偏移，十亿边图仅索引副本就接近 4~GB。静态图上的一次构造可由许多次 BFS、PageRank 或查询摊销，动态图却要让每次插入删除同时维护两份索引，否则两种方向会看到不一致拓扑。这个取舍与数据库的辅助索引相似：额外索引加速特定查询，同时增加容量、构建时间与更新一致性责任。

\discussionpoint{切换条件应估算边工作而不是只数顶点。}
一个含 1000 个度为 1 顶点的前沿只需检查约 1000 条出边，而一个超级节点可能独自拥有百万条边；把两者都记成“前沿大小 1000 或 1”会严重误判 push 成本。较可靠的启发式会估算当前前沿出边总数、剩余未访问顶点的入边量以及 pull 的近期命中率，并用两个不同阈值形成回滞。即使未访问顶点的入度总和相同，若前沿邻居常出现在第一条入边，pull 会很快早停；若几乎没有命中，它就要把候选入边扫到底，因此近期命中位置也是有效反馈。否则图在临界密度附近可能每层来回切换，转换位图和队列的成本抵消收益。这里体现了负载预测的一般原则：任务数量只是第一层信息，任务权重分布和切换成本才决定控制策略。

\discussionpoint{主机控制可能比一层遍历更昂贵。}
若每层都把新前沿长度复制到主机，再由 CPU 决定下一内核，设备同步和链路往返可能达到微秒级；对直径很小、单层工作也很短的图，这与计算本身处在同一量级。设备端队列、CUDA Graph 或持久内核可以减少启动往返，但分别引入动态控制复杂度、固定图结构约束或长期占用资源。进入多 GPU 后，远程发现去重、分区边通信和全局终止检测又会加入控制路径，单卡阈值无法原样照搬。优化 BFS 因而不只是在内核里少看几条边，还要设计谁作决策、状态在哪里以及何时必须全局一致。
\variantheading 若设备内存只容纳 CSR，禁用拉取并始终运行 \code{bfs_push_cas}，本题图仍产生新前沿规模 $2,3,2,0$ 和第 1 题相同层号；代价是不能利用未访问集合很小时的入边早停。
\practiceheading 假定 32-lane NVIDIA GPU、32 位顶点编号且 CSR/CSC 都驻留显存，用 RAPIDS cuGraph BFS 作正确性基线；在 Nsight Compute 中记录访问边数、原子吞吐和分支效率，针对真实图调节 $\alpha,\beta$ 而非沿用常数。
\sourcecheck{https://docs.nvidia.com/cuda/cuda-c-programming-guide/}{CUDA C++ Programming Guide 对原子操作、内核顺序和内存模型的官方说明；对抗检查：若推送仅先读“未访问”状态再普通写，两个父线程会重复增加新节点计数，使切换统计和终止条件失真}
\end{solutionexercise}

\begin{solutionexercise}{3}
\problemheading 修改原书图 18.17 的内核，使每层只使用一次网格范围屏障同步，而不是
三次。原内核在持久合作网格中依次清零下一前沿计数、扩展当前前沿、交换前沿并更新层号；
修改后须用线程 0 的顺序操作和一次全网格屏障同时保证本层结果对所有块可见，并避免任何块
在屏障前提前退出。

\approachheading 原实现复用一个当前前沿计数器，所以必须“完成写入、全部读完、清零完成”
三次同步。为第 $d$ 层准备独立计数器，内核启动前一次性清零，便不再需要读后清零。块私有
前沿采用有界数组：前 \code{PRIVATE_CAPACITY} 个发现留在共享内存，其余发现立即原子追加
到全局前沿；块尾再为共享部分一次性领取连续全局区间并 flush。

\answerheading 以下是可由 \code{nvcc} 编译的完整内核和主机启动函数。传入的 CSR 两个
指针已经位于当前设备，\code{rowPtr} 长度为 $V+1$，
\code{rowPtr[V]==numEdges}，所有目的顶点均小于 $V$：
\begin{lstlisting}[language=C++]
#include <cuda_runtime.h>
#include <cooperative_groups.h>
#include <algorithm>
#include <stdexcept>
#include <vector>

namespace cg = cooperative_groups;

inline void check_cuda(cudaError_t e) {
  if (e != cudaSuccess) throw std::runtime_error(cudaGetErrorString(e));
}
#define CHECK(call) check_cuda(call)

struct CSRGraphView {
  const unsigned* rowPtr;
  const unsigned* colIdx;
  unsigned numVertices;
  unsigned numEdges;
};

template<class T> struct DeviceBuffer {
  T* p = nullptr;
  explicit DeviceBuffer(size_t count) {
    CHECK(cudaMalloc(reinterpret_cast<void**>(&p), count * sizeof(T)));
  }
  ~DeviceBuffer() { if (p) cudaFree(p); }
  DeviceBuffer(const DeviceBuffer&) = delete;
  DeviceBuffer& operator=(const DeviceBuffer&) = delete;
};

constexpr unsigned UNVISITED = ~0u;
constexpr unsigned PRIVATE_CAPACITY = 256;

__global__ void bfs_one_grid_barrier(
    CSRGraphView g, unsigned* level,
    unsigned* prevFrontier, unsigned* currFrontier,
    unsigned* frontierSizes) {
  cg::grid_group grid = cg::this_grid();
  unsigned numPrev = frontierSizes[0];

  for (unsigned currLevel = 1; numPrev != 0; ++currLevel) {
    __shared__ unsigned privateFrontier[PRIVATE_CAPACITY];
    __shared__ unsigned privateCount;
    __shared__ unsigned publicBase;
    if (threadIdx.x == 0) privateCount = 0;
    __syncthreads();

    for (unsigned long long i = grid.thread_rank(); i < numPrev;
         i += grid.size()) {
      unsigned src = prevFrontier[i];
      for (unsigned e = g.rowPtr[src]; e < g.rowPtr[src + 1]; ++e) {
        unsigned dst = g.colIdx[e];
        if (atomicCAS(&level[dst], UNVISITED, currLevel) == UNVISITED) {
          unsigned slot = atomicAdd(&privateCount, 1u);
          if (slot < PRIVATE_CAPACITY) {
            privateFrontier[slot] = dst;
          } else {
            // Overflow entries bypass the full shared queue.
            unsigned out = atomicAdd(&frontierSizes[currLevel], 1u);
            currFrontier[out] = dst;
          }
        }
      }
    }
    __syncthreads();

    unsigned staged = privateCount < PRIVATE_CAPACITY
                    ? privateCount : PRIVATE_CAPACITY;
    if (threadIdx.x == 0)
      publicBase = atomicAdd(&frontierSizes[currLevel], staged);
    __syncthreads();
    for (unsigned i = threadIdx.x; i < staged; i += blockDim.x)
      currFrontier[publicBase + i] = privateFrontier[i];

    grid.sync();                 // exactly one grid-wide barrier per level
    numPrev = frontierSizes[currLevel];
    unsigned* tmp = prevFrontier;
    prevFrontier = currFrontier;
    currFrontier = tmp;
  }
}

std::vector<unsigned> cooperative_bfs(CSRGraphView g, unsigned root) {
  const unsigned V = g.numVertices;
  if (V == 0 || root >= V || !g.rowPtr || (g.numEdges && !g.colIdx))
    throw std::invalid_argument("CSR graph/root");

  int device = 0, cooperative = 0, smCount = 0, blocksPerSM = 0;
  CHECK(cudaGetDevice(&device));
  CHECK(cudaDeviceGetAttribute(&cooperative,
                               cudaDevAttrCooperativeLaunch, device));
  if (!cooperative)
    throw std::runtime_error("device has no cooperative launch support");

  constexpr int BLOCK = 256;
  CHECK(cudaDeviceGetAttribute(&smCount,
                               cudaDevAttrMultiProcessorCount, device));
  CHECK(cudaOccupancyMaxActiveBlocksPerMultiprocessor(
      &blocksPerSM, bfs_one_grid_barrier, BLOCK, 0));
  int residentLimit = smCount * blocksPerSM;
  int workBlocks = int((size_t(V) + BLOCK - 1) / BLOCK);
  int gridBlocks = std::min(workBlocks, residentLimit);
  if (gridBlocks <= 0)
    throw std::runtime_error("cooperative grid cannot be resident");

  DeviceBuffer<unsigned> level(V), prev(V), curr(V);
  DeviceBuffer<unsigned> sizes(size_t(V) + 1);
  unsigned *dLevel = level.p, *dPrev = prev.p, *dCurr = curr.p;
  unsigned *dSizes = sizes.p;
  CHECK(cudaMemset(dLevel, 0xff, size_t(V) * sizeof(unsigned)));
  unsigned zero = 0, one = 1;
  CHECK(cudaMemcpy(dLevel + root, &zero, sizeof(zero),
                   cudaMemcpyHostToDevice));
  CHECK(cudaMemcpy(dPrev, &root, sizeof(root), cudaMemcpyHostToDevice));
  CHECK(cudaMemset(dSizes, 0, (size_t(V) + 1) * sizeof(unsigned)));
  CHECK(cudaMemcpy(dSizes, &one, sizeof(one), cudaMemcpyHostToDevice));

  void* args[] = {&g, &dLevel, &dPrev, &dCurr, &dSizes};
  CHECK(cudaLaunchCooperativeKernel(
      (void*)bfs_one_grid_barrier, dim3(gridBlocks), dim3(BLOCK),
      args, 0, nullptr));
  CHECK(cudaDeviceSynchronize());

  std::vector<unsigned> result(V);
  CHECK(cudaMemcpy(result.data(), dLevel, size_t(V) * sizeof(unsigned),
                   cudaMemcpyDeviceToHost));
  return result;
}
\end{lstlisting}
每个成功的 CAS 对应一个且仅一个前沿元素。共享计数器可以超过容量，但只有槽号小于容量的
元素写共享数组；其余元素直接写入原子领取的全局槽，故不会覆盖或丢失。块尾的
\code{staged=min(privateCount,PRIVATE_CAPACITY)} 只 flush 实际存在的有界部分。CAS 使每个
顶点最多进入一次前沿，因此两个长度为 $V$ 的前沿缓冲区不会溢出。

唯一的 \code{grid.sync()} 保证所有全局前沿写入及本层计数原子更新完成；屏障后所有线程
读取同一计数并交换相同指针。下一层使用另一个预先清零、从未写过的计数项，因而不需要
“全部读完再清零”和“清零完成再写”两道屏障。长度 $V+1$ 还容纳最长简单路径末尾的空前沿。
主机既检查 cooperative launch 能力，又用 occupancy 上限乘 SM 数限制网格；若启动超过
同时驻留容量，网格屏障可能无法取得进展。块内三个 \code{__syncthreads()} 仍是共享队列
初始化、计数完成和 flush 基址发布所需，它们不属于题目要求减少的网格屏障。

\complexity{$O(V+E)$ 有效遍历工作；每层网格屏障从 3 次降至 1 次，原子争用可能增加常数}{$O(V)$ 层计数器和两个 $O(V)$ 前沿；每块另有固定 \code{PRIVATE_CAPACITY} 个共享槽}{$\min(\lceil V/256\rceil,\text{驻留上限})$ 个持久块，以 grid-stride 遍历前沿；全部块必须同时驻留}
\conclusionheading 每层独立计数器把跨层生命周期变成单赋值，因而一次“生产完成”屏障
已经足够。
\discussionheading
\discussionpoint{持久协作内核把层循环留在设备上。}
普通 BFS 每一层结束一个内核，内核边界天然提供全网格同步，也让主机有机会读取前沿长度。协作组的网格屏障可以在单个持久内核内建立同样的阶段边界，从而省去反复启动与主机往返；条件是参与屏障的所有块必须能够同时驻留，否则未调度块与已驻留块会互相等待。这个驻留上限由线程、寄存器和共享内存共同决定，不能把普通网格尺寸直接用于 cooperative launch。持久方案适合层数多而每层较短的图，但若一层本来就足以占满设备，多内核方案的弹性和可移植性可能更好。

\discussionpoint{版本化计数器是在用空间消除清零竞争。}
若上一层长度和下一层写入位置共用一个计数器，线程必须先全部读完旧值，再由一个线程清零，并在任何新写入前再次完成网格同步；少一个屏障都可能把新计数清掉。按层号保存 \code{count[level]} 相当于给状态加版本，读者只访问旧槽，写者只更新新槽，从结构上消除了生命周期重叠。最坏情况下 BFS 深度为 $V-1$，这种做法需要 $O(V)$ 计数空间；双缓冲可降到常数，但必须恢复严格的读、屏障、清零、屏障协议。版本化状态也是无锁队列和分代缓存中的常见思想：多占一点空间，换取更容易证明的时间所有权。

\discussionpoint{块私有前沿只缓解全局原子热点。}
让每个新顶点直接对全局尾指针执行原子加，会让所有块争用一个地址；先写入共享内存队列，再由块一次申请连续全局区间，可以把许多全局原子压缩成少量批量操作。典型 flush 先由一个线程以一次 \code{atomicAdd} 预留 $n$ 个连续槽，再把返回基址广播给全块并行散写；预留前后的块内屏障若缺失，就可能读取未完成的本地队列或未发布的基址。共享队列内部仍需要原子或确定分配，容量满时也必须有直接写全局等溢出路径，否则高峰前沿会静默丢顶点。容量增大能减少 flush，却提高每块共享内存并可能降低驻留块数，所以它不是越大越好。类似的分层聚合还用于直方图和日志缓冲，其共同原则是把高代价的全局争用变成本地争用，再以批量方式跨层级提交。

\discussionpoint{长期活性比单轮数值更难验证。}
协作内核即使在一个小图上给出正确层号，也可能在更大网格上因启动超过可驻留上限而失败，或因某条异常路径没有到达 \code{grid.sync()} 而死锁。测试应覆盖空图、长链、星形图、重复边、队列恰满和溢出，并用 occupancy API 根据实际寄存器与共享内存核对最大协作块数。桌面 GPU 还可能对长时间内核施加看门狗限制，使算法正确却无法在该部署环境长期运行。生产系统因此会在持久内核、多内核和 CUDA Graph 之间选择，权衡启动延迟、抢占能力、可诊断性与跨设备可移植性，而不是只比较一次内核计时。
\variantheading 若把 \code{PRIVATE_CAPACITY} 改为 128，某块一层发现 300 个唯一顶点时，128 个经共享队列批量 flush，172 个直接追加全局前沿；屏障后该块对总计数的贡献仍恰为 300。
\practiceheading 在 \code{cooperativeLaunch=true}、256 线程块且网格不超过 occupancy 驻留上限的 NVIDIA GPU 上，用 CUDA Occupancy API 核对可驻留块数，并以 Nsight Compute 的 Barrier Stall、原子吞吐和共享内存占用比较容量 64、128、256。
\sourcecheck{https://docs.nvidia.com/cuda/cuda-c-programming-guide/}{CUDA C++ Programming Guide 的 Cooperative Groups 网格同步与协作启动规范；对抗检查：若仍复用两个计数器却直接删掉读后/清零屏障，快线程可能在慢线程读取前清零，或在清零前写入下一层}
\end{solutionexercise}

\chapter{卷积神经网络}

\begin{solutionexercise}{1}
\problemheading 实现原书第 19.1 节介绍的下采样层前向传播；附录 B 对这种层作了说明。
采用本章的 $[N,C,H,W]$ 布局，每个输出通道对应同一输入通道，把每个不重叠的
$K\times K$ 区域取平均，加上该通道偏置，再通过 sigmoid 激活。应给出输出尺寸、边界
约定和 CUDA 线程到 $(n,c,h_o,w_o)$ 的映射。

\approachheading 扩展为正文采用的批处理 NCHW 布局。每个线程唯一负责
$(n,m,h_o,w_o)$，所以既不需要原子操作，也不需要块内同步。

\answerheading 若 $H_o=\lfloor H/K\rfloor$、$W_o=\lfloor W/K\rfloor$，实现为：
\begin{lstlisting}[language=C++]
__global__ void subsample_forward_nchw(
    const float* input, const float* bias, float* output,
    int N, int M, int H, int W, int K) {
  size_t Ho = H / K, Wo = W / K;
  size_t total = size_t(N) * M * Ho * Wo;
  size_t id = size_t(blockIdx.x) * blockDim.x + threadIdx.x;
  if (id >= total) return;

  int ow = id % Wo; id /= Wo;
  int oh = id % Ho; id /= Ho;
  int m  = id % M;  int n = id / M;

  float sum = 0.0f;
  for (int p = 0; p < K; ++p)
    for (int q = 0; q < K; ++q) {
      size_t x = ((size_t(n) * M + m) * H + (oh*K+p)) * W
                 + (ow*K+q);
      sum += input[x];
    }
  float a = sum / float(K*K) + bias[m];
  size_t y = ((size_t(n) * M + m) * Ho + oh) * Wo + ow;
  output[y] = 1.0f / (1.0f + expf(-a));
}

void launch_subsample(const float* x, const float* b, float* y,
                      int N, int M, int H, int W, int K,
                      cudaStream_t stream) {
  if (N < 0 || M < 0 || H < 0 || W < 0 || K <= 0)
    throw std::invalid_argument("subsampling shape");
  size_t total = size_t(N) * M * (H/K) * (W/K);
  if (total == 0) return;
  constexpr unsigned BLOCK = 256;
  unsigned grid = unsigned((total + BLOCK - 1) / BLOCK);
  subsample_forward_nchw<<<grid, BLOCK, 0, stream>>>(
      x, b, y, N, M, H, W, K);
  CHECK(cudaGetLastError());
}
\end{lstlisting}
每个输入元素在 $H,W$ 能被 $K$ 整除时恰好属于一个区域；否则该定义丢弃右侧和底部不足
$K$ 的尾部。另一种合格接口可以要求整除或显式定义 padding，但不能越界读取。$N=1$
即退化成附录中的非批处理实现。

\complexity{$O(NM H_oW_oK^2)$，整除时为 $O(NMHW)$}{$O(NMH_oW_o)$ 输出，每线程 $O(1)$ 额外状态}{$NMH_oW_o$ 个互不依赖的输出任务；相邻输出线程没有写竞态}
\conclusionheading 索引先覆盖样本和通道，再覆盖输出像素；平均、偏置、激活的先后顺序与
题设一致。
\discussionheading
\discussionpoint{算子融合首先消除中间张量的物化。}
若平均池化、加偏置和 sigmoid 分成三个内核，池化结果必须写入 HBM，后两个内核又分别读取和写回；这些字节并不属于模型的最终输入输出。融合后，线程可以让池化部分和停留在寄存器，立即加偏置并计算激活，只写一次最终值。收益主要来自减少全局流量和启动次数，而不是减少池化本身的算术。代价是原来可复用的通用算子变成形状、布局和激活函数绑定的专用代码，寄存器压力也可能提高。融合因此适合中间张量大且算术较轻的链条，却不应在所有形状上被视为无条件更快。

\discussionpoint{窗口边界是模型定义而非实现细节。}
当 $H$ 或 $W$ 不能被池化宽度整除时，可以丢弃尾部、补零到完整窗口，或保留缩小窗口并按实际元素数除；三种选择会产生不同输出形状和数值。最大池化还要定义 NaN 是否传播，平均池化要说明补零是否进入分母，偏置究竟按通道还是按输出位置广播。若训练框架与自写内核在这些约定上不同，即使内部索引完全正确，端到端模型也会发生系统偏差。接口文档和测试应把奇数尺寸、尾窗口、NaN 与广播维度列成语义用例，这揭示了高性能内核必须首先忠于上层算子契约。这类差异常在小张量单元测试中最容易暴露。

\discussionpoint{sigmoid 与池化各有自己的精度策略。}
直接用 $1/(1+\exp(-x))$ 计算 sigmoid 时，极大负 $x$ 会让 $\exp(-x)$ 溢出；可以按符号分支，正半轴使用原式，负半轴改写成 $\exp(x)/(1+\exp(x))$，使中间量保持有限。池化则涉及多项求和，FP16 或 BF16 输入常用 FP32 累加以减少舍入，再按接口要求转换输出。融合内核必须分别声明输入、累加、初等函数和输出精度，不能用“FP16 实现”掩盖内部混合精度。这个分层思路也适用于归一化和损失函数：稳定公式与高精度累加解决不同误差来源，需要分别验证。

\discussionpoint{图级优化器必须决定融合是否值得。}
在完整网络中，一层的最快融合方案可能迫使前后算子改变布局，或阻断另一条更有价值的融合链，因此局部微基准并不足以作决定。图编译器会根据具体形状、步长、布局和后继消费者生成候选，再用代价模型或 autotuning 比较；动态形状还可能触发重新编译和缓存增长。评测时应分开报告独立算子总和、融合内核时间和端到端网络时间，并计入布局转换与编译缓存成本。这个过程把手写融合扩展成编译器的全图规划问题：优化目标不是某个内核的漂亮数字，而是整段数据流最少的物化和最低的总延迟。
\variantheading 若 $N=M=1,H=W=5,K=2$，则输出为 $2\times2$，只使用 16 个输入并丢弃最右列和最底行共 9 个元素；输入全 1、偏置为 0 时四个输出均为 $\sigma(1)\approx0.7310586$。
\practiceheading 在连续 NCHW、IEEE FP32 且支持 CUDA 的 NVIDIA GPU 上，用 cuDNN average pooling 加 pointwise sigmoid 校验自写融合内核；以 Nsight Compute 的 DRAM Bytes 和 SM Throughput 判断省去中间张量是否值得。
\sourcecheck{https://docs.nvidia.com/deeplearning/cudnn/backend/latest/}{NVIDIA cuDNN Backend 官方文档的 pooling/tensor 描述；对抗检查：最大池化与平均池化的前向规则不同，尾部窗口和 NaN 传播策略也必须由接口显式约定}
\end{solutionexercise}

\begin{solutionexercise}{2}
\problemheading 本章的输入和输出特征采用 $[N,C,H,W]$（NCHW）布局。若改为
$[N,H,W,C]$（NHWC）布局，能否降低内存带宽需求？采用 $[C,H,W,N]$（CHWN）布局
可能有哪些好处？回答须区分算法必须传输的标量字节数与实际内存事务，并结合相邻线程所
遍历的连续维度说明条件。

\approachheading 区分“算法必须读取的元素数”和“硬件实际发出的内存交易数”。布局不改变
卷积数学，却会改变相邻线程看到的连续维、向量化方式和是否需要布局转换。

\answerheading 单纯把 NCHW 改为 NHWC \textbf{不会减少必须读取和写入的标量个数}，因此
不能无条件降低理论字节数。它可能降低也可能提高实际 DRAM 交易：
\begin{itemize}
  \item 原书的直接 NCHW 内核让相邻线程沿 $W$ 计算相邻输出；固定 $n,c,p,q$ 时，它们读取
  相邻输入列，天然合并。若不改变线程映射却换成 NHWC，固定通道的相邻像素相隔 $C$，反而
  可能降低合并性。
  \item NHWC 把同一像素的通道连续存放。若内核让线程或向量指令沿 $C$ 维工作，特别是
  通道数满足向量/Tensor Core 对齐时，可以合并通道加载、使用宽向量，并避免库内部转置，
  从而减少实际交易和中间布局转换。
\end{itemize}

CHWN 把批次 $N$ 设为最快变化维。CNN 的同一滤波器用于全部样本，所以若 warp 的线程对
同一 $(c,h,w)$ 处理连续样本，输入和输出都连续，权重还可在这些样本间广播/复用。这在训练
批次较大、$N$ 对 warp 或向量宽度对齐时很有吸引力。代价是常见算子和框架不一定原生支持，
$N$ 小或动态时尾部利用率差，沿空间维的内核不再自然合并，层间转换可能抵消全部收益。

可接受答案应至少给出三个评分点：理论标量字节数不因置换而自动减少；合并性取决于线程
映射；CHWN 利于跨 batch 向量化和权重复用。把某一种布局宣称为所有形状和硬件上的固定最优
不成立，最终应针对完整网络测量，并把转换成本计入。

\conclusionheading 布局是“哪一维连续”的选择，不是压缩；只有线程映射、硬件指令和相邻
算子共同配合时，才会减少实际带宽消耗。
\discussionheading
\discussionpoint{布局改变的是地址关系而非数学数据量。}
NCHW、NHWC 和 CHWN 都保存 $NCHW$ 个标量，卷积公式也不会因为字母顺序改变而少读一个逻辑元素。差别出现在相邻 lane 的地址：若一个 warp 固定像素并沿通道展开，NHWC 让 32 个 FP32 通道覆盖连续 128 字节；在 NCHW 中它们可能相隔 $HW$ 个元素，需要更多内存交易。反过来，沿空间铺线程时 NCHW 又可能更自然。因此所谓“布局带宽收益”通常来自更少的 sector、更好的对齐或缓存复用，而不是算法字节数凭空减少；分析必须先写出线程映射，再讨论哪一维应该连续。

\discussionpoint{每类算子希望连续的维度并不相同。}
逐像素通道归一化和许多 Tensor Core 卷积希望同一位置的通道连续，NHWC 便于向量打包；按特征图逐点处理时，NCHW 的空间平面连续可能更方便。CHWN 把样本索引放在最快维，早期卷积实现可让同一滤波器跨 batch 读取连续样本，但只有 batch 足够大时才能填满 lane。通道数、batch 或空间尺寸过小都会破坏这些优势，尾部还可能需要掩码或填充。布局选择因此是算子形状与线程映射的函数，不能从某个时代或某块 GPU 的成功实现推导成跨模型的永久规则。

\discussionpoint{布局必须沿整个计算图传播。}
如果某一卷积在 NHWC 下快 0.2~ms，却要求输入和输出各做一次 1~ms 的全张量转置，这个局部优化显然让网络更慢。一次转置 1~GiB 的激活至少要读入并写出各 1~GiB，产生约 2~GiB 数据移动；若不能与相邻算子融合，这笔带宽成本不会因为卷积本身更快而消失。残差连接和多分支结构还要求汇合张量拥有兼容的逻辑索引与 stride，否则转换会在多个分支重复发生。成熟框架会让布局成为张量描述符的一部分，在图上向前后传播选择，并尝试把必要转换融合进生产者或消费者。于是单层 autotuning 需要升级成全图代价问题：不仅比较内核，还要计算转换字节、分支兼容和布局在后续多层中的复用价值。

\discussionpoint{真实物理格式常超出四个字母。}
高性能内核可能把通道按 8 或 16 个元素分块，形成 NHWC8、NHWC16，也可能在共享内存或全局内存中使用 swizzle，以满足向量指令、bank 或矩阵乘片段的地址要求。逻辑通道数不是块宽倍数时要填充，因而实际 stride 与占用字节大于 $NCHW$ 的简单乘积。用户仍可通过张量描述符看到原始形状，但内核和转换必须理解物理块布局。调优报告应同时给出逻辑形状、实际 stride、填充量、转换次数和端到端延迟，这能避免把一个理想整齐尺寸上的微基准结果误推广到带尾部的真实模型。
\variantheading 若 $C=64$ 且一个 warp 的 32 个 lane 固定同一 $(n,h,w)$、分别处理 $c=0\ldots31$，NHWC 给出 32 个连续 float 的 128 字节区段；NCHW 中相邻 lane 相隔 $HW$ 个 float，不能由同一连续区段服务。
\practiceheading 假定 Ampere 或更新 NVIDIA Tensor Core GPU、FP16 输入且 $C$ 满足 16 元素对齐，用 cuDNN Frontend 和 CUTLASS Profiler 比较 NCHW/NHWC；Nsight Systems 必须把布局转换计入，Nsight Compute 再报告实际 DRAM 吞吐。
\sourcecheck{https://github.com/akrizhevsky/cuda-convnet2/blob/master/cudaconv3/src/filter_acts.cu}{Krizhevsky 的原始 cuda-convnet2 卷积实现直接把图像注释为 $(\code{numImgColors},\code{imgSizeY},\code{imgSizeX},\code{numImages})$，并让图像编号成为连续维；对抗检查：这一实现证实 CHWN 的跨样本连续访问机制，却不证明它对所有形状和算子都优于 NCHW/NHWC}
\end{solutionexercise}

\begin{solutionexercise}{3}
\problemheading 实现原书第 19.1 节介绍、附录 B 进一步说明的卷积层反向传播。采用
NCHW 输入 $X_{N\times C\times H\times W}$、滤波器
$F_{M\times C\times K\times K}$、无 padding、步长 1，并令
\[
Y_{n,m,h,w}=\sum_{c,p,q}X_{n,c,h+p,w+q}F_{m,c,p,q}.
\]
给定上游梯度 $G=\partial E/\partial Y$，实现输入梯度 $\partial E/\partial X$ 和
滤波器梯度 $\partial E/\partial F$，说明线程映射、边界、竞态处理和累加次序。

\approachheading 前向定义为
\[
Y_{n,m,h,w}=\sum_{c,p,q}X_{n,c,h+p,w+q}F_{m,c,p,q}.
\]
令上游梯度为 $G=\partial E/\partial Y$。为避免原子竞态，给每个 $dX$ 或 $dF$ 元素分配
唯一线程，在该线程内部完成全部贡献的确定次序累加。

\answerheading
\begin{align*}
dX_{n,c,i,j}
 &=\sum_{\substack{m,p,q\\0\le i-p<H_o,\ 0\le j-q<W_o}}
     G_{n,m,i-p,j-q}F_{m,c,p,q},\\
dF_{m,c,p,q}
 &=\sum_{n,h,w}X_{n,c,h+p,w+q}G_{n,m,h,w}.
\end{align*}
对应 CUDA 核心如下：
\begin{lstlisting}[language=C++]
__global__ void conv_backward_x(const float* G, const float* F,
                                float* dX, int N, int M, int C,
                                int H, int W, int K) {
  int Ho = H-K+1, Wo = W-K+1;
  size_t id = size_t(blockIdx.x)*blockDim.x + threadIdx.x;
  size_t total = size_t(N)*C*H*W;
  if (id >= total) return;
  int iw=id%W; id/=W; int ih=id%H; id/=H;
  int c=id%C; int n=id/C;
  float s=0.0f;
  for (int m=0; m<M; ++m)
    for (int p=0; p<K; ++p)
      for (int q=0; q<K; ++q) {
        int oh=ih-p, ow=iw-q;
        if (oh>=0 && oh<Ho && ow>=0 && ow<Wo) {
          size_t gy=((size_t(n)*M+m)*Ho+oh)*Wo+ow;
          size_t f=((size_t(m)*C+c)*K+p)*K+q;
          s += G[gy]*F[f];
        }
      }
  size_t x=((size_t(n)*C+c)*H+ih)*W+iw;
  dX[x]=s;
}

__global__ void conv_backward_f(const float* G, const float* X,
                                float* dF, int N, int M, int C,
                                int H, int W, int K) {
  int Ho=H-K+1, Wo=W-K+1;
  size_t id=size_t(blockIdx.x)*blockDim.x+threadIdx.x;
  size_t total=size_t(M)*C*K*K;
  if (id>=total) return;
  int q=id%K; id/=K; int p=id%K; id/=K;
  int c=id%C; int m=id/C;
  float s=0.0f;
  for (int n=0; n<N; ++n)
    for (int h=0; h<Ho; ++h)
      for (int w=0; w<Wo; ++w) {
        size_t x=((size_t(n)*C+c)*H+(h+p))*W+(w+q);
        size_t gy=((size_t(n)*M+m)*Ho+h)*Wo+w;
        s += X[x]*G[gy];
      }
  size_t f=((size_t(m)*C+c)*K+p)*K+q;
  dF[f]=s;
}

__global__ void conv_backward_b(const float* G, float* db,
                                int N, int M, int Ho, int Wo) {
  int m=blockIdx.x*blockDim.x+threadIdx.x;
  if (m>=M) return;
  float s=0.0f;
  for (int n=0;n<N;++n) for (int h=0;h<Ho;++h)
    for (int w=0;w<Wo;++w)
      s += G[((size_t(n)*M+m)*Ho+h)*Wo+w];
  db[m]=s;
}
\end{lstlisting}
主机端验证 $K>0,H\ge K,W\ge K$，按各自 \code{total} 向上取整启动网格，并在同一
stream 中依次启动。每个输出地址只有一个线程写，所以没有竞态；若改成“每个 $G$ 线程
散布贡献”，则必须原子加或做两阶段归约。更新 $F,b$ 的内核必须等待 $dF,db$ 完成，跨
stream 时应用 event 建立依赖。

\editorialnote 英文底本附录的反向传播材料存在三组相互关联的错误：公式 B.22 使用未定义
$k$ 写成 $F[k-p,k-q]$；原书图 B.10 把 $w-q$ 误写为 $w-p$ 并重复循环；原书图 B.11 给
$G$ 多写了通道下标。这里按前向式直接求偏导，使用 $F[p,q]$、$G[m,h,w]$ 和唯一循环，
与中文译稿的勘误一致。

\complexity{$dX$ 内核实际执行 $O(NMCHWK^2)$ 次边界检查，成功乘加总数为 $NMCH_oW_oK^2$；$dF$ 为 $O(NMCH_oW_oK^2)$，$db$ 为 $O(NMH_oW_o)$}{$dX,dF,db$ 输出空间；每线程 $O(1)$ 累加状态}{$dX$ 有 $NCHW$ 个任务，$dF$ 有 $MCK^2$ 个任务；基础实现的单线程归约可能成为长关键路径}
\conclusionheading 两个公式分别枚举“哪些输出受该输入影响”和“该权重影响哪些输出”，
是代码索引与有限差分梯度检查的直接依据。
\discussionheading
\discussionpoint{链式法则可以系统地产生三个梯度。}
把前向写成互相关后，每个输出 $Y$ 依赖一个输入窗口、一个滤波器和一个偏置；上游梯度 $G$ 沿这些依赖反向传播。$dX$ 把 $G$ 乘滤波器权重散回所有覆盖该输入的位置，$dF$ 则对 batch 与输出空间累加 $XG$，$db$ 是相应输出通道上全部 $G$ 的和。前向若采用数学卷积而翻转核，下标规则也必须同步改变，不能仅凭“卷积反向”名称机械套式。自动微分系统保存的正是这些局部雅可比向量积规则，而不是构造巨大的显式雅可比矩阵，这让链式法则既可组合又能高效执行。

\discussionpoint{唯一写者与归约并行度存在张力。}
让一个线程负责一个 $dF$ 元素并顺序遍历所有 $N H_o W_o$ 项，可以避免原子写并固定累加顺序，但归约很长时并行度不足，而且不同线程会重复读取相同的 $X$ 和 $G$。split-K 或分块方案让多个块计算部分和，再用原子或第二阶段归约合并，吞吐更高却增加临时写流量、workspace 和不确定的浮点加法树。库会根据形状、内存预算和确定性要求选择不同策略。这个取舍与任何大归约相同：扩大写者数量能缩短关键路径，但必须为合并、同步和数值顺序支付代价。

\discussionpoint{解析参考与有限差分承担不同验证角色。}
小张量适合用 FP64 CPU 按公式逐项计算，能精确定位某个 padding 或索引偏移；大模型则可抽查参数，用中心差分 $[L(w+h)-L(w-h)]/(2h)$ 对比反向梯度。$h$ 太小会让两个损失之差被浮点舍入淹没，太大又引入 $O(h^2)$ 截断误差，因此应扫描几个数量级并看稳定区间。测试还要覆盖 stride、dilation、奇偶尺寸和边界像素，因为内部点正确并不能证明边界下标。两类检查互补：解析参考验证实现公式，有限差分验证整个前向与反向组合的链式关系。

\discussionpoint{训练步的最优算法还受确定性与激活内存约束。}
高性能卷积反向可能使用原子、动态算法选择或不同 split-K，末位结果会随加法树变化；可复现实验需要请求确定算法，并记录 workspace 上限，因为限制临时内存可能迫使库换路径。训练系统还会用梯度检查点少存激活、在反向时重算前向，以额外 FLOP 换显存；融合和通信重叠又会改变单个内核的时间价值。于是最快的 $dX$ 或 $dF$ 内核未必带来最快训练步。完整评价应同时观察每步时间、峰值内存、数值波动和分布式通信，而不能把局部算子排行榜当成系统结论。
\variantheading 若 $N=M=C=1,H=W=3,K=2$ 且 $X,F,G$ 全为 1，则 $dF$ 的四项均为 4、$db=4$，而 $dX=\begin{bmatrix}1&2&1\\2&4&2\\1&2&1\end{bmatrix}$；这是可逐项核验的边界梯度模式。
\practiceheading 在固定同一 NVIDIA GPU、连续 NCHW 和 FP32 计算的前提下，用 cuDNN backward-data/backward-filter 对照，并以 FP64 CPU 有限差分作梯度检查；Nsight Compute 比较 workspace 方案的 DRAM 流量和 achieved occupancy。
\sourcecheck{https://docs.nvidia.com/deeplearning/cudnn/backend/latest/}{NVIDIA cuDNN Backend 官方文档的卷积数据梯度与滤波器梯度操作；高置信度}
对抗检查：前向若改成数学卷积，或改变 padding、stride、dilation，梯度下标必须同步改变，
不能照搬本题的互相关式。
\end{solutionexercise}

\begin{solutionexercise}{4}
\problemheading 分析原书图 19.11 矩阵乘法内核的内存访问模式，讨论其全局内存访问
是否合并，以及共享内存访问是否可能发生存储体冲突。该内核使用
$16\times16$ 二维线程块；一个 warp 因 \code{threadIdx.x} 最快变化而跨两个
\code{ty} 行。每一阶段以
\code{F[Row*(C*K*K)+ph*16+tx]} 加载滤波器，以
\code{u=ph*16+ty,v=Col} 把概念展开矩阵 $B[u,v]$ 映射回 NCHW 输入 $X$，再计算
\code{Fds[ty][k]*Bds[k][tx]}。须分别分析 $F$、$B/X$、$Y$ 的全局访问和两个共享 tile
的加载/乘加访问，并指出非整 tile 边界的影响。

\approachheading CUDA 一维线程线性次序中 \code{threadIdx.x} 变化最快。图中块为
$16\times16$，所以一个 warp 通常覆盖两个 \code{ty} 行，每行 16 个连续 \code{tx}。

\answerheading 全局访问逐项如下。
\begin{itemize}
  \item 加载 $F$ 时，固定 \code{ty} 的 16 个线程读取同一矩阵行的连续 16 个
  \code{float}。一个 warp 发出两个连续半行区段；地址满足正常对齐时可由少量合并交易
  服务，但它不是“整个 warp 永远读取一个连续 32 元素区段”。行距 $CK^2$ 未对齐时可能
  跨更多交易。
  \item 加载概念矩阵 $B$ 时，固定 \code{ty} 即固定 $u$，相邻 \code{tx} 使 $v$ 连续。
  在同一个输出行内，映射到 $X$ 的列也连续，因而合并；当 $v$ 跨过 $W_o$ 的行边界时，
  $X$ 地址从一行末尾窗口跳到下一行起始窗口，跳距为 $K$ 个元素，warp 会被拆成多个连续
  区段。故结论是“分段合并”，不是对任意 $W_o$ 都完全连续。
  \item 写 $Y$ 与 $F$ 类似：每个 \code{ty} 半 warp 写连续 16 列；对齐和尾部决定交易数。
\end{itemize}

共享内存加载阶段，\code{Fds[ty][tx]} 和 \code{Bds[ty][tx]} 都按连续 \code{tx}
落入连续 bank。乘加阶段读取 \code{Fds[ty][k]}：同一半 warp 的线程读取同一地址，可由
共享内存广播服务；读取 \code{Bds[k][tx]}：16 个连续 bank，各值被两个不同 \code{ty}
的 lane 读取，属于相同地址的多播而非两个不同地址落同一 bank。因此在常见 32-bank、
4 字节 bank 宽配置下没有固有 bank conflict。改变图块宽度、元素大小、bank 模式或数组
步长后必须重新分析。

原书示例还有三个前提：$M$、$H_oW_o$ 和 $CK^2$ 必须与图块边界兼容；循环条件使用
整数除法会漏掉 $CK^2$ 的余数阶段；代码也未检查 \code{Row}/\code{Col}/\code{u}
边界。通用实现应使用向上取整的阶段数，并把越界共享内存槽填 0，在最终写回前检查
\code{Row<M}、\code{Col<H_oW_o}。这些是正确性条件，不只是性能细节。

\complexity{每个样本的总乘加为 $O(MCK^2H_oW_o)$}{$2\cdot16^2$ 个共享 float，加每线程一个累加器}{每块 256 线程；两次块屏障保护共享 tile 的写后读和下一阶段覆盖}
\conclusionheading $F$、$B/X$ 和 $Y$ 在常见位置都能形成连续子交易，共享访问依靠广播/
多播避免冲突；矩阵行边界和非整图块尾部是对“完全合并、无越界”结论的反例。
\discussionheading
\discussionpoint{im2col 用数据复制换取规则矩阵乘。}
卷积的每个输出都读取一个滑动窗口，直接实现时地址包含通道、核坐标和空间边界。im2col 把所有窗口显式展开成矩阵列，同一个输入像素会因出现在多个窗口而被复制，但后续计算变成规则 GEMM，可以复用成熟矩阵乘内核。隐式 GEMM 不真正写出展开矩阵，而是在加载时根据矩阵坐标反算 $X$ 地址，省下大量中间字节，却增加整数索引和边界谓词。两者代表经典的“先物化以简化计算”与“按需生成以节省内存”取舍，数据库查询、张量编译和稀疏算子中也反复出现这种选择。

\discussionpoint{tile 形状同时支配数据复用和驻留。}
把 $16\times16$ tile 扩到 $32\times32$，一个 warp 可以完整覆盖一行，边界块比例也可能降低；但线程数从 256 增到 1024，寄存器与共享内存总量上升，设备可能只能让更少块同时驻留。较小 tile 提供更多块来隐藏延迟，却增加阶段数、同步次数和半 warp 请求。矩阵的 $M,N,K$ 比例也会让长方形 tile 优于正方形，例如输出列很多而滤波器行较少时，应把资源更多分配给 $N$ 维。最佳 tile 是复用、同步、尾块和占用共同形成的离散优化问题，不能只根据“越大复用越多”作判断。

\discussionpoint{全局合并与共享 bank 必须分别证明。}
相邻 \code{tx} 让全局地址连续，只能说明搬入阶段可能合并；数据写入共享内存后，乘加若按列读取，多个 lane 可能访问同一 bank 的不同地址而产生冲突。padding、转置或 swizzle 可以改变 bank 映射，但也可能使向量加载对齐更复杂；若多个 lane 读取完全相同地址，则通常是广播而不是冲突。Tensor Core 片段还要求特定块布局和对齐，源码二维下标的直觉不再足够。正确分析应沿实际 warp 线性 lane 写出全局源地址和共享目的地址，再分别检查尾部谓词、交易与 bank，而不是用一个“连续”结论包办两层存储。

\discussionpoint{高性能 GEMM 是多级数据流水而非单一 tile。}
现代实现会把一个 CTA tile 再划成 warp tile 和指令 tile，让全局到共享的异步复制与前一 tile 的矩阵乘重叠，并以双缓冲避免消费者读到正在覆盖的数据。共享数据进入寄存器片段后由 warp 级 MMA 累加，每一层都有对齐、同步和资源约束；增加流水级数可隐藏延迟，却消耗更多共享内存。比较自写隐式 GEMM 与 cuDNN 或 CUTLASS 时，应使用真实形状并观察有效算术、实际 DRAM 字节、bank 冲突、驻留和尾块比例。峰值 TFLOP/s 只有在操作数形状适配且数据供给不断时才有意义，这把卷积访存题延伸到了完整的层次化性能模型。
\variantheading 若 tile 从 $16\times16$ 改为 $32\times32$，每个 warp 固定一个
\code{ty}；访问 $F$、$Y$ 时是 32 个连续 float，对齐时覆盖一个 128 字节区段。访问隐式
$B/X$ tile 只有在这 32 列未跨越 $W_o$ 的输出行边界时才连续；跨界时仍分成相隔 $K$
个元素的区段。块含 1024 线程、两块共享 tile 合计 8192 字节，设备还必须允许每块
1024 线程。
\practiceheading 假定 warp 为 32 lane、共享内存有 32 个 4 字节宽的 bank，且设备允许
每块 1024 个线程。在这样的 NVIDIA GPU 上以 CUTLASS 隐式 GEMM 或 cuDNN 作基线，并用
Nsight Compute 检查每请求扇区数、共享内存 bank 冲突和占用率。
\sourcecheck{https://docs.nvidia.com/cuda/cuda-c-best-practices-guide/}{CUDA C++ Best Practices Guide 对全局合并访问、共享内存 bank 和广播的官方说明；对抗检查：只观察源码中 tx 连续而忽略 warp 跨两行、隐式展开跨输出行和尾部越界，会得出过强结论}
\end{solutionexercise}

\chapter{大语言模型}

\begin{solutionexercise}{1}
\problemheading 为原书图 20.9 中的 FlashAttention 内核补充主机端调用代码。设置输入
和输出数组的大小时，参考原书图 20.8 中所示矩阵的维度：本教学内核计算一个注意力头，
$Q,K,V,O$ 均为 $N\times d$，softmax 分母输出 $D$ 长度为 $N$；固定
$d=128$、$B_r=B_c=32$、每块 512 个线程，并要求 $N$ 是 32 的倍数。主机代码须完成
尺寸检查、设备属性/共享内存检查、分配与复制、内核启动、同步、错误检查、结果回传和释放。

\approachheading 该教学内核固定 $d=128$、$B_r=B_c=32$、每块 512 线程。$Q,K,V,O$
均为 $N\times d$，$D$ 长度为 $N$；外层 grid-stride 循环允许网格块数小于
$T_r=N/B_r$。

\answerheading 下面函数处理一个注意力头，假定主机输入已初始化：
\begin{lstlisting}[language=C++]
void run_flashattention(const float* hQ, const float* hK,
                        const float* hV, int N,
                        std::vector<float>& hO,
                        std::vector<float>& hD) {
  if (N <= 0 || N % B_r != 0 || N % B_c != 0)
    throw std::invalid_argument("teaching kernel requires N % 32 == 0");
  static_assert(d == 128 && BLOCK_SIZE == 512, "kernel constants");

  size_t elements = size_t(N) * d;
  if (elements > SIZE_MAX / sizeof(float))
    throw std::overflow_error("attention allocation");
  size_t matrixBytes = elements * sizeof(float);
  size_t vectorBytes = size_t(N) * sizeof(float);
  float *dQ=nullptr, *dK=nullptr, *dV=nullptr;
  float *dO=nullptr, *dD=nullptr;
  CHECK(cudaMalloc(&dQ, matrixBytes));
  CHECK(cudaMalloc(&dK, matrixBytes));
  CHECK(cudaMalloc(&dV, matrixBytes));
  CHECK(cudaMalloc(&dO, matrixBytes));
  CHECK(cudaMalloc(&dD, vectorBytes));
  CHECK(cudaMemcpy(dQ, hQ, matrixBytes, cudaMemcpyHostToDevice));
  CHECK(cudaMemcpy(dK, hK, matrixBytes, cudaMemcpyHostToDevice));
  CHECK(cudaMemcpy(dV, hV, matrixBytes, cudaMemcpyHostToDevice));

  cudaDeviceProp prop{};
  int device=0, activePerSM=0;
  CHECK(cudaGetDevice(&device));
  CHECK(cudaGetDeviceProperties(&prop, device));
  CHECK(cudaOccupancyMaxActiveBlocksPerMultiprocessor(
      &activePerSM, flashattention_forward_kernel, BLOCK_SIZE, 0));
  int tiles = N / B_r;
  int blocks = std::min(tiles, activePerSM * prop.multiProcessorCount);
  if (blocks <= 0) throw std::runtime_error("kernel cannot be resident");
  float scaling = 1.0f / sqrtf(float(d));
  flashattention_forward_kernel<<<blocks, BLOCK_SIZE, 0, 0>>>(
      dQ, dK, dV, N, scaling, dD, dO);
  CHECK(cudaGetLastError());
  CHECK(cudaDeviceSynchronize());

  hO.resize(elements); hD.resize(N);
  CHECK(cudaMemcpy(hO.data(), dO, matrixBytes, cudaMemcpyDeviceToHost));
  CHECK(cudaMemcpy(hD.data(), dD, vectorBytes, cudaMemcpyDeviceToHost));
  CHECK(cudaFree(dD)); CHECK(cudaFree(dO)); CHECK(cudaFree(dV));
  CHECK(cudaFree(dK)); CHECK(cudaFree(dQ));
}
\end{lstlisting}
内核不读取输出初值，所以无需预清零 \code{dO} 和 \code{dD}。
其寄存器状态必须由第 3 题的 \code{initialize()} 清零。同步既检查运行期错误，
也建立输出复制前的生命周期完成点。
异常安全的工程版本应以 RAII 管理五个设备指针。

代码只覆盖原书教学前提。支持任意 $N$ 时，$T_r,T_c$ 必须向上取整，并在加载、掩码、
写回中处理尾部；支持 batch/multi-head 时还要给指针增加批次与头偏移，或使用网格的其他
维度。静态共享内存和 512 线程还必须满足目标设备限制。

\complexity{精确注意力为 $O(N^2d)$ 算术；分块减少的是 HBM 流量而非渐近算术量}{$Q,K,V,O$ 各 $Nd$，$D$ 为 $N$；每块另用约 37 KiB 静态共享内存及寄存器 tile}{最多 $T_r=N/32$ 个独立行 tile；grid-stride 循环在块数受驻留限制时继续覆盖全部 tile}
\conclusionheading 正确尺寸为三个 $N\times128$ 输入、一个 $N\times128$ 输出和一个
$N$ 分母向量，缩放因子为 $1/\sqrt{128}$。
\discussionheading
\discussionpoint{FlashAttention 的关键是重排 I/O 而非近似 softmax。}
朴素注意力先形成 $N\times N$ 分数矩阵，再做 softmax 和乘 $V$，中间矩阵会在 HBM 中反复读写；当 $N$ 很大时，这些字节常比算术更昂贵。分块算法让 $Q$ 的一组行与 $K,V$ tile 在片上相遇，并用在线归一化立即累计输出，因此不必把完整分数矩阵物化。它仍计算精确注意力公式，只因浮点运算重排产生允许范围内的舍入差异，渐近算术仍是 $O(N^2d)$。所以性能分析不能只数 FLOP，还要问每个 tile 被复用几次、避免了多少 HBM 往返，以及额外重算是否比保存中间量更划算。

\discussionpoint{在线 softmax 保存的是可合并的行统计量。}
处理一段 logits 后，只需保留最大值 $m$、缩放后的指数和 $D=\sum_j e^{s_j-m}$，以及同一尺度下的加权输出 $O=\sum_j e^{s_j-m}V_j$。两个有效块的状态应对称合并：令 $m=\max(m_1,m_2)$，则 $D=e^{m_1-m}D_1+e^{m_2-m}D_2$，且 $O=e^{m_1-m}O_1+e^{m_2-m}O_2$；这既覆盖新最大值变大的情形，也覆盖 $m_2<m_1$ 时只缩放新块的情形。最终 $O/D$ 与一次性 softmax 等价。若两块都全掩码而使 $m_1=m_2=-\infty$，上述差值会出现 $-\infty-(-\infty)$，必须由单独的无有效元素分支处理，不能依赖公式自行得到零。这个可合并状态也适用于流式 log-sum-exp 和分布式归约，其价值在于把完整历史压缩为充分统计量。

\discussionpoint{tile 尺寸由线程、共享内存和寄存器共同约束。}
教学配置用 512 线程、$B_r=B_c=32$、$d=128$，共享的 $K^T$、$V$ 和分数 tile 合计约 37~KiB；这换来片上复用，也可能让每个 SM 同时驻留的块减少。扩大 tile 可降低外层循环和加载次数，却会增长共享内存、寄存器状态与同步成本；缩小 tile 增加并行块，却让同一数据更频繁搬运。生产内核还会针对 FP16/BF16、Tensor Core、异步复制、dropout 和掩码选择不同形状，并为非整 tile 尾部生成谓词路径。因此教材常数只是一个可审查实例，不能直接当成所有序列长度和架构的最佳配置。

\discussionpoint{训练预填充与逐 token 解码已经形成不同内核族。}
训练或长提示的 prefill 同时拥有许多 query 行，矩阵乘规模大，重点是前后向总 I/O 和 Tensor Core 利用；自回归解码每步 query 很少，却要读取不断增长的 KV cache，瓶颈更接近容量与不连续访问。paged attention 用页表管理 KV 块，grouped-query attention 让多个 query 头共享较少的 KV 头，分别处理内存碎片与容量压力。评价时应按阶段和序列长度报告输出误差、HBM 字节、有效计算率和端到端 token/s，而不是用一个峰值代表全部场景。这个分化说明算法公式相同，并不意味着部署工作负载拥有同一性能结构。
\variantheading 若 $N=96$，则有 3 个行 tile；每个 $N\times128$ 矩阵含 12,288 个 float、占 49,152 字节，$D$ 占 384 字节，缩放仍为 $1/\sqrt{128}\approx0.0883883$。
\practiceheading 假定 NVIDIA GPU 支持每块 512 线程、至少 37 KiB 共享内存且输入为 FP32、$N$ 为 32 的倍数，用 FlashAttention 官方实现或 PyTorch SDPA 校验输出；以 Nsight Compute 比较 HBM 字节、occupancy 和数值误差。
\sourcecheck{https://arxiv.org/abs/2205.14135}{FlashAttention 原始论文；对抗检查：该算法是精确重排但并不自动支持任意尾部，忽略 $N$ 的分块边界会漏算而不只是变慢}
\end{solutionexercise}

\begin{solutionexercise}{2}
\problemheading 阅读 CUB 文档，了解 \code{WarpReduce} 命名空间中数据类型的细节，
以及调用 \code{Reduce} 函数时各参数的语义。特别注意归约操作所支持的各种算术函数。
回答须说明模板数据类型和逻辑 warp 宽度、每个逻辑 warp 的临时存储、每线程输入、二元
归约操作、结果在哪些 lane 有效，以及和、最小值、最大值等操作的类型/结合性要求。

\approachheading 以当前 CCCL/CUB 接口为准，区分类模板、临时存储、每线程输入、二元
操作和结果有效 lane。

\answerheading 常见声明为
\begin{lstlisting}[language=C++]
using WR = cub::WarpReduce<T, LOGICAL_WARP_THREADS>;
__shared__ typename WR::TempStorage temp[WARPS_PER_BLOCK];
T aggregate = WR(temp[warp_id]).Reduce(thread_value, binary_op);
\end{lstlisting}
其中：
\begin{itemize}
  \item \code{T} 是每线程贡献及返回聚合值的类型；\code{LOGICAL_WARP_THREADS} 默认是
  硬件 warp 大小，也可把一个物理 warp 切成多个逻辑 warp。\code{TempStorage} 是该
  collective 的共享临时类型，每个同时工作的逻辑 warp 必须有不重叠实例。
  \item \code{thread_value} 是调用 lane 的输入，\code{binary_op} 是可调用的二元归约
  运算。\code{Reduce} 的完整聚合只保证在每个逻辑 warp 的 lane 0 有效；若所有 lane
  都需要结果，应再用 shuffle 从 lane 0 广播。
  \item 带 \code{valid_items} 的重载只纳入前若干 lane。参与的线程必须以一致控制流调用；
  临时存储在下一次复用前应遵守文档的 warp 同步要求。
\end{itemize}

\code{Reduce} 并不限于加法；它接受满足归约要求的自定义二元 functor。CUB 提供的常用
操作包括 \code{cub::Sum()}、\code{cub::Min()}、\code{cub::Max()}，以及适合键值对的
arg-min/arg-max 等操作；\code{Sum()} 还有专用便利成员。浮点加法虽然数学上结合，但在有限
精度中不严格结合，所以树形归约的结果可能与串行和有若干 ULP 差异。

\editorialnote 原书图 20.12 和图 20.15 各只声明一个共享
\code{WarpReduce::TempStorage}，而原书图 20.9 的线程块有 16 个 warp，它们会并发调用
归约。这会让多个 warp 竞态复用同一临时区。可执行实现应声明
\code{temp_store[N_WARPS]} 并以 \code{warpIdx()} 选择；图 20.4 使用的是整块唯一的
\code{BlockReduce}，不受这一问题影响。

\conclusionheading \code{WarpReduce} 是逻辑 warp collective，不是普通标量函数；
临时存储所有权、参与线程集合和 lane 0 结果语义都是正确调用的一部分。
\discussionheading
\discussionpoint{WarpReduce 是有参与契约的 collective。}
调用 \code{Reduce} 不是 32 次互不相关的标量函数，而是一组 lane 共同完成一次归约；每个 lane 提供输入，库在它们之间交换数据。完整结果通常只保证在逻辑 warp 的 lane 0 有效，其余 lane 的返回值不能未经说明直接使用，若大家都需要就应显式广播。逻辑宽度还可以是 8 或 16，使一个物理 warp 包含多个独立小组。正确使用因此要同时说明输入类型、逻辑宽度、参与 lane、结果所有者和广播步骤；CUB 封装了通信实现，却没有取消调用者对集体操作边界的责任。

\discussionpoint{临时存储的所有权必须与并发实例一一对应。}
原书教学块有 16 个 warp，若它们同时把中间值写进同一个 \code{TempStorage}，不同 collective 的生命周期重叠，就会形成共享内存竞态。最直接的设计是声明 \code{temp[WARPS_PER_BLOCK]}，用 warp 编号选择独立元素；逻辑 warp 小于 32 时还要按实际小组数分配。同一组连续两次复用临时区也应遵守 API 的同步要求，不能把硬件常见的锁步执行当成无限期可见性保证。这个问题与线程池复用工作缓冲相同：库对象的类型正确不代表并发所有权正确，资源实例必须覆盖所有同时活动的操作。

\discussionpoint{参与掩码定义了哪些值真正存在。}
手写 shuffle 归约常在尾部或分支中出错：某个 lane 已退出，却仍被掩码列为来源，其他 lane 便可能读取未定义值。逻辑 warp 设为 8 时，物理 lane 0、8、16、24 分别是四组的根，不能把物理 lane 0 当成整个 warp 唯一结果位置。CUB 可以表达这些子组，但参与线程仍须以一致控制流进入 collective，并传入与有效数据一致的数量或掩码。这个边界也出现在 ballot、scan 和协作组中，说明 SIMD/SIMT 原语的正确性不仅取决于算术，还取决于“谁同时参加”这一控制流集合。

\discussionpoint{归约操作需要代数性质与数值语义。}
树形归约假定二元操作至少具有结合性，才能在改变括号位置后保持数学结果；最大值和整数加法通常满足，字符串减法之类的非结合操作则会随树形得到不同语义。浮点加法在实数中结合，在有限精度中却不严格结合，所以结果可能与串行顺序相差若干 ULP，即使 collective 没有任何竞态。自定义 functor 还应处理单位元、NaN 和有符号零等边界。性能上可以比较 shuffle、共享访问和同步，正确性上则要用不同逻辑宽度与输入分布检查结果，这把 API 使用问题连接到了并行代数结构。
\variantheading 对 \code{WarpReduce<int,8>}，一个 32-lane 物理 warp 被分成 4 组；若物理 lane $i$ 输入 $i+1$ 并求和，则物理 lane $0,8,16,24$ 的有效聚合依次为 $36,100,164,228$。
\practiceheading 在硬件 warp 为 32、CUDA 计算能力 7.0 或更新且逻辑 warp 设为 8 的 NVIDIA GPU 上，写 CUB 微测试并用 Compute Sanitizer racecheck 排除 TempStorage 复用；Nsight Compute 再比较 shuffle/shared 指令数。
\sourcecheck{https://github.com/NVIDIA/cccl/tree/main/cub}{NVIDIA CCCL/CUB 官方源代码与文档；对抗检查：若所有 warp 共用一个 TempStorage，或未经广播就在任意 lane 使用 Reduce 返回值，教学公式正确也会被实现竞态破坏}
\end{solutionexercise}

\begin{solutionexercise}{3}
\problemheading 编写原书图 20.9 第 21 行所调用的 \code{initialize()} device
function，并解释为什么必须把 \code{O_i} 和 \code{D_i} 的元素初始化为
\code{0.f}。还应说明为什么 \code{m_i} 需要初始化为负无穷。这里每个 warp 的寄存器
状态包括 $O_i[B_{r,\mathrm{warp}}][d_{\mathrm{lane}}]$、
$D_i[B_{r,\mathrm{warp}}]$ 和 $m_i[B_{r,\mathrm{warp}}]$，后续按在线 softmax 公式逐个
$K/V$ tile 合并。

\editorialnote 所核英文第 5 版底本只要求解释 $O_i$ 与 $D_i$ 为何初始化为 0；中文译稿
另增加了 $m_i=-\infty$ 的问项。本题解保留译稿增补，并将三种状态的单位元一并说明。

\approachheading 三个状态分别参加加法累积、分母累积和最大值累积，初始化应使用各自
运算的单位元。

\answerheading
\begin{lstlisting}[language=C++]
__device__ inline void initialize(
    float O_i[B_r_warp][d_size],
    float D_i[B_r_warp],
    float m_i[B_r_warp]) {
  #pragma unroll
  for (int ii = 0; ii < B_r_warp; ++ii) {
    D_i[ii] = 0.0f;
    m_i[ii] = -CUDART_INF_F;
    #pragma unroll
    for (int dd = 0; dd < d_size; ++dd)
      O_i[ii][dd] = 0.0f;
  }
}
\end{lstlisting}
$O_i=0$ 使第一个 $PV$ tile 从零部分和开始；否则未初始化寄存器会把任意位型带入最终输出。
$D_i=0$ 是指数和的加法单位元，且在首次最大值更新时
$0\cdot\exp(-\infty-m)=0$。$m_i=-\infty$ 是最大值运算的单位元，保证第一个包含合法 key
的 tile 一定更新最大值。若误设为 0 且该行所有 logit 都为负，$m$ 不再等于已处理元素的
真实最大值，破坏 online softmax 的状态不变量；在精确实数算术且无下溢时，分子分母仍可能
因共同缩放而得到相同归一化结果，但极负 logit 会更早下溢，使分母为 0 或丢失贡献，因而
不能把 0 当作数值稳定的初始化。

这些寄存器数组是每线程私有的，无需同步；共享 $K,V,S$ tile 仍需要块屏障。若通用掩码允许
某一整行没有任何合法 key，不能只期待状态自然保持 $D=0,m=-\infty$：当旧最大值和当前 tile
最大值都是 $-\infty$ 时，缩放项会计算 $\exp[-\infty-(-\infty)]$，其中差值已经是 NaN，
分母和输出累加器也可能被污染。生产实现必须维护“已有有效项”状态或有效计数，完全无效的
tile 直接跳过，第一批有效项则直接建立最大值和累加器；最终再用 $D=0$ 检查作防御性收尾。
因果掩码的每一行至少允许对角元素，所以原书场景不会触发该边界。

\conclusionheading 0、0、$-\infty$ 分别是输出和、softmax 分母和最大值的正确单位元。
\discussionheading
\discussionpoint{负无穷是最大值递推的自然单位元。}
初始化状态要表达“尚未见过任何有效 logit”，而不是假装已经见过数值 0。对任意有限 $x$，$\max(-\infty,x)=x$，且 $e^{-\infty}=0$，所以第一项能自然建立最大值和指数和，无需另设首项分支。若把 $m$ 初始化为 0，而真实 logits 全为负，归一化在精确算术中有时仍因共同尺度抵消而得到相同比例，但极负值更容易提前下溢，使分母丢失贡献。负无穷作为 max-plus 代数的单位元也用于 Viterbi、对数概率和 masked reduction，体现了用正确代数身份简化边界状态的普遍方法。

\discussionpoint{全掩码 tile 会在最终除法之前制造 NaN。}
若旧最大值与当前 tile 最大值都是 $-\infty$，直接计算缩放因子会遇到 $-\infty-(-\infty)=\mathrm{NaN}$，随后即使旧分母是 0，$0\times\mathrm{NaN}$ 仍是 NaN。因此不能只在最终写回时检查 \code{D==0}，因为 $D$ 可能早已被污染；实现应让全无效 tile 跳过递推，或维护有效计数/谓词，只对至少一个有限元素的集合更新状态。整行最终无有效 key 时，接口还要明确定义返回零、NaN 或错误。常规 causal mask 保留对角项，原书路径不会触发整行全掩码，但通用 attention mask 不能借用这一假设。

\discussionpoint{有限的大负常数不能完全替代掩码语义。}
用 $-10^4$ 之类常数代替 $-\infty$ 看似避免了非法运算，但该值在不同精度和缩放后可能成为有限数或舍入成无穷，语义并不稳定。若真实 logit 也达到类似量级，被屏蔽项甚至可能获得非零权重；当一整行都填同一有限常数时，softmax 还会产生近似均匀分布，而不是“没有有效项”。更稳健的设计是让有效谓词参与 max 与 sum，并在无有效项时采用接口约定的收尾路径。这个区别提醒我们，数值哨兵只是编码手段，不能取代集合成员关系和异常状态的显式表示。

\discussionpoint{减最大值只是稳定递推的第一道防线。}
把 logits 减去当前最大值可避免指数上溢，但 $D$ 的长序列累计、旧输出的反复缩放和最终除法仍会舍入，低精度下小贡献也可能下溢。实践中常让 Q/K/V 使用 FP16 或 BF16，而用 FP32 保存最大值、分母和输出累加器，再在最后转换。测试应覆盖极大正负值、单个有效 key、跨 tile 最大值突增、全掩码、不同 tile 划分和长序列，并与高精度朴素参考比较。只有在这些路径下都满足误差约定，才能说明“数值稳定”；一个普通随机样例相等只证明了最容易的分支。
\variantheading 若允许某查询行屏蔽全部 key，应为每行维护 \code{hasValid} 或有效计数：空 tile 不执行 max/exp 递推，第一项有效 score 令 $m=s,D=1,O=V$；最终若仍无有效项，则按接口约定写 $O=0,D=0,m=-\infty$ 并跳过除法。只在写回前检测 $D_i=0$ 不能修复此前已经传播的 NaN。
\practiceheading 假定 IEEE FP32 \code{expf}、NVIDIA GPU 且常规因果掩码至少保留对角元素，用 FlashAttention 或 PyTorch \code{scaled_dot_product_attention} 对照随机行；随附 CPU 回归还构造全掩码行和单一有效 key，验证空 tile 跳过、零向量约定及首项初始化，再用 Compute Sanitizer 检查越界而非掩盖 NaN。
\sourcecheck{https://arxiv.org/abs/1805.02867}{Online normalizer calculation for softmax 原始论文；对抗检查：把 $m$ 初始化为最小正数或 0 都不是 max 的单位元，全掩码行还需单独处理零分母}
\end{solutionexercise}

\begin{solutionexercise}{4}
\problemheading 证明原书图 20.11 的 device function 中嵌套循环对 $K$ 和 $V$ 的
所有加载都是合并访问。代码对每个 $jj=0,\ldots,B_c-1$ 执行
\code{for (dd=threadIdx.x; dd<d; dd+=blockDim.x)}，并从两数组的同一线性下标
\[
  (B_cj+jj)d+dd
\]
分别加载到共享内存 $K_j^{\mathsf T}[dd,jj]$ 和 $V_j[jj,dd]$。证明须按 warp/lane
写出每轮全局地址，处理 $d=128$、\code{blockDim.x=512} 时的活动 warp，并区分“合并
访问”与“单个对齐内存事务”。

\approachheading 固定外层 \code{jj}，考察同一 warp 中相邻 lane 的 \code{dd} 和全局
线性地址。块是一维 512 线程，\code{threadIdx.x} 就是线性线程号。

\answerheading 对任意固定 $j,jj$，两个全局地址完全相同，只是基指针不同：
\[
 a_K(t)=a_V(t)=(B_cj+jj)d+t,
\]
其中 $t=\code{threadIdx.x}$。内层每次迭代把 $t$ 增加统一的
\code{blockDim.x}，所以某 warp 的 lane $\ell=0,\ldots,31$ 在任一迭代读取
\[
  A+\ell,\qquad A=(B_cj+jj)d+w\cdot32+k\cdot\code{blockDim.x},
\]
即 32 个连续 \code{float}。因此 $K$ 和 $V$ 都形成连续的 128 字节 warp 请求；对齐时可由
最少的内存交易服务，不对齐时至多拆成相邻交易，但仍是合并模式。

在原书常数 $d=128$、\code{blockDim.x=512} 下，每个 \code{jj} 只有线程 0--127 活跃，
恰好四个完整 warp，各自加载连续 32 个元素；其余 warp 统一不进入内层循环。\code{cudaMalloc}
提供充分的基址对齐，而每行 $128\times4=512$ 字节，所以每行也保持对齐。若 $d>512$，
下一次循环只是让每个 warp 跳到同一行中再后移 512 个元素，lane 间连续性不变。

写共享内存时，$V_j[jj][dd]$ 按行连续；$K$ 则通过
\code{addr(dd*B_c+jj)} 转置并加 padding。这个不连续的是\emph{共享内存目的地址}，不会
改变全局源地址的合并性。若 $d$ 不是 warp 倍数，最后 warp 会部分活跃；若行基址未对齐，
交易数可能增加，但不存在 lane 间大步长访问。

\complexity{每个 $K/V$ tile 各加载 $B_cd$ 个元素，合计 $O(B_cd)$}{共享 tile 各 $B_cd$，$K^T$ 另有 bank padding}{块内线程共同搬运；每个活跃 warp 在固定 jj 时访问连续区段}
\conclusionheading 合并性的关键是全局右端下标 \code{+dd}，而不是共享端的转置下标。
\discussionheading
\discussionpoint{合并访问只描述请求形状，不承诺高带宽。}
同一 warp 的 lane 若读取连续且适当对齐的 FP32 地址，硬件可以用较少 sector 服务请求，这证明了交易效率；它没有说明设备同时有多少请求在途。教学配置每个 \code{jj} 只有四个 warp 负责 $d=128$ 的搬运，其余 warp 在该循环中空闲，若网格又很小，延迟仍可能无法隐藏。加载后的 tile 是否被多次复用也决定这些字节是否值得。因而“已合并”只是必要的局部地址性质，完整吞吐还受活跃 warp、并发块、缓存命中和消费算术约束，不能把一个地址证明直接升级成性能结论。

\discussionpoint{二维线性化会让行距参与对齐。}
$d=128$ 时每行是 512 字节，四个完整 warp 恰好覆盖一行，且下一行仍保持常见的 128 字节边界相位。若 $d=130$，第五个 warp 只有两个 lane 活跃，行距变为 520 字节；下一行基址相对 128 字节边界移动 8 字节，后续完整 warp 请求可能跨更多 sector。边界谓词足以保证正确性，却不能恢复对齐与活跃 lane 效率。可以为物理行加 padding、按向量宽度分块或重新安排线程，但都要区分逻辑 $d$ 与实际 stride，这与图像 pitch 和张量 leading dimension 是同一类布局问题。

\discussionpoint{异步 copy 把加载和计算组织成生产消费流水。}
同步实现让所有 warp 搬完当前 tile、经过屏障后再开始计算，加载等待与算术没有重叠。更深的实现可让 producer warp 用异步 copy 填充下一缓冲，consumer warp 同时处理上一缓冲，并以双缓冲和阶段标记防止过早覆盖。这样可能隐藏全局延迟，却会增加共享内存、屏障状态和生命周期证明，尾 tile 还要用谓词或零填充保持消费者看到合法值。流水级数过多会降低驻留，producer 分工不当也会让计算 warp 饥饿；它是资源交换，不是打开某条指令就自动获得的重叠。

\discussionpoint{计数器验证应建立在地址推导之后。}
先按 lane、元素宽度和行距推导期望的连续区段与 sector 数，再用 profiler 比较实际 sectors/request、活跃 lane 和 stall，才能判断偏差来自对齐、尾部还是并发不足。只看 Global Load Efficiency 可能遗漏“请求很高效但数量太少”，只看 DRAM 吞吐又可能把无用流量误作成功。对注意力还应把加载字节除以 tile 的片上复用次数，并同时观察共享内存和指数计算瓶颈。理论地址分析给出可证预期，硬件计数器用于验证实现是否遵守预期；两者结合比盲目扫描指标更有解释力。
\variantheading 若 $d=130$、仍用 512 线程，则每行由 4 个完整 warp 加 1 个仅 2 lane 活跃的 warp 加载；行距为 520 字节，使相邻行基址相对 128 字节边界偏移 8 字节，后续完整请求可能拆成两个交易。
\practiceheading 假定 32-lane warp、4 字节 FP32 和 128 字节对齐的 \code{cudaMalloc} 基址，在 $d=128,130,160$ 上用 Nsight Compute 记录 Global Load Efficiency、sectors/request 与活跃 lane；再与 FlashAttention/CUTLASS 的异步 copy 路径比较。
\sourcecheck{https://docs.nvidia.com/cuda/cuda-c-programming-guide/}{CUDA C++ Programming Guide 的全局内存合并规则；对抗检查：共享目的地址的 padding 不能用来证明全局合并，证明必须直接比较同一 warp 的全局源地址}
\end{solutionexercise}

\chapter{静电势图}

\begin{solutionexercise}{1}
\problemheading 为原书图 21.6 中的聚集式 DCS 内核补全主机代码，配置线程网格并使用
全部执行配置参数调用内核。\code{cenergy} 依次接收势网格指针 \code{energygrid}、
逻辑尺寸 \code{grid}、网格间距 \code{gridspacing}、切片坐标 $z$ 和原子数
\code{numatoms}；二维线程坐标 $(i,j)$ 对应当前 $z$ 切片中的一个网格点，原子以
$(x_a,y_a,z_a,q_a)$ 四元组分块存入常量数组，每个线程累加
$\sum_a q_a/\sqrt{(x-x_a)^2+(y-y_a)^2+(z-z_a)^2}$。主机代码须给出
启动配置中的网格维度、块维度、动态共享内存字节数和 stream 四项参数，并完成常量内存
分块复制、边界检查、同步/错误检查以及输出缓冲区的生命周期管理。

\approachheading 每个线程对应切片中的一个 $(i,j)$ 网格点。原子数据按
\code{CHUNK_SIZE} 分块复制到常量符号；常量复制和内核放入同一 stream，保证下一块覆盖
符号以前，上一内核已经读完。

\answerheading 原书内核在计算 $i,j$ 后必须补上边界保护：
\begin{lstlisting}[language=C++]
if (i >= grid.x || j >= grid.y) return;
\end{lstlisting}
否则向上取整的启动网格会越界。设三个逻辑维度为 $G_x,G_y,G_z$；对已经分配为至少
$G_xG_yG_z$ 个 float 的 \code{dEnergy}，单切片主机函数为：
\begin{lstlisting}[language=C++]
void launch_cenergy_slice(float* dEnergy, dim3 logical,
                          float spacing, unsigned slice,
                          const float* hAtoms, int numAtoms,
                          cudaStream_t stream) {
  if (!dEnergy || logical.x == 0 || logical.y == 0 || logical.z == 0 ||
      spacing <= 0.0f || slice >= logical.z ||
      numAtoms < 0 || (numAtoms && !hAtoms))
    throw std::invalid_argument("cenergy arguments");

  dim3 threads(32, 8, 1); // 256 threads, x is the contiguous dimension
  dim3 blocks((logical.x + threads.x - 1) / threads.x,
              (logical.y + threads.y - 1) / threads.y, 1);
  size_t dynamicSmemBytes = 0;
  float z = float(slice) * spacing;

  size_t plane = size_t(logical.x) * logical.y;
  size_t firstEnergy = size_t(slice) * plane;
  CHECK(cudaMemsetAsync(dEnergy + firstEnergy, 0,
                        plane * sizeof(float), stream));

  for (int first = 0; first < numAtoms; first += CHUNK_SIZE) {
    int count = std::min(CHUNK_SIZE, numAtoms - first);
    CHECK(cudaMemcpyToSymbolAsync(
        atoms, hAtoms + size_t(first)*4,
        size_t(count)*4*sizeof(float), 0,
        cudaMemcpyHostToDevice, stream));
    cenergy<<<blocks, threads, dynamicSmemBytes, stream>>>(
        dEnergy, logical, spacing, z, count);
    CHECK(cudaGetLastError());
  }
}
\end{lstlisting}
\code{cudaMemsetAsync} 先把本切片置零；它与后续复制、内核同处一个 stream，所以所有
原子块的 \code{+=} 都累加到确定的零初值。若 \code{numAtoms==0}，循环不执行，但排队的
清零仍使整张切片得到正确的零势能，而不是遗留旧内容。

四项执行配置依次是 \code{blocks}、\code{threads}、动态共享内存字节数 0 和
\code{stream}。同一 stream 的有序语义使“复制第 $c$ 块、执行第 $c$ 块、复制第
$c+1$ 块”保持顺序；不能在多个 stream 中无协调地覆盖同一个 \code{__constant__ atoms}
符号。调用者在读取结果或释放输入前应同步 stream/event，并检查异步执行错误。

原书用浮点式 \code{k=z/gridspacing} 恢复切片下标；主机按
\code{z=slice*spacing} 构造可减小不一致风险，但更稳健的接口应把整数 \code{slice}
直接传给内核。$i,j,k$ 及线性乘积也应使用足够宽的无符号类型，防止大网格溢出。

题设模型严格是 $q/r$，没有 $\sqrt{r^2+\varepsilon}$ 软化项。若网格点恰与电荷位置重合且
$q\ne0$，结果按该模型发散为带符号无穷；测试数据应显式避开或专门验证这个奇点。工程上
可以把 $\varepsilon$ 做成用户可见参数以正则化，但这会改变物理模型，不能在 CPU 与 GPU
参考中同时静默加入。

\complexity{每个切片 $O(G_xG_yA)$，$A$ 为原子数}{$O(G_xG_yG_z)$ 势网格和每块最多 $4\,\code{CHUNK_SIZE}$ 个常量 float}{每次原子块启动 $\lceil G_x/32\rceil\lceil G_y/8\rceil$ 个线程块；每线程独占一个网格点，无写竞态}
\conclusionheading 二维向上取整要求内核边界判断；同一常量符号的复制、使用和覆盖必须由
stream 顺序管理。
\discussionheading
\discussionpoint{常量内存的优势来自 warp 广播。}
当一个 warp 的所有 lane 在循环第 $n$ 次都读取同一个原子坐标时，常量缓存可以把这次地址请求广播给整组线程，数据只需取一次。若各 lane 同时访问不同地址，请求会按不同地址拆分，吞吐随地址种类增加而下降；“只读”本身并不保证高效。DCS 让相邻线程处理不同网格点，却以相同原子顺序循环，正好构成广播模式。分析常量内存时因此要写出某一时刻各 lane 的地址，而不能只看到变量声明中的 \code{const}；这种统一控制流与共享参数的模式也常见于小型卷积核和查表系数。

\discussionpoint{原子分块把超容量求和变成流式归约。}
常量空间有限时，可以复制一批原子，让所有网格点把这批贡献累加到输出，再复制下一批；数学上只是把 $\sum_{a=0}^{A-1}$ 分成若干连续部分和。输出必须在第一批前清零，并在批次间保留已有值，否则每个 tile 都会覆盖前面的贡献。批次太小会增加复制和内核启动次数，太大又可能超出容量或降低缓存复用，因而存在传输固定开销与计算粒度的平衡。这个过程类似流式数据库聚合：无法一次容纳全部输入时，只要累加状态满足结合语义，就能分块处理并得到同一目标结果。跨批次误差还取决于部分和累加次序，应同时与高精度参考比较。

\discussionpoint{常量符号是设备级共享状态而非 stream 私有副本。}
把复制和内核放进同一 stream，只能保证这个 stream 内“先复制、后读取”的顺序；另一个 stream 仍可能在前一内核结束前覆盖同一个常量符号。此时两个批次读取的原子集合混杂，问题不是缓存一致性慢，而是数据生命周期发生真实竞态。需要并发处理时，可以为每项工作分配独立只读全局缓冲，通过事件串行化符号更新，或在不同设备上使用各自副本。全局常量符号使用方便，却把所有调用者绑定到同一可变资源，这与复用全局配置表或静态工作区时的可重入性问题完全相同。

\discussionpoint{大规模静电计算必须利用空间与尺度结构。}
让 $P$ 个网格点遍历 $A$ 个原子需要 $O(PA)$ 工作，常量缓存只能降低取数成本，无法改变这个渐近量。实际系统会用空间分箱和邻居表只枚举近场原子，再用多极方法、FFT 网格法或 Ewald 类分解处理远场，使工作接近线性或 $N\log N$。这些方法引入截断、插值和全局通信误差，常量分块仍可作为小规模精确参考或近场 tile。评价方案应同时观察原子与点的空间分布、每点误差和复制计算时间，因为一个缓存优化不能代替算法层面对相互作用范围的重新组织。若电荷分布随时间变化，邻居结构的重建成本也必须摊入总账。
\variantheading 若逻辑切片为 $65\times17$、线程块为 $32\times8$，网格是 $3\times3=9$ 块，共启动 2304 个线程，其中仅 1105 个通过边界检查；当 \code{numAtoms=0} 时不启动 DCS 内核但仍清零这 1105 个输出点。
\practiceheading 假定 32-lane NVIDIA GPU、\code{CHUNK_SIZE} 个 16 字节原子不超过 constant memory 容量且每个 warp 处理连续 $x$，用 Nsight Systems 核对 copy/kernel 的同 stream 顺序，再以 Nsight Compute 的 Constant Cache Hit Rate 比较全局只读版本。
\sourcecheck{https://docs.nvidia.com/cuda/cuda-c-programming-guide/}{CUDA C++ Programming Guide 对执行配置、stream 顺序和 constant memory 的官方说明；对抗检查：只配置恰好整除的示例尺寸会掩盖越界，跨 stream 覆盖常量符号则会造成数据生命周期竞态}
\end{solutionexercise}

\begin{solutionexercise}{2}
\problemheading 比较原书图 21.6 内核每次迭代执行的运算数量（内存加载、浮点运算、
分支）与把原书图 21.8 推广为粗化因子 8 后的内核。注意，后一个内核的每个线程对应前一个
内核的 8 个线程。按覆盖相同 8 个网格点、处理同一个原子的口径，分别统计四个原子字段的
加载，距离与势能表达式中的减、乘、加、平方根和除法，以及原子循环的控制分支；说明
FMA 或 \code{rsqrtf} 等编译器变换是否计入。

\editorialnote 题目文字要求把粗化因子取为 8，而其引用的原书图 21.8 代码实际定义
\code{COARSEN_FACTOR 4}。本解按题目要求把图中四点结构推广到八点后计数，并同时给出
因子 4 的变体结果；不把图中常量静默改成 8。

\approachheading 为避免把编译器的 FMA、常量折叠和 \code{rsqrt} 混入手算，采用明确
的源码级约定：只数 4 个 \code{atoms[]} 数据加载；把加、减、乘、除、\code{sqrtf} 分列；
计入 \code{energy +=}，不计地址整数运算。并列给出原书图 21.8 写法的直接八点推广，
以及再用 \code{dx[r]=dx[r-1]+spacing} 消去常数倍乘法的递推优化。

\answerheading 对一个普通线程、一个原子、一个网格点，循环体有 4 次原子数据加载，
3 减、3 乘、3 加、1 次平方根、1 次除法。因此覆盖 8 个点的 8 个线程合计为：
\begin{center}
\begin{tabular}{c|rrrrrr|c}
\toprule
方案 & 加载 & 减 & 乘 & 加 & \code{sqrtf} & 除 & 原子循环控制\\
\midrule
8 个普通线程 & 32 & 24 & 24 & 24 & 8 & 8 & 8 个 lane 各执行一次\\
图 21.8 直接推广 & 4 & 3 & 16 & 24 & 8 & 8 & 1 个 lane 执行一次\\
递推优化 & 4 & 3 & 10 & 24 & 8 & 8 & 1 个 lane 执行一次\\
\bottomrule
\end{tabular}
\end{center}
两条粗化行都有 $dx_0,dy,dz$ 共 3 减；$dy^2,dz^2$ 与 8 个 $dx_r^2$ 共 10 次必需乘法；
7 次横坐标更新、1 次公共平方和、8 次半径平方和、8 次能量累加共 24 加。原图直接推广还
按 \code{dx0+r*spacing} 计算 $r=2,\ldots,7$，故另有 6 次常数倍乘法，总数为 16；递推
版本把这 6 次乘法变成已经计入的相邻点加法，降为 10。平方根和除法是每个网格点不可复用
的，各仍为 8。表中“8 对 1”表达逻辑 lane 工作量；实际动态分支指令数还取决于点与 warp
的映射。编译器也可能融合乘加、降低除法/平方根，最终 SASS 应以 profiler/反汇编为准。

\editorialnote 英文底本相邻正文给粗化因子 4 报出“11 次加法”。按同一段对基础内核的
计数口径，$3$ 次 $dx$ 偏移、$1$ 次公共平方和、$4$ 次半径平方和及 $4$ 次能量累加共
12 次；其乘法数又取决于是否把常数倍间距移出循环。底本没有声明这一优化口径，因此本解
不沿用该数值，而显式给出可复算的源码模型。

\conclusionheading 粗化 8 倍把原子加载从 32 降到 4、减法从 24 降到 3；按原图直接推广是
16 次乘法，另作递推优化才是 10 次。每点的平方根、除法和最终累加不能由复用消除。
\discussionheading
\discussionpoint{线程粗化用较少线程换取寄存器内复用。}
普通映射让八个线程各算一个相邻网格点，它们会分别加载同一原子的四个字段，并重复计算共同的 $dy,dz$。粗化后一个线程维护八个累加器，只加载一次原子，且 $dy^2+dz^2$ 可供八个 $x$ 位置复用；代价是可独立调度的线程数降为约八分之一。这个交换与 CPU 循环展开、寄存器分块和 GEMM micro-kernel 同源，都是把跨任务共享的数据提升到线程私有快速状态。是否获益取决于原来是否受加载限制，以及剩余线程能否覆盖设备，不能从“加载减少八倍”直接推出“时间减少八倍”。

\discussionpoint{距离递推减少乘法却引入依赖与舍入传播。}
相邻位置满足 $(x+h-a_x)^2=(x-a_x)^2+2h(x-a_x)+h^2$，所以可以从前一点更新平方距离，而不是每次重新计算坐标差和乘法。这样节省算术，却让第 $r+1$ 点依赖第 $r$ 点，形成串行链；每步 FP32 舍入也可能沿长粗化区间累积。按一阶误差模型，连续 $r$ 次更新引入的舍入项上界通常随 $r u$ 增长，而从原坐标独立重算不会继承前一点的舍入误差，因此两种公式即使数学等价也未必逐位相同。对于因子 4 或 8，误差通常可由参考实现量化，若跨度很长则可每隔若干点从原坐标重算，作为数值锚点。递推优化广泛用于有限差分和光栅化，但都要在少算指令、指令级并行与累计误差之间作明确选择。

\discussionpoint{复用会把瓶颈从存储层推向执行管线。}
原子字段加载减少后，每个网格点仍要形成平方距离并执行平方根、倒数或 \code{rsqrtf}；这些特殊函数不会因相邻点共享而消失。当常量数据已命中缓存，SFU 吞吐、长依赖链或寄存器 spill 可能先成为限制，内存带宽反而不再主导。Roofline 风格分析应重新计算粗化后的字节数与算术强度，并把每点特殊函数单独看待，因为一条 \code{sqrtf} 不能简单折算成一次普通加法。优化改变瓶颈位置是常见现象，后续调优必须重新测量，而不是继续针对已经缓解的原始瓶颈。

\discussionpoint{粗化因子需要按问题形状与资源共同选择。}
较大因子提高数据复用，却为每线程增加累加器、坐标状态和尾部谓词，寄存器数上升后可能降低驻留或发生 local memory spill。总线程数也随因子下降，小网格可能无法覆盖全部 SM；因子过小则保留并行度，却失去共享原子的机会。工程实现可为 1、2、4、8 等常见因子生成特化版本，用 \code{ptxas} 报告和实测每点时间选择，并在非整宽度上验证尾部。选择过程还应与未粗化 FP64 参考比较误差，使 autotuning 不会只追求吞吐而接受未经说明的数值漂移。
\variantheading 若粗化因子改为 4，四个普通线程合计为 16 次加载、12 次减/乘/加、4 次平方根和 4 次除法；直接推广有 4 次加载、3 次减、8 次乘、12 次加，递推后乘法再降为 6 次。
\practiceheading 固定 NVIDIA \code{sm_80}、FP32、\code{--use_fast_math} 关闭和粗化因子 8，用 \code{nvdisasm} 检查实际 FMA/\code{MUFU} 指令，再由 Nsight Compute 记录寄存器数、spill、SFU 利用率和每点周期。
\sourcecheck{https://docs.nvidia.com/cuda/cuda-c-best-practices-guide/}{CUDA C++ Best Practices Guide 仅用于核验寄存器、占用率与生成指令之间的一般测量权衡，并不直接给出本题粗化算术；表中 16/10 的源码级计数由原书图 21.8 独立推广复算。对抗检查：把优化后的 10 冒充原图直接推广值，或把源码运算数直接等同于机器指令数，都会得到过强结论}
\end{solutionexercise}

\begin{solutionexercise}{3}
\problemheading 如原书第 21.3 节所示，增加每个 CUDA 线程完成的工作量可能带来哪两个
潜在缺点？回答须分别从每线程资源需求以及可供 GPU 调度的并行任务数量说明其机制，而非
只列性能现象。

\approachheading 分别考察单线程资源和整个网格可用并行任务数。

\answerheading 两个主要缺点是：
\begin{enumerate}
  \item \textbf{资源压力与占用率下降。}粗化需要同时保存更多坐标、中间距离和能量累加器，
  增加每线程寄存器数；超过寄存器预算时，同一 SM 可驻留的 warp/块减少，甚至发生 local
  memory spill。延迟隐藏能力随之下降，新增复用可能抵不过 spill 流量。
  \item \textbf{并行度与负载均衡下降。}粗化因子 $c$ 使线程数约降为原来的 $1/c$。小问题
  或边界 tile 可能没有足够线程填满 GPU；每线程运行更久，尾部无效点、数据相关分支或不同
  工作量也会拉长整个 warp/块的关键路径。
\end{enumerate}
合格的开放题还可提及指令级依赖增加、调度公平性、代码尺寸/编译器压力，但评分核心是
“寄存器/占用率”和“并行任务/负载均衡”两类不同机制。最佳粗化因子必须用目标问题尺寸和
目标 GPU 实测，不能仅凭复用次数选择。

\conclusionheading 粗化以减少冗余换取更多单线程状态和更少线程；收益存在硬件与问题规模
共同决定的拐点。
\discussionheading
\discussionpoint{occupancy 是驻留上限，不是综合性能得分。}
occupancy 表示每个 SM 实际或理论驻留 warp 数相对架构上限的比例，寄存器、共享内存、每块线程数和最大块数都可能限制它。更多驻留 warp 给调度器更多候选来隐藏延迟，却不保证内存、SFU 或发射槽已经充分利用；数据复用好的内核即使 occupancy 较低也可能更快。反过来，把寄存器强行限制到获得 100\% occupancy，若导致 spill 到 local memory，性能反而下降。因此它适合诊断“是否可能缺少并发”，不适合作为独立优化目标，更不能跨架构比较一个脱离资源配置的百分比。

\discussionpoint{粗化增加 ILP 的同时减少 TLP。}
一个线程维护多个网格点时，编译器和调度器可能交错独立的加载、平方根与累加，从而用指令级并行 ILP 隐藏单条依赖延迟。可是线程级并行 TLP 随粗化下降，尾部块和小问题可能留下空闲 SM，也减少可用于切换的 warp。理想状态不是让某一个比例达到 100\%，而是让 ILP 与 TLP 共同提供足够在途工作，同时保持数据复用。这个平衡也出现在 CPU 向量化和软件流水中：扩大单任务内部并行会减少任务数量，调度层级之间必须通过实测协调。不同问题尺寸还会改变可用并行任务数，最优粗化率不应写死。

\discussionpoint{寄存器分配以离散粒度形成驻留台阶。}
纸面上从每线程 64 个寄存器增加到 65 个似乎只增长 1.6\%，硬件却按 warp 或块的分配粒度向上取整，可能突然让一个 SM 少驻留一个完整线程块。编译器在压力过高时还会把局部值溢出到 local memory，使静态寄存器数下降但增加全局内存流量；只看一个数字会误判。应把 \code{ptxas} 的实际报告、块尺寸、动态共享内存和目标架构输入 Occupancy API，再与 profiler 的 achieved occupancy 对照。离散资源台阶说明性能空间不是光滑函数，跨过边界的小改动可能产生不成比例的结果。

\discussionpoint{诊断核心是延迟是否已经被充分隐藏。}
看到 occupancy 从 75\% 降到 50\% 时，首先应问 eligible warps 是否仍足够、发射槽是否空闲，以及 stall 来自内存依赖、SFU 还是同步。如果低 occupancy 伴随更高缓存复用和更少指令，端到端每点时间可能改善；若同时出现 spill 与 eligible warp 枯竭，才说明资源压力真正伤害吞吐。行业调优通常扫描块大小与粗化因子，寻找时间、资源和误差的 Pareto 点，并为小网格另设低延迟配置。计数器的作用是建立因果链，而不是堆成一张没有问题假设的指标清单。
\variantheading 在每 SM 有 65,536 个 32 位寄存器、块含 256 线程且暂不计分配粒度时，64 寄存器/线程可驻留 4 块、1024 线程；粗化后若升到 96 寄存器/线程，只能驻留 2 块、512 线程。
\practiceheading 假定目标 SM 正好有 65,536 个寄存器并使用 256 线程块，用 \code{nvcc --ptxas-options=-v} 和 CUDA Occupancy API 核对静态上限；Nsight Compute 同时观察 achieved occupancy、Local Memory spill 和 Eligible Warps，避免只优化单一占用率。
\sourcecheck{https://docs.nvidia.com/cuda/cuda-c-best-practices-guide/}{CUDA C++ Best Practices Guide 仅核验寄存器压力、occupancy、spill 和可用并行度的一般权衡，并非 DCS 粗化因子的直接实验依据；对抗检查：高 occupancy 不是独立性能目标，但寄存器 spill 和线程数不足是必须排除的真实反例}
\end{solutionexercise}

\begin{solutionexercise}{4}
\problemheading 修改原书图 21.8 的内核，使更新势网格点时可以使用向量存储来改善内存
访问合并。原内核的粗化因子为 4，每个线程计算同一行中 4 个连续点并执行四次标量
\code{energygrid[...] += energy}。修改后须用四元素向量完成对齐的完整组读--改--写，
同时处理 $x$ 维不足四项的尾部、地址对齐和二维网格边界。

\approachheading 要让每个线程拥有不重叠且 16 字节对齐的连续四点，起始 $i$ 必须是线性
线程号乘 4。由于原操作是 \code{+=}，需做一次 \code{float4} 向量加载、分量累加和一次
向量存储；不完整或未对齐尾部走标量路径。

\answerheading 计算 \code{energy0..energy3} 的原子循环保持不变，索引和写回改成：
\begin{lstlisting}[language=C++]
constexpr int C = 4;
int linear_x = blockIdx.x * blockDim.x + threadIdx.x;
int i = linear_x * C;
int j = blockIdx.y * blockDim.y + threadIdx.y;
if (j >= grid.y || i >= grid.x) return;

// x = i * gridspacing; compute energy0..energy3 for x + r*spacing.
size_t rowBase = (size_t(grid.x) * grid.y * k) + size_t(grid.x) * j;
size_t first = rowBase + i;
if (i + 3 < grid.x && (first & 3u) == 0) {
    float4* p = reinterpret_cast<float4*>(energygrid + first);
    float4 old = *p;
    old.x += energy0; old.y += energy1;
    old.z += energy2; old.w += energy3;
    *p = old;
} else {
    const float e[4] = {energy0, energy1, energy2, energy3};
    #pragma unroll
    for (int r = 0; r < 4; ++r)
        if (i + r < grid.x) energygrid[first + r] += e[r];
}
\end{lstlisting}
\code{cudaMalloc} 返回的基址满足向量对齐；要让每行始终走向量路径，还应令逻辑宽度
$G_x$ 是 4 的倍数，或手工把物理行步长补齐到 4 的倍数。若改用
\code{cudaMallocPitch}，内核必须接收以元素为单位的行步长
$p=\code{pitchBytes}/\code{sizeof(float)}$，并把行首改为
\[
  \code{rowBase}=(\code{size_t}(k)\,G_y+j)\,p.
\]
继续按逻辑宽度 $G_x$ 计算会索引到错误行。
相邻 warp lane 的 \code{float4} 地址依次相差
16 字节，合起来覆盖连续区域，因而能形成合并交易。每个线程的四点区间互不重叠，所以普通
读改写安全；若另一个内核或主机同时更新同一网格，仍需建立 stream/event 依赖或改变归并
策略。

\editorialnote 原书图 21.8 的索引写成：
\begin{lstlisting}[language=C++,numbers=none]
i = blockIdx.x * blockDim.x * COARSEN_FACTOR + threadIdx.x;
\end{lstlisting}
但随后访问 $i,i+1,i+2,i+3$。这样相邻线程的区间重叠，并与正文所称“相邻线程相隔
4 个元素”矛盾。正确的连续四点映射是：
\begin{lstlisting}[language=C++,numbers=none]
i = (blockIdx.x * blockDim.x + threadIdx.x) * COARSEN_FACTOR;
\end{lstlisting}
本解同时补上原图缺少的 $i,j$ 边界检查。

\complexity{算术仍为每线程 $O(4A)$，$A$ 为原子数}{$O(1)$ 每线程额外寄存器；向量路径一次 16 字节加载和一次 16 字节存储}{约 $\lceil G_x/4\rceil G_y$ 个线程；每线程独占四点，无原子竞态}
\conclusionheading 向量类型本身不会修复错误映射；先保证四点所有权不重叠和 16 字节对齐，
再使用 \code{float4} 才同时正确且合并。
\discussionheading
\discussionpoint{向量类型减少指令，却不减少逻辑数据量。}
用 \code{float4} 读改写四个连续势能值，逻辑上仍读取并写回各 16 字节，与四次标量访问处理的数据完全相同。宽指令可能减少地址计算、指令数和前端压力，但若四个标量原本已由 warp 合并进相同内存交易，DRAM 字节不会因此再降四倍。反过来，非对齐宽访问可能被拆分，甚至比标量路径更差。评价向量化应区分指令级收益和事务级收益，先问瓶颈在发射、地址计算还是带宽；这个区分也适用于 SIMD 加载和网络批量拷贝，打包操作不是数据压缩。只有硬件计数器才能判明交易数是否真的改变。

\discussionpoint{对齐契约必须覆盖基址、行首和线程起点。}
\code{cudaMalloc} 给出的整体基址通常有充分对齐，但若每行含 10 个 float，下一行首相对前一行移动 40 字节，并不总是 16 字节倍数。即使行首对齐，线程起点也必须满足元素下标为 4 的倍数，才能把 \code{float*} 合法解释为 \code{float4*}。可用 padding 或 \code{cudaMallocPitch} 保证行距，也可在运行时检查地址后选择宽路径，其他位置退回标量。对齐不是分配时一次完成的属性，而是沿每次指针偏移传播的证明义务，这正是多维数组最容易忽略的边界。

\discussionpoint{尾部策略必须保证一次且仅一次的所有权。}
行宽不是 4 的倍数时，主体可以每线程处理一个 \code{float4}，尾部再由标量循环清理；也可以为存储加 padding、使用支持掩码的宽操作，或把尾部放进单独内核。标量清理最容易证明却引入分支，padding 增加空间且要求无效槽不可影响语义，第二内核则增加启动成本。无论选择哪种方案，都要证明主体与尾部区间互斥，且它们的并集恰为 $[0,\mathrm{grid.x})$。这和并行分区中的完整覆盖不变量相同，比“样例没有越界”更强，也能直接指导 sanitizer 测试。

\discussionpoint{源码中的宽类型还要由机器指令确认。}
类型别名、潜在非对齐或复杂控制流可能让编译器把 \code{float4} 拆成多个标量指令；反过来，清晰的标量循环也可能被自动向量化。仅从 C++ 类型无法知道最终加载宽度，应检查 PTX/SASS，并用事务数与动态指令数验证。对很短或不规则的行，为对齐而增加的 padding 和尾部判断可能超过减少的指令，尤其当内核本来受平方根而非访存限制。向量化因此是一项需要编译结果和真实工作负载共同证明的优化，而不是换一个类型名就完成的性能承诺。这种证据链可以防止优化只停留在源码表象。
\variantheading 若 \code{grid.x=10} 且行首 16 字节对齐，起点 $i=0,4$ 的两个线程走 \code{float4}，$i=8$ 的线程标量处理元素 8、9；索引 0--9 恰好各更新一次，没有重叠或越界。
\practiceheading 假定 32-lane warp、\code{float4} 自然对齐为 16 字节，用两种版本比较：
当前线性分配令 \code{grid.x} 为 4 的倍数；pitched 版本用 \code{cudaMallocPitch} 并把
\code{pitchElements} 传入上述行首公式。用 Compute Sanitizer racecheck 验证所有权，再以
Nsight Compute 比较 Global Store sectors 和 scalar/vector 指令数。
\sourcecheck{https://docs.nvidia.com/cuda/cuda-c-programming-guide/}{CUDA C++ Programming Guide 对内建向量类型、自然对齐和全局内存访问的官方说明；高置信度}
对抗检查：把未对齐地址强制转换为 \code{float4*}，或让相邻线程覆盖重叠四元组，分别会
造成未定义行为和数据竞态。
\end{solutionexercise}

\begin{solutionexercise}{5}
\problemheading 使用原书图 21.13 解释：当块中的线程处理邻域列表中的一个分箱时，
控制流分歧如何产生？该图的二维示意中，一个线程块覆盖中心分箱内的一片网格点区域，并用
“截断距离加分箱半对角线”形成超级圆；真实实现对应 $8\times8\times8$ 网格点块和三维
超级球。邻域列表会保守纳入与超级球相交的分箱，各线程随后针对同一候选原子独立判断自身
网格点是否落在截断半径内。请说明哪些分箱/原子会使同一 warp 的判断结果不同。

\approachheading 块共享同一邻域分箱列表和其中的原子，但每个线程代表不同网格点；判断
条件取决于“当前线程的点到当前原子”的距离。

\answerheading 线程协作把一个候选分箱的原子载入共享内存后，对同一原子 $a$，线程
$t$ 计算
\[
 r_{t,a}^2=(x_t-x_a)^2+(y_t-y_a)^2+(z_t-z_a)^2
\]
并执行类似 \code{if (r2 < cutoff2) accumulate(...)} 的分支。原书图 21.13 中一个块覆盖
有面积/体积的区域：靠近原子的线程可能在截断球内，远端线程却在球外。同一 warp 对条件
得到混合谓词时，硬件先屏蔽 false lane 执行累加路径，再处理另一条路径，产生分歧。

分歧最严重的是只与块的截断“超级圆/球”部分相交的边界分箱；若某原子对整块所有网格点都
在截断内（或都在截断外），该原子的判断一致，不产生这处分歧。邻域列表采用保守覆盖是正确性
所需：它宁可纳入部分相交分箱，再逐线程精确判断，不能因为想消除分支就漏掉截断球内原子。

评分点包括：共同候选原子、线程坐标不同、距离谓词不同、warp 串行化两条路径，以及部分相交
分箱最易触发。可用无分支谓词乘法减少显式控制流，但会为球外原子执行昂贵的平方根/除法；
是否更快需测量，并应避免 $0\times\infty$ 等数值异常。严格小于还是小于等于截断半径也必须
统一，否则边界原子会因浮点舍入产生实现差异。

\complexity{每网格点工作为 $O(A_{\rm neighborhood})$，理想情况下与全局原子总数无关}{每块共享一个分箱原子 tile，另有每线程能量累加器}{块间独立；块内协作加载后同步，逐线程距离判断可能在 warp 内分歧}
\conclusionheading 分歧不是因为线程访问不同分箱，而是因为它们对同一候选原子拥有不同
空间距离，因而在截断谓词上分道执行。
\discussionheading
\discussionpoint{截断把全局相互作用改造成局部近似。}
原始 $q/r$ 势没有有限作用半径，每个原子理论上影响所有网格点；设置截断半径后，球外贡献被忽略，每个原子只需更新附近空间分箱，工作量才可能从 $PA$ 降到接近线性。硬截断会让势在边界处跳变，力作为导数还可能出现更明显不连续，因此物理模拟常用 switching 或 shifted potential 在一段区间内平滑过渡。半径越小计算越省，近似误差却越大。这里的分支优化只是局部实现，真正的算法选择还必须说明截断模型和允许的物理误差。切换函数的连续阶数还会影响能量漂移与积分稳定性。

\discussionpoint{分支分歧取决于线程映射的空间相干性。}
若一个 warp 处理相邻网格点，它们相对同一原子的距离通常相近，绝大多数 lane 会一起处在球内或球外；只有截断球面穿过该区域时才出现混合路径。若线程随机映射到远处分散的点，内外判断缺少相干性，分歧会更频繁。按空间 tile 分配 warp、先用包围盒粗筛或把原子按分箱排列，都能让控制流与几何局部性对齐。分支效率因此不只是 if 语句写法的属性，还由数据排序和任务映射共同塑造，这一规律同样适用于光线追踪和碰撞检测。阈值附近的数据重排成本也应计入完整内核时间。

\discussionpoint{谓词化可能隐藏分支却保留昂贵工作。}
把球外贡献乘以 0 可以让所有 lane 走同一指令序列，但它们仍可能执行平方根和倒数，既没有节省昂贵运算，又可能在 $r=0$ 时形成 $0\times\infty=\mathrm{NaN}$。本题保持物理公式 $q/r$，重合点就是有符号无穷的奇点，不能通过静默加入软化项改变模型。更好的实现通常先比较廉价的平方距离，只让活动 lane 计算倒数，同时显式约定奇点是保留、报错还是由上层模型排除。谓词化与分支各有成本，正确性语义必须先于减少控制流分歧的愿望。若活动比例很低，先压缩活动索引也可能更合算。

\discussionpoint{近场与远场方法共同承担误差预算。}
工业分子模拟不会简单丢弃所有远场：邻居表负责短程直接作用，粒子网格 Ewald、FFT 或多极方法以近似方式汇总长程贡献。截断半径不仅影响本地计算，还决定多 GPU 分区需要多宽的 halo、多久重建邻居表以及远场误差由哪一层补偿。验证应比较势能、力、长期守恒漂移与运行时间，而不是只观察分支效率；相同速度若破坏物理可接受范围就没有意义。这个多尺度分解把一道 warp 分支题连接到了科学计算的核心方法：为不同距离尺度选择不同算法，并协调它们的误差与通信成本。
\variantheading 若一个 32-lane warp 对应 $8\times4$ 网格点，而某原子的截断范围恰覆盖每行的前 4 个 $x$ 点，则 16 lane 进入累加、16 lane 跳过；该次动态判断的活动比例为 $50\%$。
\practiceheading 假定硬件 warp 为 32、一个 $8\times4$ 线程块恰映射单 warp 且距离用 FP32，在 Nsight Compute 中联合查看 Branch Efficiency、Warp Execution Efficiency 和 \code{MUFU} 利用率；用 NAMD/VMD 生成的实际原子分箱分布比较分支与谓词版本。
\sourcecheck{https://docs.nvidia.com/cuda/cuda-c-programming-guide/}{CUDA C++ Programming Guide 对 warp 控制流分歧和屏蔽执行的官方说明；对抗检查：共同遍历同一循环并不保证无分歧，数据相关的截断条件仍可让 lane 采取不同路径}
\end{solutionexercise}

\chapter{算法选择、问题分解与问题表述}
\noexercises

\chapter{多 GPU 编程}

\begin{solutionexercise}{1}
\problemheading 假设把本章的五点模板应用于一个网格，其 $x$ 维包含 64 个网格点，
$y$ 维包含 512 个网格点；计算沿 $y$ 维均匀划分到 16 个 MPI 秩。原书图 23.1 的内核
只更新 $x=1,\ldots,62$，两端列为固定边界；图 23.12 的阶段 1 先计算每个秩顶部和底部
两个边界输出行，阶段 2 在交换这两个完整的 64 元素 \code{float} 行时计算其余内部行。
分别回答：
\begin{enumerate}[label=(\alph*)]
  \item 每个进程计算多少个输出网格点值？
  \item 每个进程需要多少个 halo 网格点？
  \item 阶段 1 中，每个进程计算多少个边界网格点？
  \item 阶段 2 中，每个进程计算多少个内部网格点？
  \item 阶段 2 中，每个进程发送多少字节？
\end{enumerate}

\approachheading 一维域分解不切分 $x$ 维。每个秩拥有 $512/16=32$ 个输出行；五点模板
只需相邻分区各一行 halo。原书图 23.12 的阶段 1 先算上下两条边界输出行，阶段 2 计算其余
行并发送这两行。

\answerheading
\begin{enumerate}[label=(\alph*)]
  \item 每个进程计算
  \[
    62\times32=1984
  \]
  个输出网格点。
  \item 顶、底各一条 halo 行，共
  \[
    2\times64=128
  \]
  个 halo 槽位；所以含 halo 的本地输入数组是 $64\times34$。
  \item 阶段 1 计算顶部、底部两个边界输出行；每行只更新 $x=1,\ldots,62$，共
  $2\times62=124$ 点。
  \item 阶段 2 计算 $32-2=30$ 个内部行，同样每行更新 62 点，共
  $30\times62=1860$ 点。检查：$124+1860=1984$。
  \item 每个边界值是一个 \code{float}，逻辑发送量为
  \[
    2\times64\times4=512\ \text{字节}.
  \]
\end{enumerate}
\editorialnote 底本题面把网格写成 $64\times512$，容易把每个秩拥有的全部 64 列都算作
“计算出的输出”。但原书图 23.1 的五点模板内核明确从 $x=1$ 开始，并要求
$x<\code{nx}-1$；$x=0,63$ 是固定边界列，不由该内核更新。因此 (a)、(c)、(d) 必须按
62 个可更新列计为 1984、124、1860。图 23.14 的 halo 交换仍发送完整的 64 元素边界行，
所以 (b)、(e) 保持 128 和 512 字节。若机械地把 64 列全算作输出，才会得到底本容易诱发的
2048/128/1920 口径，但它与所指定内核的边界条件不一致。

上述数字按每个秩预留上下两侧、并执行两个方向的 halo 交换来计。若全局 $y$ 边界是非周期
边界且端点秩以 \code{MPI_PROC_NULL} 代替不存在的邻居，则中间 14 个秩实际在网络上发送
512 字节，两个端点秩各只有一个真实邻居，实际网络负载为 256 字节；它们仍可保持 128 个
halo 槽位和相同代码结构。若采用周期边界，则所有秩都是 512 字节。

\complexity{每秩 $O(62\cdot32)$ 有效点计算；总计算量与可更新网格点数成正比}{每秩 $O(64\cdot34)$ 本地输入及输出空间}{16 个秩并行；阶段 2 的 1860 点计算可与至多 512 字节通信重叠}
\conclusionheading (a)--(e) 依次是 1984、128、124、1860、512 字节；
非周期端点的实际网络字节数是需要单独说明的边界例外。
\discussionheading
\discussionpoint{halo 的宽度由数据依赖决定。}
半径为 1 的五点模板计算边界输出时，只会读取相邻分区的一层输入，所以每侧交换一行；若离散模板半径变为 $r$，或一次内核融合多个时间步，所需 halo 可能扩到 $r$ 层甚至更宽。分区方向决定哪一维的邻接变成远程通信，MPI 只是搬运这些依赖的工具。最可靠的方法是先画出每个输出点的读集合，再把跨分区部分映射成消息。这样能区分本地固定边界、halo 输入和真正由内核更新的输出，避免把含 halo 的本地数组尺寸直接当作计算点数。读集合分析还可自动生成通信表并验证角点是否重复交换。

\discussionpoint{强扩展最终受表面积与体积之比约束。}
固定全局网格而增加秩数时，每秩内部体积持续缩小，计算量近似按秩数下降；半径固定的边界表面积下降得更慢，通信占比因而上升。二维条带分解只有两个邻居却可能拥有较长边界，二维块分解能降低单秩周长，却增加邻居与角点处理；三维情况同理。局部子域很小时，一条 512 字节消息的固定延迟甚至比传输时间更重要，继续加 GPU 不再带来线性加速。表面积/体积模型因此是 stencil 强扩展的第一性解释，也能指导进程网格形状和分区维度选择。弱扩展则保持局部体积不变，瓶颈模型与此不同。

\discussionpoint{非阻塞调用只创造重叠机会。}
先发起 \code{Irecv/Isend}，再计算不依赖 halo 的 1860 个内部点，最后 \code{Waitall} 并计算 124 个边界点，在依赖图上允许通信与内部计算并行。它并不保证网络真的在后台推进：MPI 可能只在再次进入库时提供弱进展，GPU 内核也可能占满需要的复制或执行资源，设备缓冲还必须通过事件完成生产。GPU-aware MPI 又受注册、网卡与拓扑影响。是否重叠应从时间线和总时间验证，不能因为 API 名称带 \code{I} 就宣称隐藏了通信；正确表述是“建立了可重叠结构”。

\discussionpoint{物理边界会改变真实网络邻居。}
非周期 $y$ 边界下，首秩没有顶部邻居、末秩没有底部邻居，使用 \code{MPI_PROC_NULL} 可以保留统一调用结构，同时这些方向不产生传输并且接收缓冲保持原值。周期边界则让首尾秩互为邻居，所有秩都执行两个方向的实际交换。物理边界值可能由初始化或独立内核维护，所以“预留两个 halo 槽位”不等于“网络一定接收两行”。规模报告应区分逻辑 payload、实际邻居数和网络字节，并记录秩到 GPU/NUMA 节点的绑定，否则端点例外与拓扑差异会被平均值掩盖。
\variantheading 若改为沿 $y$ 均分到 8 个秩，每秩有 64 个输出行，输出点为 $62\times64=3968$，halo 仍为 128 点，阶段 1 为 124 点，阶段 2 为 $62\times62=3844$ 点，发送量仍为 512 字节。
\practiceheading 先运行 \code{solutions/code/ch17-23/ch23_ex01_partition.cpp} 核对分区
计数。随附 \code{ch23_ex01_halo_exchange_mpi.cpp} 使用非周期 $y$ 边界和
\code{MPI_PROC_NULL}：发起四个非阻塞调用后先计算 1860 个内部点，\code{MPI_Waitall} 后逐项
核验 128 个 halo 槽位（含端点保留的物理边界），再计算 124 个边界点并独立核验全部 1984 个
输出。该程序验证依赖顺序和重叠机会，不声称 MPI 一定异步推进；扩展到每秩一张 GPU 时，应把
主机向量替换为设备缓冲，明确 stream/event 所有权，并用 Nsight Systems 时间线证明 GPU-aware
MPI 与内部内核确实重叠。
\sourcecheck{https://www.mpi-forum.org/docs/mpi-4.1/mpi41-report.pdf}{MPI Forum 的 MPI 4.1 官方标准用于核验 halo 通信与 \code{MPI_PROC_NULL} 边界语义；点数另按原书图 23.1 的 $x$ 边界条件独立复算。对抗检查：把含 halo 的 34 行或固定的两列当成内核输出都会多算，把 \code{MPI_PROC_NULL} 调用算成实际网络传输则会高估端点秩}
\end{solutionexercise}

\begin{solutionexercise}{2}
\problemheading 如果以下 MPI 调用传输了 4000 字节：
\begin{center}
\small\code{MPI_Send(ptr_a, 1000, MPI_FLOAT, 2000, 4, MPI_COMM_WORLD)}
\end{center}
每个被发送数据元素的大小是多少？
\begin{enumerate}[label=(\alph*)]
  \item 1 字节
  \item 2 字节
  \item 4 字节
  \item 8 字节
\end{enumerate}

\approachheading MPI 的消息数据量是 \code{count} 乘以 datatype 的类型签名大小；题目
已经给出总字节数，无需根据指针类型猜测。

\answerheading
\[
  \frac{4000\ \text{bytes}}{1000\ \text{elements}}
  =4\ \text{bytes/element}.
\]
因此选择 (c)。在可移植程序中可用
\code{MPI_Type_size(MPI_FLOAT, \&size)} 查询该实现的类型大小；本题的已知传输量已经
把答案固定为 4。目的秩 2000 只有在 communicator 至少包含 2001 个秩时才合法，但这不改变
条件式算术。

\conclusionheading 每个被发送元素为 4 字节，答案 (c)。
\discussionheading
\discussionpoint{count 与 datatype 一起定义用户数据。}
调用 \code{MPI_Send} 时，\code{count} 说明有多少个类型实例，\code{datatype} 决定实例布局，而 \code{MPI_Type_size} 给出一个实例包含的有效基本数据字节，因此连续预定义类型的 payload 是前两项数量的乘积。目的秩选择接收者，tag 与 communicator 参与消息匹配，它们即使是数值 2000 或 4，也不能解释成字节数。发送端与接收端不必使用相同类型句柄，但类型签名中的基本数据项必须相容。把参数角色逐项分清，才能从一次函数调用正确推导数据量，也能避免把寻址元数据混进性能模型。

\discussionpoint{datatype size 与 extent 回答两个不同问题。}
派生类型可以描述结构体中的有效字段、二维数组的一列或带间隔的子数组；size 只累计真正参与消息的基本数据，extent 则覆盖从下界到上界的内存跨度，包括空洞与显式边界。连续发送多个该类型时，第 $k$ 个实例按 extent 定位，而不是紧接在前一个有效字节后。若用 size 分配本地数组跨度，可能让相邻实例重叠或越界；若用 extent 当网络 payload，又会把空洞误算成有效传输。MPI 类型系统因此同时描述“传什么”和“内存中怎样跨步”，这与序列化 schema 和数组 stride 的区别相近。

\discussionpoint{网卡观察到的字节通常大于用户 payload。}
MPI 实现可能为消息附加匹配头、序号和传输层报头，大消息还可能分片，小消息则可能采用 eager 协议先复制到接收侧缓冲。于是应用计算的 4000 字节是用户数据，不是以太网或 InfiniBand 计数器必须显示的总线上字节。常用延迟带宽模型 $T=\alpha+n/\beta$ 把每消息固定成本 $\alpha$ 与随 payload 增长的项分开，协议切换还可能形成分段曲线。性能诊断应明确测量层级，不能用网卡字节反推 datatype size，也不能用用户 payload 否认真实协议开销。

\discussionpoint{派生类型提供可移植描述，却不保证零复制。}
发送一个跨步列向量时，MPI 可以让网卡直接 gather 分散地址，也可以先把有效项 pack 到连续临时缓冲，再执行普通发送；标准允许不同实现策略。复杂嵌套类型、GPU 缓冲或未注册内存可能让打包成本占主导，而显式 pack 有时更容易与计算流水重叠。比较方案应核对 \code{Type_size}、\code{Type_get_extent} 和收发类型签名，再测实际传输与 CPU/GPU 暂存时间。派生类型的首要价值是正确、可移植地表达布局，是否获得零复制属于实现与硬件能力，需要实测而不能从 API 语法推断。
\variantheading 若调用改为发送 250 个 \code{MPI_DOUBLE}，且 \code{MPI_Type_size} 返回 8，则用户 payload 为 $250\times8=2000$ 字节；目的秩和 tag 仍不进入字节数计算。
\practiceheading 在同构 CPU 节点、\code{MPI_THREAD_SINGLE} 且不使用 GPU 缓冲的假设下，用 Open MPI 或 MPICH 调用 \code{MPI_Type_size}/\code{MPI_Type_get_extent} 对比预定义与派生类型，并用 mpiP 区分 payload 计数和通信耗时。
\sourcecheck{https://www.mpi-forum.org/docs/mpi-4.1/mpi41-report.pdf}{MPI Forum 的 MPI 4.1 官方标准对元素计数、数据类型与类型大小的定义；对抗检查：第四个实参 2000 是目的秩而非字节数，第五个实参 4 是 tag 而非元素大小}
\end{solutionexercise}

\begin{solutionexercise}{3}
\problemheading 下列哪些说法正确？
\begin{enumerate}[label=(\alph*)]
  \item \code{MPI_Send()} 默认是阻塞的。
  \item \code{MPI_Recv()} 默认是阻塞的。
  \item MPI 消息必须至少有 128 字节。
\end{enumerate}

\approachheading 区分“阻塞调用返回”和“通信对端已经完成”；再检查 MPI 是否允许
\code{count=0}。

\answerheading (a)、(b) 正确，(c) 错误。
\begin{itemize}
  \item \code{MPI_Send} 是标准模式的阻塞发送。返回时发送缓冲区可以安全复用，但实现
  可能采用 eager 缓冲，也可能等待匹配接收；“阻塞”不等于必然与接收方同步。
  \item \code{MPI_Recv} 是阻塞接收，直到匹配消息已经放入接收缓冲区并可返回状态。
  \item MPI 允许零计数消息，标准没有 128 字节的最小消息限制。传输协议内部可能有 eager
  阈值或报头，但那不是用户消息的合法最小长度。
\end{itemize}

\editorialnote 英文底本使用单数问法 “Which ... is true?”，但标准语义下 (a) 和
(b) 同时正确。中文译稿改成“哪些说法正确”是必要修正，本题解据此作多选回答。

\conclusionheading 正确选项为 (a)、(b)；(c) 没有标准依据。
\discussionheading
\discussionpoint{阻塞发送承诺的是本地缓冲安全。}
\code{MPI_Send} 返回意味着发送缓冲区可以修改或复用，数据可能已经进入接收缓冲，也可能只被复制到 MPI 内部临时区；它不承诺远端进程已经从 \code{MPI_Recv} 返回。标准模式还允许完成依赖匹配接收是否启动，因此“阻塞”也不等于纯本地操作。接收返回的含义不同：匹配数据已经放进接收缓冲，并可通过 status 查询来源和 tag。理解本地完成与远端状态的区别，才能正确安排缓冲生命周期，也能避免把一次点对点调用误作隐式屏障。非周期边界上的 \code{MPI_PROC_NULL} 仍遵循调用完成规则。

\discussionpoint{eager 缓冲可能暂时掩盖错误的等待环。}
若所有秩都先对右邻调用标准发送、再从左邻接收，小消息可能被 eager 缓冲，使发送提前完成，程序看似正常。消息增大、缓冲耗尽、实现改用 rendezvous，或将调用替换为 \code{MPI_Ssend} 后，所有秩都等待匹配接收，环形死锁便暴露。正确程序不能依赖某个 eager 阈值、当前空闲内存或机器速度，因为这些都不是标准语义保证。压力测试应跨越消息尺寸和通信模式，但最终修复必须消除等待环，而不是把阈值调大。非周期拓扑虽不形成整环，内部秩的对称等待仍应按同一原则审查。

\discussionpoint{安全交换模式在简洁、重叠和生命周期之间取舍。}
\code{MPI_Sendrecv} 把一次发送和一次接收交给实现协调，适合规则邻居交换且较容易证明不会形成阻塞环。奇偶秩分阶段也能让阻塞调用安全，却增加轮次；\code{Irecv--Isend--Waitall} 为计算重叠创造机会，但请求对象、发送缓冲和接收缓冲必须一直存活到完成。零计数消息仍参与匹配，源、目的、tag 和 communicator 仍须正确，不能把“无 payload”理解为“无通信语义”。选用模式时应先证明匹配和生命周期，再讨论性能，API 更异步并不自动意味着程序更安全。

\discussionpoint{非阻塞启动与后台进展是两个层次。}
\code{MPI_Isend} 返回只说明请求已建立，数据何时在没有后续 MPI 调用时推进，取决于实现提供强进展、弱进展、专用线程还是网卡卸载。有些系统需要周期调用 \code{MPI_Test}，另一些会在后台传输；GPU-aware 路径还涉及设备事件、内存注册和 stream 可见性。因此把内部计算放在 \code{Waitall} 前只是建立了重叠机会，真正重叠要由时间线和总时间证明。调试还应检查 status/tag、超时和缓冲所有权，把正确性、进展保证与实际并发三个问题分开。
\variantheading 若环中每个秩先对右邻执行 \code{MPI_Ssend}、再从左邻 \code{MPI_Recv}，所有秩都可能等待匹配接收而死锁；改用一次 \code{MPI_Sendrecv} 后每个秩可同时完成左右交换。
\practiceheading 假定两台同构 InfiniBand 节点、Open MPI/MPICH、\code{MPI_THREAD_SINGLE} 且无 GPU-aware 路径，用 Intel MPI Benchmarks PingPong 按 2 的幂从 0/1 字节扫描到至少数 MiB，寻找延迟或带宽曲线的协议转折，并结合该实现公开的 MPI T control variable 核验 eager/rendezvous 配置；mpiP 可做调用画像，却不能普遍直接给出阈值。阈值只解释性能，不能改写标准语义。
\sourcecheck{https://www.mpi-forum.org/docs/mpi-4.1/mpi41-report.pdf}{MPI Forum 的 MPI 4.1 官方标准关于阻塞发送、阻塞接收和零计数消息的规范；对抗检查：不能把标准发送的本地完成误写成远端完成，也不能把某实现的 eager 阈值当成 MPI 消息下限}
\end{solutionexercise}

\chapter{总结与未来展望}
\noexercises


\part{附录}
\appendix
\chapter{数值计算考量}

\begin{solutionexercise}{1}
\problemheading 为一种 6 位格式（1 位符号、3 位尾数、2 位指数）绘制与原书图 A.5
等价的图。用结果说明尾数每增加一位，会如何改变数轴上的可表示数集合。作为必要上下文，
图 A.5 的基准是 1 位符号、2 位显式尾数和 2 位指数的教学“无零”格式：指数偏置为 1，
指数编码 11 保留，不采用零或反规格化数，只绘正数，负数关于 0 对称。新图须列出全部有效
指数区间内的点及相邻间距。

\approachheading 沿用原书图 A.5 的教学格式：指数 11 保留，指数偏置为 1，不采用
反规格化数；只列正数，负数关于 0 对称。三位显式尾数 $M=k/8$，$k=0,\ldots,7$。

\answerheading 对有效指数编码 00、01、10，实际指数分别为 $-1,0,1$。正可表示数为
\[
\begin{array}{c|l|c}
E&\text{正可表示数}&\text{相邻间距}\\ \hline
-1&0.5,0.5625,0.625,0.6875,0.75,0.8125,0.875,0.9375&1/16\\
0&1,1.125,1.25,1.375,1.5,1.625,1.75,1.875&1/8\\
1&2,2.25,2.5,2.75,3,3.25,3.5,3.75&1/4
\end{array}
\]
负集合是这些数取负；该“无零”格式仍不能表示 0。一般地，每个指数区间的值是
\[
  (1+k/2^m)2^E,\qquad k=0,\ldots,2^m-1.
\]
尾数从 2 位增至 3 位后，每个区间的点数由 4 翻倍到 8，相邻间距和 round-to-nearest
在该区间内的最大绝对误差都减半；指数区间和数量级范围没有改变。换言之，增加尾数位提高
精度，不负责扩大指数范围。

\conclusionheading 与原书图 A.5 相比，在每一对原有相邻点的中点插入一个新点；正负两侧
均如此。
\discussionheading
\discussionpoint{玩具格式首先是用于隔离概念的模型。}
为什么不直接把这 6 位当作一个缩小的 IEEE 格式？本题刻意采用显式尾数，并删除零、反规格化数、无穷和 NaN，目的在于让“增加一位尾数”成为唯一变量，因而可以逐个位型枚举。IEEE binary32 则使用隐含首位、偏置指数和保留编码；套用那套解码规则会同时改变端点、点数与原点附近的结构，所得数轴已不是题目要求的实验。这样的简化模型很适合解释因果，但引用其结论时必须同时声明省略项；在数值格式设计中，先隔离变量、再恢复特殊值和异常路径，也是一种常用的分层验证方法。

\discussionpoint{尾数位决定一个 binade 内的采样密度。}
如果指数保持为 $E$，多出的一位究竟改变了哪里？对含 $p$ 个显式小数位的本题格式，区间 $[2^E,2^{E+1})$ 内的相邻间距为 $2^{E-p}$，所以 $p$ 每增加 1，点数翻倍、间距和最近舍入的最大局部绝对误差都减半。指数端点并未移动；当 $E$ 增加 1 时，绝对 ULP 又随尺度翻倍，而正规格化数的相对间距仍约为 $2^{-p}$。这一区分解释了浮点数为何能在很宽的数值范围内维持近似恒定的相对精度，也提醒工程人员不要用某个固定绝对容差评判所有数量级的数据。

\discussionpoint{零附近的表示揭示渐进下溢的作用。}
若计算结果从最小正规格化数继续趋近于零，本题格式会怎样？因为它既没有零也没有反规格化数，数轴在原点两侧留下空洞；IEEE 的反规格化数则固定指数并逐步减少有效首位，用丢失相对精度的代价保持固定的最小绝对间距，使结果能够渐进地下溢到零。若硬件采用 flush-to-zero，同一个微小差值会更早变成零，迭代停止条件、梯度传播和守恒量都可能随之改变。因此比较低精度格式时不能只列最大值，还应分别检查最小正规格化数、最小非零值、反规格化支持、零的符号及异常值语义。

\discussionpoint{低精度设计是在有限位预算内分配风险。}
同为 8 位或 6 位，为什么不同应用会选择不同的指数与尾数比例？训练梯度可能跨越许多数量级，更需要指数范围；量化后的权重若已按通道缩放，主体误差往往更受尾数密度支配，而块浮点还让一组数共享指数以节省元数据。输入格式也不能单独决定结果质量：乘积常以低精度产生，却在 FP16 或 FP32 中累加，以免长归约持续丢失低位。实际选型因而要结合张量分布、缩放粒度、饱和与非有限值计数、累加器精度和端到端任务指标；“再加一位尾数”只有在瓶颈确由量化噪声主导时才值得付出范围或存储成本。
\variantheading 若尾数再增加到 4 位，则每个指数区间有 16 个正数；$E=0$ 时依次为
$1,1+1/16,\ldots,1+15/16$，间距为 $1/16$，round-to-nearest 的最大绝对误差为
$1/32$。指数范围仍是 $-1,0,1$。
\practiceheading 设计自定义低精度格式时，应写一个穷举编码/解码测试，逐个位型验证单调性、
相邻间距、舍入边界和饱和行为。工程上还要给累加器单独选精度；低精度输入并不意味着应当
用同样低的精度累加。
\sourcecheck{https://standards.ieee.org/ieee/754/6210/}{IEEE 754-2019 官方标准页面；对抗检查：图 A.5 使用的是刻意简化的“无零”格式，不能把 IEEE 的零、反规格化数、无穷和 NaN 位模式混入本题枚举}
\end{solutionexercise}

\begin{solutionexercise}{2}
\problemheading 为另一种 6 位格式（1 位符号、2 位尾数、3 位指数）绘制与原书图 A.5
等价的图。用结果说明指数每增加一位，会如何改变数轴上的可表示数集合。沿用同一教学
“无零”约定：三位指数的偏置为 3，编码 111 保留，不采用零或反规格化数，只绘正数，负数
关于 0 对称；须列出全部有效指数区间，并区分动态范围变化和每个区间内的相对精度。

\approachheading 三位指数偏置为 3，编码 111 保留，所以实际指数为 $-3$ 到 $3$；两位
尾数给每个指数区间 4 个点。

\answerheading 正可表示数为
\[
\begin{array}{c|l}
E&\text{正可表示数}\\ \hline
-3&0.125,0.15625,0.1875,0.21875\\
-2&0.25,0.3125,0.375,0.4375\\
-1&0.5,0.625,0.75,0.875\\
0&1,1.25,1.5,1.75\\
1&2,2.5,3,3.5\\
2&4,5,6,7\\
3&8,10,12,14
\end{array}
\]
负数仍关于 0 对称，0 仍缺失。指数从 2 位增至 3 位后，有效指数层数从
$2^2-1=3$ 增至 $2^3-1=7$。相对于原书图 A.5 的 $E=-1,0,1$，新增
$E=-3,-2,2,3$ 四个区间：最小正规格化数从 $2^{-1}$ 降到 $2^{-3}$，缩小 4 倍；
最大正数从 $1.75\,2^1=3.5$ 扩到 $1.75\,2^3=14$，增大 4 倍。每个区间仍只有 4 个点，
所以相对精度不变。

\conclusionheading 指数位主要扩展动态范围；它不会增加同一二进制数量级内的采样密度。
\discussionheading
\discussionpoint{指数位以倍增编码换取指数级值域。}
给指数增加一位，为什么最大值不是简单地乘 2？指数有 $k$ 位时至多提供 $2^k$ 个编码；新增一位近似把可选指数层数翻倍，而每相邻一层又把数值尺度乘 2，于是端点按新增的指数跨度呈指数扩张。本题从三个有效层扩为七个层，最大指数由 1 到 3，故最大有限值乘 $2^{3-1}=4$；最小端点也按相反方向缩小 4 倍。偏置只决定编码区间放在数轴何处，保留编码、特殊值和反规格化规则才共同决定真实范围，这也是阅读任何硬件格式说明时必须逐项核对的边界。因此端点变化应从实际指数集合推导，而非凭位数口号估计。

\discussionpoint{固定位宽迫使范围与精度此消彼长。}
若芯片只能为每个数保留固定的总位数，多一位指数从哪里来？除符号位外，它通常要占用原本属于尾数的一位，使溢出和下溢变少，却把每个 binade 内的点数减半、相对舍入误差界加倍；本题暂时保持尾数位数不变，只是为了单独观察指数效果。具有少量离群值的数据可能需要宽范围，但若为这些离群值牺牲全部主体精度，整体均方误差反而会上升。格式选择因此是风险分配问题：既要统计极端值造成的饱和损失，也要估计高概率区间的量化噪声，并考虑误差经过后续归约或非线性算子后会怎样放大。

\discussionpoint{缩放让软件参与数值格式设计。}
硬件格式的指数不够时，软件是否只能接受溢出？损失缩放会在反向传播前放大梯度、结束后再除回去，张量或通道量化则保存 scale，把业务值域映射到编码值域；二者都等价于为一批数补充一个共享指数。设量化上界为 $q_{\max}$、观测绝对最大值为 $a$，常见初始尺度约取 $a/q_{\max}$，但单个异常值会让大多数编码挤在零附近。动态统计、异常值旁路和更细缩放粒度能缓解此问题，却增加元数据、同步与复现成本；因此工程上应把缩放策略视为数值算法的一部分，而非格式选定后的无关预处理。

\discussionpoint{FP8 变体说明格式必须匹配数据角色。}
为什么同一套模型可能同时使用不止一种 8 位浮点格式？范围宽、尖峰多的梯度更容易受溢出支配，倾向保留更多指数位；经过归一化或离线标定的权重与激活范围较窄，可能更需要尾数位来降低局部量化误差。即使输入都为 FP8，乘加的部分和通常仍提升到 FP16 或 FP32，因为数千项归约的舍入行为不同于一次乘法。部署验证应按张量角色记录饱和、下溢、非有限值及任务损失，并用真实层分布而非均匀随机数做压力测试；这种按数据角色分配格式的思想也延伸到混合精度求解器和分层存储系统。
\variantheading 若指数扩为 4 位、偏置取 7、编码 1111 保留，同时仍用两位尾数，则有效
指数为 $-7$ 到 $7$。最小正数为 $2^{-7}=1/128$，最大正数为
$1.75\times2^7=224$；每个 binade 仍只有 4 个点。
\practiceheading 混合精度训练和推理常用缩放因子把张量值域移入有限指数范围。上线前应记录
溢出、下溢和饱和计数，并把缩放策略与模型版本绑定；仅看平均误差可能掩盖少量但致命的
溢出值。
\sourcecheck{https://standards.ieee.org/ieee/754/6210/}{IEEE 754-2019 官方标准页面；对抗检查：不能简单说“范围翻倍”，新增一位指数使指数编码数翻倍，并使数值端点按 2 的新增指数跨度成倍扩展}
\end{solutionexercise}

\begin{solutionexercise}{3}
\problemheading 假设一种新处理器受技术困难限制，执行加法的浮点算术单元只能“向零
舍入”（通过向 0 截断数值来舍入）。硬件保留了足够多的位，唯一误差来自最终舍入。这台
机器上加法操作的最大 ULP 误差是多少？说明“最大值”是可达到值还是上确界，并指出溢出
等不属于题设模型的边界。

\approachheading 在同一 binade 内取两个相邻可表示数 $a<b$，间距定义为 1 ULP。向零
舍入总选择绝对值较小的一侧，而不是最近的一侧。

\answerheading 对正的精确结果 $x\in[a,b)$，结果为 $a$，故
\[
  0\le x-\operatorname{fl}_{0}(x)<b-a=1\ \mathrm{ULP}.
\]
负数情形对称。因而严格说误差的\textbf{上确界}是 1 ULP；任一有限、正常且在表示范围内
的非精确结果误差都严格小于 1 ULP，但可以任意逼近它。教材或硬件规格把上界写成“最大
1 ULP”时，指的就是这个上确界。相比之下，round-to-nearest 的对应上界是 0.5 ULP。

这个结论排除溢出；若精确和超出有限表示范围，错误不能再由局部 1 ULP 界描述。在最小
反规格化数区域，ULP 取固定的最小间距，同样适用。

\conclusionheading 向零舍入的误差界为 $<1$ ULP，通常报告为最大 1 ULP，而不是
0.5 ULP。
\discussionheading
\discussionpoint{ULP 是随数轴位置伸缩的局部刻度。}
“误差 1 ULP”能否换算成一个固定的小数？不能，因为 ULP 通常取结果附近两个相邻浮点数的间距：在同一 binade 内它保持为 $2^{E-p}$，越过二的幂边界后便翻倍；反规格化区则固定为最小间距。于是绝对误差同为 $10^{-6}$，在不同数量级上可能对应完全不同的 ULP 数，而相同的 1 ULP 也可能有不同绝对大小。ULP 特别适合检验实现是否返回附近的可表示数，但报告时必须说明参考点、格式及特殊区域；面向应用的误差预算还应同时给出绝对误差、相对误差或问题特定范数。

\discussionpoint{舍入模式同时决定误差大小和方向。}
两个相邻浮点数之间的精确结果应落向哪一端？向零舍入总选绝对值较小的一端，所以正数误差向下、负数误差向上，任一有限非精确结果的误差严格小于 1 ULP，但可无限逼近该界；向正无穷和负无穷则分别给出定向包围。区间算法会把同一表达式分别在向负无穷和向正无穷模式下求值，前者提供不大于精确值的下界，后者提供不小于精确值的上界，从而让真实结果即使经过多步运算仍被括住。最近偶数在非溢出区域把界缩为 0.5 ULP，并在恰逢中点时按末位奇偶打破平局，以避免总朝同一方向累积偏差。定向舍入因此适合区间算术和保守几何边界，而统计计算常偏好近似无偏的最近舍入，选择模式其实是在准确度、方向保证和硬件成本之间取舍。

\discussionpoint{局部舍入界不能替代问题条件分析。}
若每次加法都只错不到 1 ULP，整个算法是否就一定准确？设两个输入已分别带有微小误差，再计算 $10^8-10^8$ 一类严重抵消，结果的绝对误差可能仍小，但相对于真实的小差值却会非常大；局部正确舍入并不能恢复输入中已经丢失的信息。数值分析因而区分前向误差、后向误差与条件数：前两者描述实现偏离结果或等价输入多少，条件数描述问题本身对扰动有多敏感。工程上需要把硬件操作界沿算法传播，并用残差、缩放或稳定重写控制放大效应，不能把“每步 1 ULP”误读成端到端质量承诺。

\discussionpoint{舍入测试必须精确命中边界。}
随机生成一百万组加数，为什么仍可能漏掉舍入错误？最难的路径往往只出现在两个相邻浮点数的中点附近、binade 交界、反规格化边界或溢出前一格，随机实数几乎不会精确命中这些位型。可靠测试可用整数位模式构造相邻数，以多精度算术形成中点及其两侧的值，并分别覆盖正负号、精确结果、进位和 tie-to-even。还要防止编译器把乘加收缩为 FMA 或重排表达式，从而改变中间舍入点；把编译选项、实际指令和参考舍入模式纳入测试记录，才能区分算术单元错误与工具链改变语义。
\variantheading 若改为向 $+\infty$ 舍入，有限非精确结果的绝对误差上确界仍是 1 ULP，
但误差方向总为非负：正数向绝对值更大处移动，负数向绝对值更小处移动。若改为最近偶数
舍入，上确界降为 0.5 ULP。
\practiceheading 定向舍入常用于区间算术构造保守上下界。实现时必须确认编译器没有重排或
融合破坏舍入模式假设；CUDA 代码可使用带明确舍入后缀的算术内建函数，并用相邻浮点数测试
每个边界。
\sourcecheck{https://docs.nvidia.com/cuda/floating-point/index.html}{NVIDIA CUDA Floating Point 官方资料对 IEEE 舍入模式和 ULP 的说明；对抗检查：把“1 ULP”说成每个非精确结果都能恰好达到，或忽略溢出，均不严谨}
\end{solutionexercise}

\begin{solutionexercise}{4}
\problemheading 一名研究生编写了 CUDA 内核，把一个大型浮点数组归约为所有元素之和。
数组总是按数值从小到大排序。为避免分支分歧，他决定实现原书图 10.7 的算法。该低分歧
分配在首轮令前半数组的线程 $i$ 计算 $x_i+x_{i+N/2}$，以后继续折半。请解释这为何可能
降低结果的准确度，并把浮点加法次序、ULP 和结合律纳入说明。

\approachheading 原书图 10.7 首轮把前半数组与后半数组逐项配对，而不是让相邻元素配对。
对已排序的非负数据，这会过早地把小数加到大数上。

\answerheading 设数组长 $2n$。该分配首轮计算
\[
  x_0+x_n,\ x_1+x_{n+1},\ldots,x_{n-1}+x_{2n-1}.
\]
若 $x_i$ 与 $x_{n+i}$ 的指数差很大，较小加数在对齐尾数后可能完全落到保留位之外；当它
小于大数处的半个 ULP 时，round-to-nearest 甚至直接返回大数。大量小贡献会在尚未彼此
聚合时分别丢失。若先把数量级接近的相邻小数做平衡树归约，小数先形成较大的部分和，之后
再与大数相加，通常能保留更多有效位。补偿求和或按绝对值分桶还能进一步改善。

这里的“降低”是可能性而不是逐个输入的定理：若数值同量级、加法恰好可表示，或正负抵消
模式不同，两棵归约树可能相同或各有优劣。浮点加法不满足结合律，所以改变配对次序足以改变
舍入路径和最终 ULP；算法仍在实数算术意义下计算同一个和。

\editorialnote 英文底本把本题写成“图 A.4 的算法”，但原书图 A.4 是可表示数表，并非
归约算法。结合“避免分支分歧”和第 10 章内容，实际指向原书图 10.7；中文译稿已作同样修正。

\complexity{两种树形归约都是 $O(N)$ 工作、$O(\log N)$ 深度}{$O(N)$ 原地数据或 $O(1)$ 每线程额外状态}{低分歧映射提高 SIMT 利用率，但线程映射并不保证数值稳定性}
\conclusionheading 低控制分歧方案用数值上相距较远的操作数换取连续活跃线程；对升序
非负数据，这一配对通常不利于保留小量贡献。
\discussionheading
\discussionpoint{归约树同时改变关键路径与舍入路径。}
同一批数只是换一棵加法树，为何性能和答案都会变化？顺序累加有 $O(N)$ 的依赖深度，平衡树将其降为 $O(\log N)$，并常把一阶误差增长从约 $O(uN)$ 降到 $O(u\log N)$；这里的 $u$ 是单位舍入误差，但这只是典型上界而非逐输入保证。原书的低分歧映射让更多线程连续活跃，却在首轮把已排序数组的前后半段配对，使小量可能分别消失在大数的 ULP 下。并行数值算法因此不能只优化工作量和占用率，还要把数据到叶节点的映射视为算法语义，并用高精度参考比较不同树的误差分布。

\discussionpoint{稳定求和方法形成精度与成本的谱系。}
如果普通平衡树仍不够准确，还有哪些选择？成对求和保留并行结构；Kahan 或 Neumaier 求和为被舍掉的低位维护补偿量；按指数分桶、超累加器或长定点累加则能把大量项延迟到更精确的表示中合并，甚至支持正确舍入。代价依次可能表现为更多指令、寄存器和共享内存、更复杂的合并状态，或较低的 SIMD/SIMT 利用率，而且补偿更新本身也必须按规定次序实现。数据库聚合、深度学习梯度和科学守恒量对准确度预算不同，合理方案是在给定吞吐、存储和确定性约束下选择谱系中的一点，而非默认最高精度总是最优。

\discussionpoint{排序只有结合符号和条件数才有意义。}
“从小到大相加更准确”是否能无条件套用？对全非负数据，按绝对值从小到大通常让小量先形成可见的部分和；含正负数时，按代数值升序会先聚合大负数，再聚合大正数，可能产生很大的中间量，最后一次抵消仍会暴露输入误差。按绝对值排序也不是万能定理，因为有意让相反符号的近等量先抵消有时反而减少中间范围，效果取决于和的条件数与分布。排序还需额外的 $O(N\log N)$ 工作和数据搬移，所以工程上更常采用分桶或局部重排，并以真实数据、高精度参考和残差共同决定其收益。

\discussionpoint{可复现性与高准确度是两个独立目标。}
同一归约多跑一次得到末位不同，是否就说明某次更不准确？原子更新、动态任务调度或不同分区可能改变合法的加法顺序，从而得到多个都接近参考值的结果；固定树可保证逐位复现，却可能稳定地重复一个误差较大的配对。反过来，补偿或高精度累加能改善误差，但若合并次序仍不固定，结果依旧未必 bitwise 相同。审计、回归定位和某些迭代求解器重视确定性，探索性训练可能更重视吞吐与统计等价；接口应分别声明顺序契约、误差界和性能，而实验记录还应固定分区、工具链及设备拓扑。
\variantheading 在 binary32 中取 $[2^{-24},2^{-24},1]$。若先分别把两个小数加到 1，
每次都落在 tie 并按偶数舍回 1；若先算两个小数，得到 $2^{-23}$，再加 1 得到可表示的
$1+2^{-23}$。两种次序相差 1 ULP，直接展示结合律失效。
\practiceheading 需要跨运行或跨 GPU 可复现时，应固定归约树、块大小和分区顺序，并把
确定性作为接口契约；只要求统计准确时，可在性能预算内选择分桶、Kahan/Neumaier 补偿或
更高精度累加器，并用高精度参考计算 ULP 分布。
\sourcecheck{https://docs.nvidia.com/cuda/floating-point/index.html}{NVIDIA CUDA Floating Point 官方资料对加法非结合性和并行归约差异的说明；对抗检查：排序应按绝对值理解数值稳定性，含正负数据时仅按代数值升序并不能保证相邻项同量级}
\end{solutionexercise}

\begin{solutionexercise}{5}
\problemheading 假设某算术单元以迭代近似算法实现 \code{sin()} 函数，每个时钟周期为
结果再产生 2 个准确的尾数位。架构师决定让该算术函数迭代 9 个时钟周期。假设硬件把剩余
尾数位全部填为 0。对于 IEEE 单精度数，这一 \code{sin()} 硬件实现的最大 ULP 误差是
多少？假设被省略的 ``$1.$'' 尾数位也必须由算术单元产生；亦即 binary32 的有效精度按
隐含首位加 23 个显式 fraction 位共 24 位计。

\approachheading IEEE binary32 的有效数是 24 位：1 个隐含首位和 23 个显式 fraction 位。
九周期只确定 $2\times9=18$ 个最高有效位。

\answerheading 未确定并清零的低位数为
\[
  24-18=6.
\]
硬件结果的相邻量化级相隔 $2^6=64$ 个 binary32 末位单位。因此，以精确实数为参照，
向下清零造成的误差满足
\[
  0\le |\widehat{s}-s|<64\ \mathrm{ULP},
\]
上确界为 64 ULP。若参照“正确舍入到 binary32”的结果，精确值尾部接近下一个粗量化级时，
正确舍入会进位而该硬件仍截断到当前级，二者可相差整整 64 ULP；所以题目通常要求的最大值
是 \textbf{64 ULP}。只观察一个未进位的正确舍入尾数字段会得到 63 ULP，但那漏掉了跨粗
量化边界的进位情形。

该推导假设正常有限结果，并以结果所在 binade 的 binary32 间距定义 ULP。零点附近的
反规格化处理、精确的 $\sin$ 参数归约误差以及 NaN 都不在题设的“唯一误差来自迭代截断”
模型内。

\conclusionheading 18 个准确有效位比 binary32 的 24 位少 6 位，故最大误差规格为
$2^6=64$ ULP（对实数误差是严格小于该值的上确界）。
\discussionheading
\discussionpoint{初等函数的误差来自一整条实现管线。}
把低 6 位清零得到 64 ULP 后，能否把它直接当作任意 \code{sin} 实现的总误差？实际函数通常先判别特殊值，再做周期和象限的参数归约，在小区间使用查表、多项式或迭代近似，随后重构符号并舍入；近似系数、有限精度运算与末端舍入都会增加误差。本题特意假定前 18 位完全准确，因此只隔离了最后截断这一环，不能涵盖前面任何误差。数学库设计常为各阶段分配独立预算，再用组合界或穷举验证总误差；这种分层方法也适用于指数、对数和倒数平方根等硬件函数。如果前级预算未受控，末端多迭代一轮也无法给总误差兜底。

\discussionpoint{巨大输入让参数归约成为主导难题。}
为什么一个看似简单的周期函数会为巨大 $x$ 准备慢速路径？计算 \code{sin(x)} 前必须确定 $x$ 属于哪个 $\pi/2$ 区间并求出很小的余量；当 $x$ 很大时，普通精度的 $x-k\pi/2$ 是两个大数相减，常数低位或整数 $k$ 的一位错误都可能选错象限，其误差远大于末端 64 ULP。高质量实现会保存更多位的 $2/\pi$，以整数式乘法提取象限，并仅在小参数上使用快速近似。快速路径与精确慢速路径的分界需要兼顾延迟、代码体积和最坏误差，也说明输入范围是函数准确度契约不可省略的条件。

\discussionpoint{数学库必须声明准确度契约。}
用户看到同名的 \code{sin}，能否假定不同编译选项都给出相同末位？正确舍入要求结果等于精确实数按目标模式舍入后的唯一浮点数；faithful rounding 通常只保证落在包围精确值的相邻候选之一，而快速数学路径可能仅承诺更宽的 ULP 或相对误差，并放宽 NaN、反规格化或特殊值行为。前者需要昂贵的难例处理，后者可换取吞吐和更短延迟。金融定价、可证明科学模拟与机器学习推理承担的风险不同，API、编译标志和测试报告都应明确契约，调用方再决定是否以性能交换准确度。

\discussionpoint{最坏误差验证不能依赖随机平均。}
随机输入的平均 ULP 很小，是否足以证明 64 ULP 上界没有被突破？高误差点可能集中在零点、极值、binade 边界、舍入中点和参数归约换路处，随机采样命中这些狭窄集合的概率极低；应使用位型枚举、邻接浮点数及多精度参考主动搜索，并覆盖反规格化与非有限值。对“只截取精确展开的前 24 位”这一变体，必须构造正确舍入向上进位的例子，才能验证它与正确舍入结果至多相差 1 ULP、且该界确实可达到。最后还要分别测量指令延迟和流水吞吐，因为九周期延迟并不意味着执行单元每九周期才接受一个新操作。
\variantheading 若迭代 10 个周期，则得到 20 个准确有效位，还缺 $24-20=4$ 位，最大
规格变为 $2^4=16$ ULP；若迭代 12 个周期得到的是已经正确舍入的 24 位 binary32 结果，
则末端清零误差为 0。若硬件只是截取精确展开的前 24 位而没有 guard/sticky 位，仍可能与
正确舍入结果相差至多 1 ULP，并且在正确舍入选择相邻上端点时可以恰差 1 ULP；“得到 24 位”
本身不足以推出零 ULP。
\practiceheading 选择快速数学函数时，应在业务输入域上同时测延迟、吞吐量和最大/分位 ULP，
并专门覆盖零点、极值、巨大参数与 NaN/无穷。编译选项和 GPU 架构会改变实现路径，验证结果
应与工具链及目标架构一起归档。
\sourcecheck{https://standards.ieee.org/ieee/754/6210/}{IEEE 754-2019 官方标准页面；对抗检查：若忘记隐含首位会误算为少 5 位并得到 32 ULP，若把截断尾部最大位型机械写成 63 又会漏掉正确舍入的进位}
\end{solutionexercise}

\chapter{深度学习基础}
\noexercises

\chapter{CUDA 内存、地址空间与指针}
\noexercises


\backmatter
\chapter*{验证说明}
\addcontentsline{toc}{chapter}{验证说明}
数值题使用独立算式复核；数组与图算法题使用小规模逐步状态复核；CUDA 题检查索引范围、
同步点、竞态条件、分配与释放生命周期，并给出 CPU 参考路径或可手算样例。当前构建环境用
CUDA Toolkit 13.3 的 \code{nvcc} 与 MSVC 14.44 编译链接全部 27 个 CUDA 程序，并实际运行
全部 27 个 \code{--cpu-only} 参考路径；普通 C++ 与两秩 MPI 的 CTest 为 6/6。另以
\code{sm_90} 编译 Smith--Waterman，检查 DPX 条件分支。本机未检测到 NVIDIA GPU，因此不把
CPU 参考运行误报为设备端运行。

\end{document}
